% Generated by build_chapter.py from the GY1980 files. Do not edit by hand; edit the files and rebuild.
\chapter{The academic year 1980--81}\label{ch:1980}
\noindent\textbf{Scope and source order.}
This chapter covers all 36 records in the supplied 1980--81 manifest and restores three omitted records, 510--512: the archive catalog dates them 23 February 1991, while the autograph P000960 reads 23 February 1981. The resulting 39 records comprise 38 talks or interventions and one catalog year divider. They are ordered by first appearance in the manuscript, including Gelfand's introductions before Beilinson and the Gelfand--Cherednik sequence on P000960. Overlapping catalog page ranges have been replaced by observed footprints. P000916 is a preceding-year continuation already indexed as 1979-488.1--.7; it is accounted for without being counted again (\citealp{mccme2026gelfand}).

\noindent\textbf{What the year contains.}
The 246 mathematical occurrences range from classroom exercises to announcements. Autumn's intersection homology, Segre families and regular holonomic systems lead into Beilinson's sustained exposition of localization, transfer bimodules and derived direct image. The chronology matters: an announcement circulating before a talk is distinguished from a later full proof, and a later geometric implementation does not automatically make an old formula new. The categorical affinity question is followed from its November introduction to its answer in the continuation.

Novikov's December lecture combines already published Pauli zero modes with the topological discussion appearing in a paper received the following day. The physical line bundle has degree \(N\); the reciprocal-torus zero-mode bundle has rank \(N\) and Chern number of absolute value one. The notebook conflates these numbers. MacPherson's Grassmannian constructions lead to positive torus orbits, harmonic forms and convexity; Khovanskii and Kazarnovskii connect Newton polyhedra, toric geometry and zero counting. Nondegeneracy, regularity and normalization conditions are restored explicitly.

Spring brings Yang--Baxter equations, commuting differential operators, algebraic geometry and \(K\)-theory, Hamiltonian kinetic theory, eigenvalue inequalities, gas dynamics and Lie semigroups. The transfer-matrix argument needs the correct auxiliary trace legs; the Weyl-operator examples need noncommutative calculation; the kinetic constructions need their actual functional and coefficient classes. The classification records the fate of the qualified statement and separately assesses the literal notebook claim.

\begin{center}
\begin{tabular}{lrrrrrr}
\toprule
Class & A & B & C & D & E & F\\
\midrule
Occurrences & 42 & 13 & 1 & 0 & 165 & 25\\
\bottomrule
\end{tabular}
\end{center}
A denotes a subsequent speaker publication; B a later development by others; C a partly developed precise question; D an undeveloped identified question; E prior knowledge; F an insufficiently identified statement. There are 16 problem identities, 94 bibliography entries and 90 correction rows. Counts describe occurrences, so a theorem repeated in several talks is counted several times. No D classification is justified by the search evidence.

\noindent\textbf{The bounded development and its limits.}
The remaining class-C item is Gelfand's conjecture that every complex algebraic curve's function field embeds in the first Weyl division ring. All genus-one and genus-two fields are positive cases: in particular Oganesyan's 2016 preprint, published in 2018, realizes every standard genus-two quintic by rational-coefficient commuting operators (\citealp[Theorem~3]{oganesyan2016rational}). The genus-two item is therefore B. The development proves an elementary obstruction for a prescribed positive-genus affine model with a nonconstant unit, supplies a complete exact certificate for a known genus-two family, and proves the injectivity step from operators to a function field. These are bounded clarifications, not a solution of the universal conjecture or a novelty claim for the known family. Independent checks caught both a notebook formal-expression error and a scalar normalization in the checked Mironov preprint.

\noindent\textbf{Continuity and evidence.}
The unextended Gelfand--Kirillov conjecture recurs under its existing identity GP-1967-21; its general form is false, with surviving simple-type cases stated cautiously. The preceding year's collective-field remarks connect to the Vlasov and hydrodynamic material without being equated with its later precise brackets. The field-embedding question is kept separate from automorphism and different-rank Weyl-field problems. The report and research ledgers identify full-text, metadata-only and unresolved evidence; a verified bibliographic link alone is never presented as a verified proof.

\GStalk{490}{1980-09-22}{D. B. Fuchs}{P000917-P000918}{Intersection homology, duality, and the L2 comparison}
The dated Fuchs heading begins on P000917; the first line of P000918 finishes the discussion of Deligne’s construction. The talk explains allowable chains, Poincar\'e duality, the two-sphere nodal example, a proposed comparison with \(L^2\)-cohomology, and isolated-singularity truncation. This continues the subject indexed in 1979-484.1--4. The catalogue’s P000916 is an undated continuation of Zakharov’s Benney-equation discussion already indexed as 1979-488.1--7: hydrostatic equations, height conservation, additional conservation laws, a Hamiltonian layer formulation, a two-horizontal-dimensional extension, kinetic interpretation, and Gelfand’s collective-equation remark. It is covered here by this explicit boundary cross-reference, without creating a second occurrence or attributing it to Fuchs.

\begin{GSitem}{1980-490.1}{E}{Stratifications and middle-perversity allowable chains}
\GSfield{Kind; source}{definition; \GSpage{P000917}}
\GSfield{The notes}{The notes give \(X=X_{2n}\supset X_{2n-2}\supset X_{n-3}\supset\cdots\), locally constant stratum geometry, and \(\dim(\sigma_j\cap X_{2n-i})=j-2i\).}
\GSfield{Statement}{Let \(X\) be a pure complex \(n\)-dimensional variety with a conically locally trivial stratification by even real codimension. For middle perversity \(\bar m(c)=\lfloor(c-2)/2\rfloor\), a PL \(j\)-chain \(\xi\) is allowable if \(\dim(|\xi|\cap X_{2n-c})\le j-c+\bar m(c)\), and its boundary is allowable with \(j\) replaced by \(j-1\). Their chain complex defines \(IH_j^{\bar m}(X)\).}
\GSfield{Reconstruction}{The inequality, boundary condition, dimension convention, and perversity are supplied explicitly; the inconsistent displayed formula is not silently repaired.}
\GSfield{Fate}{Already announced in 1977 and published before this talk: \citet[\S\S1.1,1.3]{goreskymacpherson1980intersection}. The same construction occurs in 1979-484.1.}
\GSfield{Present status}{As of November 2026, the construction is standard; the displayed equality and codimension indices are not a correct definition.}
\GSfield{Assessment}{literal statement: false as written; publication match: exact; verification: read the result; development: none; confidence: high}
\end{GSitem}
\begin{GSitem}{1980-490.2}{E}{Poincare duality for intersection homology}
\GSfield{Kind; source}{theorem; \GSpage{P000917}}
\GSfield{The notes}{\ru{В этих когом. есть двойств. Пуанкаре.}}
\GSfield{Statement}{For a compact oriented real \(d\)-dimensional stratified pseudomanifold and complementary perversities \(\bar p(c)+\bar q(c)=c-2\), the intersection pairing \(IH_j^{\bar p}(X;\mathbb Q)\otimes IH_{d-j}^{\bar q}(X;\mathbb Q)\to\mathbb Q\) is perfect. On complex varieties the even-codimension stratification makes middle perversity self-complementary on the relevant strata.}
\GSfield{Fate}{The duality theorem is \citet[\S3.3]{goreskymacpherson1980intersection}; no new 1980 conjecture is being posed.}
\GSfield{Present status}{As of November 2026, duality holds with these compactness, orientation, coefficient and perversity qualifications.}
\GSfield{Assessment}{literal statement: correct with missing hypotheses; publication match: exact; verification: read the result; development: none; confidence: high}
\end{GSitem}
\begin{GSitem}{1980-490.3}{E}{The nodal union of two spheres and failure of homotopy invariance}
\GSfield{Kind; source}{construction; \GSpage{P000917}}
\GSfield{The notes}{Two \(2\)-spheres are joined at one point. Ordinary groups in degrees \(0,1,2\) are \(\mathbb Z,0,\mathbb Z^2\); the MacPherson groups are \(\mathbb Z^2,0,\mathbb Z^2\). A path through the node is forbidden.}
\GSfield{Statement}{Stratify \(X=S^2\vee S^2\) by its node and smooth complement. Middle-perversity intersection homology is the homology of the disjoint normalization \(S^2\sqcup S^2\), hence the displayed groups. Intersection homology is not invariant under arbitrary homotopy equivalence. It is invariant under homeomorphism of suitable pseudomanifolds; the two statements are compatible.}
\GSfield{Fate}{This is the elementary nonnormal-curve model of the allowable-chain theory; topological invariance is \citet[\S\S3.2,6.2]{goreskymacpherson1980intersection}. The group calculation follows directly from the chain restriction.}
\GSfield{Present status}{As of November 2026, the displayed groups are correct; the example demonstrates failure of ordinary duality, while general failure of homotopy invariance requires the distinction between homotopy and homeomorphism.}
\GSfield{Assessment}{literal statement: correct; publication match: exact; verification: read the result; development: annotated; confidence: high}
\end{GSitem}
\begin{GSitem}{1980-490.4}{B}{Intersection cohomology versus L2 cohomology}
\GSfield{Kind; source}{conjecture; \GSpage{P000917}}
\GSfield{The notes}{\ru{Гипотеза: совпадают с \(L_2\)-когомологиями.} Borel’s interest is mentioned, but no metric, domain of the differential, compactification or arithmetic hypothesis is specified.}
\GSfield{Statement}{A precise later specialization is the Zucker comparison: for an arithmetic quotient \(X^\circ=\Gamma\backslash D\) of a Hermitian symmetric domain, its Baily--Borel compactification \(X\), and the metrized local system attached to a finite-dimensional algebraic representation, the sheaf of locally square-integrable forms with square-integrable differential represents the middle intersection complex on \(X\). This is a formulation from the literature, not an assertion that the notebook supplies those hypotheses.}
\GSfield{Reconstruction}{The arithmetic setting and differential domain are explicitly later hypotheses.}
\GSfield{Fate}{\citet[theorem on p.~3]{looijenga1988l2} proves this arithmetic specialization; independent work of Saper--Stern was announced in February 1987 and published in 1990. The missing metric/domain data prevent an exact match of the notebook sentence to that theorem.}
\GSfield{Present status}{As of November 2026, the cited specialization is proved. The sentence in the notebook is underdetermined as a general mathematical conjecture and is not entered as a newly open problem.}
\GSfield{Assessment}{literal statement: correct with missing hypotheses; publication match: partial; verification: read the result; development: none; confidence: medium}
\end{GSitem}
\begin{GSitem}{1980-490.5}{E}{Isolated singularities and Deligne truncation}
\GSfield{Kind; source}{method; \GSpage{P000917-P000918}}
\GSfield{The notes}{For one isolated singularity the notes say that above half the dimension the groups agree with ordinary groups, and below it they follow by duality; the next page calls Deligne’s general definition inductive.}
\GSfield{Statement}{For a compact pure complex \(n\)-fold with isolated singularities, rational middle intersection homology satisfies \(IH_j(X)=H_j(X\setminus\Sigma)\) for \(j<n\), \(IH_n(X)=\operatorname{im}[H_n(X\setminus\Sigma)\to H_n(X)]\), and \(IH_j(X)=H_j(X)\) for \(j>n\). In sheaf language the intersection complex is obtained from the constant sheaf on the regular stratum by successive derived direct images and perversity truncations.}
\GSfield{Reconstruction}{The precise isolated-singularity formula and the middle-degree clause are editorial reconstruction.}
\GSfield{Fate}{The sheaf formulation was proposed by Deligne in 1976 and is formalized in \citet[\S\S1--3]{goreskymacpherson1983intersection}; the isolated-singularity formula is the allowable-chain calculation. See also \citet[\S6.1, pp.~158--159]{goreskymacpherson1980intersection} for the isolated-singularity chain computation.}
\GSfield{Present status}{As of November 2026, the middle-degree image is essential. The notebook alternates homology and cohomology terminology, so the displayed high-degree assertion is read homologically.}
\GSfield{Assessment}{literal statement: correct with missing hypotheses; publication match: partial; verification: indirect; development: none; confidence: high}
\end{GSitem}

\GStalk{491}{1980-09-22}{M. I. Graev}{P000918}{Admissible families of curves and local integral geometry}
Graev’s actual heading is on P000918. He discusses weighted curve transforms, differential forms on curves through a point, the distinction between closure and point-supported inversion when that incidence fiber is one-dimensional, and quotienting trivial exact forms. The two explicit FA references are identified separately: Gelfand--Gindikin--Shapiro, 13(2), 1979, and Gelfand--Graev, 14(4), 1980; the second does not have the same author list. The preceding Fuchs and Zakharov material is not reassigned to Graev.

\begin{GSitem}{1980-491.1}{E}{The weighted curve transform and incidence fiber}
\GSfield{Kind; source}{construction; \GSpage{P000918}}
\GSfield{The notes}{An \((n+k-1)\)-parameter family is written \(x_i=\varphi_i(x_1,\xi)\), with positive weight \(\chi\), and \(I_\chi f(\xi)=\int f(x_1,\varphi_1,\ldots,\varphi_n)\chi\,dx_1\). The set of curves through \(x^0\) is \(\Gamma_{x^0}\), of dimension \(k\).}
\GSfield{Statement}{On a domain \(D\subset\mathbb R^n\), let a smooth incidence submanifold \(Z\subset D\times\Gamma\), \(\dim\Gamma=n+k-1\), have submersive projections and one-dimensional fibers over \(\Gamma\). Locally write \(x_i=\varphi_i(x_1,\xi)\), \(2\le i\le n\). For \(f\in C_c^\infty(D)\), set \(I_\chi f(\xi)=\int f(x_1,\varphi_2,\ldots,\varphi_n)\chi(x_1,\xi)\,dx_1\). Transversality gives \(\dim\Gamma_x=k\).}
\GSfield{Fate}{The general local framework is in \citet[\S1]{gelfandgindikinshapiro1979local}; the \(k=1\) setting is written explicitly in \citet[\S1, equations (1)--(4)]{gelfandgraev1980admissible}, received 16 June 1980 and published in issue 14(4).}
\GSfield{Present status}{As of November 2026, this is a standard double-fibration transform. The redundant coordinate and parameter indices in the notebook require correction.}
\GSfield{Assessment}{literal statement: false as written; publication match: exact; verification: read the result; development: none; confidence: high}
\end{GSitem}
\begin{GSitem}{1980-491.2}{E}{Local inversion from curves through a point}
\GSfield{Kind; source}{problem; \GSpage{P000918}}
\GSfield{The notes}{Find families permitting recovery using integrals over curves through the point or infinitesimally nearby; these are called admissible.}
\GSfield{Statement}{For the preceding incidence family, seek a differential operator on transform data and cycles in \(\Gamma_x\) whose periods recover \(f(x)\), or a nonzero distribution supported at \(x\). Closure of a form on \(\Gamma_x\) alone is insufficient: its relevant period must be nonzero and normalized.}
\GSfield{Fate}{\citet[introduction and \S1]{gelfandgindikinshapiro1979local} explicitly distinguishes closed range forms from nonvanishing inversion periods; it also explains the real/complex distinction. This is the published research program being explained, not evidence of an unsolved theorem with stated quantifiers.}
\GSfield{Present status}{As of November 2026, the construction is established in the cited admissible settings. No blanket inversion theorem for arbitrary curve families is asserted.}
\GSfield{Assessment}{literal statement: correct with missing hypotheses; publication match: partial; verification: read the result; development: none; confidence: high}
\end{GSitem}
\begin{GSitem}{1980-491.3}{E}{Admissible first-order differential forms}
\GSfield{Kind; source}{definition; \GSpage{P000918}}
\GSfield{The notes}{\(\mathcal H_xF=\sum a_{ij}(x_1,\xi)\partial_{\xi_j}F\,d\xi_i+F\sum b_i(x_1,\xi)d\xi_i\). Admissibility means that \(\mathcal H_xI_\chi f\) is closed for every \(f\); for \(k>1\) this is said to be the complete definition.}
\GSfield{Statement}{An operator taking data functions to differential one-forms restricted to \(\Gamma_x\) is admissible when \(d_{\Gamma_x}(\mathcal H_x I_\chi f)=0\) for all test functions \(f\). Nontriviality is a separate condition after quotienting exact operators; local recovery also needs a nonzero period.}
\GSfield{Fate}{This is the closed-form approach of \citet[\S1]{gelfandgindikinshapiro1979local}. The notebook’s closure criterion is meaningful in dimensions \(k>1\), but does not itself imply an inversion formula.}
\GSfield{Present status}{As of November 2026, the definition is mathematically meaningful once the omitted target \(\Gamma_x\) and the distinction from inversion are restored.}
\GSfield{Assessment}{literal statement: correct with missing hypotheses; publication match: exact; verification: read the result; development: none; confidence: high}
\end{GSitem}
\begin{GSitem}{1980-491.4}{A}{The one-dimensional incidence case and support at a point}
\GSfield{Kind; source}{definition; \GSpage{P000918}}
\GSfield{The notes}{For \(k=1\), every one-form is closed; require \(f\mapsto\int\mathcal H_xI_\chi f\) to be concentrated at \(x\).}
\GSfield{Statement}{For \(\dim\Gamma_x=1\), the published definition requires the period functional to vanish for test functions vanishing near \(x\). In the specific first-order incidence setting of Gelfand--Graev this implies \(L_xI_\chi f=c_x f(x)\). The authors further point out that for real curves the local inversion constant is necessarily zero; the support condition is therefore not itself a nonzero inversion formula.}
\GSfield{Fate}{The exactly contemporary treatment is \citet[\S2, Definition 1 and the following remark]{gelfandgraev1980admissible}; received 16 June 1980, published in issue 14(4). Its point-evaluation conclusion is stronger than the general distribution-theoretic fact that a point-supported distribution can contain delta derivatives.}
\GSfield{Present status}{As of November 2026, the precise published support condition is established; nonvanishing, the real/complex setting and the choice of cycle remain essential to an inversion assertion.}
\GSfield{Assessment}{literal statement: correct with missing hypotheses; publication match: exact; verification: read the result; development: none; confidence: high}
\end{GSitem}
\begin{GSitem}{1980-491.5}{E}{Quotient by trivial exact forms}
\GSfield{Kind; source}{construction; \GSpage{P000918}}
\GSfield{The notes}{\(\mathcal H_xF=d(\mu F)\) is called trivial; the quotient by such functionals is important.}
\GSfield{Statement}{For a closed one-cycle \(\gamma\subset\Gamma_x\) and a single-valued smooth \(\mu F\), \(\int_\gamma d(\mu F)=0\). Hence one quotients admissible operators by these exact operators when studying nonzero period functionals.}
\GSfield{Fate}{The exact-form quotient is part of the local integral-geometry construction in \citet[\S1]{gelfandgindikinshapiro1979local}; the zero-period assertion is Stokes' theorem.}
\GSfield{Present status}{As of November 2026, the quotient removes operators that cannot contribute a period inversion formula.}
\GSfield{Assessment}{literal statement: correct; publication match: exact; verification: read the result; development: none; confidence: high}
\end{GSitem}

\GStalk{492}{1980-10-13}{I. M. Gelfand}{P000919-P000921}{Double fibrations, differential equations, and Segre structures}
The actual talk occupies P000919--P000921. It passes from a two-parameter family of plane curves and its dual differential equation to incidence double fibrations, homogeneous examples, separate-coordinate equivalence, and the coframe definition of Segre structures. The explicit Inventiones reference is Faran’s 1980 paper. P000922 begins the separately dated 20 October session, so its Goncharov, Fuchs and Beilinson interventions belong to record 493. The second-order-equation discussion recurs from 1976-403.1--5.

\begin{GSitem}{1980-492.1}{E}{A second-order equation and its dual family}
\GSfield{Kind; source}{construction; \GSpage{P000919}}
\GSfield{The notes}{Starting with \(y''=F(x,y,y')\), the notes write a general integral \(y=\varphi(x,\xi,\eta)\), then exchange \((x,y)\) with \((\xi,\eta)\).}
\GSfield{Statement}{A regular scalar second-order ODE determines locally a two-parameter family of solution graphs. Conversely, if \(\det\partial(\varphi,\varphi_x)/\partial(\xi,\eta)\ne0\), the parameters can be eliminated from \(y=\varphi\), \(p=\varphi_x\), producing \(y''=F(x,y,p)\). A dual scalar ODE is obtained by regarding the incident parameters as curves, provided the dual projection has the corresponding nondegeneracy and graph property.}
\GSfield{Fate}{This classical incidence interpretation and duality are treated in \citet[\S\S22--23]{cartan1924projective}. It is an exact subject recurrence of 1976-403.1 and 1976-403.3, not a new 1980 proposal.}
\GSfield{Present status}{As of November 2026, the correspondence is local; arbitrary two-parameter families can fail the determinant condition.}
\GSfield{Assessment}{literal statement: correct with missing hypotheses; publication match: exact; verification: read the result; development: none; confidence: high}
\end{GSitem}
\begin{GSitem}{1980-492.2}{E}{Straightening a family of curves by point transformations}
\GSfield{Kind; source}{problem; \GSpage{P000919}}
\GSfield{The notes}{Can the family be reduced to \(\xi x+\eta y=1\) by independent changes in the point and parameter planes?}
\GSfield{Statement}{For a nondegenerate local path geometry, ask whether a local diffeomorphism of the point plane takes its unparametrized paths to projective lines; a separate parameter change merely relabels the paths. The criterion is vanishing of Cartan’s projective curvature invariants for the corresponding second-order ODE. The displayed affine line chart omits lines outside that chart.}
\GSfield{Fate}{Solved in the classical equivalence theory: \citet[\S\S22--24]{cartan1924projective}. The same question is 1976-403.1; the notebook supplies no new unrestricted linearization theorem.}
\GSfield{Present status}{As of November 2026, the local criterion is established. A general equation is not point-equivalent to the flat line family.}
\GSfield{Assessment}{literal statement: correct with missing hypotheses; publication match: exact; verification: read the result; development: none; confidence: high}
\end{GSitem}
\begin{GSitem}{1980-492.3}{E}{Incidence manifolds and double fibrations}
\GSfield{Kind; source}{definition; \GSpage{P000919-P000920}}
\GSfield{The notes}{There are maps \(A\to B\), \(A\to\Gamma\), and a pair of projections is required to specify a point of \(A\) uniquely. A hypersurface \(f(z,w)=0\) is considered modulo separate transformations in \(z,w\).}
\GSfield{Statement}{A regular incidence geometry is a smooth submanifold \(Z\subset B\times\Gamma\) with submersive projections. Joint injectivity of two maps \(A\to B,\Gamma\) identifies their points set-theoretically; to identify \(A\) smoothly with a submanifold one also needs an immersion and appropriate embedding condition. Equivalence consists of local product maps \((b,\gamma)\mapsto(F(b),G(\gamma))\) preserving \(Z\).}
\GSfield{Reconstruction}{The embedding hypotheses distinguish the smooth formulation from the set-theoretic uniqueness in the notes.}
\GSfield{Fate}{The differential-equation case is classical in \citet[\S\S22--23]{cartan1924projective}; local integral geometry uses the same incidence organization in \citet[\S1]{gelfandgindikinshapiro1979local}.}
\GSfield{Present status}{As of November 2026, the formulation is standard with the regularity conditions made explicit.}
\GSfield{Assessment}{literal statement: correct with missing hypotheses; publication match: exact; verification: read the result; development: none; confidence: high}
\end{GSitem}
\begin{GSitem}{1980-492.4}{E}{Homogeneous double fibrations and the intersection stabilizer}
\GSfield{Kind; source}{construction; \GSpage{P000920}}
\GSfield{The notes}{Take \(A=G\), \(B=H\backslash G\), and \(\Gamma=K\backslash G\); matrix examples use unitary and triangular subgroups.}
\GSfield{Statement}{For closed subgroups \(H,K\) of a Lie group \(G\), the natural joint map from \(G\) has fibers the left cosets of \(H\cap K\). The effective incidence manifold is \((H\cap K)\backslash G\), with its two quotient maps; the map from \(G\) itself is injective only if \(H\cap K=\{1\}\). Smooth embeddedness of its image needs a further regularity hypothesis.}
\GSfield{Fate}{This is the standard homogeneous incidence construction underlying the local integral-geometric framework of \citet[\S1]{gelfandgindikinshapiro1979local}; the stabilizer correction follows directly from equality of the two cosets.}
\GSfield{Present status}{As of November 2026, the construction is valid after accounting for the common stabilizer.}
\GSfield{Assessment}{literal statement: false as written; publication match: partial; verification: read the result; development: annotated; confidence: high}
\end{GSitem}
\begin{GSitem}{1980-492.5}{E}{The hyperbolic and horospherical matrix example}
\GSfield{Kind; source}{remark; \GSpage{P000920}}
\GSfield{The notes}{For \(GL(n,\mathbb C)\), unitary and triangular subgroups are proposed; the \(n=2\) case is described as a Lobachevsky plane with horospheres. The choice of subgroup is tentative.}
\GSfield{Statement}{A precise standard model is hyperbolic three-space \(SL(2,\mathbb C)/SU(2)\), with oriented horospheres parametrized by \(SL(2,\mathbb C)/(MN)\), where \(M\) is the compact diagonal subgroup and \(N\) upper unipotent. A full complex Borel subgroup instead parametrizes the sphere at infinity. \(GL(2,\mathbb C)/U(2)\) has real dimension four, so it is not a hyperbolic plane.}
\GSfield{Reconstruction}{The SL(2,C), SU(2), MN formulation is editorial reconstruction of the example, not a transcription.}
\GSfield{Fate}{The group/quotient calculation identifies the classical intended example; it does not identify the notebook’s tentative triangular subgroup with the horosphere stabilizer. Its integral-geometric context is \citet[\S1]{gelfandgindikinshapiro1979local}.}
\GSfield{Present status}{As of November 2026, the exact example requires the subgroup and dimension corrections stated here.}
\GSfield{Assessment}{literal statement: false as written; publication match: thematic only; verification: indirect; development: none; confidence: medium}
\end{GSitem}
\begin{GSitem}{1980-492.6}{E}{Segre families and separate-coordinate equivalence}
\GSfield{Kind; source}{construction; \GSpage{P000920-P000921}}
\GSfield{The notes}{The notes cite an Inventiones paper on Segre families, use \(\tau(z,\zeta)=0\), and allow independent holomorphic changes in its two sets of variables and multiplication of \(\tau\) by a nonvanishing factor.}
\GSfield{Statement}{A real-analytic hypersurface in \(\mathbb C^{n+1}\) has a complexified defining equation \(\tau(z,\zeta)=0\). Its fibers over \(\zeta\) form the Segre family. Near a regular nondegenerate point, consider the equivalence of this complex incidence hypersurface under independent biholomorphisms in \(z\) and \(\zeta\). Multiplication by a nonzero function preserves the hypersurface and rescales the contact form \(d_z\tau\) on it.}
\GSfield{Fate}{The exact bibliographic identification is \citet[pp.~135--172]{faran1980segre}. The paper was received 15 November 1979, revised 28 April 1980 and published in June 1980. Publisher metadata, rather than the full proof, was checked.}
\GSfield{Present status}{As of November 2026, Faran’s article is an exact subject match; no assertion that its detailed normal-form theorem has been independently verified is made.}
\GSfield{Assessment}{literal statement: correct with missing hypotheses; publication match: exact; verification: publisher metadata checked; development: none; confidence: high}
\end{GSitem}
\begin{GSitem}{1980-492.7}{E}{Adapted coframes and the Segre structure group}
\GSfield{Kind; source}{definition; \GSpage{P000920-P000921}}
\GSfield{The notes}{An adapted coframe \((\theta,\theta^\alpha,\theta_\alpha)\) changes by a block matrix with diagonal blocks \(u,U,V\) and independent \(\theta\)-terms in both families. Several displayed indices and stars are inconsistent.}
\GSfield{Statement}{On a complex manifold of dimension \(2n+1\), an adapted coframe changes as \(\widehat\theta=a\theta\), \(\widehat\theta^\alpha=A^\alpha_\beta\theta^\beta+b^\alpha\theta\), and \(\widehat\theta_\alpha=B_\alpha^\beta\theta_\beta+c_\alpha\theta\), where \(a\ne0\) and \(A,B\) are invertible. This preserves the common hyperplane and its two distinguished rank-\(n\) subbundles. If a normalized Levi pairing is imposed, \(A,B,a\) satisfy its additional compatibility relation; the unnormalized block display alone does not impose it.}
\GSfield{Fate}{The notebook explicitly associates this structure with \citet[pp.~135--172]{faran1980segre}. The corrected coframe action is elementary matrix bookkeeping; the exact normalization used in the article was not read.}
\GSfield{Present status}{As of November 2026, the unnormalized reduction is well-defined. The defective displayed transformation laws are not literal identities.}
\GSfield{Assessment}{literal statement: false as written; publication match: partial; verification: publisher metadata checked; development: none; confidence: medium}
\end{GSitem}
\begin{GSitem}{1980-492.8}{E}{Frobenius conditions for the two Segre foliations}
\GSfield{Kind; source}{claim; \GSpage{P000921}}
\GSfield{The notes}{Both ideals generated by \((\theta,\theta^1,\ldots,\theta^n)\) and by \((\theta,\theta_1,\ldots,\theta_n)\) must be closed under exterior differentiation; invariants of curve and surface families are requested.}
\GSfield{Statement}{For an adapted coframe, require \(d\theta,d\theta^\alpha\equiv0\pmod{\theta,\theta^1,\ldots,\theta^n}\), and the analogous condition for the lower-index forms. Frobenius' theorem then gives two local rank-\(n\) foliations. Their transversality and, for a nondegenerate Segre geometry, nondegeneracy of the Levi pairing must also be required. The request for invariants is an equivalence problem, not an assertion that every such structure is flat.}
\GSfield{Fate}{The integrability implication is the classical Frobenius theorem; the intended equivalence problem is the published subject of \citet[pp.~135--172]{faran1980segre}, with the ODE antecedent in \citet[\S\S22--24]{cartan1924projective}.}
\GSfield{Present status}{As of November 2026, the local foliation criterion is established; no separate precise unresolved invariant problem is specified.}
\GSfield{Assessment}{literal statement: correct with missing hypotheses; publication match: partial; verification: indirect; development: none; confidence: high}
\end{GSitem}

\GStalk{493}{1980-10-20}{A. B. Goncharov}{P000922-P000924}{G-structures, intersection complexes, and regular holonomic systems}
The session begins on P000922 with Goncharov’s definition of a G-structure and a problem about geometries on the sphere. Fuchs then revisits allowable chains, and Beilinson explains intersection complexes through differential operators, minimal extension and the regular Riemann--Hilbert correspondence, continuing through P000924. These interventions are expressly headed in the manuscript although absent from the short catalogue label. P000925 begins 27 October and belongs to records 494--495.

\begin{GSitem}{1980-493.1}{E}{A G-structure as a reduction of the frame bundle}
\GSfield{Kind; source}{definition; \GSpage{P000922}}
\GSfield{The notes}{For a smooth \(n\)-manifold \(M\), its tangent frame bundle is a principal \(GL(n,\mathbb R)\)-bundle; a \(G\)-subbundle is called a \(G\)-structure.}
\GSfield{Statement}{For a closed Lie subgroup \(G\subset GL(n,\mathbb R)\), a smooth \(G\)-structure is a principal \(G\)-subbundle of the frame bundle \(F(M)\), equivalently a reduction of its structure group. Integrability or flatness is an additional condition.}
\GSfield{Fate}{This is classical structure-group language underlying Cartan’s equivalence method; compare \citet[\S\S22--24]{cartan1924projective} for the projective-equation example discussed the previous week.}
\GSfield{Present status}{As of November 2026, the definition is standard and supplies no new classification theorem.}
\GSfield{Assessment}{literal statement: correct; publication match: thematic only; verification: read the result; development: none; confidence: high}
\end{GSitem}
\begin{GSitem}{1980-493.2}{F}{Geometries on the two-sphere associated with subgroups}
\GSfield{Kind; source}{problem; \GSpage{P000922}}
\GSfield{The notes}{The problem is to consider all geometries on \(S^2\) associated with different subgroups \(G\subset GL(n,\mathbb R)\).}
\GSfield{Statement}{For \(S^2\), a reduction question must use \(G\subset GL(2,\mathbb R)\). One can ask which such reductions exist, classify the reductions up to bundle isomorphism or diffeomorphism, or classify locally homogeneous integrable structures; these are different problems. The notes do not choose one, restrict the subgroups, or state an equivalence relation.}
\GSfield{Fate}{No determinate classification problem can be extracted from the single sentence; the neighboring definition identifies the subject but not its quantifiers.}
\GSfield{Present status}{As of November 2026, this remains an underspecified program in the source, not evidence of an open or undeveloped theorem.}
\GSfield{Assessment}{literal statement: not determinable; publication match: none found; verification: unresolved lead; development: none; confidence: low}
\end{GSitem}
\begin{GSitem}{1980-493.3}{E}{Allowability revisited with complex-dimensional skeleta}
\GSfield{Kind; source}{definition; \GSpage{P000922}; D. B. Fuchs}
\GSfield{The notes}{Fuchs writes \(\dim(\sigma^q\cap X_{n-i})\le q-i\), explicitly noting real codimension \(2i\). He contrasts this with bounds \(q\) and \(q-2i\).}
\GSfield{Statement}{For middle perversity and a stratum of real codimension \(2i\), the bound is \(q-2i+\bar m(2i)=q-i-1\), with the analogous boundary condition. The unrestricted \(q\) condition gives ordinary chains; the transverse \(q-2i\) condition is zero perversity. Its relation to ordinary cohomology uses duality and the normality and compact-support conventions, rather than a literal equality of arbitrary chain complexes.}
\GSfield{Fate}{Compare \citet[\S\S1.1,1.3,3.3]{goreskymacpherson1980intersection} and the earlier occurrences 1979-484.1 and 1980-490.1. The repeated formula is another faithfully recorded missing-index error.}
\GSfield{Present status}{As of November 2026, the theory is established, but the displayed middle bound differs by one from middle perversity.}
\GSfield{Assessment}{literal statement: false as written; publication match: exact; verification: read the result; development: none; confidence: high}
\end{GSitem}
\begin{GSitem}{1980-493.4}{E}{Sheaf-local intersection cohomology and its arithmetic motivation}
\GSfield{Kind; source}{method; \GSpage{P000922}; D. B. Fuchs and A. A. Beilinson}
\GSfield{The notes}{Deligne wants a definition valid over finite fields; a local complex of sheaves should compute MacPherson’s groups. The comparison with \(L^2\) forms on the smooth locus is again mentioned.}
\GSfield{Statement}{The intersection complex is specified locally by extension and truncation from the regular stratum. This sheaf formulation admits an \(\ell\)-adic analogue over a field of characteristic different from \(\ell\). A proposed \(L^2\) identification is a separate analytic problem and requires a metric and a differential domain.}
\GSfield{Fate}{Deligne’s sheaf construction predates the seminar and is treated in \citet[\S\S1--3]{goreskymacpherson1983intersection}. The precise arithmetic analytic specialization is \citet[p.~3]{looijenga1988l2}. This repeats 1980-490.4--5 rather than supplying a new complete conjecture.}
\GSfield{Present status}{As of November 2026, the sheaf formulation is established. The notebook’s unrestricted analytic comparison remains underdetermined.}
\GSfield{Assessment}{literal statement: correct with missing hypotheses; publication match: partial; verification: indirect; development: none; confidence: high}
\end{GSitem}
\begin{GSitem}{1980-493.5}{E}{Embedding a singular variety and the delta-module picture}
\GSfield{Kind; source}{construction; \GSpage{P000923}; A. A. Beilinson}
\GSfield{The notes}{Embed \(X\) in smooth \(Y\), regard functions on \(X_{\rm reg}\) as differential equations, and extend them transversely by delta functions. A cross in the plane illustrates the construction.}
\GSfield{Statement}{If \(i:U\hookrightarrow V\) is a smooth closed embedding, the direct image \(i_+\mathcal O_U\) is a holonomic \(\mathcal D_V\)-module supported on \(U\); locally it includes all normal derivatives of a delta generator. For \(U=X_{\rm reg}\), this is first formed on \(V=Y\setminus X_{\rm sing}\), then extended across the omitted locus.}
\GSfield{Fate}{This is the direct-image and support formalism of holonomic systems used in \citet[\S5 and the references there]{kashiwara1984riemann}; its regular-system formulation was already circulated before the talk, as \citet[pp.~XIX.1--2]{kashiwara1980constructible} records on 22 April 1980.}
\GSfield{Present status}{As of November 2026, the construction is established; a single informal delta function must be read as its generated D-module.}
\GSfield{Assessment}{literal statement: correct with missing hypotheses; publication match: exact; verification: read the result; development: none; confidence: high}
\end{GSitem}
\begin{GSitem}{1980-493.6}{E}{Minimal extension and the irreducibility hypothesis}
\GSfield{Kind; source}{theorem; \GSpage{P000923}; A. A. Beilinson}
\GSfield{The notes}{After maximal direct image the notes assert that there is a unique irreducible submodule; this is the desired extension across the singularities.}
\GSfield{Statement}{For a connected smooth dense stratum \(U\) of an irreducible variety \(X\), an irreducible regular connection has a unique intermediate extension with no nonzero submodule or quotient supported on \(X\setminus U\). For the trivial rank-one connection this is the intersection-cohomology D-module. If \(X\) is reducible, the direct sum of the componentwise extensions need not be irreducible; the cross drawn in the notes has two components.}
\GSfield{Fate}{The correspondence between simple regular holonomic modules and intersection complexes belongs to the regular Riemann--Hilbert theory announced in \citet[p.~XIX.1]{kashiwara1980constructible} and proved in detail in \citet[\S\S5--7]{kashiwara1984riemann}.}
\GSfield{Present status}{As of November 2026, the intermediate extension is canonical. The notebook’s unqualified irreducibility is false for its reducible example.}
\GSfield{Assessment}{literal statement: false as written; publication match: exact; verification: read the result; development: none; confidence: high}
\end{GSitem}
\begin{GSitem}{1980-493.7}{E}{Derived solutions from a locally free resolution}
\GSfield{Kind; source}{method; \GSpage{P000923-P000924}; A. A. Beilinson}
\GSfield{The notes}{Resolve the D-module by free modules \(F_i\) and apply \(\operatorname{Hom}_{\mathcal D_Y}(F_i,\mathcal O_Y)\); the resulting cohomology is called MacPherson cohomology after a codimension shift.}
\GSfield{Statement}{The resulting sheaf complex computes \(\operatorname{Sol}(\mathcal M)=R\mathcal Hom_{\mathcal D_Y}(\mathcal M,\mathcal O_Y)\), using local resolutions. For the intersection-cohomology D-module supported on a pure \(n\)-dimensional \(X\subset Y\), with \(\dim Y=m\), one has \(\operatorname{Sol}(\mathcal M)[m]\simeq IC_X\), where \(IC_X|_{X_{\rm reg}}=\mathbb C[n]\), for the self-dual trivial local system. Thus unshifted solutions restrict in degree \(m-n\); hypercohomology, not stalk cohomology alone, gives the global groups.}
\GSfield{Fate}{The solution functor and derived equivalence are stated in \citet[p.~XIX.1]{kashiwara1980constructible} and \citet[Introduction and \S6]{kashiwara1984riemann}. The explicit shifts reconcile the source’s codimension remark.}
\GSfield{Present status}{As of November 2026, the construction is standard with the local/global and shift conventions stated.}
\GSfield{Assessment}{literal statement: correct with missing hypotheses; publication match: exact; verification: read the result; development: none; confidence: high}
\end{GSitem}
\begin{GSitem}{1980-493.8}{E}{Flat connections and local systems}
\GSfield{Kind; source}{theorem; \GSpage{P000924}; A. A. Beilinson}
\GSfield{The notes}{The simplest systems are locally equivalent to \(\partial u/\partial x_i=0\); their solution sheaves are locally constant, while global solutions may not exist.}
\GSfield{Statement}{On a complex manifold a finite-rank holomorphic vector bundle with integrable connection is locally a constant connection and has a finite-rank local system of horizontal sections. Conversely, tensoring such a local system with \(\mathcal O\) gives the connection. Nontrivial monodromy can make the space of global horizontal sections zero.}
\GSfield{Fate}{This is the nonsingular starting case of \citet[Introduction]{kashiwara1984riemann}; the regular-singular extension was already announced in \citet[p.~XIX.1]{kashiwara1980constructible}.}
\GSfield{Present status}{As of November 2026, the local and global statements are compatible.}
\GSfield{Assessment}{literal statement: correct; publication match: exact; verification: read the result; development: none; confidence: high}
\end{GSitem}
\begin{GSitem}{1980-493.9}{E}{The regular Riemann--Hilbert correspondence}
\GSfield{Kind; source}{theorem; \GSpage{P000924}; A. A. Beilinson}
\GSfield{The notes}{The notes assert equivalence between regular singular holonomic systems and MacPherson complexes, and attribute the regularity definition to Kashiwara.}
\GSfield{Statement}{On a complex manifold, the solution functor gives an anti-equivalence between the bounded derived category of regular holonomic analytic D-modules and the bounded derived category of complex-constructible sheaves. Under the standard dimension shift, the abelian category corresponds to perverse sheaves. Intersection complexes are the simple objects obtained from irreducible local systems, not all constructible complexes.}
\GSfield{Reconstruction}{The distinction between derived constructible complexes, perverse sheaves and simple intersection complexes is supplied.}
\GSfield{Fate}{An explicit theorem and an inverse-functor construction were announced on 22 April 1980 in \citet[pp.~XIX.1--2]{kashiwara1980constructible}, which also acknowledges Mebkhout’s independent proof. The full paper was received 4 March 1983 and published in 1984: \citet[Main Theorem, \S6]{kashiwara1984riemann}.}
\GSfield{Present status}{As of November 2026, the theorem is established; identifying every regular system with a single intersection complex would be too restrictive.}
\GSfield{Assessment}{literal statement: correct with missing hypotheses; publication match: exact; verification: read the result; development: none; confidence: high}
\end{GSitem}
\begin{GSitem}{1980-493.10}{E}{Torsion-freeness does not characterize connections}
\GSfield{Kind; source}{claim; \GSpage{P000924}; A. A. Beilinson}
\GSfield{The notes}{The source defines no torsion by \(f x\ne0\) for every nonzero \(f\in\mathcal D\) and nonzero \(x\in\mathcal M\), then says the system reduces to the locally constant case.}
\GSfield{Statement}{A coherent D-module which is coherent as an \(\mathcal O\)-module on a smooth complex manifold is a finite-rank vector bundle with integrable connection. D-torsion-freeness is not this hypothesis: \(\mathcal O\) itself satisfies \(\partial\cdot1=0\), whereas the rank-one free module \(\mathcal D\) is D-torsion-free and is not a finite-rank connection when the manifold has positive dimension.}
\GSfield{Fate}{The proper finite-rank condition is part of the holonomic-system framework of \citet[\S5]{kashiwara1984riemann}. The two counterexamples directly refute the literal criterion; the manuscript visibly writes \(f\in\mathcal D\), not \(f\in\mathcal O\).}
\GSfield{Present status}{As of November 2026, the source criterion is false; changing only D-torsion to O-torsion would still not supply O-coherence.}
\GSfield{Assessment}{literal statement: false as written; publication match: exact; verification: read the result; development: annotated; confidence: high}
\end{GSitem}
\begin{GSitem}{1980-493.11}{E}{Regular singular equations and logarithmic lattices}
\GSfield{Kind; source}{definition; \GSpage{P000924}; A. A. Beilinson}
\GSfield{The notes}{A Fuchsian equation is written \(dx/dt=t^{-1}A(t)x\).}
\GSfield{Statement}{A meromorphic connection on a complex curve is regular singular at t=0 when it admits a logarithmic lattice. In a suitable meromorphic frame its system has coefficient \(t^{-1}A(t)\) with A holomorphic near zero. A particular displayed system requires A to have this regularity; an arbitrary meromorphic A does not suffice.}
\GSfield{Fate}{This is classical regular-singular ordinary differential equation theory, part of the regularity framework recalled in \citet[\S5]{kashiwara1984riemann}.}
\GSfield{Present status}{As of November 2026, the logarithmic-lattice criterion is established; the holomorphic coefficient qualification is required.}
\GSfield{Assessment}{literal statement: correct with missing hypotheses; publication match: partial; verification: read the result; development: none; confidence: high}
\end{GSitem}
\begin{GSitem}{1980-493.12}{E}{Testing higher-dimensional regularity on curves}
\GSfield{Kind; source}{claim; \GSpage{P000924}; A. A. Beilinson}
\GSfield{The notes}{The notes say that it is sufficient to restrict to a sufficiently regular curve.}
\GSfield{Statement}{Regularity of a meromorphic connection can be tested by requiring regular singular pullbacks on all appropriate smooth analytic curve germs through the boundary. This is a quantified curve criterion; checking one unspecified generic curve is not such a criterion. For general holonomic modules, derived inverse image and the support stratification replace a naive restriction of equations.}
\GSfield{Fate}{The curve-wise formulation is part of the regular holonomic theory summarized in \citet[\S5]{kashiwara1984riemann}; the regularity preprint was already cited by \citet[bibliography, item~2]{kashiwara1980constructible}. The notebook’s quantifier is incomplete.}
\GSfield{Present status}{As of November 2026, the qualified regularity criterion is established. No unrestricted single-curve test is certified.}
\GSfield{Assessment}{literal statement: correct with missing hypotheses; publication match: partial; verification: read the result; development: none; confidence: high}
\end{GSitem}
\begin{GSitem}{1980-493.13}{E}{Finite length and generically flat simple factors}
\GSfield{Kind; source}{theorem; \GSpage{P000924}; A. A. Beilinson}
\GSfield{The notes}{In the general case every module is said to have a finite filtration with irreducible factors, each a locally constant system away from some subvariety.}
\GSfield{Statement}{Coherent holonomic D-modules have finite length locally in the analytic setting and globally on smooth algebraic varieties of finite type. Each simple factor has irreducible support and is, on a dense smooth open subset of that support, an irreducible connection, interpreted by direct image from the support. Arbitrary D-modules need not have finite length. Restricting to regular holonomic factors gives regular connections and their intermediate extensions.}
\GSfield{Fate}{These are the holonomic-category qualifications underlying \citet[\S5]{kashiwara1984riemann} and the regular derived equivalence announced in \citet[p.~XIX.1]{kashiwara1980constructible}. Infinite direct sums refute the literal assertion for every module.}
\GSfield{Present status}{As of November 2026, the finite-length theorem is established with coherence and holonomicity; its unrestricted formulation is false.}
\GSfield{Assessment}{literal statement: false as written; publication match: partial; verification: read the result; development: none; confidence: high}
\end{GSitem}

\GStalk{494}{1980-10-27}{I. M. Gelfand}{P000925-P000926}{Elementary group problems and projective structures on the line}
The dated heading is on P000925. The problems concern finite symmetry groups, the topology and group law of the universal cover of SL(2,R), a triple-line cocycle, and Kirillov’s classification of locally projective structures on the real line. The material ends above Gaveau’s heading on P000926. The catalogue’s preceding P000924 still belongs to the 20 October D-module discussion.

\begin{GSitem}{1980-494.1}{E}{Subgroups and normal subgroups of finite symmetry groups}
\GSfield{Kind; source}{problem; \GSpage{P000925}}
\GSfield{The notes}{Find all subgroups and normal subgroups of symmetry groups of regular polyhedra and of the symmetric group; these are described as problems for first-year students.}
\GSfield{Statement}{For a specified finite group \(G\), enumerate its subgroups and test normality. For \(S_n\), \(n\ge5\), the normal subgroups are \(1,A_n,S_n\); for \(S_4\), add the normal Klein four group. The small cases must be treated separately. For the Platonic solids one must distinguish rotation groups \(A_4,S_4,A_5\) from full reflection-containing symmetry groups. Enumerating all subgroups of \(S_n\) is a finite problem for each specified \(n\), not a short uniform list.}
\GSfield{Reconstruction}{The explicit normal-subgroup answer and enumeration method are expository annotations; the notebook only sets the exercise.}
\GSfield{Fate}{These are classical finite-group exercises. A complete finite enumeration can start at the trivial subgroup and repeatedly adjoin elements until no new subgroup is obtained; every subgroup is eventually reached by its finite generating list. The source supplies neither a value of n nor the rotation/full-symmetry convention.}
\GSfield{Present status}{As of November 2026, the normal-subgroup assertion and finite enumeration are established; no new finite-group conjecture is posed.}
\GSfield{Assessment}{literal statement: correct with missing hypotheses; publication match: none found; verification: indirect; development: annotated; confidence: high}
\end{GSitem}
\begin{GSitem}{1980-494.2}{E}{Topology of SL(2,R) and its universal covering manifold}
\GSfield{Kind; source}{question; \GSpage{P000925}}
\GSfield{The notes}{Is \(SL(2,\mathbb R)\) simply connected? A quadratic change of coordinates gives \(SL(2,\mathbb R)\simeq S^1\times\mathbb R^2\).}
\GSfield{Statement}{The Iwasawa decomposition \(SL(2,\mathbb R)=KAN\), with \(K=SO(2)\), positive diagonal \(A\simeq\mathbb R\), and upper unipotent \(N\simeq\mathbb R\), is a diffeomorphism of manifolds. Hence \(\pi_1SL(2,\mathbb R)\simeq\mathbb Z\), and its universal cover is diffeomorphic to \(\mathbb R^3\). The product decomposition is not a direct-product decomposition of Lie groups.}
\GSfield{Fate}{This classical covering structure is used explicitly in \citet[\S4.2]{belgunmoroianu2025projective}; the notebook already supplies a direct quadratic-coordinate proof of the manifold decomposition.}
\GSfield{Present status}{As of November 2026, the question has the stated negative answer, with infinite cyclic fundamental group.}
\GSfield{Assessment}{literal statement: correct; publication match: exact; verification: read the result; development: none; confidence: high}
\end{GSitem}
\begin{GSitem}{1980-494.3}{E}{The universal covering group and a triple-line cocycle}
\GSfield{Kind; source}{construction; \GSpage{P000925-P000926}}
\GSfield{The notes}{Pairs \((g,n)\) are multiplied with a triple-line index correction. The display writes \(g\in SL(2,\mathbb C)\), then defines a phase using signs of real matrix entries; the kernel of a phase condition is said to have four components, from which the connected one is taken.}
\GSfield{Statement}{The universal covering group of \(SL(2,\mathbb R)\) can be defined by homotopy classes relative endpoints of paths from the identity, with pointwise path multiplication. Its kernel is \(\mathbb Z\). Choosing a section gives an integer cocycle \(c(g,h)\), and \((g,m)(h,n)=(gh,m+n+c(g,h))\). A triple-line index can express a cohomologous cocycle only after degenerate triples, normalization, coboundary and topology have been specified.}
\GSfield{Reconstruction}{The path-group construction is supplied to separate the identified object from an unverified explicit cocycle.}
\GSfield{Fate}{The lift-action description is classical and is made precise in \citet[\S4.2]{belgunmoroianu2025projective}. The path formulation reconstructs the intended covering group; it does not certify the notebook’s incomplete phase/cocycle normalization.}
\GSfield{Present status}{As of November 2026, the covering group is established. The literal SL(2,C), real sign formulas and component assertions do not define that group without correction.}
\GSfield{Assessment}{literal statement: false as written; publication match: partial; verification: read the result; development: none; confidence: medium}
\end{GSitem}
\begin{GSitem}{1980-494.4}{E}{Countably many projective structures on the real line up to isomorphism}
\GSfield{Kind; source}{problem; \GSpage{P000926}; A. A. Kirillov}
\GSfield{The notes}{Kirillov asks for all atlases on \(\mathbb R\) with fractional-linear transitions and says there are countably many structures. The affine line and coverings of the projective line are examples.}
\GSfield{Statement}{Up to projective diffeomorphism, connected open one-dimensional real projective manifolds have countably many isomorphism classes. A developing map identifies one with an interval in \(\widetilde{\mathbb RP^1}\simeq\mathbb R\). Up to lifted projective transformations and orientation reversal, the classes are the full cover, a half-line, and bounded intervals indexed by positive half-integral winding number. If equivalence instead means compatible atlases on a fixed underlying line, without quotienting by diffeomorphisms, the countability statement is false.}
\GSfield{Fate}{The classification is traced to Kuiper (1954) in \citet[Proposition 4.15 and Corollary 4.16]{belgunmoroianu2025projective}; that primary 2025 preprint explicitly recovers the countable open-curve classification while correcting separate closed-curve issues. This differs from the circle problem in 1976-403.6.}
\GSfield{Present status}{As of November 2026, the countability assertion is proved under projective-isomorphism equivalence; this is a resolved classification problem.}
\GSfield{Assessment}{literal statement: correct with missing hypotheses; publication match: exact; verification: read the result; development: none; confidence: high; problem \hyperref[prob:GP-1980-01]{GP-1980-01}}
\end{GSitem}

\GStalk{495}{1980-10-27}{Bernard Gaveau}{P000926-P000928}{Levi-flat fillings and gauge-covariant stochastic constructions}
Gaveau’s heading begins near the bottom of P000926 and continues through P000928. The first topic concerns envelopes of holomorphy and disc fillings of spheres in complex dimension two; the second constructs gauge-covariant diffusion on connection spaces with a fixed heat-kernel cutoff; the final formulas concern Feynman--Kac, spectral growth and exponential moment estimates. His brief list of possible topics is not evidence that a full scattering theorem was stated. P000929--P000930 begin the dated 3 November session and are not part of this talk.

\begin{GSitem}{1980-495.1}{A}{Levi-flat filling of an embedded two-sphere}
\GSfield{Kind; source}{problem; \GSpage{P000926}}
\GSfield{The notes}{Find a Levi-flat three-dimensional manifold \(M\subset\mathbb C^2\) with \(\partial M=S^2\), foliated by complex curves. The drawing shows a one-parameter family of discs with two end points.}
\GSfield{Statement}{A precise proved version assumes a smooth embedded sphere in the boundary of a strictly pseudoconvex domain in \(\mathbb C^2\), exactly two elliptic complex points, and total reality elsewhere. It bounds a Levi-flat hypersurface foliated by analytic discs, smooth up to the sphere away from the elliptic points. The notebook does not state these hypotheses; arbitrary embedded two-spheres need not admit such a filling.}
\GSfield{Reconstruction}{Strict pseudoconvexity and the two elliptic-point assumptions come from the later published theorem, not the manuscript.}
\GSfield{Fate}{The closely matching speaker publication is Bedford--Gaveau (1983), identified in the primary account \citet[Introduction, p.~516]{gaussiersukhov2011filling}, which states its hypotheses and theorem. The same paper’s Theorem 1.1 gives a later almost-complex extension and notes counterexamples without pseudoconvexity. The 1983 original was identified bibliographically rather than read in full.}
\GSfield{Present status}{As of November 2026, the specified sphere theorem is established; a universal assertion for arbitrary embedded spheres is false.}
\GSfield{Assessment}{literal statement: correct with missing hypotheses; publication match: partial; verification: indirect; development: none; confidence: high; problem \hyperref[prob:GP-1980-02]{GP-1980-02}}
\end{GSitem}
\begin{GSitem}{1980-495.2}{E}{The asserted obstruction for a torus}
\GSfield{Kind; source}{claim; \GSpage{P000926}}
\GSfield{The notes}{The parenthetical sentence says that this problem cannot be solved for a torus.}
\GSfield{Statement}{There is no obstruction based only on torus topology: \(M=\{(z,w):|z|\le1,\ |w|=1\}\subset\mathbb C^2\) is a smooth solid torus with boundary the Clifford torus. Its interior is Levi-flat and is foliated by the discs \(w=\mathrm{constant}\). A prohibition would need additional embedding, hull or boundary conditions absent from the notes.}
\GSfield{Fate}{The explicit product example is an elementary expository counterexample. The sphere-filling theorem in \citet[Introduction]{gaussiersukhov2011filling} does not imply a universal nonfilling theorem for tori.}
\GSfield{Present status}{As of November 2026, the literal torus assertion is false.}
\GSfield{Assessment}{literal statement: false as written; publication match: partial; verification: read the result; development: annotated; confidence: high}
\end{GSitem}
\begin{GSitem}{1980-495.3}{F}{A degenerate boundary problem of index zero}
\GSfield{Kind; source}{claim; \GSpage{P000927}}
\GSfield{The notes}{An isolated foliation singularity is said to lead to a degenerate Riemann problem of index zero, hence solvability; an integro-differential equation is mentioned.}
\GSfield{Statement}{Index zero alone states equality of kernel and cokernel dimensions for a Fredholm operator; it does not imply surjectivity. The notes do not specify the operator, domain, boundary regularity, ellipticity or a separate uniqueness argument, so no definite solvability theorem can be recovered.}
\GSfield{Fate}{The adjacent Levi-flat discussion identifies the area, but the displayed implication has no recoverable hypotheses or operator. Searches for the Bedford--Gaveau filling work did not identify this exact analytic step.}
\GSfield{Present status}{As of November 2026, this is an unrecoverable assertion as recorded; index zero without a kernel calculation is insufficient.}
\GSfield{Assessment}{literal statement: not determinable; publication match: none found; verification: unresolved lead; development: none; confidence: low}
\end{GSitem}
\begin{GSitem}{1980-495.4}{A}{Gauge-invariant functional integration as a regularized program}
\GSfield{Kind; source}{method; \GSpage{P000927}}
\GSfield{The notes}{The formal density is \(\exp(-\int\operatorname{Tr}(\Omega\Omega))D\omega\), and gauge transformations act by \(g^{-1}dg+g^{-1}\omega g\). A functional integral on trajectories is proposed.}
\GSfield{Statement}{For a principal compact-group bundle over a compact Riemannian manifold, the gauge action on connections is \(A^g=g^{-1}Ag+g^{-1}dg\), and the Yang--Mills energy is \(\int|F_A|^2\,d\mathrm{vol}\), using an invariant positive inner product. The symbol \(DA\) is formal. The identified construction replaces it by a stochastic process with fixed positive smoothing parameter; it is not a construction of a translation-invariant infinite-dimensional Lebesgue measure or of a cutoff-free four-dimensional Yang--Mills theory.}
\GSfield{Fate}{Gaveau--Trauber’s connection-space construction appeared in JFA 38 (1980), 324--341, after a 1979 C. R. announcement. The author’s primary follow-up \citet[Introduction and references 2,5]{gaveaumazet1981elliptic} describes that result and the separate 1980 cutoff quantization announcement. It explicitly acknowledges the Moscow and Leningrad autumn-1980 seminars.}
\GSfield{Present status}{As of November 2026, the regularized construction is identified; no claim about removing the cutoff follows from these notes.}
\GSfield{Assessment}{literal statement: correct with missing hypotheses; publication match: partial; verification: indirect; development: none; confidence: high}
\end{GSitem}
\begin{GSitem}{1980-495.5}{A}{Covariant Hodge heat kernels as diffusion covariance}
\GSfield{Kind; source}{construction; \GSpage{P000927}}
\GSfield{The notes}{For a connection \(\Pi\), define \(\Box_\Pi=d_\Pi^*d_\Pi+d_\Pi d_\Pi^*\), use its heat kernel at fixed \(\tau\), and contract it with the second functional derivative of a functional.}
\GSfield{Statement}{On adjoint-bundle-valued one-forms over a compact Riemannian base, \(\Box_A\) is a nonnegative self-adjoint elliptic operator. For \(\tau>0\), \(Q_A=e^{-\tau\Box_A}\) is smoothing and trace class. For smooth cylindrical \(\Phi\), \(L\Phi(A)=\tfrac12\operatorname{Tr}(Q_A D^2\Phi(A))\) is a well-defined second-order expression. Gauge covariance of \(\Box_A\) implies covariance of its heat kernel and invariance of this expression.}
\GSfield{Fate}{The exact construction is Gaveau--Trauber (1980), described by the authors in \citet[pp.~105--106 and \S2]{gaveaumazet1981elliptic}; the latter explicitly says that the connection-space operator uses the heat kernel as its symbol.}
\GSfield{Present status}{As of November 2026, the fixed-cutoff cylindrical generator is established. Its stochastic realization needs the regularity and lifetime hypotheses considered separately.}
\GSfield{Assessment}{literal statement: correct with missing hypotheses; publication match: exact; verification: indirect; development: none; confidence: high}
\end{GSitem}
\begin{GSitem}{1980-495.6}{A}{Heat-kernel regularity and the stochastic differential equation}
\GSfield{Kind; source}{claim; \GSpage{P000927-P000928}}
\GSfield{The notes}{The heat operators satisfy a Hilbert--Schmidt Lipschitz estimate in a \(C^{k,\alpha}\) connection norm. The notes write \(d\Pi_t=e^{-\tau\Box_\tau}d\beta_t\), using white noise; Sobolev spaces are said not to work.}
\GSfield{Statement}{For compact base, fixed positive \(\tau\), sufficient coefficient regularity, and a bounded neighborhood of connections, parabolic estimates give local Lipschitz dependence of smoothing heat operators in suitable operator norms. A local stochastic equation uses the current connection \(A_t\), and a noise coefficient \(e^{-\tau\Box_{A_t}/2}\) has covariance \(e^{-\tau\Box_{A_t}}\). With coefficient \(e^{-\tau\Box_{A_t}}\), the covariance is instead \(e^{-2\tau\Box_{A_t}}\); the It\^o generator has the factor \(1/2\). The construction does not imply that every possible Sobolev formulation is impossible.}
\GSfield{Fate}{The method is described in \citet[Introduction, \S3, Theorem 2 and its proof]{gaveaumazet1981elliptic}, referring back to Gaveau--Trauber (1980), \S\S2--3. The opened follow-up proves its analogous loop-space diffusion up to its lifetime, not an unconditional global theorem for the notebook equation.}
\GSfield{Present status}{As of November 2026, the regularized local construction is identified; global lifetime, the exact Banach-space norm and the cutoff normalization are not supplied by the manuscript.}
\GSfield{Assessment}{literal statement: false as written; publication match: partial; verification: indirect; development: none; confidence: medium}
\end{GSitem}
\begin{GSitem}{1980-495.7}{E}{The Feynman--Kac representation of a Schroedinger semigroup}
\GSfield{Kind; source}{theorem; \GSpage{P000928}}
\GSfield{The notes}{The expectation \(\mathbb E_x[f(B_t)\exp(-\int_0^tV(B_s)ds)]\) is the evolution operator for \(-\tfrac12\Delta+V\).}
\GSfield{Statement}{For Brownian motion with generator \(\tfrac12\Delta\) and bounded real V, \(e^{-tH}f(x)=\mathbb E_x[f(B_t)\exp(-\int_0^tV(B_s)ds)]\), where \(H=-\tfrac12\Delta+V\). The equality extends to suitable Kato-class potentials, with the operator defined by its closed semibounded quadratic form. The bounded-potential version already makes the displayed formula precise.}
\GSfield{Fate}{This is the classical Feynman--Kac formula; \citet[formula (A26), Theorem A.2.7]{simon1982semigroups} treats it systematically. The later survey is a verification source, not a post-1980 priority claim.}
\GSfield{Present status}{As of November 2026, the formula is established with the specified operator sign and potential hypotheses.}
\GSfield{Assessment}{literal statement: correct with missing hypotheses; publication match: partial; verification: read the result; development: none; confidence: high}
\end{GSitem}
\begin{GSitem}{1980-495.8}{E}{Khasminskii exponential bound and the missing negative part}
\GSfield{Kind; source}{theorem; \GSpage{P000928}}
\GSfield{The notes}{The displayed assumption is \(|\Delta^{-1}V|\le C<1\), followed by \[\mathbb E_x\exp(-\int_0^\infty V(B_s)ds)\le(1-C)^{-1}.\]}
\GSfield{Statement}{For a Markov process and \(W\ge0\), if \(\sup_x\mathbb E_x\int_0^T W(B_s)ds\le C<1\), then \(\sup_x\mathbb E_x\exp(\int_0^TW(B_s)ds)\le(1-C)^{-1}\). The infinite-time version follows when the same occupation bound holds with \(T=\infty\). For the notebook’s sign take \(W=V_-\). A bound on the absolute value of the Green potential of signed \(V\) permits cancellations and is not this hypothesis; the Green operator must also match the Brownian normalization.}
\GSfield{Fate}{The estimate is the classical Khasminskii lemma; see \citet[Lemma B.1.2, p.~461]{simon1982semigroups}. Expanding the exponential and applying the Markov property bounds its ordered-integral term of order \(n\) by \(C^n\), yielding the geometric series.}
\GSfield{Present status}{As of November 2026, the corrected bound is established. The signed Green-potential condition as written is insufficient.}
\GSfield{Assessment}{literal statement: false as written; publication match: exact; verification: read the result; development: annotated; confidence: high}
\end{GSitem}
\begin{GSitem}{1980-495.9}{E}{Large-time logarithmic growth and the bottom of the spectrum}
\GSfield{Kind; source}{claim; \GSpage{P000928}}
\GSfield{The notes}{The logarithmic expectation divided by t is asserted to approach \(\pm\lambda_1\).}
\GSfield{Statement}{On a compact connected Riemannian manifold, for continuous real V and \(H=-\tfrac12\Delta+V\), the positive ground state gives \(\lim_{t\to\infty}t^{-1}\log\mathbb E_x\exp(-\int_0^tV(B_s)ds)=-\lambda_0(H)\). A noncompact version needs additional potential, spectral and pointwise-growth hypotheses; the compact formulation is an explicit safe specialization, not a claim that the notes stated compactness.}
\GSfield{Reconstruction}{Compactness and connectedness are supplied as an explicit sufficient setting; no claim is made that they were present in the notebook.}
\GSfield{Fate}{The spectral-semigroup mechanism is classical; compare \citet[\S B.5]{simon1982semigroups}. The sign follows from \(e^{-tH}\), and is not an independent plus-or-minus choice.}
\GSfield{Present status}{As of November 2026, the compact-case limit holds. The source’s unspecified noncompact generality is not certified.}
\GSfield{Assessment}{literal statement: correct with missing hypotheses; publication match: partial; verification: read the result; development: none; confidence: high}
\end{GSitem}

\GStalk{497}{1980-11-03}{I. M. Gelfand}{P000929}{Representation theory, noncommutative birational geometry, and sheaves}
The actual intervention occupies P000929, headed I.M.G. and History of representation theory. Its historical list runs through characters, Haar integration, operator algebras and Lorentz-group representations; these names are contextual remarks, not newly stated theorems or a reliable exclusive priority list. The mathematical content concerns enveloping algebras, the Gelfand--Kirillov birational conjecture, a poorly specified extension beyond polynomial coefficients, and the BGG sheaf method on projective space. The subsequent Beilinson heading on P000930 starts record 496. The Gelfand--Kirillov problem is the earlier GP-1967-21, not a new 1980 problem.

\begin{GSitem}{1980-497.1}{E}{Enveloping algebras and their quotient division rings}
\GSfield{Kind; source}{method; \GSpage{P000929}}
\GSfield{The notes}{The analogue of the group algebra for a Lie algebra is its enveloping algebra; there are many such algebras, and no Wedderburn theorem is available in the intended sense.}
\GSfield{Statement}{For a finite-dimensional Lie algebra \(\mathfrak g\) over a characteristic-zero field, \(U(\mathfrak g)\) is a noetherian domain and has an Ore quotient division ring. The usual Wedderburn structure theorem for finite-dimensional semisimple associative algebras does not classify these infinite-dimensional algebras. Comparing their quotient division rings is a noncommutative birational classification problem.}
\GSfield{Reconstruction}{The finite-dimensional and characteristic-zero hypotheses and the intended meaning of the Wedderburn comparison are supplied.}
\GSfield{Fate}{This is the framework of \citet[Introduction, pp.~5--6]{gelfandkirillov1966corps}, already discussed in 1967-144.5--6. The statement that there is no Wedderburn theorem is interpreted as a limitation of that finite-dimensional theorem, not a claim that no structure theorem can exist.}
\GSfield{Present status}{As of November 2026, the Ore construction and birational formulation are established.}
\GSfield{Assessment}{literal statement: correct with missing hypotheses; publication match: exact; verification: read the result; development: none; confidence: high}
\end{GSitem}
\begin{GSitem}{1980-497.2}{B}{The unextended Gelfand--Kirillov conjecture revisited}
\GSfield{Kind; source}{conjecture; \GSpage{P000929}}
\GSfield{The notes}{The skew field of fractions of an enveloping algebra is identified with that of a ring of differential operators; Gelfand calls this unfinished.}
\GSfield{Statement}{For a finite-dimensional algebraic Lie algebra \(\mathfrak g\) over an algebraically closed field \(k\) of characteristic zero, the original conjecture asks whether \(\operatorname{Frac}U(\mathfrak g)\simeq\operatorname{Frac}A_m(k(t_1,\ldots,t_r))\) for some \(m,r\ge0\), where \(A_m\) is a Weyl algebra and the \(t_i\) are central algebraically independent variables. An isomorphism after a finite extension of the center is a different, weaker assertion.}
\GSfield{Fate}{The exact recurrence is GP-1967-21 (1967-144.7); see \citet[Introduction]{gelfandkirillov1966corps}. The finite-central-extension result of \citet{gelfandkirillov1969structure} must be kept separate. \citet[Theorem~1.1]{premet2010modular} excludes simple types \(B_n\) for \(n\ge3\), \(D_n\) for \(n\ge4\), \(E_6,E_7,E_8,F_4\). Type A is positive. \citet[Introduction, pp.~2--3 of the author version]{marianoschwarz2024gelfand} still lists types C and G as unresolved; its version of record appeared in August 2025.}
\GSfield{Present status}{As of November 2026, the unrestricted conjecture is false. The checked literature through the present search leaves the simple cases \(C_n\), \(n\ge2\), and \(G_2\) open; the necessary-type restriction A, C, G is not a sufficiency theorem. No claimed 2026 resolution of these enveloping-algebra cases was found.}
\GSfield{Assessment}{literal statement: correct with missing hypotheses; publication match: exact; verification: read the result; development: none; confidence: high; problem \hyperref[prob:GP-1967-21]{GP-1967-21}}
\end{GSitem}
\begin{GSitem}{1980-497.3}{F}{Differential operators beyond polynomial coefficients}
\GSfield{Kind; source}{remark; \GSpage{P000929}}
\GSfield{The notes}{Polynomial coefficients are said to be insufficient, with partial understanding attributed to Novikov, Krichever and Shabat; an additional name is cancelled.}
\GSfield{Statement}{No class of coefficient functions, differential operators, equivalence relation, or proposed replacement for the Weyl division field is specified. The named researchers suggest the theory of commuting differential operators and integrable systems, but the sentence does not determine a theorem or conjecture in that theory.}
\GSfield{Fate}{The names and the neighboring birational statement are recoverable; a precise publication match cannot be inferred from them. The statement is retained as an identifiable research direction without assigning it an invented finite-gap or spectral-curve theorem.}
\GSfield{Present status}{As of November 2026, the intended mathematical assertion remains unrecoverable from this source.}
\GSfield{Assessment}{literal statement: not determinable; publication match: none found; verification: unresolved lead; development: none; confidence: low}
\end{GSitem}
\begin{GSitem}{1980-497.4}{E}{The BGG sheaf approach on projective space}
\GSfield{Kind; source}{method; \GSpage{P000929}}
\GSfield{The notes}{Bernstein, S. Gelfand and I.M. Gelfand are named in connection with projective space, followed by the instruction to work with sheaves.}
\GSfield{Statement}{The identifiable construction associates to a finite-dimensional graded module over the exterior algebra on an \((n+1)\)-dimensional space a linear complex of vector bundles on \(\mathbb P^n\). In the BGG theorem, the stable category modulo graded free modules is equivalent to \(D^b(\operatorname{Coh}\mathbb P^n)\). This precise statement identifies the named method from its publication; the notebook itself states no categorical equivalence.}
\GSfield{Reconstruction}{The exterior-algebra and derived-category formulation comes from the identified paper, not from the notebook sentence.}
\GSfield{Fate}{\citet[Theorems~1--2, pp.~212--213]{bernsteingelfandgelfand1978bundles}, submitted 10 January 1978, is the exact author/subject identification. It is also discussed in the previous companion’s records 1977-437 and 1977-438. The November 1980 sentence is a retrospective pointer.}
\GSfield{Present status}{As of November 2026, the construction and equivalence are established. The short pointer cannot be used as evidence that every detail of the published theorem was stated in this intervention.}
\GSfield{Assessment}{literal statement: correct with missing hypotheses; publication match: partial; verification: read the result; development: none; confidence: high}
\end{GSitem}

\GStalk{496}{1980-11-03}{A. A. Beilinson}{P000930-P000932}{D-affinity and localization on the flag variety}
The Beilinson heading is on P000930, following Gelfand’s separately catalogued historical introduction on P000929. His talk continues through P000932, outside this record’s short catalogue range. It moves from affine schemes and quasi-coherent sheaves to D-affinity, induced modules, left/right operator issues and the Beilinson--Bernstein localization theorem at the trivial infinitesimal character. The uncertainty about the converse categorical criterion is resolved in his 10 November continuation, record 498.

\begin{GSitem}{1980-496.1}{E}{Affine objects and their rings of functions}
\GSfield{Kind; source}{theorem; \GSpage{P000930}}
\GSfield{The notes}{Affine varieties are uniquely determined by their function rings; the notes then identify a commutative ring with an affine algebraic variety.}
\GSfield{Statement}{Affine schemes are contravariantly equivalent to commutative unital rings via \(A\mapsto\operatorname{Spec}A\). Classical reduced affine varieties over \(\mathbb C\) correspond to finitely generated reduced \(\mathbb C\)-algebras (and to domains if irreducibility is part of the definition). An arbitrary commutative ring is not the coordinate ring of a classical complex affine variety.}
\GSfield{Fate}{The classical affine-variety and coherent-sheaf foundations are \citet[Chapter II]{serre1955coherent}; the notebook moves from that category to the later scheme formulation without marking the change.}
\GSfield{Present status}{As of November 2026, the ring equivalence is correct in the scheme category; finite generation and reducedness cannot be omitted in the classical variety formulation.}
\GSfield{Assessment}{literal statement: false as written; publication match: partial; verification: read the result; development: none; confidence: high}
\end{GSitem}
\begin{GSitem}{1980-496.2}{E}{Localization of modules and quasi-coherent sheaves}
\GSfield{Kind; source}{construction; \GSpage{P000930}}
\GSfield{The notes}{For \(X=\operatorname{Spec}A\), the sheaf associated with \(M\) has \(\mathcal M(X_f)=M_f=M\otimes_A A_f\). The affine-open base-change condition is used as the definition of quasi-coherence.}
\GSfield{Statement}{For any \(A\)-module \(M\), the localization rule on principal opens defines a quasi-coherent sheaf \(\widetilde M\). Global sections and sheafification are inverse equivalences between \(A\)-modules and quasi-coherent sheaves on \(\operatorname{Spec}A\). For affine opens \(V\subset U\), a quasi-coherent sheaf satisfies \(\mathcal M(V)=\mathcal O(V)\otimes_{\mathcal O(U)}\mathcal M(U)\).}
\GSfield{Fate}{The coherent finite-type version is established in \citet[Chapter II, \S3]{serre1955coherent}; the arbitrary-module quasi-coherent version is its standard affine-scheme extension.}
\GSfield{Present status}{As of November 2026, the localization construction and equivalence are standard.}
\GSfield{Assessment}{literal statement: correct; publication match: partial; verification: read the result; development: none; confidence: high}
\end{GSitem}
\begin{GSitem}{1980-496.3}{E}{Serre’s cohomological criterion for affineness}
\GSfield{Kind; source}{theorem; \GSpage{P000930}}
\GSfield{The notes}{Affineness is characterized by vanishing of higher cohomology and generation of every quasi-coherent sheaf by global sections.}
\GSfield{Statement}{For a separated algebraic variety of finite type, affineness is equivalent to \(H^i(X,\mathcal F)=0\) for every quasi-coherent \(\mathcal F\) and \(i>0\). On an affine variety, the evaluation map \(\mathcal O_X\otimes_{\Gamma\mathcal O_X}\Gamma\mathcal F\to\mathcal F\) is also surjective. Thus the two conditions stated together in the notes are equivalent to affineness in this setting.}
\GSfield{Fate}{This is the affine-cohomology criterion introduced in \citet[Chapter II, \S3]{serre1955coherent}, with the quasi-coherent formulation standard in the scheme setting.}
\GSfield{Present status}{As of November 2026, the statement is established with the finite-type and separation hypotheses.}
\GSfield{Assessment}{literal statement: correct with missing hypotheses; publication match: exact; verification: read the result; development: none; confidence: high}
\end{GSitem}
\begin{GSitem}{1980-496.4}{A}{Quasi-coherent D-modules and D-affinity}
\GSfield{Kind; source}{definition; \GSpage{P000930-P000931}}
\GSfield{The notes}{A D-module is called quasi-coherent if its underlying O-module is quasi-coherent. D-affinity requires higher-cohomology vanishing and generation by global sections.}
\GSfield{Statement}{For smooth \(X/\mathbb C\), let \(\mathcal D_X\) be finite-order algebraic differential operators and consider left D-modules quasi-coherent over \(\mathcal O_X\). D-affinity means that every such \(\mathcal M\) is acyclic for \(\Gamma\) and the evaluation \(\mathcal D_X\otimes_{\Gamma\mathcal D_X}\Gamma\mathcal M\to\mathcal M\) is surjective. The order of factors records the left-module convention.}
\GSfield{Fate}{This is the localization setting of \citet[pp.~15--18]{beilinsonbernstein1981localisation}. Contemporary independent testimony in \citet[p.~44]{brylinski1981differential} explicitly records Beilinson--Bernstein’s stronger vanishing theorem for every quasi-coherent D-module. The November 1980 notes are an early seminar exposition.}
\GSfield{Present status}{As of November 2026, the definition is standard and supports the categorical equivalence in the sequel.}
\GSfield{Assessment}{literal statement: correct with missing hypotheses; publication match: exact; verification: read the result; development: none; confidence: high}
\end{GSitem}
\begin{GSitem}{1980-496.5}{A}{Affine varieties, projective spaces, and the elliptic counterexample}
\GSfield{Kind; source}{theorem; \GSpage{P000931}}
\GSfield{The notes}{Every affine variety and every projective space is D-affine; an elliptic curve is not.}
\GSfield{Statement}{A smooth affine complex variety is D-affine by ordinary affine quasi-coherent cohomology and generation. Complex \(\mathbb P^N\) is D-affine by localization (or its projective-space case). An elliptic curve \(E\) is not: \(\mathcal O_E\) is a left D-module and \(H^1(E,\mathcal O_E)\ne0\).}
\GSfield{Fate}{The projective assertion is a special case of \citet[pp.~15--18]{beilinsonbernstein1981localisation}; its all-module vanishing is independently recorded in \citet[p.~44]{brylinski1981differential}. The affine and elliptic observations are elementary consequences of the definition and earlier coherent cohomology.}
\GSfield{Present status}{As of November 2026, all three examples have the stated outcomes; the positive-dimensional projective example is the new localization content.}
\GSfield{Assessment}{literal statement: correct; publication match: exact; verification: read the result; development: none; confidence: high}
\end{GSitem}
\begin{GSitem}{1980-496.6}{A}{Induced D-modules from negative line bundles}
\GSfield{Kind; source}{construction; \GSpage{P000931}}
\GSfield{The notes}{Even for \(\mathcal L=\mathcal O(-i)\) on projective space, \(\mathcal D_X\otimes_{\mathcal O_X}\mathcal L\) is generated by global sections and acyclic. The notes say \(\mathcal O(-i)\) itself has no sections and has higher cohomology.}
\GSfield{Statement}{For quasi-coherent \(\mathcal L\), induction uses the right O-action on \(\mathcal D_X\), producing a left D-module. On D-affine \(X\) it is acyclic and D-generated by its global sections. On \(\mathbb P^N\), \(H^0\mathcal O(-i)=0\) for \(i>0\), while nonzero higher cohomology occurs only for \(i\ge N+1\); for \(1\le i\le N\) all cohomology vanishes. For an invertible \(\mathcal L\), the induced module is locally free of rank one over D.}
\GSfield{Fate}{The induced-module consequence follows directly from the all-module localization theorem recorded in \citet[p.~44]{brylinski1981differential} and \citet{beilinsonbernstein1981localisation}; the line-bundle cohomology is classical.}
\GSfield{Present status}{As of November 2026, the induced-module assertion holds, but the blanket higher-cohomology assertion for every negative twist does not.}
\GSfield{Assessment}{literal statement: false as written; publication match: exact; verification: read the result; development: none; confidence: high}
\end{GSitem}
\begin{GSitem}{1980-496.7}{E}{Reversing the tensor product and formal transposition}
\GSfield{Kind; source}{claim; \GSpage{P000931-P000932}}
\GSfield{The notes}{The reversed \(\mathcal L\otimes_{\mathcal O}\mathcal D\) has \(\mathcal L\) as a direct summand. An involution of D fixing functions would identify the two constructions. On an elliptic curve formal adjoint is written \(f\partial_x\mapsto\partial_xf\).}
\GSfield{Statement}{The reversed tensor is naturally a right D-module, with the left O-action on D used in forming it. As an O-module, evaluation \(P\mapsto P(1)\) splits \(\mathcal O\hookrightarrow\mathcal D\), so \(\mathcal L\) is a direct summand. Formal transposition reverses order and sends \(\partial_x\mapsto-\partial_x\) in a constant-volume coordinate; globally it exchanges left and right modules via the canonical bundle. It is not an unrestricted canonical involution of the same untwisted ring.}
\GSfield{Fate}{This is the classical left/right duality of differential operators used in \citet[Introduction, formulae for duality]{kashiwara1984riemann}. The source’s missing minus sign already contradicts \([\partial_x,x]=1\).}
\GSfield{Present status}{As of November 2026, the splitting is correct; the side convention, canonical-bundle twist and adjoint sign are essential.}
\GSfield{Assessment}{literal statement: false as written; publication match: partial; verification: read the result; development: annotated; confidence: high}
\end{GSitem}
\begin{GSitem}{1980-496.8}{A}{Does equivalence by global sections imply D-affinity?}
\GSfield{Kind; source}{question; \GSpage{P000932}}
\GSfield{The notes}{D-affinity implies that \(\Gamma:\mathcal D_X\text{-mod}\to\Gamma(\mathcal D_X)\text{-mod}\) is an equivalence. The speaker says he does not know the converse.}
\GSfield{Statement}{For the category of O-quasi-coherent left D-modules on a smooth separated finite-type variety, D-affinity is equivalent to the natural adjunction between global sections and localization being an equivalence. In the reverse implication, equivalence gives exactness and detection of nonzero objects; an affine \v{C}ech resolution computes the higher sheaf cohomology.}
\GSfield{Fate}{The notebook itself supplies the converse one week later, in record 498, by the compact projective generator argument. This abstract mechanism is the standard module-category criterion, stated with proof outline in \citet[\S1.2]{keller1997categories}; the localization application is \citet{beilinsonbernstein1981localisation}.}
\GSfield{Present status}{As of November 2026, the question is answered affirmatively under these hypotheses, including in the next recorded talk.}
\GSfield{Assessment}{literal statement: correct with missing hypotheses; publication match: exact; verification: read the result; development: none; confidence: high; problem \hyperref[prob:GP-1980-03]{GP-1980-03}}
\end{GSitem}
\begin{GSitem}{1980-496.9}{E}{The canonical class obstruction and the limits of Chern-class arguments}
\GSfield{Kind; source}{claim; \GSpage{P000932}}
\GSfield{The notes}{Vanishing of the first Chern class of the canonical or tangent bundle is said to imply non-D-affinity and existence of an involution. An insertion identifies line-bundle isomorphism with equality of first Chern classes.}
\GSfield{Statement}{Without properness the non-D-affinity claim is false: affine space has trivial canonical bundle and is D-affine. A safe precise special case is that a smooth connected projective variety of positive dimension with trivial canonical bundle is not D-affine, since \(H^{\dim X}(X,\mathcal O_X)\simeq\mathbb C\) by Serre duality. Equality of first Chern classes does not imply algebraic or holomorphic isomorphism of line bundles: nontrivial degree-zero line bundles on an elliptic curve have the same Chern class as \(\mathcal O\).}
\GSfield{Fate}{The obstruction is an elementary application of coherent duality; the global transpose naturally involves the canonical bundle, as in \citet[Introduction]{kashiwara1984riemann}. The source does not specify whether its Chern class is integral, rational, de Rham or an Atiyah class, so the stronger claimed equivalence is not certified.}
\GSfield{Present status}{As of November 2026, the unqualified implications are false. The stated projective trivial-canonical case is established; it is not claimed to exhaust the intended theorem.}
\GSfield{Assessment}{literal statement: false as written; publication match: partial; verification: read the result; development: annotated; confidence: medium}
\end{GSitem}
\begin{GSitem}{1980-496.10}{A}{Global operators on the flag variety and localization}
\GSfield{Kind; source}{theorem; \GSpage{P000932}}
\GSfield{The notes}{For semisimple \(G\) and Borel \(B\), \(X=G/B\) is D-affine and \(\Gamma\mathcal D_X=U(\mathfrak g)/U(\mathfrak g)I_0\), where \(I_0\) is the kernel of the trivial central character. Hence D-modules correspond to representations with that infinitesimal character.}
\GSfield{Statement}{For a connected complex semisimple algebraic group \(G\), put \(X=G/B\) and let \(\chi_{\rm triv}:Z(U\mathfrak g)\to\mathbb C\) be the character on the trivial representation. The infinitesimal action induces \(U(\mathfrak g)/U(\mathfrak g)\ker\chi_{\rm triv}\simeq\Gamma(X,\mathcal D_X)\). Global sections and localization give inverse equivalences between O-quasi-coherent left \(\mathcal D_X\)-modules and modules over this central reduction. This uses the exact central character, not merely generalized central-character torsion.}
\GSfield{Fate}{\citet[pp.~15--18]{beilinsonbernstein1981localisation} is the direct speaker publication (1981; seminar exposition November 1980). \citet[theorem on p.~43 and addendum p.~44]{brylinski1981differential}, from the August--September 1980 Toru\'n conference, independently proves the global-ring assertion and records the Beilinson--Bernstein all-module vanishing.}
\GSfield{Present status}{As of November 2026, the untwisted localization theorem is established. Its regular dominant character convention is fixed here by the trivial representation.}
\GSfield{Assessment}{literal statement: correct; publication match: exact; verification: read the result; development: none; confidence: high}
\end{GSitem}
\begin{GSitem}{1980-496.11}{E}{The projective-line operator algebra and its central relation}
\GSfield{Kind; source}{construction; \GSpage{P000932}}
\GSfield{The notes}{On \(\mathbb P^1\) the operators \(\partial_x,x\partial_x,x^2\partial_x\) generate the algebra; the notes say there is a unique operator commuting with all of them.}
\GSfield{Statement}{The three vector fields span \(\mathfrak{sl}_2\). Their universal enveloping algebra maps onto \(\Gamma(\mathbb P^1,\mathcal D)\), with kernel generated by the Casimir minus its trivial-character scalar. The center of \(U(\mathfrak{sl}_2)\) is a polynomial algebra in one Casimir; this means one algebraically independent generator, not one commuting operator. In the global differential-operator quotient the center consists of scalar operators.}
\GSfield{Fate}{The rank-one example already occurs in 1967-144.13. The general quotient description is proved in \citet[p.~43]{brylinski1981differential}; the Harish--Chandra center statement is recalled on p.~45. The generation and relation can also be checked directly in one variable.}
\GSfield{Present status}{As of November 2026, the example is classical; uniqueness must be read as uniqueness of a central generator up to normalization, not literally a single commuting operator.}
\GSfield{Assessment}{literal statement: false as written; publication match: partial; verification: read the result; development: none; confidence: high}
\end{GSitem}

\GStalk{499}{1980-11-10}{I. M. Gelfand}{P000933}{Cayley numbers, absolute-valued algebras, and projective planes}
The actual intervention is the upper part of P000933, before Beilinson’s heading. Its questions concern multiplicative norms, the finite-dimensional restriction, a possible infinite-dimensional example, projective-plane coordinatization and a geometric approach to exceptional Lie theory. The notebook does not state a unit axiom. That omission changes both the finite classification and the answer to the infinite-dimensional question, and is recorded explicitly. No content from Beilinson’s following pages is assigned to this record.

\begin{GSitem}{1980-499.1}{E}{The Cayley norm and multiplicative composition}
\GSfield{Kind; source}{definition; \GSpage{P000933}}
\GSfield{The notes}{In a basis of eight Cayley units, \(\lvert\sum x_je_j\rvert^2=\sum x_j^2\), and \(|ab|=|a||b|\).}
\GSfield{Statement}{The real octonion algebra \(\mathbb O\) is an eight-dimensional unital, nonassociative alternative algebra. Conjugation gives \(N(a)=a\bar a=\sum_{j=1}^8x_j^2\), and \(N(ab)=N(a)N(b)\). Thus its Euclidean norm is multiplicative.}
\GSfield{Fate}{This is the classical Cayley composition algebra; see \citet[\S\S2.1--2.2]{baez2002octonions}. The note asks for a review of its classification, not a newly announced composition law.}
\GSfield{Present status}{As of November 2026, the norm identity is correct.}
\GSfield{Assessment}{literal statement: correct; publication match: exact; verification: read the result; development: none; confidence: high}
\end{GSitem}
\begin{GSitem}{1980-499.2}{E}{Dimensions of finite-dimensional absolute-valued algebras}
\GSfield{Kind; source}{theorem; \GSpage{P000933}}
\GSfield{The notes}{There are none of dimension greater than eight, and in finite dimensions only the real numbers, complex numbers, quaternions and Cayley numbers occur.}
\GSfield{Statement}{A nonzero finite-dimensional real algebra with a bilinear product and a multiplicative norm has dimension \(1,2,4\), or \(8\). If it has a two-sided identity, it is isometrically isomorphic to \(\mathbb R,\mathbb C,\mathbb H\), or \(\mathbb O\). Without an identity the dimension assertion remains true, but there are additional isotopes, so the four-algebra isomorphism list is incomplete.}
\GSfield{Fate}{The unital composition classification and its representation-theoretic proof are reviewed in \citet[Theorem~1 and \S2.3]{baez2002octonions}. The reduction of a finite-dimensional absolute-valued algebra to a unital isotope is explicitly proved before Proposition~1 in \citet[pp.~1--2]{kaidiramirezrodriguez1999absolute}.}
\GSfield{Present status}{As of November 2026, the dimension restriction is established. The unqualified isomorphism classification needs the missing unit hypothesis.}
\GSfield{Assessment}{literal statement: correct with missing hypotheses; publication match: exact; verification: read the result; development: none; confidence: high}
\end{GSitem}
\begin{GSitem}{1980-499.3}{E}{The Clifford-algebra route to the dimension restriction}
\GSfield{Kind; source}{method; \GSpage{P000933}}
\GSfield{The notes}{The norm identity gives a representation-theoretic problem and an explicit construction.}
\GSfield{Statement}{For a unital finite-dimensional Euclidean composition algebra A, left multiplication by vectors orthogonal to 1 satisfies \(L_uL_v+L_vL_u=-2\langle u,v\rangle\operatorname{Id}\). Thus A is a module for the real Clifford algebra of its imaginary subspace. Comparing possible module dimensions forces \(\dim A\in\{1,2,4,8\}\); the four usual algebras realize these possibilities. For a nonunital finite-dimensional algebra, pass first to a unital isotope.}
\GSfield{Reconstruction}{The Clifford action and dimension comparison are supplied from the cited account; the source names neither the algebra nor the precise relation.}
\GSfield{Fate}{The explicit representation-theoretic proof is \citet[\S2.3]{baez2002octonions}; the identity is obtained by polarizing the multiplicative Euclidean norm. This is an identification and expansion of the source’s one-line method, not a newly supplied classification proof.}
\GSfield{Present status}{As of November 2026, the Clifford argument is established and explains the speaker’s reference to representation theory.}
\GSfield{Assessment}{literal statement: correct with missing hypotheses; publication match: exact; verification: read the result; development: none; confidence: high}
\end{GSitem}
\begin{GSitem}{1980-499.4}{E}{Can a multiplicative norm exist in infinite dimension?}
\GSfield{Kind; source}{question; \GSpage{P000933}; S. G. Gindikin (affirmative answer to Gelfand)}
\GSfield{The notes}{Gelfand asks whether the multiplicative-norm identity can hold in infinite dimension; Gindikin says yes.}
\GSfield{Statement}{Without a unit, yes: take \(H=\ell^2(\mathbb N,\mathbb R)\), choose a bijection \(b:\mathbb N^2\to\mathbb N\), and define \((x*y)_{b(i,j)}=x_iy_j\). This is a continuous bilinear product with \(\|x*y\|^2=\sum_{i,j}|x_i|^2|y_j|^2=\|x\|^2\|y\|^2\). It has no two-sided identity: under the induced unitary \(H\mathbin{\widehat\otimes}_2H\to H\), the range of multiplication by a nonzero fixed x is the proper subspace corresponding to \(\mathbb Rx\otimes H\). With a two-sided identity, no infinite-dimensional example exists: the Urbanik--Wright theorem gives precisely \(\mathbb R,\mathbb C,\mathbb H,\mathbb O\), without assuming finite dimension or completeness.}
\GSfield{Reconstruction}{The explicit example is supplied from the identified literature, not attributed historically to Gindikin.}
\GSfield{Fate}{The Urbanik--Wright theorem (1960) is stated explicitly in \citet[p.~2, preceding Corollary~1]{kaidiramirezrodriguez1999absolute}. The tensor-coordinate example is precisely the Hilbert-space case of \citet[Proposition~2.2 and Corollary~2.5]{becerramorenorodriguez2005absolute}. Its elementary verification here is an annotation and is not asserted to reproduce Gindikin’s unpublished argument.}
\GSfield{Present status}{As of November 2026, the question is resolved under either explicit convention. Gindikin’s affirmative answer is consistent with the nonunital reading and inconsistent with the unital one.}
\GSfield{Assessment}{literal statement: not determinable; publication match: partial; verification: read the result; development: annotated; confidence: high; problem \hyperref[prob:GP-1980-04]{GP-1980-04}}
\end{GSitem}
\begin{GSitem}{1980-499.5}{E}{Associativity, Desargues, and the Cayley projective plane}
\GSfield{Kind; source}{theorem; \GSpage{P000933}}
\GSfield{The notes}{After coordinatization, associativity is identified with Desargues; the Cayley projective plane is called the simplest non-Desarguesian example.}
\GSfield{Statement}{For projective planes coordinatized by an alternative division algebra, the plane is Desarguesian precisely when that algebra is associative. The octonionic Moufang plane \(\mathbb OP^2\) is therefore non-Desarguesian. For an arbitrary projective plane, coordinatization uses a planar ternary ring and the bare word associativity is not a complete formulation. Calling the Cayley plane simplest is qualitative, not a minimality theorem.}
\GSfield{Fate}{The construction of \(\mathbb OP^2\), its incidence geometry and its relation to the failure of associativity are reviewed in \citet[\S3]{baez2002octonions}.}
\GSfield{Present status}{As of November 2026, the division-algebra formulation is established; the source’s shorthand requires the coordinatization category.}
\GSfield{Assessment}{literal statement: correct with missing hypotheses; publication match: exact; verification: read the result; development: none; confidence: high}
\end{GSitem}
\begin{GSitem}{1980-499.6}{F}{Differential and locally homogeneous projective planes}
\GSfield{Kind; source}{remark; \GSpage{P000933}}
\GSfield{The notes}{It is interesting to study differential and locally homogeneous projective planes.}
\GSfield{Statement}{One could impose smooth point and line manifolds with smooth incidence and join/intersection maps, then ask for transitive or locally transitive collineation groups. The notes do not fix these regularity axioms, the dimension, completeness, the acting group or an equivalence relation, so they specify no unique classification question.}
\GSfield{Fate}{The program is retained because its geometric object is identifiable, but no precise assertion can be reconstructed from this single sentence. The modern study of smooth projective planes does not by itself answer an unspecified classification request.}
\GSfield{Present status}{As of November 2026, the source remains underdetermined; no new open problem or claim of nondevelopment is inferred.}
\GSfield{Assessment}{literal statement: not determinable; publication match: none found; verification: unresolved lead; development: none; confidence: low}
\end{GSitem}
\begin{GSitem}{1980-499.7}{F}{A unified geometric explanation of exceptional Lie types}
\GSfield{Kind; source}{remark; \GSpage{P000933}}
\GSfield{The notes}{Classifying Lie algebras is called too particular; the exceptional types \(G_2,E_6,E_7,E_8\) should not appear as isolated fragments.}
\GSfield{Statement}{The request is for a conceptual geometric organization of exceptional Lie theory. Classical octonionic and Jordan-algebra constructions do provide relationships among exceptional groups, but the notebook gives no proposed object, theorem or classification target whose completion can be assessed.}
\GSfield{Fate}{\citet[\S4]{baez2002octonions} supplies a precise later account of octonionic constructions of exceptional Lie algebras. It is a thematic comparison, not evidence for a uniquely recoverable theorem in this source.}
\GSfield{Present status}{As of November 2026, the mathematical motivation is clear, but an exact assertion is not recoverable. This item is not classified as an undeveloped conjecture.}
\GSfield{Assessment}{literal statement: not determinable; publication match: thematic only; verification: indirect; development: none; confidence: low}
\end{GSitem}

\GStalk{498}{1980-11-10}{A. A. Beilinson}{P000933-P000936}{A-affinity, projective generators, localization and transfer bimodules}
The actual heading is in the lower part of P000933; the talk continues through P000936. It proves the categorical affinity criterion left uncertain on 3 November, formulates the compact projective generator lemma, repeats projective-space and flag-variety localization, and begins the transfer-bimodule construction. The preceding Cayley-algebra intervention on the same page is record 499. The short catalogue range stops before most of the proof and must not determine the content boundary.

\begin{GSitem}{1980-498.1}{E}{Global sections and localization for a sheaf of rings}
\GSfield{Kind; source}{construction; \GSpage{P000933-P000934}}
\GSfield{The notes}{Let \(\mathcal A\) contain \(\mathcal O_X\) and be quasi-coherent over it. For \(A=\Gamma(X,\mathcal A)\), the functors are \(\Gamma\) and \(\Delta(M)=\mathcal A\otimes_A M\); the source calls the constant sheaf of M a sheaf of \(\mathcal A\)-modules.}
\GSfield{Statement}{Let X be a separated finite-type complex variety and \(\mathcal A\) a sheaf of unital rings with a given inclusion of \(\mathcal O_X\), quasi-coherent for its left O-action. In the category of left A-sheaves whose underlying O-module is quasi-coherent, localization \(\Delta(M)=\mathcal A\otimes_{\underline A}\underline M\) is left adjoint to global sections. Here \(\underline A\) and \(\underline M\) are constant sheaves, with the natural map \(\underline A\to\mathcal A\). The O-subring need not be central; differential operators are the motivating example.}
\GSfield{Fate}{This is the standard tensor--Hom adjunction used by the projective-generator proof, cf. \citet[\S1.2]{keller1997categories}. It generalizes the affine sheaf/module construction of \citet[Chapter II, \S3]{serre1955coherent}.}
\GSfield{Present status}{As of November 2026, the adjunction is valid. A general left A-module does not already carry a canonical constant-sheaf action by the varying ring \(\mathcal A\).}
\GSfield{Assessment}{literal statement: false as written; publication match: partial; verification: read the result; development: none; confidence: high}
\end{GSitem}
\begin{GSitem}{1980-498.2}{A}{Cohomological affinity and equivalence of module categories}
\GSfield{Kind; source}{theorem; \GSpage{P000934-P000935}}
\GSfield{The notes}{Two conditions are stated to be equivalent: every quasi-coherent A-module has no higher cohomology and every nonzero one has a nonzero global section; or \(\Gamma,\Delta\) are inverse equivalences. The reverse direction is assigned to the \v{C}ech resolution.}
\GSfield{Statement}{For X and \(\mathcal A\) as in the preceding item, assume the category is closed under the indicated quasi-coherent operations. Vanishing of \(H^{i>0}(X,\mathcal M)\) for every object and detection of nonzero objects by \(\Gamma\) are equivalent to the natural localization/global-sections adjunction being an equivalence. They also imply D-style generation by global sections: the cokernel of the evaluation map is zero. For the converse, use an affine \v{C}ech resolution whose terms and kernels remain in the category.}
\GSfield{Fate}{This explicitly resolves 1980-496.8 and GP-1980-03 in the source itself. The general category-theoretic mechanism is \citet[\S1.2]{keller1997categories}; the D-module application is the setting of \citet{beilinsonbernstein1981localisation}.}
\GSfield{Present status}{As of November 2026, the criterion is established under the stated quasi-coherent and finite-cover hypotheses. Detection of nonzero objects and exactness, taken together, suffice to recover the generation condition.}
\GSfield{Assessment}{literal statement: correct with missing hypotheses; publication match: exact; verification: read the result; development: none; confidence: high; problem \hyperref[prob:GP-1980-03]{GP-1980-03}}
\end{GSitem}
\begin{GSitem}{1980-498.3}{E}{A compact projective generator identifies an abelian category with modules}
\GSfield{Kind; source}{theorem; \GSpage{P000934-P000935}}
\GSfield{The notes}{An abelian category has arbitrary direct sums and a projective generator P with \(\operatorname{Hom}(P,-)\) preserving them. The notes identify its target as left \(\operatorname{End}(P)\)-modules and write \(\operatorname{End}_{\mathcal A}(\mathcal A)=A\).}
\GSfield{Statement}{For a locally small abelian category with arbitrary coproducts and a compact projective generator P, \(\operatorname{Hom}(P,-)\) is an equivalence to right modules over \(\operatorname{End}(P)\), or left modules over its opposite ring. Compact means preservation of arbitrary coproducts. For left \(\mathcal A\)-modules, \(\operatorname{End}_{\mathcal A}(\mathcal A)\simeq\Gamma(X,\mathcal A)^{\mathrm{op}}\), so the resulting target is indeed left A-modules. On a quasi-compact separated variety, global sections of quasi-coherent sheaves commute with coproducts.}
\GSfield{Fate}{This classical module-category criterion is stated with its historical attribution and proof outline in \citet[\S1.2, p.~4]{keller1997categories}. The proof takes a presentation by coproducts of P and applies its exact Hom functor. The source has the compactness hypothesis but omits the opposite-ring convention.}
\GSfield{Present status}{As of November 2026, the lemma is established. The missing opposite ring is a literal error under the explicitly chosen left-module convention.}
\GSfield{Assessment}{literal statement: false as written; publication match: exact; verification: read the result; development: annotated; confidence: high; problem \hyperref[prob:GP-1980-03]{GP-1980-03}}
\end{GSitem}
\begin{GSitem}{1980-498.4}{A}{Projective space is D-affine}
\GSfield{Kind; source}{theorem; \GSpage{P000935}}
\GSfield{The notes}{Theorem (13): \(\mathbb P^N\) is D-affine.}
\GSfield{Statement}{For complex projective space and the sheaf of finite-order algebraic differential operators, every O-quasi-coherent left D-module is acyclic for global sections and is generated as a D-module by them. Consequently the localization/global-sections adjunction is an equivalence.}
\GSfield{Fate}{This repeats the projective-space assertion in 1980-496.5. It follows from the localization theorem of \citet{beilinsonbernstein1981localisation}; the all-module vanishing is explicitly recorded in \citet[p.~44]{brylinski1981differential}.}
\GSfield{Present status}{As of November 2026, the theorem is established; the source means algebraic D-affinity, not ordinary affineness of projective space.}
\GSfield{Assessment}{literal statement: correct; publication match: exact; verification: read the result; development: none; confidence: high}
\end{GSitem}
\begin{GSitem}{1980-498.5}{A}{Flag-variety localization for a reductive group}
\GSfield{Kind; source}{theorem; \GSpage{P000935}}
\GSfield{The notes}{For a reductive group G and Borel B, the Lie action gives \(U(\mathfrak g)\to\Gamma\mathcal D_{G/B}\); quotienting by the central annihilator of the trivial representation gives an isomorphism and the expected category equivalence.}
\GSfield{Statement}{For a connected complex reductive algebraic group G, let \(X=G/B\) and let \(I\subset Z(U\mathfrak g)\) be the kernel of the trivial central character. Then X is D-affine and \(U(\mathfrak g)/U(\mathfrak g)I\simeq\Gamma(X,\mathcal D_X)\). Thus O-quasi-coherent left D-modules correspond to all left modules over that central reduction. The central Lie algebra of G acts trivially on X and is killed by I; the reductive formulation reduces to the semisimple one.}
\GSfield{Fate}{The exact theorem is \citet[pp.~15--18]{beilinsonbernstein1981localisation}; compare \citet[p.~43 and addendum p.~44]{brylinski1981differential}. This is a substantive second occurrence of 1980-496.10, now preceded by the categorical proof.}
\GSfield{Present status}{As of November 2026, the theorem is established with connectedness and the exact trivial infinitesimal character specified.}
\GSfield{Assessment}{literal statement: correct with missing hypotheses; publication match: exact; verification: read the result; development: none; confidence: high}
\end{GSitem}
\begin{GSitem}{1980-498.6}{E}{The transfer bimodule and its two commuting operator actions}
\GSfield{Kind; source}{construction; \GSpage{P000936}}
\GSfield{The notes}{The source defines \(\mathcal D_{X\to Y}=\mathcal O_X\otimes_{f^{-1}\mathcal O_Y}f^{-1}\mathcal D_Y\), with a left \(D_X\) action by Leibniz and a right pulled-back \(D_Y\) action. The last sentence says df(v) acts on \(f^{-1}\mathcal D_Y\) itself.}
\GSfield{Statement}{For a morphism f of smooth complex algebraic varieties, the displayed tensor is a \((\mathcal D_X,f^{-1}\mathcal D_Y)\)-bimodule. Locally, if \(df(v)=\sum_i a_i\otimes w_i\), then \[v\cdot(b\otimes P)=v(b)\otimes P+\sum_i ba_i\otimes w_iP.\] The right action is multiplication on P. The coefficients \(a_i\) are functions on X, so the formula acts on the tensor product, not generally on the unextended inverse-image sheaf alone.}
\GSfield{Fate}{The transfer-module formula is standard D-module functoriality; see \citet[\S4.3, p.~343]{kashiwara1984riemann}. The notebook captures the correct bimodule but abbreviates the vector-field action too far.}
\GSfield{Present status}{As of November 2026, the construction is established; \(O_X\) scalar extension is essential to the claimed left action.}
\GSfield{Assessment}{literal statement: correct with missing hypotheses; publication match: exact; verification: read the result; development: annotated; confidence: high}
\end{GSitem}
\begin{GSitem}{1980-498.7}{E}{Derived direct and inverse image in place of a ring pullback}
\GSfield{Kind; source}{method; \GSpage{P000936}}
\GSfield{The notes}{There is no corresponding morphism of rings of differential operators for a general map, but natural functors exist in derived categories; the notes say there are four.}
\GSfield{Statement}{For a morphism of smooth varieties, differential operators do not pull back by a functorial homomorphism of sheaves of rings in the way functions do. Derived inverse and direct image use transfer bimodules and derived tensor products, with the canonical-bundle twist for the reverse transfer module. On the appropriate holonomic derived categories, duality produces the complementary extraordinary operations. Preservation of regularity and comparison with constructible sheaves require the regular-holonomic hypotheses.}
\GSfield{Fate}{These operations are part of the holonomic-system formalism in \citet[\S4.3 and \S8]{kashiwara1984riemann}; the regular derived correspondence was already announced in \citet[p.~XIX.1]{kashiwara1980constructible}. The source does not define four functors on every unbounded module category.}
\GSfield{Present status}{As of November 2026, the derived functoriality is established, with category and duality restrictions; the notebook’s count of four is a programmatic abbreviation.}
\GSfield{Assessment}{literal statement: correct with missing hypotheses; publication match: partial; verification: read the result; development: none; confidence: high}
\end{GSitem}

\GStalk{501}{1980-11-17}{I. M. Gelfand}{P000937}{Elementary analysis exercises and rapidly oscillating integrals}
Four numbered exercises precede Beilinson's heading: Thomae's function and its powers, a double-limit representation of the Dirichlet function, the obstruction to a single limit of continuous functions, and Fourier integrals with quadratic, quartic, and sextic decay. These are definite questions with classical solutions, not evidence of newly announced unresolved conjectures.

\begin{GSitem}{1980-501.1}{E}{Continuity and differentiability of Thomae’s function}
\GSfield{Kind; source}{problem; \GSpage{P000937}}
\GSfield{The notes}{Exercise 1 asks for the continuity and differentiability of \(f(p/q)=1/q\) for reduced fractions and \(f(x)=0\) for irrational \(x\).}
\GSfield{Statement}{On \(\mathbb R\), with \(q>0\) and \(\gcd(p,q)=1\), this function is continuous exactly at the irrational points and differentiable nowhere. At an irrational point, any derivative would be zero along irrational increments, whereas the convergents \(p_k/q_k\) satisfy \(|x-p_k/q_k|<q_k^{-2}\), so the absolute difference quotients exceed \(q_k\).}
\GSfield{Reconstruction}{The domain, reduced denominators, and the short differentiability proof are explicit completion of the exercise.}
\GSfield{Fate}{This is the classical Thomae example. \citet[pp.353--354]{varona2009differentiability} reproduces the exact function and proof; the 2009 article is a later exposition and extension, not the origin of this elementary result.}
\GSfield{Present status}{As of November 2026, the statement with the displayed hypotheses is established; the notebook occurrence is expository.}
\GSfield{Assessment}{literal statement: correct; publication match: exact; verification: read the result; development: annotated; confidence: high; problem \hyperref[prob:GP-1980-20]{GP-1980-20}}
\end{GSitem}
\begin{GSitem}{1980-501.2}{E}{Differentiability of powers of Thomae’s function}
\GSfield{Kind; source}{problem; \GSpage{P000937}}
\GSfield{The notes}{The same exercise asks what changes when \(1/q\) is replaced by \(1/q^\alpha\).}
\GSfield{Statement}{For real \(\alpha\), let \(f_\alpha(p/q)=q^{-\alpha}\), \(f_\alpha(x)=0\) at irrationals. If \(\alpha\le0\), it is nowhere continuous; if \(\alpha>0\), continuity holds exactly at irrationals. It is nowhere differentiable for \(\alpha\le2\). For \(\alpha>2\) and irrational \(x=[a_0;a_1,\ldots]\) with convergent denominators \(q_k\), differentiability is equivalent to \(q_{k+1}/q_k^{\alpha-1}\to0\), and then \(f_\alpha^{\prime}(x)=0\). Equivalently, \(q^\alpha|x-p/q|\to\infty\) along all reduced rationals tending to \(x\).}
\GSfield{Reconstruction}{The exact continued-fraction criterion is our explicit reformulation of the difference quotient. Necessity uses \((q_k(q_k+q_{k+1}))^{-1}<|x-p_k/q_k|<(q_kq_{k+1})^{-1}\); sufficiency also uses that a nonconvergent approximation has error at least \(1/(2q^2)\). Statements using only the irrationality exponent require care at equality.}
\GSfield{Fate}{The family had already been studied by Nymann in 1969, as documented in \citet[introduction, \S\S1--3, reference 7]{varona2009differentiability}; that article gives further exact differentiability and measure statements. Thus the exercise is not classified as an unresolved 1980 conjecture. No assertion that the present denominator formulation first appeared in 1969 is made.}
\GSfield{Present status}{As of November 2026, the statement with the displayed hypotheses is established; the notebook occurrence is expository.}
\GSfield{Assessment}{literal statement: correct; publication match: exact; verification: read the result; development: annotated; confidence: high; problem \hyperref[prob:GP-1980-20]{GP-1980-20}}
\end{GSitem}
\begin{GSitem}{1980-501.3}{E}{A double continuous limit for the Dirichlet function}
\GSfield{Kind; source}{construction; \GSpage{P000937}}
\GSfield{The notes}{Exercise 2 gives \(\lim_{n\to\infty}\lim_{m\to\infty}\cos^{2m}(2\pi n!x)\).}
\GSfield{Statement}{For every real \(x\), this iterated limit equals \(\mathbf1_{\mathbb Q}(x)\). For fixed \(n\), the inner limit is one exactly when \(2n!x\in\mathbb Z\), and zero otherwise. Every rational satisfies the condition for all sufficiently large \(n\), whereas an irrational never does.}
\GSfield{Reconstruction}{The order of limits is essential and is retained. This elementary calculation proves the formula without interchanging limits.}
\GSfield{Fate}{This is the classical distinction between first and second Baire class. The underlying classification was established before 1900; \citet[p.173]{baire1900nouvelle} states the continuous-limit criterion. The displayed factorial formula is checked directly here, not attributed to that article.}
\GSfield{Present status}{As of November 2026, the statement with the displayed hypotheses is established; the notebook occurrence is expository.}
\GSfield{Assessment}{literal statement: correct; publication match: exact; verification: read the result; development: annotated; confidence: high}
\end{GSitem}
\begin{GSitem}{1980-501.4}{E}{Can the Dirichlet function be a pointwise limit of continuous functions?}
\GSfield{Kind; source}{problem; \GSpage{P000937}}
\GSfield{The notes}{Exercise 3 asks whether a single sequence of continuous functions can tend pointwise to the Dirichlet function.}
\GSfield{Statement}{No sequence \(g_j\in C(\mathbb R)\) converges at every real point to \(\mathbf1_{\mathbb Q}\). A real-valued Baire-class-one function has a dense \(G_\delta\) set of continuity points on every nonempty closed interval, while \(\mathbf1_{\mathbb Q}\) is nowhere continuous.}
\GSfield{Reconstruction}{Pointwise convergence at every point is the required meaning; convergence almost everywhere would give a different, trivial answer.}
\GSfield{Fate}{The negative answer is an immediate classical consequence of Baire's criterion; see \citet[p.173, opening theorem and reference to the 1899 thesis]{baire1900nouvelle}.}
\GSfield{Present status}{As of November 2026, the statement with the displayed hypotheses is established; the notebook occurrence is expository.}
\GSfield{Assessment}{literal statement: correct; publication match: exact; verification: read the result; development: annotated; confidence: high; problem \hyperref[prob:GP-1980-21]{GP-1980-21}}
\end{GSitem}
\begin{GSitem}{1980-501.5}{E}{Fourier integrals with quadratic, quartic, and sextic decay}
\GSfield{Kind; source}{problem; \GSpage{P000937}}
\GSfield{The notes}{Exercise 4 asks for approximate values of \(\int_{-\infty}^{\infty}e^{-x^2}\cos(1000x)\,dx\), and the analogues with \(x^4\) and \(x^6\).}
\GSfield{Statement}{Write \(I_p(\lambda)=\int_{\mathbb R}e^{-x^p}\cos(\lambda x)\,dx\). Then \(I_2(1000)=\sqrt\pi e^{-250000}\). For \(p=4,6\), put \(r=(\lambda/p)^{1/(p-1)}\), \(\theta=\pi/(2(p-1))\), \(A=(p-1)r^p\). As \(\lambda\to+\infty\), the two dominant conjugate saddle contributions give \[I_p(\lambda)=2\sqrt{\frac{2\pi}{p(p-1)}}r^{-(p-2)/2}e^{-A\sin\theta}\left\{\cos\!\left(A\cos\theta-\frac{p-2}{2}\theta\right)+O(A^{-1})\right\}.\] The remainder is absolute relative to the displayed envelope, not relative to the cosine, which can vanish.}
\GSfield{Reconstruction}{The parameter \(\lambda\), saddle computation, and error convention are supplied; the notebook asks only for approximations. Saddles solve \(pz^{p-1}=i\lambda\). The quartic and sextic formulas are explicit applications of the classical method, not a claimed new asymptotic theorem.}
\GSfield{Fate}{Gaussian Fourier transformation and steepest descent predate the seminar. The simple-saddle expansion and branch convention are checked in \citet[\S2.4(iv), equations 2.4.15--16]{nist2026dlmf}; its cited pre-1980 sources include Olver and Bleistein--Handelsman. The displayed specialization is our calculation.}
\GSfield{Present status}{As of November 2026, the statement with the displayed hypotheses is established; the notebook occurrence is expository.}
\GSfield{Assessment}{literal statement: correct; publication match: exact; verification: read the result; development: annotated; confidence: medium; problem \hyperref[prob:GP-1980-22]{GP-1980-22}}
\end{GSitem}

\GStalk{500}{1980-11-17}{A. A. Beilinson}{P000937-P000939}{Inverse and direct images of D-modules}
The manuscript heading on the lower part of P000937 begins this talk, following I. M. Gelfand's exercises. Beilinson constructs ordinary inverse images of left modules and transfer bimodules for direct images of right modules, tests them on point inclusions and delta modules, and introduces derived integration and its de Rham interpretation. The actual continuation reaches P000939, beyond the catalogue range. P000936 belongs to the preceding meeting.

\begin{GSitem}{1980-500.1}{E}{Pullback of a left D-module}
\GSfield{Kind; source}{construction; \GSpage{P000937-P000938}}
\GSfield{The notes}{For \(f:X\to Y\), the notes put \(f^*M=\mathcal O_X\otimes_{f^{-1}\mathcal O_Y}f^{-1}M\) and differentiate both factors through \(df\).}
\GSfield{Statement}{Let \(X,Y\) be smooth varieties over a characteristic-zero field and \(M\) a left \(\mathcal D_Y\)-module quasicoherent over \(\mathcal O_Y\). If \(df(v)=\sum_i a_i\otimes w_i\), the pullback left action is \(v(b\otimes m)=v(b)\otimes m+\sum_i ba_i\otimes w_i m\). The resulting module is independent of the local expression for \(df(v)\).}
\GSfield{Reconstruction}{Smoothness, characteristic zero, and the inverse-image sheaves are written explicitly. This is the ordinary tensor pullback; it need not preserve coherent \(\mathcal D\)-modules.}
\GSfield{Fate}{This is the standard inverse-image construction of algebraic \(\mathcal D\)-module theory. The exact bimodule and action are read in \citet[\S7 and \S13]{gourevitch2025dmodules}; the contemporary framework is already represented by \citet[\S1]{bernstein1971modules}. The 2025 course is used as a precise later exposition, not a priority date.}
\GSfield{Present status}{As of November 2026, the statement with the displayed hypotheses is established; the notebook occurrence is expository.}
\GSfield{Assessment}{literal statement: correct with missing hypotheses; publication match: exact; verification: read the result; development: annotated; confidence: high}
\end{GSitem}
\begin{GSitem}{1980-500.2}{E}{Pullback of connections and of the structure sheaf}
\GSfield{Kind; source}{theorem; \GSpage{P000938}}
\GSfield{The notes}{The pullback agrees with the usual pullback of a bundle with flat connection, and \(f^*\mathcal O_Y=\mathcal O_X\).}
\GSfield{Statement}{For a finite-rank algebraic vector bundle \((E,\nabla)\) with integrable connection on smooth \(Y\), the displayed module action gives the integrable pullback connection on \(f^*E\); the trivial connection on \(\mathcal O_Y\) pulls back to the trivial connection on \(\mathcal O_X\).}
\GSfield{Fate}{The compatibility is part of the classical connection--\(\mathcal D\)-module correspondence; compare \citet[Lemma 9.11 and inverse-image construction]{gourevitch2025dmodules}.}
\GSfield{Present status}{As of November 2026, the statement with the displayed hypotheses is established; the notebook occurrence is expository.}
\GSfield{Assessment}{literal statement: correct; publication match: exact; verification: read the result; development: annotated; confidence: high}
\end{GSitem}
\begin{GSitem}{1980-500.3}{E}{Ordinary pullback can destroy D-module coherence}
\GSfield{Kind; source}{claim; \GSpage{P000938}}
\GSfield{The notes}{For a point inclusion \(i:x\hookrightarrow Y\), the fibre of the free module \(\mathcal D_Y\) is the infinite-dimensional space of constant-coefficient differential operators.}
\GSfield{Statement}{If \(\dim Y=d>0\), then \(i^*\mathcal D_Y\simeq k[\partial_1,\ldots,\partial_d]\) as a module over \(\mathcal D_{\{x\}}=k\). Thus a coherent \(\mathcal D_Y\)-module can have an ordinary inverse image that is not coherent over \(k\).}
\GSfield{Reconstruction}{The polynomial description uses local smooth coordinates and the positive-dimensional assumption.}
\GSfield{Fate}{The example follows directly from the Weyl-algebra normal form; see \citet[\S1.2 and inverse images]{gourevitch2025dmodules}. It is a classical warning about inverse images, not a later counterexample to holonomic stability.}
\GSfield{Present status}{As of November 2026, the statement with the displayed hypotheses is established; the notebook occurrence is expository.}
\GSfield{Assessment}{literal statement: correct with missing hypotheses; publication match: exact; verification: read the result; development: annotated; confidence: high}
\end{GSitem}
\begin{GSitem}{1980-500.4}{E}{The ordinary fibre of a delta module vanishes}
\GSfield{Kind; source}{claim; \GSpage{P000938}}
\GSfield{The notes}{The notes assert \(i^*\delta_x=0\), justifying it with the written identity \(\delta(y)=y\delta^{\prime}(y)\).}
\GSfield{Statement}{For the left delta module at the origin of \(\mathbb A^d\), \(\delta_0\simeq k[\partial_1,\ldots,\partial_d]\), one has \(x_i\partial^\alpha\delta_0=-\alpha_i\partial^{\alpha-e_i}\delta_0\). Consequently each coordinate acts surjectively and \(k\otimes_{\mathcal O_{\mathbb A^d}}\delta_0=0\). The derived fibre is nonzero; in the unshifted tensor convention its only cohomology is \(k\) in degree \(-d\).}
\GSfield{Reconstruction}{The sign is repaired: \(y\delta^{\prime}=-\delta\). Derived and ordinary pullbacks are distinguished; the Koszul-complex calculation supplies the derived degree.}
\GSfield{Fate}{The delta-module and closed-embedding calculations are classical; see \citet[Example 9.5 and Theorem 9.1]{gourevitch2025dmodules}. The sign follows from differentiating \(y\delta=0\).}
\GSfield{Present status}{As of November 2026, the statement with the displayed hypotheses is established; the notebook occurrence is expository.}
\GSfield{Assessment}{literal statement: false as written; publication match: exact; verification: read the result; development: annotated; confidence: high}
\end{GSitem}
\begin{GSitem}{1980-500.5}{E}{The transfer bimodule and ordinary right direct image}
\GSfield{Kind; source}{construction; \GSpage{P000938}}
\GSfield{The notes}{The notebook defines \(\mathcal D_{X\to Y}=f^*\mathcal D_Y\) and \(f_*M=f_\bullet(M\otimes_{\mathcal D_X}\mathcal D_{X\to Y})\) for right modules.}
\GSfield{Statement}{For a morphism of smooth varieties, \(\mathcal D_{X\to Y}=\mathcal O_X\otimes_{f^{-1}\mathcal O_Y}f^{-1}\mathcal D_Y\) is a \((\mathcal D_X,f^{-1}\mathcal D_Y)\)-bimodule. Its underived direct image sends a right module \(M\) to \(f_\bullet(M\otimes_{\mathcal D_X}\mathcal D_{X\to Y})\).}
\GSfield{Reconstruction}{The subscript \(\bullet\) is reserved for ordinary sheaf direct image, avoiding confusion with the differential-module direct image.}
\GSfield{Fate}{This transfer-module construction is classical and appears with the derived version in \citet[\S13]{gourevitch2025dmodules}.}
\GSfield{Present status}{As of November 2026, the statement with the displayed hypotheses is established; the notebook occurrence is expository.}
\GSfield{Assessment}{literal statement: correct; publication match: exact; verification: read the result; development: annotated; confidence: high}
\end{GSitem}
\begin{GSitem}{1980-500.6}{E}{Closed-embedding direct images and normal delta derivatives}
\GSfield{Kind; source}{construction; \GSpage{P000938-P000939}}
\GSfield{The notes}{Pushing a scalar at a point gives a delta module; a closed embedding of a subvariety adds delta functions and derivatives in transverse directions.}
\GSfield{Statement}{If \(i:X\hookrightarrow Y\) is a closed embedding of smooth varieties, the right direct image is exact and locally adjoins freely the normal derivatives to a right \(\mathcal D_X\)-module. At a coordinate origin it sends \(k\) to \(\mathcal D_Y/\sum_j y_j\mathcal D_Y\) as a right module.}
\GSfield{Reconstruction}{The closed and smooth hypotheses are essential. The quotient is a right-module quotient by the right ideal generated by multiplication by the coordinates, not its left-module analogue.}
\GSfield{Fate}{The exact identification is the closed-embedding calculation underlying Kashiwara's equivalence; \citet[Theorem 9.1, Example 9.5]{gourevitch2025dmodules} gives the read modern formulation.}
\GSfield{Present status}{As of November 2026, the statement with the displayed hypotheses is established; the notebook occurrence is expository.}
\GSfield{Assessment}{literal statement: correct with missing hypotheses; publication match: exact; verification: read the result; development: annotated; confidence: high}
\end{GSitem}
\begin{GSitem}{1980-500.7}{E}{Canonical right module and conversion from left modules}
\GSfield{Kind; source}{construction; \GSpage{P000939}}
\GSfield{The notes}{Top-degree forms, and top-degree forms valued in a flat bundle, are presented as natural right modules.}
\GSfield{Statement}{For smooth \(X\), \(\omega_X=\Omega_X^{\dim X}\) has the right action \(\eta\cdot v=-\mathcal L_v\eta\). Tensoring a left module \(E\) with \(\omega_X\) gives a right module, with \((\eta\otimes e)\cdot v=(\eta\cdot v)\otimes e-\eta\otimes ve\). This is an equivalence between left and right \(\mathcal D_X\)-modules.}
\GSfield{Reconstruction}{The signs and the full equivalence are supplied to make the note about volume forms precise.}
\GSfield{Fate}{The standard left/right conversion is read in \citet[Exercise 7.16]{gourevitch2025dmodules}; it belongs to the established theory used in the talk.}
\GSfield{Present status}{As of November 2026, the statement with the displayed hypotheses is established; the notebook occurrence is expository.}
\GSfield{Assessment}{literal statement: correct with missing hypotheses; publication match: exact; verification: read the result; development: annotated; confidence: high}
\end{GSitem}
\begin{GSitem}{1980-500.8}{E}{Derived pullback and derived integration}
\GSfield{Kind; source}{construction; \GSpage{P000939}}
\GSfield{The notes}{The notes replace pullback by \(Lf^*\) and define \(\int M=Rf_\bullet(M\otimes^L_{\mathcal D_X}\mathcal D_{X\to Y})\), with its cohomology objects \(\int^iM\).}
\GSfield{Statement}{For bounded complexes of quasicoherent right \(\mathcal D_X\)-modules, define \(f_+M=Rf_\bullet(M\otimes^L_{\mathcal D_X}\mathcal D_{X\to Y})\). For left modules, the unshifted inverse image is \(\mathcal D_{X\to Y}\otimes^L_{f^{-1}\mathcal D_Y}f^{-1}N\). These derived functors retain the higher Tor and sheaf-cohomology terms lost by ordinary operations.}
\GSfield{Reconstruction}{A bounded setting and the distinction from dimension-shifted extraordinary inverse image are stated explicitly.}
\GSfield{Fate}{These are established derived operations; see \citet[\S13]{gourevitch2025dmodules} and the derived-functor formalism in \citet[\S10.5]{weibel1994homological}.}
\GSfield{Present status}{As of November 2026, the statement with the displayed hypotheses is established; the notebook occurrence is expository.}
\GSfield{Assessment}{literal statement: correct; publication match: exact; verification: read the result; development: annotated; confidence: high}
\end{GSitem}
\begin{GSitem}{1980-500.9}{E}{Integration of the canonical module computes shifted de Rham cohomology}
\GSfield{Kind; source}{theorem; \GSpage{P000939}}
\GSfield{The notes}{For smooth complex \(X\), the final formula is \(\int^i\omega_X=H^{\dim X+i}(X,\mathbb C)\).}
\GSfield{Statement}{If \(X\) is a smooth complex algebraic variety of dimension \(d\), and \(p:X\to\mathrm{pt}\), then \(p_+\omega_X\simeq R\Gamma(X,\Omega_X^\bullet)[d]\). Hence \(H^i(p_+\omega_X)\simeq H^{i+d}(X^{\mathrm{an}},\mathbb C)\) by algebraic de Rham comparison. This is ordinary cohomology, not cohomology with compact support.}
\GSfield{Reconstruction}{The algebraic/analytic comparison and the cohomological shift convention are made explicit. Properness is unnecessary for this ordinary algebraic de Rham comparison.}
\GSfield{Fate}{The de Rham description of direct image is classical; see \citet[\S13, de Rham direct image]{gourevitch2025dmodules}. The notebook states an already established comparison, rather than a new 1980 theorem.}
\GSfield{Present status}{As of November 2026, the statement with the displayed hypotheses is established; the notebook occurrence is expository.}
\GSfield{Assessment}{literal statement: correct; publication match: exact; verification: read the result; development: annotated; confidence: high}
\end{GSitem}

\GStalk{502}{1980-11-17}{S. I. Gelfand}{P000940 (upper part)}{Two exercises on derived categories}
The upper portion of P000940, before the explicit 24 November heading, gives an exercise separating isomorphic cohomology from a chain map and an incompletely specified dual-sphere equivalence. The latter does not determine a kernel, an equivariance convention, or a precise category; its resemblance to a Radon transform is insufficient for an exact identification.

\begin{GSitem}{1980-502.1}{E}{Isomorphic cohomology need not lift to a chain map}
\GSfield{Kind; source}{problem; \GSpage{P000940}}
\GSfield{The notes}{Exercise 1 says that an isomorphism of the cohomology of two complexes need not extend to a chain map inducing it.}
\GSfield{Statement}{In complexes of abelian groups, let \(A=(\mathbb Z/2)[0]\), and let \(B\) be \(\mathbb Z\xrightarrow{2}\mathbb Z\) in degrees \(-1,0\). Both have only \(H^0\simeq\mathbb Z/2\). Every chain map \(A\to B\) is zero because \(\operatorname{Hom}(\mathbb Z/2,\mathbb Z)=0\); hence none induces the prescribed identity on \(H^0\). They are nevertheless isomorphic in the derived category through the quasi-isomorphism \(B\to A\).}
\GSfield{Reconstruction}{The explicit example and the direction of the attempted map are supplied. Over a field, complexes split and this particular obstruction disappears; the ambient abelian category cannot be omitted.}
\GSfield{Fate}{This is established homological algebra, explained by localization at quasi-isomorphisms; see \citet[\S\S10.1--10.3]{weibel1994homological}. It continues earlier seminar discussions of the derived-category universal property, especially 1971-266.3.}
\GSfield{Present status}{As of November 2026, the statement with the displayed hypotheses is established; the notebook occurrence is expository.}
\GSfield{Assessment}{literal statement: correct; publication match: exact; verification: read the result; development: annotated; confidence: high; problem \hyperref[prob:GP-1980-23]{GP-1980-23}}
\end{GSitem}
\begin{GSitem}{1980-502.2}{F}{An unspecified derived equivalence between a sphere and its dual}
\GSfield{Kind; source}{claim; \GSpage{P000940}}
\GSfield{The notes}{The notes write sheaf categories \(C,C_1\) on \(S^n,S^{*n}\), an \(O(n)\)-action, \(C_1\not\simeq C\) because the action differs, and \(D(S)\sim D(S^*)\).}
\GSfield{Statement}{An exact equivalence cannot be specified from the record: the definition of the dual sphere, the action on each category, the coefficients and boundedness or constructibility restrictions, and the transforming kernel are absent. With the standard Euclidean duality and orthogonal action, the spheres are already equivariantly isomorphic, so their ordinary sheaf categories are equivalent; that literal reading does not produce the claimed distinction.}
\GSfield{Reconstruction}{No Radon-transform equivalence is substituted for the missing assertion. Such a reconstruction would require extra geometry not fixed by the manuscript.}
\GSfield{Fate}{Searches for S. I. Gelfand, dual spheres, equivariant sheaves, and derived/Radon transformations did not establish an exact publication match. \citet[Chapter 10]{weibel1994homological} supplies only the general derived-category language, not this claimed equivalence.}
\GSfield{Present status}{As of November 2026, the exact assertion and its fate remain unidentifiable from the surviving formulation; this is a source ambiguity, not a well-posed open problem.}
\GSfield{Assessment}{literal statement: not determinable; publication match: none found; verification: unresolved lead; development: annotated; confidence: low}
\end{GSitem}

\GStalk{504}{1980-11-24}{D. B. Fuchs}{P000940 (middle)}{Whitehead products and the Jacobi identity}
A short intervention gives the boundary-of-a-product construction of the Whitehead product and its graded Jacobi identity. The target homotopy degree is one too high in the manuscript itself; this is a faithful mathematical error, not an English translation error.

\begin{GSitem}{1980-504.1}{E}{Whitehead product from the boundary of a product of cubes}
\GSfield{Kind; source}{construction; \GSpage{P000940}}
\GSfield{The notes}{The notes write \(\pi_n(X)\times\pi_m(X)\to\pi_{n+m}(X)\) and decompose \(\partial(I^n\times I^m)\).}
\GSfield{Statement}{For a pointed space and \(n,m\ge2\), the Whitehead product is \([\ ,\ ]:\pi_n(X)\times\pi_m(X)\to\pi_{n+m-1}(X)\). Representatives on \(I^n/\partial I^n\) and \(I^m/\partial I^m\) define the map on the two pieces of the boundary of \(I^n\times I^m\), agreeing at the basepoint on their intersection.}
\GSfield{Reconstruction}{The target is corrected by one degree; the boundary has dimension \(n+m-1\). Low-dimensional nonabelian variants are not needed for the next graded statement.}
\GSfield{Fate}{This is Whitehead's classical construction; the attaching-map formulation is read in \citet[\S4.2, p.381]{hatcher2002algebraic}.}
\GSfield{Present status}{As of November 2026, the statement with the displayed hypotheses is established; the notebook occurrence is expository.}
\GSfield{Assessment}{literal statement: false as written; publication match: exact; verification: read the result; development: annotated; confidence: high}
\end{GSitem}
\begin{GSitem}{1980-504.2}{E}{Shifted graded Jacobi identity for Whitehead products}
\GSfield{Kind; source}{theorem; \GSpage{P000940}}
\GSfield{The notes}{The note compares the product to a Lie bracket and says that the Jacobi identity holds, with a “super” sign convention.}
\GSfield{Statement}{After the usual sign adjustment identifying Whitehead products with Samelson products in \(\Omega X\), the graded group \(L_d=\pi_{d+1}(X)\), \(d\ge1\), satisfies \([a,b]=-(-1)^{|a||b|}[b,a]\) and \((-1)^{|a||c|}[a,[b,c]]+(-1)^{|b||a|}[b,[c,a]]+(-1)^{|c||b|}[c,[a,b]]=0\). Over \(\mathbb Z\), torsion can prevent the alternating axiom \([a,a]=0\) in even degree; after tensoring with \(\mathbb Q\), this is a graded Lie algebra in the usual sense.}
\GSfield{Reconstruction}{The degree shift and the integral torsion caveat are supplied. The notebook's analogy should not be read as an unrestricted ordinary Lie algebra on unshifted homotopy groups.}
\GSfield{Fate}{The identity predates the talk by decades. \citet[p.180, publisher extract]{hilton1961jacobi} explicitly describes the already-known identity; the 1961 print date is distinct from the 2008 digital posting. Only the extract, not the full three-page proof, was read.}
\GSfield{Present status}{As of November 2026, the statement with the displayed hypotheses is established; the notebook occurrence is expository.}
\GSfield{Assessment}{literal statement: correct with missing hypotheses; publication match: exact; verification: publisher metadata checked; development: annotated; confidence: high}
\end{GSitem}

\GStalk{505}{1980-11-24}{V. I. Lomonosov}{P000940 (middle)}{A compressed perturbation statement for commuting operators}
The manuscript records a trace-class difference for two operators, a compact difference for another pair, two commutation assertions, and a claimed unitary equivalence of absolutely continuous restrictions. The primes and the operators used in the spectral subspaces do not yield a valid identifiable theorem as written. A standard scattering-theoretic repair can be proved, but is explicitly separated from any attribution to Lomonosov.

\begin{GSitem}{1980-505.1}{F}{An ambiguous four-operator perturbation assertion}
\GSfield{Kind; source}{claim; \GSpage{P000940}}
\GSfield{The notes}{The manuscript has \(A-A^{\prime}\) nuclear, \([B,A]=[B,A^{\prime}]=0\), \(B-B^{\prime}\) completely continuous, and concludes \(B|_{\mathcal H_{ac}(B)}\simeq B^{\prime}|_{\mathcal H_{ac}(B^{\prime})}\). It then proposes treating commuting families directly.}
\GSfield{Statement}{The exact intended theorem is undetermined. A valid nearby statement is: for bounded self-adjoint \(A,A^{\prime},B,B^{\prime}\), if \(A-A^{\prime}\) is trace class, \([A,B]=[A^{\prime},B^{\prime}]=0\), and \(B-B^{\prime}\) is compact, the complete wave operators \(W_\pm(A^{\prime},A)\) intertwine \(B|_{\mathcal H_{ac}(A)}\) and \(B^{\prime}|_{\mathcal H_{ac}(A^{\prime})}\). They are unitary between those two spectral subspaces.}
\GSfield{Reconstruction}{This repair changes both the second commutator and the operators defining the absolutely continuous subspaces; it is ours, not a corrected reading. Indeed \(e^{-itA}u\rightharpoonup0\) on \(\mathcal H_{ac}(A)\), so \((B^{\prime}-B)e^{-itA}u\to0\) in norm. Passing to the wave-operator limit proves the extra intertwining. The literal conditions, with \(A=A^{\prime}=0\), would impose no control beyond compact perturbation on \(B,B^{\prime}\), which is insufficient to preserve absolutely continuous spectrum.}
\GSfield{Fate}{\citet[Theorem 1, pp.260--261]{kato1957trace} proves the required complete wave operators, already in 1957. The compact-perturbation counterexample follows from the classical Weyl--von Neumann theorem, whose exact statement is read in \citet[Theorem 1.1]{andomatsuzawa2014weyl}. Lomonosov's \citet[pp.54--55]{lomonosov1980intertwiner} is a bibliographically verified contemporary lead, received 11 December 1978; its full text was not retrieved and it is not claimed as a match. The commuting-family remark has no recoverable additional hypotheses.}
\GSfield{Present status}{As of November 2026, the repaired statement is a standard consequence of trace-class scattering theory, but the exact Lomonosov assertion and its publication fate remain unresolved at the source-identification level. No open mathematical problem is inferred.}
\GSfield{Assessment}{literal statement: false as written; publication match: thematic only; verification: unresolved lead; development: annotated; confidence: low}
\end{GSitem}

\GStalk{503}{1980-11-24}{A. A. Beilinson}{P000940-P000943}{Derived direct images and holonomic modules}
The talk starts below the Fuchs and Lomonosov interventions on P000940 and continues through P000943. It treats composition of derived direct images, Spencer and relative de Rham complexes, Gauss--Manin connections, closed embeddings, point-fibre criteria for holonomicity, finite length, and simple modules as minimal extensions of irreducible connections. Several assertions need the distinction between a closed embedding and an arbitrary locally closed embedding.

\begin{GSitem}{1980-503.1}{E}{Pulling solutions back and integrating solutions along fibres}
\GSfield{Kind; source}{method; \GSpage{P000940-P000941}}
\GSfield{The notes}{The transfer-module formulas are repeated with the explanation that pullback carries solutions to solutions and direct image supplies equations for integrals, including improper integrals.}
\GSfield{Statement}{For a morphism of smooth characteristic-zero varieties, inverse-image modules encode the chain rule and derived direct-image modules encode relative de Rham classes. In an analytic realization, an integral of a differential form solving a \(\mathcal D_X\)-system gives a solution of the appropriate direct-image system when differentiation under the integral and integration by parts are valid, for example along a proper smoothly varying closed cycle or with controlled compact support.}
\GSfield{Reconstruction}{The analytic convergence, boundary, and differentiation conditions are supplied; the bare module formula does not justify arbitrary improper integration.}
\GSfield{Fate}{This is the established direct-image interpretation of differential equations and periods; see \citet[\S13, especially the de Rham and Gauss--Manin discussion]{gourevitch2025dmodules}. It repeats the construction of 1980-500.5 and develops its analytic meaning.}
\GSfield{Present status}{As of November 2026, the statement with the displayed hypotheses is established; the notebook occurrence is expository.}
\GSfield{Assessment}{literal statement: correct with missing hypotheses; publication match: exact; verification: read the result; development: annotated; confidence: high}
\end{GSitem}
\begin{GSitem}{1980-503.2}{E}{Derived direct images compose}
\GSfield{Kind; source}{theorem; \GSpage{P000941}}
\GSfield{The notes}{The underived operation is said not to satisfy composition; equation (7) states \(Rg_*Rf_*=R(g\circ f)_*\).}
\GSfield{Statement}{For morphisms \(X\xrightarrow fY\xrightarrow gZ\) of smooth varieties in characteristic zero, the derived right-module direct images satisfy a natural isomorphism \(g_+f_+\simeq(gf)_+\) on the bounded quasicoherent derived categories where these operations are defined. The ordinary tensor-and-sheaf-direct-image construction can lose higher cohomology and does not in general have this property.}
\GSfield{Reconstruction}{The notation \(f_+\) distinguishes this operation from the ordinary derived direct image of sheaves. The assertion concerns natural derived isomorphisms, not equality of chosen complexes.}
\GSfield{Fate}{The exact composition theorem is read in \citet[Theorem 13.22]{gourevitch2025dmodules}. Its mechanism is associativity of derived tensor products and the derived projection formula, established before this expository lecture.}
\GSfield{Present status}{As of November 2026, the statement with the displayed hypotheses is established; the notebook occurrence is expository.}
\GSfield{Assessment}{literal statement: correct; publication match: exact; verification: read the result; development: annotated; confidence: high}
\end{GSitem}
\begin{GSitem}{1980-503.3}{E}{Projection of a plane and the degree minus one direct image}
\GSfield{Kind; source}{construction; \GSpage{P000941}}
\GSfield{The notes}{For the projection of a plane to a line, the direct image of volume forms has zero degree-zero component and volume forms on the line in degree \(-1\).}
\GSfield{Statement}{For the algebraic projection \(f:\mathbb A^2\to\mathbb A^1\) over a characteristic-zero field, \(f_+\omega_{\mathbb A^2}\simeq\omega_{\mathbb A^1}[1]\). Its only cohomology is \(H^{-1}=\omega_{\mathbb A^1}\). Relative de Rham cohomology of the affine-line fibre is concentrated in unshifted degree zero.}
\GSfield{Reconstruction}{The notes' word “component” is read as a cohomology object, not a canonical zero/nonzero pattern for every representative complex. The affine algebraic setting fixes the cohomology convention.}
\GSfield{Fate}{This elementary instance of the relative de Rham description is verified from \citet[\S13]{gourevitch2025dmodules}.}
\GSfield{Present status}{As of November 2026, the statement with the displayed hypotheses is established; the notebook occurrence is expository.}
\GSfield{Assessment}{literal statement: correct with missing hypotheses; publication match: exact; verification: read the result; development: annotated; confidence: high}
\end{GSitem}
\begin{GSitem}{1980-503.4}{E}{The right-module Spencer complex and de Rham shift}
\GSfield{Kind; source}{construction; \GSpage{P000941}}
\GSfield{The notes}{The complex \(\cdots\to M\otimes\Lambda^2T_X\to M\otimes T_X\to M\) computes direct image to a point. For volume forms, it is called the de Rham complex shifted to the left by a symbol transcribed as \(M\).}
\GSfield{Statement}{For a right module \(M\) on smooth \(d\)-dimensional \(X\), the Spencer differential uses the right vector-field action and Lie brackets on \(M\otimes_{\mathcal O_X}\Lambda^jT_X\), placed in degree \(-j\). This complex computes \(M\otimes^L_{\mathcal D_X}\mathcal O_X\); applying \(R\Gamma\) computes \(p_+M\). For \(M=\omega_X\), contraction identifies it with \(\Omega_X^\bullet[d]\).}
\GSfield{Reconstruction}{The shift is \(d=\dim X\), as confirmed by the following displayed formula \(R^ip_*\omega_X=H^{d+i}(X)\); the ambiguous handwritten symbol cannot serve as a shift index.}
\GSfield{Fate}{The Spencer resolution and de Rham interpretation are standard; see \citet[\S13]{gourevitch2025dmodules}. This makes explicit the calculation already announced in 1980-500.9.}
\GSfield{Present status}{As of November 2026, the statement with the displayed hypotheses is established; the notebook occurrence is expository.}
\GSfield{Assessment}{literal statement: correct with missing hypotheses; publication match: exact; verification: read the result; development: annotated; confidence: high}
\end{GSitem}
\begin{GSitem}{1980-503.5}{E}{Relative de Rham cohomology and the Gauss–Manin action}
\GSfield{Kind; source}{construction; \GSpage{P000941-P000942}}
\GSfield{The notes}{For a smooth morphism, the relative tangent complex is pushed forward and receives a \(\mathcal D_Y\)-action called the Gauss--Manin connection.}
\GSfield{Statement}{If \(f:X\to Y\) is smooth between smooth varieties, the relative Spencer complex computes direct image after the canonical left/right conversion and the appropriate relative-dimension shift. For a smooth proper morphism and a finite-rank integrable connection on \(X\), its relative de Rham cohomology sheaves are locally free and carry integrable Gauss--Manin connections, compatible with the \(\mathcal D_Y\)-module direct image.}
\GSfield{Reconstruction}{The canonical-bundle conversion is necessary for a right-module formula: locally \(\omega_X\simeq\omega_{X/Y}\otimes f^*\omega_Y\). Properness is added only for the stated finite-rank bundle conclusion; without it the direct image is still a derived \(\mathcal D\)-module and need not be a vector bundle.}
\GSfield{Fate}{The exact smooth-proper connection statement and its qualifications are read in \citet[\S13, Gauss--Manin discussion]{gourevitch2025dmodules}.}
\GSfield{Present status}{As of November 2026, the statement with the displayed hypotheses is established; the notebook occurrence is expository.}
\GSfield{Assessment}{literal statement: correct with missing hypotheses; publication match: exact; verification: read the result; development: annotated; confidence: high}
\end{GSitem}
\begin{GSitem}{1980-503.6}{E}{Closed-embedding exactness and Kashiwara equivalence}
\GSfield{Kind; source}{theorem; \GSpage{P000942}}
\GSfield{The notes}{For an “embedding”, both direct-image steps are said to be exact; volume forms become delta forms in normal directions.}
\GSfield{Statement}{For a closed embedding \(i:X\hookrightarrow Y\) of smooth varieties, the transfer bimodule is flat for the required left \(\mathcal D_X\)-action, and \(i_+\) is exact on modules. It identifies \(\mathcal D_X\)-modules with \(\mathcal D_Y\)-modules supported on \(X\), with an inverse given by the appropriate degree of extraordinary restriction.}
\GSfield{Reconstruction}{“Closed” must be supplied; an arbitrary open or locally closed embedding does not make ordinary sheaf direct image exact in the same fashion.}
\GSfield{Fate}{This is Kashiwara's equivalence, read in \citet[Theorem 9.1]{gourevitch2025dmodules} and independently in \citet[\S8]{bernstein1982dmodules}. The later lecture notes are verification sources for an already established theorem.}
\GSfield{Present status}{As of November 2026, the statement with the displayed hypotheses is established; the notebook occurrence is expository.}
\GSfield{Assessment}{literal statement: correct with missing hypotheses; publication match: exact; verification: read the result; development: annotated; confidence: high}
\end{GSitem}
\begin{GSitem}{1980-503.7}{E}{Holonomicity tested by all derived point fibres}
\GSfield{Kind; source}{definition; \GSpage{P000942}}
\GSfield{The notes}{A finitely generated module is called holonomic if the cohomology of \(Li_x^*M\) is finite-dimensional for every point. The free module is explicitly excluded.}
\GSfield{Statement}{For a coherent \(\mathcal D_X\)-module on a smooth complex algebraic variety, holonomicity is equivalent to finite-dimensional cohomology of its derived inverse image at every closed point. Equivalently, its nonempty characteristic variety has dimension \(\dim X\). A nonzero locally free \(\mathcal D_X\)-module is not holonomic if \(\dim X>0\). Dimension shifts between \(i^*\) and \(i^!\) do not affect this finiteness test.}
\GSfield{Reconstruction}{Coherence, the complex algebraic category, and derived rather than ordinary fibres are essential. For example, the ordinary fibre of a delta module vanishes, while its derived fibre does not.}
\GSfield{Fate}{The exact converse as well as the forward implication are read in \citet[Proposition 16.6]{gourevitch2025dmodules}. The foundational dimension framework is in \citet[\S\S1--3]{bernstein1971modules}.}
\GSfield{Present status}{As of November 2026, the statement with the displayed hypotheses is established; the notebook occurrence is expository.}
\GSfield{Assessment}{literal statement: correct with missing hypotheses; publication match: exact; verification: read the result; development: annotated; confidence: high}
\end{GSitem}
\begin{GSitem}{1980-503.8}{E}{Generic connections and gluing holonomic modules along strata}
\GSfield{Kind; source}{claim; \GSpage{P000942}}
\GSfield{The notes}{The examples say that a holonomic module is a flat bundle away from exceptional points or curves, and ask whether all such modules are glued from transverse delta modules.}
\GSfield{Statement}{A holonomic module has a finite algebraic stratification of its support such that the appropriate restrictions to smooth strata are finite-rank integrable connections. Its simple constituents are minimal extensions of irreducible connections on dense smooth locally closed strata; general holonomic modules are finite successive extensions of these. This does not assert a direct-sum decomposition, regular singularities, or a classification by arbitrary gluing matrices.}
\GSfield{Reconstruction}{The circled question is resolved in the precise finite-extension sense supplied by the following finite-length and simple-object statements. “Away from curves” on a surface is not an assertion that every module has full support.}
\GSfield{Fate}{The reconstruction follows the structure results read in \citet[\S\S16--17, especially Lemma 16.18 and Theorem 17.2]{gourevitch2025dmodules}; the next notebook page itself supplies these two ingredients.}
\GSfield{Present status}{As of November 2026, the statement with the displayed hypotheses is established; the notebook occurrence is expository.}
\GSfield{Assessment}{literal statement: correct with missing hypotheses; publication match: partial; verification: read the result; development: annotated; confidence: high}
\end{GSitem}
\begin{GSitem}{1980-503.9}{E}{Finite length of holonomic modules}
\GSfield{Kind; source}{theorem; \GSpage{P000943}}
\GSfield{The notes}{Theorem (15)(1) states finite length and equality of the lengths of filtrations with irreducible factors.}
\GSfield{Statement}{The category of holonomic \(\mathcal D_X\)-modules on a smooth characteristic-zero algebraic variety is both noetherian and artinian. Every object admits a finite composition series; the multiset of simple factors, with multiplicities, is independent of the series.}
\GSfield{Reconstruction}{The Jordan--Hölder conclusion is stated in its stronger usual form, including multiplicities rather than only length.}
\GSfield{Fate}{Finite length is classical holonomic-module theory. The exact statement and its proof using duality are read in \citet[Lemma 16.18]{gourevitch2025dmodules}; compare the early finiteness framework in \citet[\S3]{bernstein1971modules}.}
\GSfield{Present status}{As of November 2026, the statement with the displayed hypotheses is established; the notebook occurrence is expository.}
\GSfield{Assessment}{literal statement: correct; publication match: exact; verification: read the result; development: annotated; confidence: high}
\end{GSitem}
\begin{GSitem}{1980-503.10}{E}{Simple holonomic modules as minimal extensions}
\GSfield{Kind; source}{theorem; \GSpage{P000943}}
\GSfield{The notes}{The notes classify irreducibles by an irreducible variety and irreducible flat bundle, identify pairs after shrinking, and claim a unique irreducible submodule of \(i_*\mathcal L_Y\).}
\GSfield{Statement}{Simple holonomic modules on smooth \(X\) correspond to pairs consisting of an irreducible closed support \(Z\subset X\) and an irreducible integrable connection on a nonempty smooth open subset \(U\subset Z\), modulo isomorphism after further shrinking. The corresponding module is the intermediate extension to \(X\); equivalently, it is the unique simple submodule of \(H^0j_+\mathcal L\) for the locally closed embedding \(j:U\hookrightarrow X\).}
\GSfield{Reconstruction}{The connection lives on a smooth dense open subset, not necessarily all of a singular support. Equivalence is generated by restriction to a common dense open, rather than the asymmetric single inclusion written in the notes.}
\GSfield{Fate}{The exact minimal-extension classification is read in \citet[Theorem 17.2]{gourevitch2025dmodules}. It is standard holonomic theory and does not require regular singularities unless one also invokes a constructible-sheaf Riemann--Hilbert equivalence.}
\GSfield{Present status}{As of November 2026, the statement with the displayed hypotheses is established; the notebook occurrence is expository.}
\GSfield{Assessment}{literal statement: correct with missing hypotheses; publication match: exact; verification: read the result; development: annotated; confidence: high}
\end{GSitem}
\begin{GSitem}{1980-503.11}{E}{Smoothness alone does not make a direct image simple}
\GSfield{Kind; source}{claim; \GSpage{P000943}}
\GSfield{The notes}{The final line says that if \(Y\) is smooth, the direct image is already irreducible.}
\GSfield{Statement}{This is valid for a closed embedding of smooth varieties and a simple source module, by Kashiwara equivalence. It is false for a general locally closed embedding: for \(j:\mathbb G_m\hookrightarrow\mathbb A^1\), the simple connection \(\mathcal O_{\mathbb G_m}\) has \(j_+\mathcal O_{\mathbb G_m}=\mathcal O_{\mathbb A^1}(*0)\), which contains the proper nonzero submodule \(\mathcal O_{\mathbb A^1}\). Its quotient is supported at the origin.}
\GSfield{Reconstruction}{The counterexample distinguishes smooth source from closed support and explains why intermediate extension is needed.}
\GSfield{Fate}{The exact closed-embedding and intermediate-extension theorems are read in \citet[Theorems 9.1 and 17.2]{gourevitch2025dmodules}; the localization counterexample is a direct calculation in the first Weyl algebra.}
\GSfield{Present status}{As of November 2026, the statement with the displayed hypotheses is established; the notebook occurrence is expository.}
\GSfield{Assessment}{literal statement: false as written; publication match: exact; verification: read the result; development: annotated; confidence: high}
\end{GSitem}

\GStalk{506}{1980-12-01}{S. P. Novikov}{P000944-P000948}{Magnetic translations, Pauli ground states, and Bloch bundles}
The actual heading is on P000944; the earlier pages assigned to this record by the supplied catalog are the end of Beilinson's lectures. Novikov discusses a magnetic Schr\"odinger operator with periodic curvature and potential, the flux obstruction to ordinary Bloch theory, the explicit zero modes of the two-dimensional Pauli operator, and the topology and perturbation of their bundle over the reciprocal torus. The zero-mode results were already published in September 1980; the perturbation and topological discussion matches a paper received on 2 December 1980, the day after this seminar (\citealp[Secs.~II--IV]{dubrovinnovikov1980ground}; \citealp[Lemmas~1--2, Theorem~2]{novikov1981magnetic}).

The notebook's physical line bundle and its bundle of Bloch zero modes must be distinguished: the former has degree \(N\), whereas the latter has rank \(N\) and first Chern number \(1\) in Novikov's orientation. Gelfand's concluding suggestion of a complex or parabolic regularization in the irrational-flux case is a program, not a specified spectral theorem. The final reading assignment on \(G\)-structures, \(\mathcal D\)-structures and the Spencer resolution is recorded as an assignment, not converted into an additional conjecture.

\begin{GSitem}{1980-506.1}{E}{Periodic magnetic Schrödinger and Pauli operators}
\GSfield{Kind; source}{definition; \GSpage{P000944}}
\GSfield{The notes}{Equations (1)--(7) give \(\widehat H=\frac12\sum_\alpha(i\partial_\alpha-eA_\alpha)^2+cV\), \(H_{\alpha\beta}=\partial_\beta A_\alpha-\partial_\alpha A_\beta\), and a spin term \(e\sigma_jH^j\). The curvature and electric potential are periodic; the vector potential need not be. The text first calls \(V\) the electric field but then explicitly distinguishes field from potential.}
\GSfield{Statement}{Let \(\Lambda\subset\mathbb R^2\) be a lattice, let \(B,V\) be smooth real \(\Lambda\)-periodic functions, and choose a smooth real potential \(A\) with \(\partial_xA_y-\partial_yA_x=B\). Put \(\Pi_j=-i\partial_j-eA_j\). The magnetic Schr\"odinger operator is the self-adjoint operator associated with the semibounded quadratic form \(\sum\|\Pi_jf\|^2+\int V|f|^2\). In a consistent Pauli normalization, \(P=\Pi_1^2+\Pi_2^2-e\sigma_3B\) splits into two scalar spin sectors. The notebook uses the opposite curvature sign and suppresses mass and spin-normalization conventions.}
\GSfield{Reconstruction}{The smoothness, real coefficients, form realization, dimension restriction for the displayed Pauli operator, and a single consistent normalization are supplied here. They are not a claim that the notebook fixes all physical units.}
\GSfield{Fate}{The model and its two spin sectors appear in \citet[Sec.~I, equations (1)--(8)]{dubrovinnovikov1980ground}, received 15 April 1980 and published before this talk.}
\GSfield{Present status}{As of November 2026, the precisely qualified statement above is established; the historical shorthand does not add a further open assertion.}
\GSfield{Assessment}{literal statement: correct with missing hypotheses; publication match: exact; verification: read the result; development: annotated; confidence: high}
\end{GSitem}
\begin{GSitem}{1980-506.2}{F}{The aside on a constant electric field}
\GSfield{Kind; source}{remark; \GSpage{P000944}}
\GSfield{The notes}{The parenthesis following (4) says that a constant electric field gives a potential of the form \(\vec E\cdot\vec x\), followed by \ru{не изучена даже 1-мерн. задача}.}
\GSfield{Statement}{A constant nonzero electric field requires an affine potential, which is not periodic on a full lattice. The phrase about an unstudied one-dimensional problem does not specify the periodic background, spectral property, or limiting regime, so no unique open problem can be extracted.}
\GSfield{Fate}{The observation about nonperiodicity is elementary. Searches of the Novikov bibliography and the cited 1980 papers identify no exact one-dimensional question corresponding to this aside; a general reference to Stark operators would be only thematic.}
\GSfield{Present status}{As of November 2026, the intended historical question remains unidentified; this is not a claim that one-dimensional Stark spectral theory is undeveloped.}
\GSfield{Assessment}{literal statement: not determinable; publication match: none found; verification: unresolved lead; development: annotated; confidence: low}
\end{GSitem}
\begin{GSitem}{1980-506.3}{E}{Integral magnetic flux as the degree of a physical line bundle}
\GSfield{Kind; source}{construction; \GSpage{P000944}}
\GSfield{The notes}{Equations (8)--(10) impose \(e\Phi_{\alpha\beta}/(2\pi)\in\mathbb Z\) and identify these integers with \(\langle c_1(\xi),Z_{\alpha\beta}\rangle\). The operator on Euclidean space covers an operator in a bundle over a torus.}
\GSfield{Statement}{A smooth closed real two-form \(B\) on \(T^n=\mathbb R^n/\Lambda\) is the curvature, divided by \(-ie\), of a unitary connection on a complex line bundle precisely when \([eB/(2\pi)]\) is integral. On \(T^2\), its degree is \(N=e\int_{T^2}B/(2\pi)\). Pullback to \(\mathbb R^n\) trivializes the bundle; lattice transitions yield the magnetic boundary conditions. Changing the sign of the curvature convention changes the sign of this degree.}
\GSfield{Reconstruction}{Closedness, a unitary connection and the sign convention are explicit here. This is the bundle on physical space, not the reciprocal-torus eigenbundle in item 13.}
\GSfield{Fate}{The integrality interpretation is already used in \citet[Sec.~I, equation (3)]{dubrovinnovikov1980ground}; the talk presents it as background.}
\GSfield{Present status}{As of November 2026, the precisely qualified statement above is established; the historical shorthand does not add a further open assertion.}
\GSfield{Assessment}{literal statement: correct with missing hypotheses; publication match: exact; verification: read the result; development: annotated; confidence: high}
\end{GSitem}
\begin{GSitem}{1980-506.4}{E}{Reduction of rational flux to integral flux}
\GSfield{Kind; source}{method; \GSpage{P000944-P000945}}
\GSfield{The notes}{The last line of P000944 says \ru{Рац. свод. к нему укрупнением решётки}, after mentioning the irrational case. The next page imposes integral flux to obtain commuting translations.}
\GSfield{Statement}{If \(e\Phi/(2\pi)=p/q\) with coprime integers \(p,q\), replacing one period by \(q\) times that period gives cell flux \(p\). Replacing both periods by \(q\) gives flux \(pq\). A finite-index enlargement cannot make an irrational flux integral: multiplying an irrational number by a nonzero integer remains irrational.}
\GSfield{Reconstruction}{The antecedent must be the integral case, not the irrational case. The arithmetic argument is an explanatory reconstruction, not a new result.}
\GSfield{Fate}{The intended rational-to-integral reduction is explicit in \citet[Sec.~I and Sec.~IV]{dubrovinnovikov1980ground}. The immediately preceding pronoun in the notes is misleading.}
\GSfield{Present status}{As of November 2026, the precisely qualified statement above is established; the historical shorthand does not add a further open assertion.}
\GSfield{Assessment}{literal statement: false as written; publication match: exact; verification: read the result; development: annotated; confidence: high}
\end{GSitem}
\begin{GSitem}{1980-506.5}{E}{Flux quantization depends on the charge being modeled}
\GSfield{Kind; source}{remark; \GSpage{P000945}}
\GSfield{The notes}{Two situations are contrasted: a prescribed crystal lattice in an external field and a superconducting lattice whose field is said to be automatically integral. The notebook assigns cell areas \(10^{-10}\,\mathrm{cm}^2\) to the former and \(10^{-16}\,\mathrm{cm}^2\) to the latter.}
\GSfield{Statement}{Flux quantization refers to the charge in the effective order parameter. If one superconducting flux quantum is \(2\pi/(2e_0)\), it is an integer quantum for charge \(2e_0\), but a half quantum for an electron of charge \(e_0\). Integrality therefore cannot be asserted independently of the chosen charge and model. The two numerical cell-area examples in the notes are reversed relative to the cited paper.}
\GSfield{Reconstruction}{This is a correction of a faithfully transcribed physical example, not a correction inferred from handwriting alone.}
\GSfield{Fate}{\citet[Sec.~I, Remark]{dubrovinnovikov1980ground} explicitly makes the charge distinction and gives the crystal area as about \(10^{-16}\,\mathrm{cm}^2\).}
\GSfield{Present status}{As of November 2026, the precisely qualified statement above is established; the historical shorthand does not add a further open assertion.}
\GSfield{Assessment}{literal statement: false as written; publication match: exact; verification: read the result; development: annotated; confidence: high}
\end{GSitem}
\begin{GSitem}{1980-506.6}{E}{Magnetic translations form a projective lattice representation}
\GSfield{Kind; source}{theorem; \GSpage{P000945}}
\GSfield{The notes}{If \(A_\alpha(x+T_j)=A_\alpha(x)+\partial_\alpha f_j\), the notes set \(\widehat T_j^*\psi=\psi(x+T_j)e^{-ief_j(x)}\) and state \(\widehat T_j^*\widehat T_k^*=\widehat T_k^*\widehat T_j^*e^{-ie\Phi_{kj}}\), crediting Zak and Brown.}
\GSfield{Statement}{For a smooth real connection potential on \(\mathbb R^n\) with lattice-periodic curvature, the difference between a translate of \(A\) and \(A\) is exact. The corresponding gauge-corrected translations are unitary, commute with a magnetic Hamiltonian whose scalar potential is periodic, and commute with one another up to the exponential of the flux through the spanning parallelogram. They commute exactly on a pair of lattice generators when that normalized flux is integral.}
\GSfield{Reconstruction}{The phase depends on orientation and on whether active or passive translations are used. A missing hat on the last translation is merely typographical.}
\GSfield{Fate}{Brown and Zak published magnetic-translation theory in 1964 (\citealp{brown1964bloch,zak1964magnetic}); the periodic-field formula is read in \citet[Sec.~I, equations (4)--(6)]{dubrovinnovikov1980ground}.}
\GSfield{Present status}{As of November 2026, the precisely qualified statement above is established; the historical shorthand does not add a further open assertion.}
\GSfield{Assessment}{literal statement: correct with missing hypotheses; publication match: exact; verification: read the result; development: annotated; confidence: high}
\end{GSitem}
\begin{GSitem}{1980-506.7}{E}{Irrational flux prevents finite-dimensional magnetic translation fibers}
\GSfield{Kind; source}{claim; \GSpage{P000945; P000948}}
\GSfield{The notes}{The notes say irrational flux gives groups \ru{не типа 1}, as subgroups of the Heisenberg group; P000948 again writes the noncommuting two-translation relation and mentions factors and ergodicity.}
\GSfield{Statement}{For \(\theta\notin\mathbb Q\), the universal rotation algebra generated by unitaries \(U,V\) with \(UV=e^{2\pi i\theta}VU\) has no nonzero finite-dimensional unital representation: determinants in dimension \(d\) would require \(e^{2\pi id\theta}=1\). Its operator-algebraic representation theory replaces a finite-dimensional magnetic Bloch fiber. A particular representation or von Neumann completion must be specified before assigning a factor type.}
\GSfield{Reconstruction}{The determinant proof and the precise rotation-algebra formulation are supplied here. The notebook does not assert a classification of all Heisenberg subgroups.}
\GSfield{Fate}{The irrational-representation obstruction is already described in \citet[Sec.~I]{dubrovinnovikov1980ground}; later observable-algebra treatments are developed in \citet[Secs.~3.1, 3.5--3.6]{bellissard1994quantumhall}.}
\GSfield{Present status}{As of November 2026, the precisely qualified statement above is established; the historical shorthand does not add a further open assertion.}
\GSfield{Assessment}{literal statement: correct with missing hypotheses; publication match: partial; verification: read the result; development: annotated; confidence: high}
\end{GSitem}
\begin{GSitem}{1980-506.8}{F}{Describe the spectrum: a program without a specified invariant}
\GSfield{Kind; source}{question; \GSpage{P000945}}
\GSfield{The notes}{An underlined \ru{Задача: описать спектр} precedes the magnetic Bloch construction.}
\GSfield{Statement}{For the smooth periodic-field operators of item 1, the spectrum depends on the entire functions \(B,V\), on the flux and on the spin sector. The notes do not specify whether the desired description concerns explicit bands, spectral type, gap labels, a density of states, or continuity in the field. Thus no single complete yes-or-no problem is determined by this line.}
\GSfield{Fate}{The talk supplies the special Pauli zero-mode solution. Later noncommutative spectral geometry gives substantial invariant information (\citealp[Sec.~3.4, Theorem~1]{bellissard1994quantumhall}), but is not an exact solution to an unspecified request for the whole spectrum.}
\GSfield{Present status}{As of November 2026, this is retained as an identifiable research program with an unspecified target, not promoted to a newly discovered open conjecture.}
\GSfield{Assessment}{literal statement: not determinable; publication match: thematic only; verification: read the result; development: annotated; confidence: medium}
\end{GSitem}
\begin{GSitem}{1980-506.9}{E}{Magnetic Bloch decomposition for commuting lattice translations}
\GSfield{Kind; source}{method; \GSpage{P000945}}
\GSfield{The notes}{Equations (16)--(18) seek simultaneous eigenfunctions of \(\widehat H\) and the translations, with \(\widehat T_j^*\psi_p=e^{ip_jT_j}\psi_p\); \(L^2\) is a direct integral of torus problems and \(\varepsilon_j(p_1,p_2)\) are dispersion laws.}
\GSfield{Statement}{For integral flux and smooth periodic coefficients, the commuting magnetic lattice action decomposes the Hilbert space as a direct integral over its character torus. The fiber operator acts on sections of a line bundle over the physical torus twisted by a flat unitary character, and has compact resolvent. Its ordered eigenvalues, counted with multiplicity, give the fiberwise bands. Individual Bloch waves need not belong to \(L^2(\mathbb R^2)\).}
\GSfield{Reconstruction}{Integrality and generalized eigenfunctions are essential qualifications to the notebook's simultaneous-eigenfunction wording.}
\GSfield{Fate}{This is the framework used before the seminar in \citet[Sec.~I, equation (7), and Sec.~III]{dubrovinnovikov1980ground}.}
\GSfield{Present status}{As of November 2026, the precisely qualified statement above is established; the historical shorthand does not add a further open assertion.}
\GSfield{Assessment}{literal statement: correct with missing hypotheses; publication match: exact; verification: read the result; development: annotated; confidence: high}
\end{GSitem}
\begin{GSitem}{1980-506.10}{E}{Explicit Pauli zero modes in a periodic magnetic field}
\GSfield{Kind; source}{theorem; \GSpage{P000946}}
\GSfield{The notes}{For \(V=0\), spin \(1/2\) and \(N=e\Phi_{12}/(2\pi)\), equation (20) gives \(\Psi_A=e^{-e\varphi}\prod_{j=1}^N\sigma(z-a_j)e^{-az}\), with \(\Delta\varphi=H\) and a logarithmic sigma integral in (22). The margin cites Dubrovin--Novikov, JETP, September 1980.}
\GSfield{Statement}{Let a smooth real magnetic field on a two-torus have positive integral normalized flux \(N>0\), and take the Pauli operator with no electric potential. In the spin sector whose first-order factor is a Dolbeault operator, the zero modes are holomorphic sections of a degree-\(N\) line bundle, expressed on the universal cover as a gauge factor times an entire theta function with \(N\) zeros per cell. Sigma products provide such functions; their exponential factor must satisfy the magnetic boundary multipliers. For negative flux exchange the spin sector; at zero flux the statement requires a separate analysis.}
\GSfield{Reconstruction}{Positivity of N, the spin sector, and a consistent curvature normalization have been supplied. The sign of a can be changed by renaming that parameter; the curvature inconsistency between (2) and (27) cannot.}
\GSfield{Fate}{This is \citet[Sec.~III, Theorem~(a)]{dubrovinnovikov1980ground}, not a post-seminar discovery. The JETP account appeared in September 1980 (\citealp[Sec.~2]{dubrovinnovikov1980electron}).}
\GSfield{Present status}{As of November 2026, the precisely qualified statement above is established; the historical shorthand does not add a further open assertion.}
\GSfield{Assessment}{literal statement: correct with missing hypotheses; publication match: exact; verification: read the result; development: annotated; confidence: high}
\end{GSitem}
\begin{GSitem}{1980-506.11}{E}{Divisors, sigma multipliers, and real quasimomenta}
\GSfield{Kind; source}{construction; \GSpage{P000946}}
\GSfield{The notes}{Equations (21)--(24) relate \(p_1+ip_2\) to \(\sum a_j\); \(a=a(a_1,\ldots,a_N)\) is chosen to select real quasimomenta. The parameters are symmetric and defined modulo the lattice.}
\GSfield{Statement}{For an elliptic curve \(E=\mathbb C/\Lambda\), an effective degree-\(N\) divisor determines a degree-\(N\) line bundle. Fixing the Abel sum modulo \(\Lambda\) fixes its isomorphism class. The quasi-periodicity of the sigma product fixes its two multipliers; an exponential correction normalizes their absolute values to one after the gauge factor is included. The scalar multiplier and unordered zeros parametrize nonzero sections with the appropriate lattice identifications, not a globally free Cartesian coordinate system.}
\GSfield{Reconstruction}{No correction is made solely because the two publications use different signs for the complex quasimomentum; orientation and exponential normalization must first be fixed.}
\GSfield{Fate}{The exact multiplier equations are (14)--(16) of \citet[Sec.~III]{dubrovinnovikov1980ground}; alternative signs occur in \citet[equation (6)]{novikov1981magnetic}.}
\GSfield{Present status}{As of November 2026, the precisely qualified statement above is established; the historical shorthand does not add a further open assertion.}
\GSfield{Assessment}{literal statement: correct with missing hypotheses; publication match: partial; verification: read the result; development: annotated; confidence: high}
\end{GSitem}
\begin{GSitem}{1980-506.12}{E}{The N-dimensional zero-mode fiber}
\GSfield{Kind; source}{theorem; \GSpage{P000946-P000947}}
\GSfield{The notes}{At fixed \(p\), the notes identify the fiber with \(\mathcal L(\mathcal D)\), state \(\dim\mathcal L(\mathcal D)=N\), and invoke Atiyah--Singer with zero cokernel. All functions in one fiber are obtained from one by meromorphic elliptic multipliers.}
\GSfield{Statement}{If \(L\) is a holomorphic line bundle of degree \(N>0\) on a smooth compact elliptic curve and \(F_p\) is any flat line bundle, then \(H^1(E,L\otimes F_p)=0\) and \(h^0(E,L\otimes F_p)=N\). Indeed Serre duality identifies \(H^1\) with the dual of sections of a negative-degree line bundle, hence zero; Riemann--Roch gives the dimension. A quotient of two nonzero sections is a meromorphic elliptic function with poles canceled by the zeros of the denominator section.}
\GSfield{Reconstruction}{The elementary curve-theoretic proof clarifies the zero-cokernel hypothesis; an index identity alone never proves the kernel dimension without it.}
\GSfield{Fate}{\citet[Sec.~III, Theorem~(b) and its proof]{dubrovinnovikov1980ground} applies this index calculation to the Pauli kernel.}
\GSfield{Present status}{As of November 2026, the precisely qualified statement above is established; the historical shorthand does not add a further open assertion.}
\GSfield{Assessment}{literal statement: correct with missing hypotheses; publication match: exact; verification: read the result; development: annotated; confidence: high}
\end{GSitem}
\begin{GSitem}{1980-506.13}{A}{Rank and Chern number of the Bloch bundle}
\GSfield{Kind; source}{claim; \GSpage{P000946-P000947}}
\GSfield{The notes}{Equation (25) states that the zero modes form an \(N\)-dimensional holomorphic vector bundle over the reciprocal torus with \(c_1=N\), identifying this with the number of magnetic-flux quanta.}
\GSfield{Statement}{The family \(\mathcal E_p=H^0(E,L\otimes F_p)\) has rank \(N\). Its first Chern number has absolute value \(1\), and equals \(1\) in Novikov's reciprocal-torus convention; it is not \(N\) for general positive \(N\). In the standard algebraic Poincar\'e-bundle orientation, the Fourier--Mukai transform of a degree-\(N\) line bundle has rank \(N\) and degree \(-1\); reversing the base orientation gives Novikov's sign.}
\GSfield{Reconstruction}{The scan clearly reads \(c_1=N\); this is a faithful notebook error, not an OCR correction. The Fourier--Mukai terminology is a modern reconstruction; the cited source proves the topological normalization directly.}
\GSfield{Fate}{\citet[Sec.~II, Lemma~1]{novikov1981magnetic} explicitly gives \(c_1=1\). This paper was received 2 December 1980 and published in 1981.}
\GSfield{Present status}{As of November 2026, the distinction is settled: physical line bundle degree N; reciprocal-torus zero-mode bundle rank N and Chern number of absolute value 1.}
\GSfield{Assessment}{literal statement: false as written; publication match: exact; verification: read the result; development: annotated; confidence: high}
\end{GSitem}
\begin{GSitem}{1980-506.14}{E}{The same fiber multiplicity at every constant-field Landau level}
\GSfield{Kind; source}{theorem; \GSpage{P000947}}
\GSfield{The notes}{The notes extend the \(N\)-dimensionality at fixed quasimomentum to all energy levels when the magnetic field is uniform. A question about the origin of the degeneracy is crossed out.}
\GSfield{Statement}{For a constant nonzero magnetic field with positive integral flux \(N\) per cell, each Landau eigenspace in a magnetic Bloch fiber in a fixed scalar spin sector has dimension \(N\). Creation operators identify the higher-level eigenbundles with the lowest-level bundle. This extension concerns constant field; an arbitrary periodic magnetic perturbation does not preserve all higher Landau levels. Combining both Pauli spin sectors can give an additional spin multiplicity at positive energies.}
\GSfield{Reconstruction}{The crossed-out question is not counted as an independent unresolved problem.}
\GSfield{Fate}{The constant-field comparison is made in \citet[Sec.~3]{dubrovinnovikov1980electron}; \citet[Sec.~II, Remark after Lemma~1]{novikov1981magnetic} also records the eigenbundle identification.}
\GSfield{Present status}{As of November 2026, the precisely qualified statement above is established; the historical shorthand does not add a further open assertion.}
\GSfield{Assessment}{literal statement: correct with missing hypotheses; publication match: exact; verification: read the result; development: annotated; confidence: high}
\end{GSitem}
\begin{GSitem}{1980-506.15}{E}{Compressing a weak electric perturbation to the zero-mode bundle}
\GSfield{Kind; source}{method; \GSpage{P000947}}
\GSfield{The notes}{For \(\widehat H=\widehat H_0+ev\), equation (30) restricts the Hermitian form \(\int_{K_{12}}\overline{\Psi_A}v\Psi_A\) to the unperturbed kernel at fixed \((p_1,p_2)\).}
\GSfield{Statement}{Let \(v\) be a smooth real periodic function and assume the zero-mode cluster of \(H_0(p)\) is isolated by a uniform gap. If \(P_0(p)\) is its orthogonal projection, the first-order coefficients of its eigenvalues under \(H_0+\varepsilon v\) are the eigenvalues of \(P_0(p)vP_0(p)\) on \(\ker H_0(p)\). This is an \(N\times N\) Hermitian endomorphism of the zero-mode bundle; it is not the exact full perturbed Hamiltonian.}
\GSfield{Reconstruction}{The small perturbation parameter and isolated-cluster qualification are made explicit.}
\GSfield{Fate}{The compression is already in \citet[Sec.~IV]{dubrovinnovikov1980ground}; its topological use is sharpened in \citet[Sec.~III]{novikov1981magnetic}.}
\GSfield{Present status}{As of November 2026, the precisely qualified statement above is established; the historical shorthand does not add a further open assertion.}
\GSfield{Assessment}{literal statement: correct with missing hypotheses; publication match: exact; verification: read the result; development: annotated; confidence: high}
\end{GSitem}
\begin{GSitem}{1980-506.16}{A}{Codimension three and generic splitting into line bundles}
\GSfield{Kind; source}{theorem; \GSpage{P000947}}
\GSfield{The notes}{Theorem 2 states that a generic two-parameter family of complex Hermitian forms has no eigenvalue coincidences because the coincidence condition has codimension three; in three parameters coincidences are isolated. A corollary splits the Bloch bundle into lines.}
\GSfield{Statement}{In the real vector space of Hermitian \(N\times N\) matrices, the stratum with exactly one double eigenvalue and all other eigenvalues simple has codimension three. Globally, consider a smooth section of the bundle of Hermitian endomorphisms of the rank-\(N\) Bloch bundle, transverse to its fiberwise multiplicity strata. On a two-dimensional compact base it avoids all multiple eigenvalues and splits the bundle into smooth eigenlines. A transverse three-parameter family meets the simple-double stratum in isolated points. Applying this to scalar periodic potentials additionally requires transversality of the induced family of compressions. No global trivialization of the nontrivial Bloch bundle is assumed.}
\GSfield{Reconstruction}{The notebook's three equations concern coincidence at fixed quasimomentum. Allowing the two quasimomenta to vary gives codimension one in the space of potentials; the 1981 paper distinguishes these counts.}
\GSfield{Fate}{The potential-dependent transversality and resulting splitting are stated in \citet[Sec.~III, Lemma~2 and Theorem~2]{novikov1981magnetic}, received the day after the talk.}
\GSfield{Present status}{As of November 2026, the corrected genericity statement is established. The global Chern classes of the eigenlines need not vanish and sum to the total Bloch Chern class.}
\GSfield{Assessment}{literal statement: correct with missing hypotheses; publication match: exact; verification: read the result; development: annotated; confidence: high}
\end{GSitem}
\begin{GSitem}{1980-506.17}{F}{The aside about one-dimensional second-order operators}
\GSfield{Kind; source}{remark; \GSpage{P000947}}
\GSfield{The notes}{An uncancelled parenthesis says that in the one-dimensional case general position never occurs, because the operator is second order.}
\GSfield{Statement}{For unrestricted smooth one-parameter Hermitian families, second-order coincidence is not forced: for example the constant diagonal family with distinct diagonal entries has simple spectrum. The notebook must therefore refer to a more restricted family of scalar differential operators, boundary conditions, or a different notion of general position; those data are not specified in this parenthesis.}
\GSfield{Reconstruction}{The diagonal-matrix observation only shows why the preceding finite-dimensional genericity formulation cannot be transferred literally to this aside. It does not reconstruct a missing Sturm--Liouville theorem.}
\GSfield{Fate}{The codimension statements in \citet[Sec.~III, Lemma~2]{novikov1981magnetic} concern the two-dimensional periodic-field construction. They do not identify the separate one-dimensional assertion or its intended quantifiers.}
\GSfield{Present status}{As of November 2026, the intended one-dimensional comparison remains unidentified; no open spectral problem is inferred from it.}
\GSfield{Assessment}{literal statement: not determinable; publication match: none found; verification: unresolved lead; development: annotated; confidence: low}
\end{GSitem}
\begin{GSitem}{1980-506.18}{F}{Complex or parabolic regularization in the irrational case}
\GSfield{Kind; source}{method; \GSpage{P000947-P000948}; I. M. Gelfand}
\GSfield{The notes}{Gelfand says \ru{Интересен иррац. случай!} and suggests entering the complex domain, \ru{внутри круга}, in a non-self-adjoint setting, or a parabolic regularization. The next page recalls the irrational magnetic-translation relation and von Neumann factors.}
\GSfield{Statement}{No regularized operator, analytic parameter, topology of convergence, or spectral invariant is specified. The strongest recoverable content is a proposal to investigate irrational-flux spectral theory through a non-self-adjoint or heat-type deformation. The relation \(UV=e^{2\pi i\theta}VU\) survives, but by itself does not define the proposed regularization.}
\GSfield{Fate}{Noncommutative observable algebras and trace/Chern invariants later gave rigorous tools (\citealp[Secs.~3--4]{bellissard1994quantumhall}). No exact implementation of this notebook prescription was identified in the checked Novikov and quantum-Hall sources.}
\GSfield{Present status}{As of November 2026, the precise suggested regularization remains unidentifiable. No open existence theorem is manufactured from this incomplete method proposal.}
\GSfield{Assessment}{literal statement: not determinable; publication match: thematic only; verification: unresolved lead; development: annotated; confidence: low}
\end{GSitem}

\GStalk{507}{}{-}{Catalogue range P000948-P000950}{YEAR 1981}
This catalogue record is a calendar-year divider before the 9 February MacPherson talk. No independent mathematical assertion belongs to it, so it has no occurrence-index items. The actual P000949 scan begins with the date and MacPherson heading; the catalogue divider is not represented as a separate manuscript theorem or talk. Adjacent Novikov and MacPherson material is assigned to their own manuscript headings.


\GStalk{508}{1981-02-09}{Robert D. MacPherson}{P000949-P000953}{Grassmannian simplices, the dilogarithm, and convexity}
The talk, described as joint work with I. M. Gelfand, develops positive diagonal-torus orbit closures in real Grassmannians, hypersimplex combinatorics, cross-ratios, Pontryagin-form integration, and a five-term dilogarithm identity. It then reports Atiyah's extensions to Kähler and symplectic geometry. The unrestricted harmonic-form assertion, the use of an orbit in place of an orbit closure, and the final Hamiltonian-orbit diagram each need distinct qualifications. The construction continues the 1979 occurrences 1979-486.12--15 and 1979-487.1--2.

\begin{GSitem}{1980-508.1}{E}{Grassmannian dimension notation and coordinate inclusions}
\GSfield{Kind; source}{definition; \GSpage{P000949}}
\GSfield{The notes}{The talk uses \(G_n^m\) for \(n\)-planes in \(\mathbb R^{n+m}\) and arranges coordinate inclusions in a two-dimensional array.}
\GSfield{Statement}{Write \(G_n^m=\operatorname{Gr}_n(\mathbb R^{n+m})\), of dimension \(nm\). The horizontal inclusion adds an unused coordinate; the vertical inclusion sends \(V\) to \(V\oplus\mathbb R e_{n+m+1}\). Equivalently \(G_n^m\) parametrizes projective \((n-1)\)-planes in \(\mathbb{RP}^{n+m-1}\).}
\GSfield{Reconstruction}{The third entry in the second row must be \(G_2^3\), not the handwritten \(G_1^3\). These coordinate choices specify which smaller Grassmannians contain faces.}
\GSfield{Fate}{The underlying coordinate Grassmannian geometry is classical; compare the explicit Grassmannian setting in \citet[pp.13--14]{atiyah1982convexity}. The table is checked directly from its dimensions.}
\GSfield{Present status}{As of November 2026, the statement with the displayed hypotheses is established; the notebook occurrence is expository.}
\GSfield{Assessment}{literal statement: false as written; publication match: exact; verification: read the result; development: annotated; confidence: high}
\end{GSitem}
\begin{GSitem}{1980-508.2}{F}{The inadequately specified geodesic-simplex construction}
\GSfield{Kind; source}{remark; \GSpage{P000949}}
\GSfield{The notes}{The opening asks for polyhedra whose faces lie in smaller Grassmannians; it says that filling a triangle with geodesic segments fails when the rank is at least two.}
\GSfield{Statement}{The record does not specify a metric, a canonical filling rule, or an exact compatibility condition that is alleged to fail. Thus it records motivation for using torus orbit closures, rather than a quantified nonexistence theorem for geodesic simplices.}
\GSfield{Reconstruction}{No universal claim that higher-rank Grassmannians lack totally geodesic simplices is inferred. Totally geodesic subspaces and small geodesic triangles certainly exist.}
\GSfield{Fate}{The subsequent construction has a precise publication match in \citet{gelfandmacpherson1982geometry}, but no exact theorem corresponding to this preliminary objection was identified. The previous problem GP-1979-04 concerns a Grassmannian counterpart of a projective-space linear model; title-level resemblance does not establish that its exact question recurs here.}
\GSfield{Present status}{As of November 2026, this motivating remark does not determine an independently assessable mathematical problem.}
\GSfield{Assessment}{literal statement: not determinable; publication match: thematic only; verification: indirect; development: annotated; confidence: low}
\end{GSitem}
\begin{GSitem}{1980-508.3}{E}{The hypersimplex}
\GSfield{Kind; source}{definition; \GSpage{P000949}}
\GSfield{The notes}{\(\Delta_n^m\) is defined by \(0\le x_i\le1\) and \(\sum x_i=n\).}
\GSfield{Statement}{For positive integers \(n,m\), put \(N=n+m\) and \(\Delta(n,N)=\{x\in[0,1]^N:\sum_i x_i=n\}\). This \((N-1)\)-dimensional polytope is the convex hull of the characteristic vectors of all \(n\)-element subsets of \(\{1,\ldots,N\}\).}
\GSfield{Reconstruction}{The vertex description and dimension are explicit consequences of the inequalities; \(\Delta_n^m=\Delta(n,n+m)\).}
\GSfield{Fate}{The combinatorial object is classical and already occurs in 1979-486.2. Its Grassmannian realization is part of \citet{gelfandmacpherson1982geometry}; see also the read orbit-polytope theorem \citet[Theorem 2′, p.13]{atiyah1982convexity}.}
\GSfield{Present status}{As of November 2026, the statement with the displayed hypotheses is established; the notebook occurrence is expository.}
\GSfield{Assessment}{literal statement: correct; publication match: exact; verification: read the result; development: annotated; confidence: high}
\end{GSitem}
\begin{GSitem}{1980-508.4}{A}{Positive torus orbit closures and the genericity condition}
\GSfield{Kind; source}{construction; \GSpage{P000949}}
\GSfield{The notes}{A projective frame determines the identity component \(H^0\) of its stabilizer. A Grassmannian simplex is the closure of a generic orbit; generic means transverse to every coordinate subspace, not merely maximal orbit dimension.}
\GSfield{Statement}{In \(\operatorname{Gr}_n(\mathbb R^N)\), let \(H=(\mathbb R_{>0})^N/\mathbb R_{>0}\) act diagonally. If every Plücker coordinate \(p_I(V)\) is nonzero, define the Grassmannian simplex to be \(\overline{H\cdot V}\). The open orbit has dimension \(N-1\). Nonvanishing of all Plücker coordinates is stronger than maximal orbit dimension and is equivalent to the stated coordinate general-position condition.}
\GSfield{Reconstruction}{The central scalar quotient, all-minor condition, and orbit closure are explicit. Non-generic supports generally yield smaller matroid polytopes, not the full hypersimplex.}
\GSfield{Fate}{This is the announced Gelfand--MacPherson construction, explicitly identified by \citet[p.13]{atiyah1982convexity}. The work was published by Gelfand and MacPherson in June 1982; this 9 February 1981 seminar is a verified earlier oral occurrence. No earlier public preprint date was independently established.}
\GSfield{Present status}{As of November 2026, this construction is established. The distinction between full Plücker support and maximal orbit dimension remains essential.}
\GSfield{Assessment}{literal statement: correct; publication match: exact; verification: read the result; development: annotated; confidence: high}
\end{GSitem}
\begin{GSitem}{1980-508.5}{A}{The orbit closure is a hypersimplex with its face stratification}
\GSfield{Kind; source}{theorem; \GSpage{P000950; P000952-P000953}}
\GSfield{The notes}{P000950 calls the two objects diffeomorphic; P000952 identifies face interiors with positive-torus orbits, and P000953 corrects the global conclusion to homeomorphism with smoothness on individual orbits.}
\GSfield{Statement}{For the generic \(V\) above, the Grassmannian moment map restricts to a homeomorphism \(\overline{H\cdot V}\to\Delta(n,N)\). One convenient expression is \(\mu(V)=\sum_I |p_I(V)|^2\mathbf1_I/\sum_I|p_I(V)|^2\). Each positive orbit maps diffeomorphically to a relative face interior. A global diffeomorphism of embedded orbit closures is not asserted; boundary strata can be singular.}
\GSfield{Reconstruction}{The corrected conclusion follows the notebook's own later qualification. The formula fixes the normalization and makes the polytope correspondence concrete.}
\GSfield{Fate}{\citet[Theorems 2 and 2′, pp.3 and 13]{atiyah1982convexity} proves the exact orbit-closure statement and explicitly identifies the real-Grassmannian case with Gelfand--MacPherson. The work was published by Gelfand and MacPherson in June 1982; this 9 February 1981 seminar is a verified earlier oral occurrence. No earlier public preprint date was independently established.}
\GSfield{Present status}{As of November 2026, the homeomorphism and orbit-by-face smoothness are established; the initial unrestricted diffeomorphism wording requires the recorded correction.}
\GSfield{Assessment}{literal statement: false as written; publication match: exact; verification: read the result; development: annotated; confidence: high}
\end{GSitem}
\begin{GSitem}{1980-508.6}{A}{The projective triangle and its dual Grassmannian triangle}
\GSfield{Kind; source}{construction; \GSpage{P000950}}
\GSfield{The notes}{The examples \(G_1^2\) and \(G_2^1\) give four triangular chambers in the projective plane and a dual triangle of lines, with sides in \(G_1^1\).}
\GSfield{Statement}{The complement of the three coordinate lines in \(\mathbb{RP}^2\) has four sign chambers, each a positive diagonal orbit; its closure is a triangle. Dually, two-planes in \(\mathbb R^3\) are kernels of nonzero covectors, and each covector-sign chamber similarly gives a triangle. Its edges consist of planes containing a fixed coordinate axis, and its vertices are coordinate two-planes.}
\GSfield{Reconstruction}{The dual description depends on the selected sign chamber. The informal phrase about lines avoiding a hatched region is interpreted through these explicit coordinate conditions.}
\GSfield{Fate}{These elementary models are exact special cases of the orbit-closure construction; \citet[Theorem 2′]{atiyah1982convexity}. They are already implicit in the earlier configuration descriptions 1979-487.1--2.}
\GSfield{Present status}{As of November 2026, the coordinate and dual models are established; the drawings specify a chosen chamber, not the whole dual projective plane.}
\GSfield{Assessment}{literal statement: correct with missing hypotheses; publication match: exact; verification: read the result; development: annotated; confidence: high}
\end{GSitem}
\begin{GSitem}{1980-508.7}{A}{The octahedron and the cross-ratio in Gr(2,4)}
\GSfield{Kind; source}{construction; \GSpage{P000951}}
\GSfield{The notes}{\(G_2^2\) is four-dimensional, its positive orbit is three-dimensional, and the cross-ratio of four planes through a projective line is constant along the orbit. The six vertices are labelled by pairs.}
\GSfield{Statement}{For a generic projective line \(\ell\subset\mathbb{RP}^3\), the planes spanned by \(\ell\) and the four coordinate points determine four distinct points in the projective line of planes through \(\ell\). Their cross-ratio is invariant under diagonal transformations. The positive-orbit closure is \(\Delta(2,4)\), an octahedron with six pair-labelled vertices and eight triangular facets, four of each coordinate-facet type.}
\GSfield{Reconstruction}{Cross-ratio conventions and sign chambers must be fixed to identify a particular connected orbit in a level set; the notes state only containment in a level set.}
\GSfield{Fate}{This is the same four-ray/cross-ratio geometry discussed in 1979-486.12--13. The hypersimplex identification follows the exact theorem read in \citet[Theorem 2′]{atiyah1982convexity}; its use in the dilogarithm construction was published in \citet{gelfandmacpherson1982geometry}.}
\GSfield{Present status}{As of November 2026, the octahedral model and cross-ratio invariant are established.}
\GSfield{Assessment}{literal statement: correct; publication match: exact; verification: read the result; development: annotated; confidence: high}
\end{GSitem}
\begin{GSitem}{1980-508.8}{A}{Dilogarithmic volume between two Grassmannian simplices}
\GSfield{Kind; source}{claim; \GSpage{P000951}}
\GSfield{The notes}{The recorded volume between \(K=c_1\) and \(K=c_2\), with \(0<c_1<c_2<1\), is \(\Phi(c_1)-\Phi(c_2)\), where \(\Phi(x)=\operatorname{Li}_2(x)-\operatorname{Li}_2(1-x)\).}
\GSfield{Statement}{The formula can describe an oriented, normalized characteristic-form integral, but it cannot literally be an unsigned positive volume with the stated order: \(\Phi^{\prime}(x)=-\log(1-x)/x-\log x/(1-x)>0\). The natural increasing primitive satisfies \(\Phi=2L-\pi^2/6\), where \(L(x)=\operatorname{Li}_2(x)+\tfrac12\log x\log(1-x)\) is the Rogers dilogarithm. An unsigned volume requires the opposite difference after fixing the form normalization.}
\GSfield{Reconstruction}{The sign test is direct. The exact normalization of the invariant four-form, orientation of the slab, and choice of cross-ratio are not recoverable from this line alone, so no numerical constant is silently supplied.}
\GSfield{Fate}{The geometric dilogarithm construction is the subject of \citet{gelfandmacpherson1982geometry}; publisher metadata is verified, but its full text was not retrieved. The precise Rogers identity and normalization are read in \citet[Chapter II, \S1.A]{zagier2007dilogarithm}. The work was published by Gelfand and MacPherson in June 1982; this 9 February 1981 seminar is a verified earlier oral occurrence. No earlier public preprint date was independently established.}
\GSfield{Present status}{As of November 2026, the characteristic-form/dilogarithm construction is established, but the notebook's unsigned-volume wording has a sign defect and omits normalization.}
\GSfield{Assessment}{literal statement: false as written; publication match: partial; verification: indirect; development: annotated; confidence: medium}
\end{GSitem}
\begin{GSitem}{1980-508.9}{E}{The two families of hypersimplex boundary faces}
\GSfield{Kind; source}{construction; \GSpage{P000952}}
\GSfield{The notes}{Equation (12) starts an alternating boundary sum over \(x_i=0\) faces and then breaks off in the second family; its second lower limit is written \(-1\).}
\GSfield{Statement}{For \(1<n<N-1\), the facets of \(\Delta(n,N)\) are \(x_i=0\), isomorphic to \(\Delta(n,N-1)\), and \(x_i=1\), isomorphic to \(\Delta(n-1,N-1)\). Giving the ambient affine hyperplane and each facet compatible orientations yields its ordinary signed cellular boundary. At \(n=1\) or \(n=N-1\), the degenerate family must be omitted as a family of facets.}
\GSfield{Reconstruction}{The missing \(x_i=1\) condition and index range \(i=1,\ldots,N\) are reconstructed from the defining inequalities. The incomplete manuscript does not determine a fully checked orientation-sign formula.}
\GSfield{Fate}{The face geometry is elementary and underlies \citet{gelfandmacpherson1982geometry}; the face/orbit compatibility is independently read in \citet[Theorem 2]{atiyah1982convexity}. No unverified completion of the alternating signs is represented as a transcription.}
\GSfield{Present status}{As of November 2026, the statement with the displayed hypotheses is established; the notebook occurrence is expository.}
\GSfield{Assessment}{literal statement: not determinable; publication match: partial; verification: read the result; development: annotated; confidence: medium}
\end{GSitem}
\begin{GSitem}{1980-508.10}{A}{Vanishing of harmonic forms on real Grassmannian simplices}
\GSfield{Kind; source}{theorem; \GSpage{P000952}}
\GSfield{The notes}{Equation (13) claims that every harmonic form on \(G_n^m\), for a coordinate-compatible metric, vanishes on every Grassmannian simplex.}
\GSfield{Statement}{For the unoriented real Grassmannian \(G_n^m\) with \(nm\) even and its orthogonally invariant metric, every harmonic differential form of positive degree restricts to zero on each positive diagonal orbit. The same assertion holds for the corresponding invariant positive-degree forms. Constants are excluded; the unrestricted assertion fails in general for \(nm\) odd and for oriented Grassmannians.}
\GSfield{Reconstruction}{Positive degree, the unoriented space, and the parity hypothesis are supplied from the primary theorem. Restriction is understood on the smooth orbit strata, avoiding a fictitious smooth structure at singular boundary points.}
\GSfield{Fate}{\citet[pp.13--14, discussion before and after Theorem 4]{atiyah1982convexity} states this Gelfand--MacPherson result, gives a new proof, and explicitly records the odd-dimensional and oriented failures. The work was published by Gelfand and MacPherson in June 1982; this 9 February 1981 seminar is a verified earlier oral occurrence. No earlier public preprint date was independently established.}
\GSfield{Present status}{As of November 2026, the qualified vanishing theorem is established; the literal all-degrees, all-Grassmannians formulation is false.}
\GSfield{Assessment}{literal statement: false as written; publication match: exact; verification: read the result; development: annotated; confidence: high}
\end{GSitem}
\begin{GSitem}{1980-508.11}{E}{Five configurations and the differential five-term identity}
\GSfield{Kind; source}{theorem; \GSpage{P000952}}
\GSfield{The notes}{Five points on \(\mathbb P^1\), modulo projectivities, form a two-dimensional space. Omitting each point gives a cross-ratio \(k_i\), and equation (14) alternates the differentials \(d\Phi(k_i)\).}
\GSfield{Statement}{The moduli space of five ordered distinct points on \(\mathbb P^1\) modulo \(\operatorname{PGL}_2\) has dimension two. On a real chamber with consistent cross-ratio conventions, the five-term relation is equivalent to the Rogers identity, for \(0<x,y<1\): \[L(x)+L(y)=L(xy)+L\!\left(\frac{x(1-y)}{1-xy}\right)+L\!\left(\frac{y(1-x)}{1-xy}\right).\] Since \(\Phi=2L-\pi^2/6\), differentiating gives a five-term zero relation; its displayed signs depend on the ordered cross-ratios.}
\GSfield{Reconstruction}{A concrete branch-free chamber and the exact five arguments replace the notebook's unspecified cross-ratio convention. The constant in the corresponding undifferentiated \(\Phi\)-relation is \(\pi^2/6\).}
\GSfield{Fate}{The analytic identity is classical, not new in 1981; its exact Rogers form is read in \citet[Chapter II, \S1.A]{zagier2007dilogarithm}. The geometric derivation is the Gelfand--MacPherson contribution discussed separately below.}
\GSfield{Present status}{As of November 2026, the statement with the displayed hypotheses is established; the notebook occurrence is expository.}
\GSfield{Assessment}{literal statement: correct with missing hypotheses; publication match: exact; verification: read the result; development: annotated; confidence: high}
\end{GSitem}
\begin{GSitem}{1980-508.12}{A}{Pontryagin-form fibre integration and Stokes’ formula}
\GSfield{Kind; source}{method; \GSpage{P000952-P000953}}
\GSfield{The notes}{The notes use the first Pontryagin form on \(G_3^2\), say its harmonic degree-four space is one-dimensional, and apply fibre integration with boundary terms over the space of five-point configurations.}
\GSfield{Statement}{The invariant degree-four cohomology of \(\operatorname{Gr}_3(\mathbb R^5)\) is one-dimensional, generated by the first Pontryagin class of the tautological bundle (equivalently the negative of the quotient bundle class). On the suitable oriented family of compactified positive orbits, fibre integration of its closed invariant representative satisfies the Stokes formula with oriented facet terms. The vanishing theorem and the two facet families turn those terms into the differential five-term identity.}
\GSfield{Reconstruction}{The tautological rank-three and complementary rank-two bundles have opposite first Pontryagin classes because their sum is trivial. “Exactly one form” means a one-dimensional vector space. Boundary orientations and the family compactification must be fixed; the abbreviated diagram does not supply them.}
\GSfield{Fate}{This is a recurrence of the characteristic-form construction in 1979-486.14--15 and 1979-487.2, published in \citet{gelfandmacpherson1982geometry}. The exact vanishing input is independently read in \citet[Theorem 4 and its Grassmannian application]{atiyah1982convexity}. Full verification of every 1982 normalization was unavailable.}
\GSfield{Present status}{As of November 2026, the geometric mechanism is established; this occurrence does not fix a unique sign and normalization for the primitive.}
\GSfield{Assessment}{literal statement: correct with missing hypotheses; publication match: partial; verification: indirect; development: annotated; confidence: medium}
\end{GSitem}
\begin{GSitem}{1980-508.13}{B}{Positive orbits in compact Kähler manifolds and their polytope closures}
\GSfield{Kind; source}{theorem; \GSpage{P000953}}
\GSfield{The notes}{The generalization says that algebraic torus actions with a fixed point have positive-real orbits homeomorphic to polyhedra, with only stratumwise differentiability.}
\GSfield{Statement}{Let a compact torus \(T\) act by Kähler automorphisms on a compact connected Kähler manifold \(M\), with a fixed point. Let \(T^{\mathbb C}=T\cdot H\), where \(\operatorname{Lie}H=J\operatorname{Lie}T\), and choose a moment map. For each \(y\), the moment map restricts to a homeomorphism \(\overline{H\cdot y}\to P_y\), a convex polytope; \(H\cdot y\) maps to its relative interior, and lower orbits to face interiors.}
\GSfield{Reconstruction}{Compactness, a Kähler-preserving action, and an orbit closure are supplied. A positive orbit itself is generally noncompact and cannot be homeomorphic to a compact polytope.}
\GSfield{Fate}{This is the exact general theorem of \citet[Theorems 2 and 2′]{atiyah1982convexity}, published in 1982 after receipt on 27 March 1981. The notebook's 9 February report is earlier evidence of oral circulation; receipt is not treated as a public preprint date.}
\GSfield{Present status}{As of November 2026, the orbit-closure theorem is established with the compact Kähler and Hamiltonian hypotheses.}
\GSfield{Assessment}{literal statement: false as written; publication match: exact; verification: read the result; development: annotated; confidence: high}
\end{GSitem}
\begin{GSitem}{1980-508.14}{B}{Kähler harmonic-form vanishing on the positive torus part}
\GSfield{Kind; source}{theorem; \GSpage{P000953}}
\GSfield{The notes}{A torus-invariant harmonic form, with respect to a Kähler metric, is said to vanish on every positive-real orbit.}
\GSfield{Statement}{For compact connected Kähler \(M\), a torus \(T\) preserving its Kähler structure and having a fixed point, every harmonic form of degree \(\ell>0\) restricts to zero on the orbits of the positive subgroup \(H\subset T^{\mathbb C}\). A torus orbit is null-homotopic through the fixed point; its invariant exact pullback vanishes, and Hodge-type decomposition transfers this conclusion from \(T\)-directions to \(J T\)-directions.}
\GSfield{Reconstruction}{The degree-zero exception and the compact connected hypotheses are explicit. A \(T\)-invariant Kähler metric ensures the form-invariance used in the proof.}
\GSfield{Fate}{The statement and proof are read in \citet[Theorem 4, pp.13--14]{atiyah1982convexity}. The paper appeared in 1982; its recorded receipt was 27 March 1981, after this seminar report.}
\GSfield{Present status}{As of November 2026, the positive-degree theorem is established; nonzero harmonic constants are the immediate exception to the literal wording.}
\GSfield{Assessment}{literal statement: correct with missing hypotheses; publication match: exact; verification: read the result; development: annotated; confidence: high}
\end{GSitem}
\begin{GSitem}{1980-508.15}{B}{Convexity for commuting Hamiltonians and the orbit-image conflation}
\GSfield{Kind; source}{theorem; \GSpage{P000953}}
\GSfield{The notes}{Theorem (15), attributed to Atiyah, assumes Poisson commutation and a common critical point, then depicts an orbit of the Hamiltonian fields mapping to a convex polyhedron.}
\GSfield{Statement}{Let \(M\) be compact connected symplectic and \(f_1,\ldots,f_r\) commuting Hamiltonians whose flows generate a compact torus action (Atiyah allows almost-periodic fields). The map \(f:M\to\mathbb R^r\) has connected nonempty fibres and convex-polytope image, the convex hull of its common critical values. Its restriction to an orbit of the Hamiltonian fields is constant, since \(H_{f_i}(f_j)=\{f_i,f_j\}=0\). Nontrivial orbit polytopes instead concern the positive complexified action in the Kähler version.}
\GSfield{Reconstruction}{The missing compactness and almost-periodicity hypotheses and the correct domain of the moment-image assertion are supplied. Mere commutation plus a common critical point is insufficient: on a symplectic two-torus, \((\cos x,\cos^2x)\) has these properties and a nonconvex parabolic image.}
\GSfield{Fate}{The full theorem is read in \citet[Theorem 1 and following remarks, pp.2--3]{atiyah1982convexity}; the orbit theorem is its distinct Theorem 2′. Publication was 1982, receipt 27 March 1981. The counterexample and constant-orbit observation are direct checks here.}
\GSfield{Present status}{As of November 2026, moment-map convexity is established. The notebook conflates the whole-manifold moment image, positive complexified orbits, and Hamiltonian orbits; the latter have point images.}
\GSfield{Assessment}{literal statement: false as written; publication match: exact; verification: read the result; development: annotated; confidence: high}
\end{GSitem}
\begin{GSitem}{1980-508.16}{F}{The unspecified relation with integration by quadratures}
\GSfield{Kind; source}{question; \GSpage{P000953}}
\GSfield{The notes}{The final line asks for the relation with integration by quadratures.}
\GSfield{Statement}{No precise implication, class of differential equations, or meaning of global quadratures is specified. The nearby theorem concerns global moment-map geometry; a regular Liouville integrable system would additionally require an appropriate number of independent commuting first integrals. These are related frameworks but are not interchangeable assertions.}
\GSfield{Reconstruction}{No new conjecture equating convexity and integrability is inferred. The final sentence is retained as a research-direction question, not assigned a fabricated definite problem.}
\GSfield{Fate}{\citet[Theorem 1 and introduction]{atiyah1982convexity} provides the exact geometric result being discussed, but does not identify this unspecified requested relation as a separate theorem. Searches found no exact formulation attributable to this sentence.}
\GSfield{Present status}{As of November 2026, no independently assessable open question can be recovered from the one-line formulation.}
\GSfield{Assessment}{literal statement: not determinable; publication match: thematic only; verification: unresolved lead; development: annotated; confidence: low}
\end{GSitem}

\GStalk{509}{1981-02-16}{A. G. Khovanskii}{P000954-P000959}{Newton polygons, genus, mixed volumes, and toric compactification}
Khovanskii explains how Newton polygons determine the genus and punctures of a nondegenerate curve, derives root counts through boundary valuations, introduces mixed area, and connects Pick's formula and Minkowski's inequality with toric intersection theory. The notes cite his 1977 and 1978 papers and Danilov's 1978 survey. P000960 starts the 23 February meeting, so it is not part of this lecture despite the catalogue endpoint.

\begin{GSitem}{1980-509.1}{E}{Genus of a hyperelliptic curve}
\GSfield{Kind; source}{theorem; \GSpage{P000955}}
\GSfield{The notes}{The motivating example is \(y^2=P_n(x)\), with \(P_n\) squarefree, and genus \(\lfloor(n-1)/2\rfloor\); degrees three and four give one handle and degrees five and six give two.}
\GSfield{Statement}{Over \(\mathbb C\), let \(P\) be a squarefree polynomial of degree \(n\geq1\). The smooth projective model of \(y^2=P(x)\) has genus \(\lfloor(n-1)/2\rfloor\). Indeed, its degree-two map to \(\mathbb P^1\) has \(n\) finite simple branch points and one further branch point at infinity exactly when \(n\) is odd; Riemann--Hurwitz gives the formula.}
\GSfield{Reconstruction}{The genus belongs to the smooth compact model. The branch count supplies an elementary proof and includes the point at infinity.}
\GSfield{Fate}{This is a classical hyperelliptic genus calculation. It also follows from the Newton-polygon genus theorem read in \citet[\S2]{khovanskii1978genus}; that publication predates this explanatory talk.}
\GSfield{Present status}{As of November 2026, the statement with the displayed hypotheses is established; the notebook occurrence is expository.}
\GSfield{Assessment}{literal statement: correct; publication match: exact; verification: read the result; development: annotated; confidence: high}
\end{GSitem}
\begin{GSitem}{1980-509.2}{E}{Degree genus formula and why the hyperelliptic plane model is special}
\GSfield{Kind; source}{theorem; \GSpage{P000955}}
\GSfield{The notes}{The notes give \((n-1)(n-2)/2\) for a general degree-\(n\) equation, then say that the hyperelliptic equation is “not an equation of degree \(n\)”. }
\GSfield{Statement}{A nonsingular plane projective curve of degree \(d\) has genus \((d-1)(d-2)/2\). For \(n\geq2\), \(y^2-P_n(x)\) does have total degree \(n\); its projective closure is a special degree-\(n\) model, generally singular at infinity, so the nonsingular plane-curve formula cannot simply be applied to that model.}
\GSfield{Reconstruction}{The intended distinction is generic versus special support, not total degree. Passing to normalization explains the discrepancy in genera.}
\GSfield{Fate}{This is the classical plane genus formula. Its realization as the count of interior lattice points of the degree-\(d\) triangle is covered by \citet[\S2]{khovanskii1978genus} and the speaker's later explicit comparison in \citet[slides 6--8]{khovanskii2025chern}.}
\GSfield{Present status}{As of November 2026, the statement with the displayed hypotheses is established; the notebook occurrence is expository.}
\GSfield{Assessment}{literal statement: false as written; publication match: exact; verification: read the result; development: annotated; confidence: high}
\end{GSitem}
\begin{GSitem}{1980-509.3}{E}{Newton polygons and Laurent-monomial translation}
\GSfield{Kind; source}{definition; \GSpage{P000955-P000956}}
\GSfield{The notes}{The Newton polygon is the convex hull of exponent pairs with nonzero coefficients. Multiplication by \(x^ay^b\) translates it; the curve is considered away from the coordinate cross.}
\GSfield{Statement}{For a nonzero Laurent polynomial \(P=\sum_{m\in A}c_mx^m\) on \(T=(\mathbb C^*)^n\), define \(\Delta(P)=\operatorname{conv}\{m:c_m\ne0\}\subset\mathbb R^n\). Since a Laurent monomial is a unit on \(T\), multiplying by it leaves the torus zero scheme unchanged and translates \(\Delta(P)\) by the exponent vector.}
\GSfield{Reconstruction}{The torus restriction is essential: multiplication by a polynomial monomial can add coordinate-axis components in affine space.}
\GSfield{Fate}{The definition and its character-lattice interpretation are read in \citet[\S\S1--2]{khovanskii1977toroidal}, published in 1977. These are already established tools in the 1981 lecture.}
\GSfield{Present status}{As of November 2026, the statement with the displayed hypotheses is established; the notebook occurrence is expository.}
\GSfield{Assessment}{literal statement: correct; publication match: exact; verification: read the result; development: annotated; confidence: high}
\end{GSitem}
\begin{GSitem}{1980-509.4}{E}{The Newton-polygon genus theorem}
\GSfield{Kind; source}{theorem; \GSpage{P000955}}
\GSfield{The notes}{The answer to the first motivating question is the number of interior lattice points of the Newton polygon, for coefficients in general position; the notebook attributes the earlier result to Baker and its rediscovery to Khovanskii.}
\GSfield{Statement}{Let \(P\in\mathbb C[x^{\pm1},y^{\pm1}]\) have a two-dimensional Newton polygon \(\Delta\), and assume its torus curve is irreducible and Newton nondegenerate: for every face \(F\), including \(\Delta\), the equations \(P_F=x\partial_xP_F=y\partial_yP_F=0\) have no common solution in \((\mathbb C^*)^2\). The genus of its smooth projective model is \(\#(\operatorname{int}\Delta\cap\mathbb Z^2)\). Generic coefficients with prescribed full Newton polygon satisfy these hypotheses.}
\GSfield{Reconstruction}{Face nondegeneracy and the smooth projective model make “general position” precise. The historical teacher--student anecdote is not needed for the theorem and was not independently verified.}
\GSfield{Fate}{The exact hypersurface genus result is read in \citet[\S2]{khovanskii1978genus}; its companion paper \citet[\S2]{khovanskii1977toroidal} defines nondegeneracy and the good toric compactification. Thus the talk is exposition of published 1977--1978 work, not a new unresolved genus question.}
\GSfield{Present status}{As of November 2026, the statement with the displayed hypotheses is established; the notebook occurrence is expository.}
\GSfield{Assessment}{literal statement: correct with missing hypotheses; publication match: exact; verification: read the result; development: annotated; confidence: high; problem \hyperref[prob:GP-1980-24]{GP-1980-24}}
\end{GSitem}
\begin{GSitem}{1980-509.5}{E}{Boundary valuations and unimodular monomial changes}
\GSfield{Kind; source}{construction; \GSpage{P000955-P000956}}
\GSfield{The notes}{At an omitted point, the orders of \(x,y\) give an integral vector. A monomial substitution with determinant one changes this vector by the corresponding integral matrix; primitive vectors can be made coordinate vectors.}
\GSfield{Statement}{On the smooth compactification of an irreducible torus curve, each boundary point \(p\) has \(v_p=(\operatorname{ord}_p x,\operatorname{ord}_p y)\in\mathbb Z^2\setminus\{0\}\). Under \(x=u^{a_{11}}v^{a_{12}},\ y=u^{a_{21}}v^{a_{22}}\), the order column is multiplied by \(A\). For \(A\in\operatorname{GL}_2(\mathbb Z)\), the map is an isomorphism of tori. Any primitive lattice vector is carried to a coordinate vector by such a matrix. In the nondegenerate compactification the boundary order vectors are primitive inward edge normals.}
\GSfield{Reconstruction}{The determinant may be either \(1\) or \(-1\); the written determinant-one subcase is sufficient. Primitivity of boundary vectors uses transverse intersection with the toric boundary and cannot be assumed for arbitrary singular curves.}
\GSfield{Fate}{The lattice duality and admissible smooth toric charts are read in \citet[\S1]{khovanskii1977toroidal}; its \S2 gives the transverse compactification for nondegenerate systems.}
\GSfield{Present status}{As of November 2026, the statement with the displayed hypotheses is established; the notebook occurrence is expository.}
\GSfield{Assessment}{literal statement: correct with missing hypotheses; publication match: exact; verification: read the result; development: annotated; confidence: high}
\end{GSitem}
\begin{GSitem}{1980-509.6}{E}{Edge lengths count punctures}
\GSfield{Kind; source}{theorem; \GSpage{P000956-P000957}}
\GSfield{The notes}{After translating an edge to a coordinate axis, its polynomial has as many nonzero roots as the length of the edge, with distinct roots in general position. A later sentence instead says “number of lattice points on the segment”; the total puncture count is the number of boundary lattice points.}
\GSfield{Statement}{Under the nondegeneracy hypotheses of 1980-509.4, the compactified curve meets the toric divisor associated to an edge \(E\) in exactly its lattice length \(\ell(E)=\#(E\cap\mathbb Z^2)-1\) distinct transverse points. Consequently the number of punctures of the curve in \((\mathbb C^*)^2\) is \(\sum_E\ell(E)=\#(\partial\Delta\cap\mathbb Z^2)\). Vertex initial polynomials are monomials and contribute no roots.}
\GSfield{Reconstruction}{The local edge count and the global boundary-point count differ by one per closed edge because vertices are counted twice in the latter edgewise listing. Simple roots supply the transverse local parameter.}
\GSfield{Fate}{This is the elementary boundary calculation in the nondegenerate toric compactification of \citet[\S\S1--2]{khovanskii1977toroidal}; the speaker states the exact curve count again in \citet[slide 6]{khovanskii2025chern}.}
\GSfield{Present status}{As of November 2026, the statement with the displayed hypotheses is established; the notebook occurrence is expository.}
\GSfield{Assessment}{literal statement: false as written; publication match: exact; verification: read the result; development: annotated; confidence: high}
\end{GSitem}
\begin{GSitem}{1980-509.7}{E}{Divisor degree and the signed boundary root count}
\GSfield{Kind; source}{method; \GSpage{P000957}}
\GSfield{The notes}{The notes restrict \(Q\) to \(P=0\), say its divisor is zero, substitute \(x\sim t^n,y\sim t^m\), and sum minimum scalar products over punctures to count the torus roots.}
\GSfield{Statement}{Let \(C\) be the smooth projective model of \(P=0\) and let \(Q|_C\) be a nonzero meromorphic function. Then \(\deg\operatorname{div}(Q|_C)=0\), not generally \(\operatorname{div}(Q|_C)=0\). If the initial terms of \(Q\) do not cancel at a boundary point \(p\), its order is \(\min_{a\in\Delta(Q)}\langle a,v_p\rangle\). As \(Q\) has no poles in the torus, the number of its torus zeros, with multiplicity, is \[-\sum_{p\in C\setminus T}\operatorname{ord}_p(Q)=-\sum_{p\in C\setminus T}\min_{a\in\Delta(Q)}\langle a,v_p\rangle.\]}
\GSfield{Reconstruction}{The missing minus sign and the difference between a divisor and its degree are explicit corrections. Local expansions have nonvanishing unit factors and the right-hand side of equation (11) requires the full sum; the minimum formula needs noncancellation of its leading coefficient.}
\GSfield{Fate}{The valuation and intersection mechanism is already present in \citet[\S\S1--2]{khovanskii1977toroidal} and the root-count application in \citet[\S1]{khovanskii1978genus}. The sign is checked directly from the degree-zero identity.}
\GSfield{Present status}{As of November 2026, the statement with the displayed hypotheses is established; the notebook occurrence is expository.}
\GSfield{Assessment}{literal statement: false as written; publication match: exact; verification: read the result; development: annotated; confidence: high}
\end{GSitem}
\begin{GSitem}{1980-509.8}{E}{Mixed area and the two incompatible normalizations}
\GSfield{Kind; source}{definition; \GSpage{P000957-P000958}}
\GSfield{The notes}{The planar coefficient \(S(K,L)\) is defined by \(\operatorname{Area}(aK+bL)=a^2\operatorname{Area}(K)+b^2\operatorname{Area}(L)+abS(K,L)\). The following general definition instead sets \(V(K,\ldots,K)=\operatorname{Vol}(K)\).}
\GSfield{Statement}{The standard mixed volume is the symmetric Minkowski-multilinear functional with \(V(K,\ldots,K)=\operatorname{Vol}_n(K)\); equivalently, the coefficient of \(t_1\cdots t_n\) in \(\operatorname{Vol}(\sum t_iK_i)\) is \(n!V(K_1,\ldots,K_n)\). Thus \(S(K,L)=2V(K,L)=\operatorname{Area}(K+L)-\operatorname{Area}(K)-\operatorname{Area}(L)\), and \(S(K,K)=2\operatorname{Area}(K)\).}
\GSfield{Reconstruction}{Symmetry and multilinearity in every slot are included in the characterization. The notebook shifts from the coefficient normalization \(S\) to the diagonal-volume normalization \(V\), so they cannot be substituted into root counts without the factorial.}
\GSfield{Fate}{The exact polarization normalization is read in \citet[\S1, mixed-volume definition]{khovanskii1978genus} and \citet[slide 12]{khovanskii2025chern}. This is classical mixed-volume theory; the later lecture is an explicit normalization check.}
\GSfield{Present status}{As of November 2026, the statement with the displayed hypotheses is established; the notebook occurrence is expository.}
\GSfield{Assessment}{literal statement: correct with missing hypotheses; publication match: exact; verification: read the result; development: annotated; confidence: high}
\end{GSitem}
\begin{GSitem}{1980-509.9}{E}{Bernshtein’s torus root count, with multiplicities and nondegeneracy}
\GSfield{Kind; source}{theorem; \GSpage{P000957-P000958}}
\GSfield{The notes}{The second motivating question is answered by mixed area, and theorem (14), attributed to D. Bernshtein, states the generic system root count as the mixed volume of its Newton polytopes.}
\GSfield{Statement}{For finite supports \(A_i\subset\mathbb Z^n\), put \(\Delta_i=\operatorname{conv}(A_i)\). A generic system of Laurent polynomials \(f_i\) with these Newton polytopes has \(n!V(\Delta_1,\ldots,\Delta_n)\) isolated roots in \((\mathbb C^*)^n\), counted with multiplicity. A sufficient precise hypothesis for equality is that for every nonzero weight \(w\in\mathbb R^n\), the simultaneous initial system \(f_{1,w}=\cdots=f_{n,w}=0\) has no torus root and that the original system is zero-dimensional. In two variables the count is \(S(\Delta(P),\Delta(Q))\).}
\GSfield{Reconstruction}{The torus, multiplicities, finite intersection, and nonzero-weight conditions are explicit. With the normalization immediately preceding theorem (14), its right-hand side lacks \(n!\), or lacks \(2\) in the stated planar case.}
\GSfield{Fate}{\citet{bernshtein1975roots} was published in 1975 (received 7 April 1975; receipt is not a public-circulation date). The exact normalized theorem is read in \citet[\S1, number-of-solutions consequence]{khovanskii1978genus}. It is therefore known before this talk. Earlier 1974 Newton-polyhedron occurrences concern local singularity invariants and are not automatically the same problem.}
\GSfield{Present status}{As of November 2026, the statement with the displayed hypotheses is established; the notebook occurrence is expository.}
\GSfield{Assessment}{literal statement: false as written; publication match: exact; verification: read the result; development: annotated; confidence: high; problem \hyperref[prob:GP-1980-25]{GP-1980-25}}
\end{GSitem}
\begin{GSitem}{1980-509.10}{E}{Equal polygons and the Bézout specialization}
\GSfield{Kind; source}{theorem; \GSpage{P000957}}
\GSfield{The notes}{When both polygons coincide, a support-distance sum equals twice their area; for the degree triangle this gives \(n^2\), identified with Bézout.}
\GSfield{Statement}{Two generic torus polynomials with the same Newton polygon \(\Delta\) have \(2\operatorname{Area}(\Delta)\) common roots. For generic polynomials of total degrees \(d,e\), the polygons are \(d\Sigma,e\Sigma\), where \(\Sigma=\operatorname{conv}(0,e_1,e_2)\); hence the mixed-area root count is \(de\). Genericity ensures that the projective Bézout intersections do not escape to coordinate divisors.}
\GSfield{Reconstruction}{The displayed \(n^2\) is the equal-degree subcase. The support-distance formula is invariant under translation because the length-weighted outward normals sum to zero.}
\GSfield{Fate}{This is the classical Bézout count recovered from the already published Bernshtein theorem; see \citet[\S1]{khovanskii1978genus}.}
\GSfield{Present status}{As of November 2026, the statement with the displayed hypotheses is established; the notebook occurrence is expository.}
\GSfield{Assessment}{literal statement: correct; publication match: exact; verification: read the result; development: annotated; confidence: high}
\end{GSitem}
\begin{GSitem}{1980-509.11}{E}{Euler characteristic of a nondegenerate torus hypersurface}
\GSfield{Kind; source}{theorem; \GSpage{P000958}}
\GSfield{The notes}{The word “curve” is crossed out in favor of “hypersurface”, while equation (15) retains the constant negative sign \(-n!\operatorname{Vol}(\Delta)\).}
\GSfield{Statement}{Let \(f\) have full-dimensional Newton polytope \(\Delta\subset\mathbb R^n\), and assume every face truncation, including the full polynomial, defines a smooth hypersurface in \((\mathbb C^*)^n\) when nonempty. Then \[\chi\{f=0\text{ in }(\mathbb C^*)^n\}=(-1)^{n-1}n!\operatorname{Vol}_n(\Delta).\] For complex algebraic varieties the ordinary and compactly supported Euler characteristics agree, although their cohomology groups do not generally agree.}
\GSfield{Reconstruction}{The parity-dependent sign and the nondegeneracy assumption are supplied. The negative sign is correct for plane torus curves, but already fails for one-variable polynomials, whose zero set has positive cardinality.}
\GSfield{Fate}{The exact result is read in \citet[\S3, Theorem 2 and hypersurface example]{khovanskii1978genus}. It was published three years before the seminar occurrence.}
\GSfield{Present status}{As of November 2026, the statement with the displayed hypotheses is established; the notebook occurrence is expository.}
\GSfield{Assessment}{literal statement: false as written; publication match: exact; verification: read the result; development: annotated; confidence: high}
\end{GSitem}
\begin{GSitem}{1980-509.12}{E}{Pick’s formula from genus, punctures, and Euler characteristic}
\GSfield{Kind; source}{theorem; \GSpage{P000958}}
\GSfield{The notes}{Combining \(\chi=2-2B-\Gamma=-2S\), with \(B\) interior and \(\Gamma\) boundary lattice points, gives \(B+\Gamma/2=S+1\).}
\GSfield{Statement}{For a two-dimensional convex lattice polygon \(\Delta\), write \(I=\#(\operatorname{int}\Delta\cap\mathbb Z^2)\) and \(B=\#(\partial\Delta\cap\mathbb Z^2)\). Then \(\operatorname{Area}(\Delta)=I+B/2-1\). Choosing a generic nondegenerate torus curve gives genus \(I\), exactly \(B\) punctures, and Euler characteristic \(-2\operatorname{Area}(\Delta)\), which proves the identity.}
\GSfield{Reconstruction}{Notation distinguishes the notebook's area \(S\) from its mixed-area coefficient. A generic curve exists for any such polygon; the formula is for the polygon itself.}
\GSfield{Fate}{Pick's identity is classical. The same algebraic derivation is explicitly presented by the speaker in \citet[slides 6--7]{khovanskii2025chern}, and its ingredients were already published in \citet{khovanskii1978genus}.}
\GSfield{Present status}{As of November 2026, the statement with the displayed hypotheses is established; the notebook occurrence is expository.}
\GSfield{Assessment}{literal statement: correct; publication match: exact; verification: read the result; development: annotated; confidence: high}
\end{GSitem}
\begin{GSitem}{1980-509.13}{E}{Mixed area with a disk and the isoperimetric inequality}
\GSfield{Kind; source}{theorem; \GSpage{P000958}}
\GSfield{The notes}{The mixed area of a polygon and the unit disk is its perimeter. The following inequality is written in the diagonal-volume normalization, \(V(K,L)^2\geq\operatorname{Area}(K)\operatorname{Area}(L)\), and equality for a disk is mentioned.}
\GSfield{Statement}{For compact convex planar bodies \(K,L\) with nonempty interior, \(V(K,L)^2\geq\operatorname{Area}(K)\operatorname{Area}(L)\). If \(B\) is the unit disk, \(2V(K,B)=\operatorname{Perimeter}(K)\), because \(\operatorname{Area}(K+tB)=\operatorname{Area}(K)+t\operatorname{Perimeter}(K)+\pi t^2\). Hence \(\operatorname{Perimeter}(K)^2\geq4\pi\operatorname{Area}(K)\), with equality for disks.}
\GSfield{Reconstruction}{The notebook's perimeter statement uses its coefficient normalization \(S=2V\). This accounts for the factor of four in the isoperimetric inequality; convexity and full planar dimension are explicit.}
\GSfield{Fate}{These are established Minkowski and isoperimetric inequalities, not questions left open in 1981. The precise normalization and disk consequence are read in \citet[slide 12]{khovanskii2025chern}; the general Alexandrov--Fenchel statement is read in \citet[Theorem 2.1]{kavehkhovanskii2012equations}.}
\GSfield{Present status}{As of November 2026, the statement with the displayed hypotheses is established; the notebook occurrence is expository.}
\GSfield{Assessment}{literal statement: correct with missing hypotheses; publication match: exact; verification: read the result; development: annotated; confidence: high}
\end{GSitem}
\begin{GSitem}{1980-509.14}{E}{The surface intersection inequality requires positivity}
\GSfield{Kind; source}{theorem; \GSpage{P000958-P000959}}
\GSfield{The notes}{For any two curves on any algebraic surface the notes claim \((\Gamma_1\cdot\Gamma_2)^2\geq\Gamma_1^2\Gamma_2^2\), citing the Hodge structure and relating it to mixed area.}
\GSfield{Statement}{On a smooth projective complex surface, nef real divisor classes \(D,E\) satisfy \((D\cdot E)^2\geq D^2E^2\), by the Hodge index theorem and approximation by ample classes. The assertion for arbitrary curves is false: the two disjoint exceptional curves on the blow-up of \(\mathbb P^2\) at two distinct points have squares \(-1\) and mutual intersection zero. For globally generated toric divisors associated to polygons, \(D_K\cdot D_L=2V(K,L)\), which gives the mixed-area inequality.}
\GSfield{Reconstruction}{Smoothness, projectivity, and positivity are substantive hypotheses. The exceptional-curve example is a direct diagnostic counterexample, not a historical attribution. The last toric identity requires a common smooth projective refinement of the polygons' normal fans.}
\GSfield{Fate}{The surface Hodge inequality is classical. The exact positive-form version and its proof are read in \citet[Corollary 7.11]{kavehkhovanskii2010mixed}; the same paper explains the toric intersection/mixed-volume correspondence in \S2. The later work broadens the intersection theory but does not make the unrestricted notebook assertion true.}
\GSfield{Present status}{As of November 2026, the statement with the displayed hypotheses is established; the notebook occurrence is expository.}
\GSfield{Assessment}{literal statement: false as written; publication match: exact; verification: read the result; development: annotated; confidence: high}
\end{GSitem}
\begin{GSitem}{1980-509.15}{E}{Toric compactification from compatible monomial charts}
\GSfield{Kind; source}{construction; \GSpage{P000959}}
\GSfield{The notes}{The end of the talk begins a compactification by monomial changes of coordinates and cones, explaining how a curve or divisor becomes visible at the boundary.}
\GSfield{Statement}{A finite fan \(\Sigma\) of strongly convex rational polyhedral cones in the cocharacter lattice defines a normal toric variety by gluing \(U_\sigma=\operatorname{Spec}\mathbb C[\sigma^\vee\cap M]\) along their face charts. When every cone is generated by part of a lattice basis, the variety is smooth. A complete fan gives a compact complex toric variety. Refining the normal fan of the Newton polytopes supplies a compactification in which nondegenerate complete intersections meet torus strata transversely.}
\GSfield{Reconstruction}{Compatibility means that every cone intersection is a common face and all faces are included; mere disjointness of maximal-cone interiors does not suffice. The later 23 February occurrence supplies the fuller chart and completeness discussion.}
\GSfield{Fate}{The exact construction and admissibility/completeness theorems are read in \citet[\S1]{khovanskii1977toroidal}, and the systematic fan construction in \citet[\S5]{danilov1978toric}. Both sources were explicitly available before this lecture.}
\GSfield{Present status}{As of November 2026, the statement with the displayed hypotheses is established; the notebook occurrence is expository.}
\GSfield{Assessment}{literal statement: correct with missing hypotheses; publication match: exact; verification: read the result; development: annotated; confidence: high}
\end{GSitem}

\GStalk{511}{1981-02-23}{I. M. Gelfand}{P000960}{Dixmier operators and Grinevich’s construction}
The opening of the manuscript dated 23 February 1981 reports Dixmier’s commuting operators of orders four and six, then Grinevich’s extension using Tyurin parameters. The catalogue places this record in 1991; the autograph date and the surrounding 1981 sequence establish a catalogue decade error. This recovered entry precedes Cherednik’s talk in manuscript order.

\begin{GSitem}{1980-511.1}{E}{Dixmier’s commuting operators of orders four and six}
\GSfield{Kind; source}{theorem; \GSpage{P000960}}
\GSfield{The notes}{Two monic polynomial-coefficient operators of orders four and six commute but are not both functions of a single operator of order two; the note writes a square-root relation.}
\GSfield{Statement}{In the first Weyl algebra over \(\mathbb C\), there exist commuting ordinary differential operators of orders \(4,6\), rank \(2\), whose joint spectral curve is an elliptic cubic. The rank-two qualification concerns the generic common eigenspace, not merely the greatest common divisor of the two displayed orders.}
\GSfield{Fate}{This is the 1968 Dixmier example, revisited explicitly on P000970; see \citet{dixmier1968weyl}. The corrected operators and their algebraic identity are checked in the development section.}
\GSfield{Present status}{As of November 2026, the qualified statement above is established; no new open problem is inferred from the omissions in the notebook.}
\GSfield{Assessment}{literal statement: correct with missing hypotheses; publication match: exact; verification: publisher metadata checked; development: annotated; confidence: high}
\end{GSitem}
\begin{GSitem}{1980-511.2}{A}{Grinevich’s rational rank-two generalization}
\GSfield{Kind; source}{claim; \GSpage{P000960}}
\GSfield{The notes}{Gelfand says Grinevich obtained many examples using Tyurin parameters.}
\GSfield{Statement}{The identifiable content is a family of commuting rational-coefficient rank-two operators associated with an elliptic spectral curve, developed in the detailed talk on P000969--P000970. No numerical count or completeness assertion is recorded here.}
\GSfield{Fate}{Compare the rational-coefficient construction in \citet{grinevich1982rational}; the fuller occurrence is record 518.}
\GSfield{Present status}{As of November 2026, the qualified statement above is established; no new open problem is inferred from the omissions in the notebook.}
\GSfield{Assessment}{literal statement: correct with missing hypotheses; publication match: partial; verification: publisher metadata checked; development: annotated; confidence: high}
\end{GSitem}

\GStalk{510}{1981-02-23}{I. V. Cherednik}{P000960-P000961}{Inverse scattering, elliptic zero curvature, and Yang–Baxter structures}
Cherednik presents three stages of soliton theory, from scalar inverse scattering through matrix zero curvature to the Yang--Baxter equation, with elliptic examples and Landau--Lifshitz dynamics. Only the upper portion of P000961 belongs here; Khovanskii begins below it. The manuscript dates this material 23 February 1981, correcting the catalogue’s 1991 date.

\begin{GSitem}{1980-510.1}{E}{KdV inverse scattering and the recorded Lax pair}
\GSfield{Kind; source}{method; \GSpage{P000960}}
\GSfield{The notes}{Stage I displays \(u_t=6uu_x+u_{xxx}\), \(L=D^2-u\), \(A=D^3-uD-Du\), and \([L,A]=L_t\).}
\GSfield{Statement}{For a smooth scalar function \(u\), a consistent normalization is \(L=D^2+u\), \(A=-4D^3-3(uD+Du)\), and \(L_t=[L,A]\), which yields \(u_t=6uu_x+u_{xxx}\). For sufficiently decaying real potentials the associated one-dimensional scattering data supply the inverse-scattering variables.}
\GSfield{Reconstruction}{The signs and factors in \(L,A\) are an explicit emendation. The decay conditions belong to the scattering interpretation, not to the purely differential-operator identity.}
\GSfield{Fate}{The isospectral method precedes the talk; see \citet[pp.~467--490]{lax1968integrals}. Direct use of \(Df=fD+f^{\prime}\) shows that the notebook’s three displayed formulas do not form this Lax pair.}
\GSfield{Present status}{As of November 2026, the qualified statement above is established; no new open problem is inferred from the omissions in the notebook.}
\GSfield{Assessment}{literal statement: false as written; publication match: partial; verification: read the result; development: annotated; confidence: high}
\end{GSitem}
\begin{GSitem}{1980-510.2}{E}{Matrix zero-curvature equations}
\GSfield{Kind; source}{method; \GSpage{P000960}}
\GSfield{The notes}{Stage II uses \([D_x-U,D_t-V]=0\), with \(U,V\) rational in the spectral parameter; an example has degrees two and one.}
\GSfield{Statement}{For matrix-valued \(U(x,t,\lambda),V(x,t,\lambda)\), compatibility of \(\psi_x=U\psi\), \(\psi_t=V\psi\) is \(U_t-V_x+[U,V]=0\) at each regular value of \(\lambda\). A spectral-parameter representation alone does not establish a Hamiltonian structure or completeness of inverse scattering.}
\GSfield{Fate}{This is the established zero-curvature formulation; the relativistic elliptic models discussed next belong to the programme in \citet{cherednik1981quasiclassical}, submitted 8 August 1980 and published in 1981.}
\GSfield{Present status}{As of November 2026, the qualified statement above is established; no new open problem is inferred from the omissions in the notebook.}
\GSfield{Assessment}{literal statement: correct with missing hypotheses; publication match: partial; verification: publisher metadata checked; development: annotated; confidence: high}
\end{GSitem}
\begin{GSitem}{1980-510.3}{F}{The third-stage unification slogan}
\GSfield{Kind; source}{remark; \GSpage{P000960-P000961}}
\GSfield{The notes}{The Yang--Baxter equation is presented as the third stage, encompassing all or almost all integrable systems and connecting scattering with statistical mechanics.}
\GSfield{Statement}{The diagram records a programme of constructions and limiting procedures between specified models. It is not a theorem that every system called integrable arises from a Yang--Baxter equation.}
\GSfield{Fate}{\citet[Introduction]{belavin1980discrete} gives concrete scattering/discrete-symmetry constructions; it does not prove the unqualified universality slogan.}
\GSfield{Present status}{As of November 2026, this broad slogan has no single precise mathematical proposition whose truth or resolution could be assigned.}
\GSfield{Assessment}{literal statement: not determinable; publication match: thematic only; verification: read the result; development: annotated; confidence: high}
\end{GSitem}
\begin{GSitem}{1980-510.4}{A}{Elliptic linear systems for anisotropic vector equations}
\GSfield{Kind; source}{construction; \GSpage{P000961}}
\GSfield{The notes}{The auxiliary vector \(E\) satisfies \(E_x=(J(\lambda-\alpha)V)\times E\), \(E_t=(J(\lambda)U)\times E\). The ensuing nonlinear system has \(U_x=(JV)\times U\) and a second line written with \(U_t\).}
\GSfield{Statement}{An elliptic spectral parameter can enter an \(\mathfrak{so}_3\)-valued flat connection by identifying vectors with cross-product matrices. Compatibility requires the coefficient of each independent elliptic function in its curvature to vanish. The exact two-field system cannot be recovered from an unspecified elliptic matrix \(J(\lambda)\) and the repeated \(U\) in the second line alone.}
\GSfield{Reconstruction}{The missing equality sign in the \(E_t\) line is syntactic; replacing the second \(U_t\) by \(V_t\) would be a mathematical conjecture about the intended system and is not silently adopted.}
\GSfield{Fate}{The subject matches Cherednik’s relativistic quasiclassical models \citep[abstract and pp.~422--425]{cherednik1981quasiclassical}; this is a partial, not equation-by-equation, identification.}
\GSfield{Present status}{As of November 2026, the construction is identifiable at the level of elliptic zero-curvature models, but the exact notebook normalization is unresolved.}
\GSfield{Assessment}{literal statement: not determinable; publication match: partial; verification: publisher metadata checked; development: annotated; confidence: medium}
\end{GSitem}
\begin{GSitem}{1980-510.5}{E}{Landau–Lifshitz dynamics and scattering brackets}
\GSfield{Kind; source}{claim; \GSpage{P000961}}
\GSfield{The notes}{The notes write \(U_t=U\times U_{xx}+(JU)\times U\) and say that Yang--Baxter theory supplies Hamiltonian and scattering brackets.}
\GSfield{Statement}{For a smooth spin field of fixed length and a constant diagonal anisotropy matrix \(J\), this is an anisotropic Landau--Lifshitz equation, with sign determined by the spin Poisson bracket convention. A statement about brackets of scattering data additionally requires a chosen Lax representation and boundary conditions.}
\GSfield{Fate}{Classical and quantum inverse-scattering constructions already precede this lecture; compare \citet{sklyaninfaddeev1978quantum}. The notebook does not specify the scattering-data bracket to be identified.}
\GSfield{Present status}{As of November 2026, the qualified statement above is established; no new open problem is inferred from the omissions in the notebook.}
\GSfield{Assessment}{literal statement: correct with missing hypotheses; publication match: thematic only; verification: publisher metadata checked; development: annotated; confidence: medium}
\end{GSitem}

\GStalk{512}{1981-02-23}{A. G. Khovanskii}{P000961 (lower part)-P000962 (upper part)}{Continuation: monomial charts and fans}
This recovered catalogue record continues Khovanskii after the Cherednik and Gelfand interventions of 23 February. The live catalogue's decade typo had excluded it from the supplied academic-year list. Character--cocharacter pairings, cones of one-parameter subgroups, a modification chart, the completeness criterion, and the fan of the projective plane precede the separate Kazarnovskii heading. The actual manuscript date is 1981, not 1991.

\begin{GSitem}{1980-512.1}{E}{Unimodular torus charts and character--cocharacter duality}
\GSfield{Kind; source}{construction; \GSpage{P000961-P000962}}
\GSfield{The notes}{The continuation writes \(x_i=\prod_j u_j^{a_{ij}}\), one-parameter groups \(x_i=t^{\lambda_i}\), and the value \(t^{\langle m,\lambda\rangle}\) of a monomial on such a group.}
\GSfield{Statement}{For \(T=(\mathbb C^*)^n\), the character lattice \(M=\operatorname{Hom}(T,\mathbb C^*)\) and cocharacter lattice \(N=\operatorname{Hom}(\mathbb C^*,T)\) are dual. In chosen bases, \(\chi^m(t^\lambda)=t^{\langle m,\lambda\rangle}\). The monomial change \(x_i=\prod_j u_j^{a_{ij}}\) is an isomorphism of tori precisely when \(A\in\operatorname{GL}_n(\mathbb Z)\). An \(\mathbb A^n\) chart in a smooth toric compactification corresponds to a maximal cone generated by its cocharacter-basis columns.}
\GSfield{Reconstruction}{Unimodularity distinguishes a chart isomorphism from a finite monomial covering. The statement concerns a chosen torus with its action, not arbitrary analytic compactifications.}
\GSfield{Fate}{The exact pairing, monomial coordinate change, and associated cones are read in \citet[\S1]{khovanskii1977toroidal}, published in 1977, and \citet[\S5]{danilov1978toric}.}
\GSfield{Present status}{As of November 2026, the statement with the displayed hypotheses is established; the notebook occurrence is expository.}
\GSfield{Assessment}{literal statement: correct with missing hypotheses; publication match: exact; verification: read the result; development: annotated; confidence: high}
\end{GSitem}
\begin{GSitem}{1980-512.2}{E}{The fan compatibility condition}
\GSfield{Kind; source}{theorem; \GSpage{P000962}}
\GSfield{The notes}{Cones are generated by the columns of the chart matrix. Overlapping interiors force equality; a second statement applies the same requirement to faces.}
\GSfield{Statement}{A finite collection of strongly convex rational cones defines a separated toric gluing when it is closed under faces and the intersection of every two cones is a face of each. Equivalently, relative interiors of all its cones are disjoint unless the cones coincide. Smoothness requires every cone to be generated by part of a lattice basis.}
\GSfield{Reconstruction}{The second, facewise sentence is essential: disjoint interiors of maximal cones alone are insufficient. Including all faces makes the notebook's criterion the standard fan condition.}
\GSfield{Fate}{This is the admissibility theorem in \citet[\S1]{khovanskii1977toroidal}; compare the systematic fan gluing in \citet[\S5]{danilov1978toric}. It is established before the lecture.}
\GSfield{Present status}{As of November 2026, the statement with the displayed hypotheses is established; the notebook occurrence is expository.}
\GSfield{Assessment}{literal statement: correct with missing hypotheses; publication match: exact; verification: read the result; development: annotated; confidence: high}
\end{GSitem}
\begin{GSitem}{1980-512.3}{E}{The sigma-process example is one blow-up chart}
\GSfield{Kind; source}{construction; \GSpage{P000962}}
\GSfield{The notes}{The substitution \(x_1=u_1u_2^{-1},\ x_2=u_2\), or \(u_1=x_1x_2,\ u_2=x_2\), is called a sigma-process.}
\GSfield{Statement}{The map \((x_1,x_2)\mapsto(u_1,u_2)=(x_1x_2,x_2)\) is one affine chart of the blow-up of the origin in the \(u\)-plane; the other standard chart is needed for the complete blow-up. Torically, the modification subdivides the positive-quadrant cone by the ray \(e_1+e_2\). On the dense torus the displayed inverse is the stated unimodular monomial transformation.}
\GSfield{Reconstruction}{The direction of the morphism and the need for the second chart are explicit; the inverse formula is not regular along \(u_2=0\) in the original plane.}
\GSfield{Fate}{The standard toric modification construction is read in \citet[\S5 and the discussion of subdivisions]{danilov1978toric}; its unimodular chart interpretation is also in \citet[\S1]{khovanskii1977toroidal}.}
\GSfield{Present status}{As of November 2026, the statement with the displayed hypotheses is established; the notebook occurrence is expository.}
\GSfield{Assessment}{literal statement: correct with missing hypotheses; publication match: exact; verification: read the result; development: annotated; confidence: high}
\end{GSitem}
\begin{GSitem}{1980-512.4}{E}{Completeness, the projective-plane fan, and toric classification}
\GSfield{Kind; source}{theorem; \GSpage{P000962}}
\GSfield{The notes}{The compactness criterion is that the cones cover the whole vector space. A three-ray drawing gives the projective plane, and the compactification is said to be determined up to an integral automorphism.}
\GSfield{Statement}{The complex toric variety \(X_\Sigma\) is complete (equivalently compact analytically) exactly when \(|\Sigma|=N_\mathbb R\). The fan with rays \((1,0),(0,1),(-1,-1)\) and the three intervening two-dimensional cones gives \(\mathbb P^2\). Smooth complete toric varieties with their dense torus action are classified by smooth complete fans; allowing an automorphism of the torus identifies fans under \(\operatorname{GL}(N)\). Completeness alone does not assert projectivity.}
\GSfield{Reconstruction}{The classification is equivariant and concerns valid fans. The projective-plane ray generators are reconstructed from the standard three-chart model shown schematically.}
\GSfield{Fate}{The completeness equivalence is read in \citet[\S1, compactness theorem]{khovanskii1977toroidal}; \citet[\S5]{danilov1978toric} supplies the fan classification. Both precede 1981.}
\GSfield{Present status}{As of November 2026, the statement with the displayed hypotheses is established; the notebook occurrence is expository.}
\GSfield{Assessment}{literal statement: correct with missing hypotheses; publication match: exact; verification: read the result; development: annotated; confidence: high}
\end{GSitem}

\GStalk{513}{1981-02-23}{B. Ya. Kazarnovskii}{P000962 (lower part)-P000963; follow-up P000964}{Zeros of exponential sums and mixed pseudovolume}
Kazarnovskii extends Newton-polytope root counting to exponential sums with complex spectra, explains mixed pseudovolume through faces and complex tangent geometry, and mentions its support-function formulation. A separately dated 2 March question by I. M. Gelfand asks for the connection with the growth indicatrix; that follow-up is included with its own speaker and date, while the neighbouring Gaussian-tail exercise is assigned to record 514.

\begin{GSitem}{1980-513.1}{A}{Exponential sums and complex Newton polytopes}
\GSfield{Kind; source}{definition; \GSpage{P000962}}
\GSfield{The notes}{The proposed extension replaces Laurent polynomials by finite sums \(f(z)=\sum_{\lambda\in\Lambda}c_\lambda e^{\langle z,\lambda\rangle}\), with \(\Lambda\subset\mathbb C^n\), and calls \(\operatorname{conv}(\Lambda)\) the Newton polyhedron.}
\GSfield{Statement}{For a finite spectrum \(\Lambda\subset(\mathbb C^n)^*\), with nonzero coefficients on its support, the expression defines an entire exponential sum. Its complex Newton polytope is the real convex hull \(\Gamma\subset\mathbb R^{2n}\) of that spectrum. The pairing in the exponential is complex bilinear; the associated real support function is \(h_\Gamma(z)=\max_{\lambda\in\Gamma}\operatorname{Re}\langle z,\lambda\rangle\). Integer real spectra recover Laurent polynomials after \(x_j=e^{z_j}\).}
\GSfield{Reconstruction}{The real convex hull and the complex-bilinear pairing are distinguished. The logarithmic substitution is a covering map, so it changes finite torus counts into densities of infinitely many lifted roots.}
\GSfield{Fate}{The definition is read at the opening of \citet[p.347]{kazarnovskii1981zeros}. The paper appeared in 1981 and records receipt on 6 November 1980; receipt is not treated as evidence of public circulation. The seminar report on 23 February 1981 is the verified oral occurrence.}
\GSfield{Present status}{As of November 2026, the complex-spectrum formulation is established; it is the precise setting of the announced extension.}
\GSfield{Assessment}{literal statement: correct; publication match: exact; verification: read the result; development: annotated; confidence: high}
\end{GSitem}
\begin{GSitem}{1980-513.2}{A}{Asymptotic zero count for regular exponential systems}
\GSfield{Kind; source}{theorem; \GSpage{P000962}}
\GSfield{The notes}{For \(F=(f_1,\ldots,f_n)\), the notebook states \(N_F(R)=cV(\Gamma_1,\ldots,\Gamma_n)R^n+O(R^{n-1})\), without recording the regularity condition or exact constant.}
\GSfield{Statement}{Put \(\Gamma=\sum_i\Gamma_i\). For each face \(\Delta\) of \(\Gamma\), write \(\Delta=\sum_i\Delta_i\) and form the face sums \(f_i^\Delta\). Assume there is \(\varepsilon>0\) such that for every face whose complex affine span has dimension less than \(n\), and every \(z\in\mathbb C^n\), \[\sum_i\frac{|f_i^\Delta(z)|}{\max_{\lambda\in\Lambda_i\cap\Delta_i}|e^{\langle z,\lambda\rangle}|}\geq\varepsilon.\] This is Kazarnovskii's regularity condition. Counting isolated common roots with multiplicity in the radius-\(R\) Euclidean ball, \[N_F(R)=\frac{\kappa_n n!}{(2\pi)^n}P_n(\Gamma_1,\ldots,\Gamma_n)R^n+O(R^{n-1}),\] where \(\kappa_n\) is the volume of the real unit \(n\)-ball and \(P_n\) agrees with ordinary mixed volume for bodies in \(\mathbb R^n\).}
\GSfield{Reconstruction}{The condition and normalization are imported from the publication. A system with common positive-dimensional zero components cannot be counted by the literal bare formula. For \(e^z-1\), \(P_1([0,1])=1\) gives \(N(R)=R/\pi+O(1)\), checking the constant. The actual notebook scan clearly has uppercase \(O\); the retrieved 1981 OCR renders that glyph ambiguously.}
\GSfield{Fate}{This matches \citet[Definition 3 and Theorem 1, pp.348--349]{kazarnovskii1981zeros}, specialized to \(r=n\). The paper appeared in 1981 and records receipt on 6 November 1980; receipt is not treated as evidence of public circulation. The seminar report on 23 February 1981 is the verified oral occurrence.}
\GSfield{Present status}{As of November 2026, the stated regular-system asymptotic is established. No claim is made for an arbitrary fixed system with the same spectra.}
\GSfield{Assessment}{literal statement: correct with missing hypotheses; publication match: exact; verification: read the result; development: annotated; confidence: high; problem \hyperref[prob:GP-1980-26]{GP-1980-26}}
\end{GSitem}
\begin{GSitem}{1980-513.3}{A}{The face formula for pseudovolume}
\GSfield{Kind; source}{definition; \GSpage{P000963}}
\GSfield{The notes}{The diagonal functional is a sum over real \(n\)-faces of volume times external angle times the area distortion of \(E^\perp\to iE\).}
\GSfield{Statement}{For a convex polytope \(\Gamma\subset\mathbb C^n\), \[P_n(\Gamma)=\sum_{\dim_{\mathbb R} F=n}\operatorname{Vol}_n(F)\,\Omega(F)\,g(E_F),\qquad g(E)=\left|\det\bigl(\operatorname{pr}_{iE}|_{E^\perp}\bigr)\right|.\] Here \(E_F\) is the direction space of the face, and \(\Omega(F)\) is the solid angle of its exterior normal cone, normalized as a fraction of the unit ball in \(E_F^\perp\). In particular \(g(\mathbb R^n)=1\), while \(g(E)=0\) when \(E\) contains a complex line. The formula is unitarily invariant; for \(n>1\) it is not generally monotone under inclusion.}
\GSfield{Reconstruction}{All face dimensions are real; the determinant is an absolute Euclidean Jacobian between equal-dimensional real spaces. For a lower-dimensional body its whole body can be a face and its normal cone can be a full orthogonal space.}
\GSfield{Fate}{The exact face construction, complex-angle factor, and nonmonotonicity warning are read in \citet[\S1, Definition 1]{kazarnovskii1981zeros}. The notebook's projection is the transpose version of the same equal-dimensional projection determinant. The paper appeared in 1981 and records receipt on 6 November 1980; receipt is not treated as evidence of public circulation. The seminar report on 23 February 1981 is the verified oral occurrence.}
\GSfield{Present status}{As of November 2026, the face formula and its Hermitian dependence are established; ordinary real volume is only a special case.}
\GSfield{Assessment}{literal statement: correct with missing hypotheses; publication match: exact; verification: read the result; development: annotated; confidence: high}
\end{GSitem}
\begin{GSitem}{1980-513.4}{A}{Polarization and the real-spectrum specialization}
\GSfield{Kind; source}{theorem; \GSpage{P000962-P000963}}
\GSfield{The notes}{The multilinear \(V(\Gamma_1,\ldots,\Gamma_n)\) is described through the single-polytope functional; it is ordinary mixed volume when all spectra are real. In one variable a perimeter appears.}
\GSfield{Statement}{There is a symmetric Minkowski-multilinear polarization \(P_n(\Gamma_1,\ldots,\Gamma_n)\) whose diagonal is the preceding face functional. For bodies in the same real form \(\mathbb R^n\subset\mathbb C^n\), it is ordinary mixed volume in the normalization \(V(K,\ldots,K)=\operatorname{Vol}_n(K)\). For \(n=1\), \(P_1\) is half the planar perimeter, with a line segment assigned perimeter twice its length. Thus the one-variable leading zero count is \(\operatorname{Perimeter}(\Gamma)R/(2\pi)\).}
\GSfield{Reconstruction}{The polarization and semiperimeter normalization explain the unspecified constant in the notes. The statement is not ordinary \(2n\)-dimensional mixed volume of complex polytopes.}
\GSfield{Fate}{The exact polarization is \citet[Propositions 1--2 and opening comparison, pp.347--348]{kazarnovskii1981zeros}. The paper appeared in 1981 and records receipt on 6 November 1980; receipt is not treated as evidence of public circulation. The seminar report on 23 February 1981 is the verified oral occurrence.}
\GSfield{Present status}{As of November 2026, the mixed pseudovolume and real-spectrum reduction are established.}
\GSfield{Assessment}{literal statement: correct with missing hypotheses; publication match: exact; verification: read the result; development: annotated; confidence: high}
\end{GSitem}
\begin{GSitem}{1980-513.5}{A}{Smooth-boundary and Levi-form descriptions of pseudovolume}
\GSfield{Kind; source}{construction; \GSpage{P000963}}
\GSfield{The notes}{The notes extend the polytope functional to bodies with smooth boundary, introduce the complex tangent hyperplane, symmetrize the second fundamental form, and write an integral of the product of Levi eigenvalues, labelled \(S(\Delta)\).}
\GSfield{Statement}{For a bounded domain \(K\subset\mathbb C^n\) with smooth boundary, let \(\nu\) be its outward unit normal and \(\alpha(v)=\langle v,i\nu\rangle\) on \(\partial K\). The normalized extension is \[P_n(K)=\frac1{\kappa_n n!}\int_{\partial K}\alpha\wedge(d\alpha)^{n-1}.\] For convex bodies it agrees with the polytope functional by approximation. On the complex tangent space \(H=T\partial K\cap iT\partial K\), the Hermitian part of the second fundamental form is the Levi form; the integral can equivalently be expressed through its determinant with the convention-dependent numerical factor.}
\GSfield{Reconstruction}{The notebook's final label conflates the dimensionless face-angle factor \(S(\Delta)\) with the integrated body functional. Its division by two in defining the Hermitian part changes the Levi eigenvalue normalization, so no unverified factor is attached to its eigenvalue-product formula. The contact-form expression fixes a consistent normalization independently.}
\GSfield{Fate}{All three constructions are read in \citet[\S1(a)--(c), p.348]{kazarnovskii1981zeros}; the paper explicitly allows nonconvex smooth domains for the first two boundary expressions. The paper appeared in 1981 and records receipt on 6 November 1980; receipt is not treated as evidence of public circulation. The seminar report on 23 February 1981 is the verified oral occurrence.}
\GSfield{Present status}{As of November 2026, the normalized boundary construction is established; the notebook's notation and unsupplied eigenvalue factor require qualification.}
\GSfield{Assessment}{literal statement: false as written; publication match: exact; verification: read the result; development: annotated; confidence: medium}
\end{GSitem}
\begin{GSitem}{1980-513.6}{A}{Support functions and the growth-indicatrix question of 2 March}
\GSfield{Kind; source}{question; \GSpage{P000963-P000964}; I. M. Gelfand (follow-up, 1981-03-02)}
\GSfield{The notes}{P000963 ends by mentioning a support-function formulation. On 2 March 1981, I. M. Gelfand asks on P000964 to relate Kazarnovskii's result to the growth indicatrix.}
\GSfield{Statement}{With \(d^c=i(\bar\partial-\partial)\) and \(h_j=h_{\Gamma_j}\), the precise bridge is \[P_n(\Gamma_1,\ldots,\Gamma_n)=\frac1{\kappa_n n!}\int_{B_1}dd^ch_1\wedge\cdots\wedge dd^ch_n.\] These positive currents are well defined for the continuous convex support functions, by smooth approximation. For a finite exponential sum, \(\log|f(tz)|/t\leq h_\Gamma(z)+O(1/t)\); if one spectral exponent uniquely maximizes the real pairing at a direction, the limit equals the support function at that direction. Exceptional directions require the usual upper regularization of the indicator. Poincaré--Lelong converts logarithmic growth to zero currents. Passing to a product of zero currents or to a fixed-system remainder still requires appropriate regularity; individual pointwise bounds alone do not justify that passage.}
\GSfield{Reconstruction}{The statement answers the requested structural relation without turning it into an unrestricted convergence conjecture. The date and intervening speaker are explicit. The support-function integral already occurs in the 1981 publication, so the connection is not credited solely to much later work.}
\GSfield{Fate}{\citet[\S1(c), p.348]{kazarnovskii1981zeros} gives the exact support-function expression. \citet[Definition 5 and Theorems 1--2]{kazarnovskii2024distribution} later treats averaged zero currents for finite-dimensional spaces whose normalized evaluation norms converge locally uniformly to the support function; that is a precise extension with averaging, published 11 July 2024 after submissions in 2023 and 2024. Its displayed HTML normalizations are not imported here. The factor above uses the 1981 real-volume normalization and is checked on an interval.}
\GSfield{Present status}{As of November 2026, the finite-spectrum support-function connection is established, and later averaged generalizations have precise regularity assumptions. The notebook sentence supplies no broader independently stated open conjecture.}
\GSfield{Assessment}{literal statement: correct with missing hypotheses; publication match: exact; verification: read the result; development: annotated; confidence: high; problem \hyperref[prob:GP-1980-27]{GP-1980-27}}
\end{GSitem}

\GStalk{514}{1981-03-02}{I. V. Cherednik}{P000964-P000965}{A diagram relating Yang–Baxter structures}
A short diagrammatic reprise links factorized scattering, two-dimensional statistical mechanics, quantum one-dimensional systems, and integrable equations. Belavin adds unitarity. Gelfand’s preceding Gaussian-tail exercise is retained here as an explicitly attributed intervention; his adjacent Kazarnovskii question is indexed with record 513. The detailed three-stage account is in recovered record 510, not this March entry.

\begin{GSitem}{1980-514.1}{E}{Borel summation of the Gaussian tail}
\GSfield{Kind; source}{problem; \GSpage{P000964}; I. M. Gelfand}
\GSfield{The notes}{Gelfand asks for the divergent asymptotic expansion of \(\int_x^\infty e^{-t^2}\,dt\) and its summation by Borel’s method.}
\GSfield{Statement}{For \(x>0\), put \(z=x^{-2}\) and \(F(z)=2xe^{x^2}\int_x^\infty e^{-t^2}\,dt\). Then \(F(z)\sim\sum_{n\ge0}(-1)^n(1/2)_n z^n\), whose Borel transform is \((1+s)^{-1/2}\). Its positive-direction Laplace sum \(z^{-1}\int_0^\infty e^{-s/z}(1+s)^{-1/2}\,ds\) equals \(F(z)\).}
\GSfield{Fate}{The asymptotic expansion is classical; \citet[§7.12(i), Eq.~7.12.1]{nist2026dlmf} records it and remainder bounds. The elementary substitution \(t^2=x^2+s x^2\) proves the stated Borel integral.}
\GSfield{Present status}{As of November 2026, the qualified statement above is established; no new open problem is inferred from the omissions in the notebook.}
\GSfield{Assessment}{literal statement: correct; publication match: exact; verification: read the result; development: annotated; confidence: high}
\end{GSitem}
\begin{GSitem}{1980-514.2}{E}{Yang–Baxter connections and unitarity}
\GSfield{Kind; source}{method; \GSpage{P000965}}
\GSfield{The notes}{The diagram connects factorized scattering, two-dimensional statistical mechanics, integrable systems, and quantum one-dimensional systems. Belavin adds unitarity conditions.}
\GSfield{Statement}{In a specified factorized relativistic theory, a two-body matrix must satisfy the Yang--Baxter consistency relation; physical unitarity and crossing impose additional independent restrictions. Transfer-matrix and continuum constructions relate certain such matrices to lattice and differential models.}
\GSfield{Fate}{These relationships already occur in \citet[§§1--2]{belavin1980discrete}; the more detailed construction follows in record 516. They recur from 1978-453.19, 1978-454.3, and 1978-460.14 without creating a new problem.}
\GSfield{Present status}{As of November 2026, the qualified statement above is established; no new open problem is inferred from the omissions in the notebook.}
\GSfield{Assessment}{literal statement: correct with missing hypotheses; publication match: partial; verification: read the result; development: annotated; confidence: high}
\end{GSitem}

\GStalk{515}{1981-03-02}{A. A. Belavin}{P000965-P000966}{Inversion, crossing, and Baxter invariant curves}
Belavin studies matrix inversion and partial transposition, then the invariant quantities of the symmetric eight-vertex matrix. Elliptic curves appear after projectivizing and imposing generic smoothness. The notebook’s symmetry sign and the distinction between generic and exceptional orbits require qualification.

\begin{GSitem}{1980-515.1}{B}{Inversion and partial-transpose dynamics}
\GSfield{Kind; source}{construction; \GSpage{P000965}}
\GSfield{The notes}{\(\Omega R=R^{-1}\), crossing \(C\) transposes the block indices, \(\Omega^2=C^2=1\), and the product is said to have no finite order.}
\GSfield{Statement}{On the Zariski-open set where the required inverses exist, inversion and partial transpose are birational involutions of projective matrix space. Their product generally has infinite order; special points and invariant subvarieties can have periodic orbits.}
\GSfield{Reconstruction}{Projectivization, the domain of the rational maps, and the distinction between generic and periodic behaviour are supplied.}
\GSfield{Fate}{The corresponding birational-involution approach was developed systematically in \citet[§2]{bellonmaillardviallet1991higher}; the paper was received 25 February 1991 and published in October. The match is to the method, not a universal classification.}
\GSfield{Present status}{As of November 2026, the qualified statement above is established; no new open problem is inferred from the omissions in the notebook.}
\GSfield{Assessment}{literal statement: correct with missing hypotheses; publication match: partial; verification: read the result; development: annotated; confidence: high}
\end{GSitem}
\begin{GSitem}{1980-515.2}{E}{The symmetric eight-vertex subspace}
\GSfield{Kind; source}{construction; \GSpage{P000965-P000966}}
\GSfield{The notes}{The displayed \(4\times4\) matrix has entries \(a,b,c,d\) in the symmetric eight-vertex pattern and is described by two simultaneous \(\mathbb Z_2\) conjugations.}
\GSfield{Statement}{For \(V=\mathbb C^2\), the matrices commuting with both \(\sigma_x\otimes\sigma_x\) and \(\sigma_z\otimes\sigma_z\) have the form \(\left(\begin{smallmatrix}a&0&0&d\\0&b&c&0\\0&c&b&0\\d&0&0&a\end{smallmatrix}\right)\). This four-dimensional linear space is preserved by inversion where defined and by partial transpose.}
\GSfield{Fate}{This is the zero-field eight-vertex family associated with \citet{baxter1971eight}. The minus sign preceding simultaneous conjugation in the notebook is incompatible with literal fixed-point invariance, although a global sign disappears in projective coordinates.}
\GSfield{Present status}{As of November 2026, the qualified statement above is established; no new open problem is inferred from the omissions in the notebook.}
\GSfield{Assessment}{literal statement: false as written; publication match: exact; verification: publisher metadata checked; development: annotated; confidence: high}
\end{GSitem}
\begin{GSitem}{1980-515.3}{E}{Baxter invariants and the generic elliptic fibre}
\GSfield{Kind; source}{theorem; \GSpage{P000966}}
\GSfield{The notes}{Two invariants are given as \(I_1=(a^2+b^2-c^2-d^2)/(ab)\), \(I_2=cd/(ab)\), and their level sets are called elliptic curves.}
\GSfield{Statement}{On \(ab\ne0\), these functions are invariant under the projective inversion/crossing dynamics. In \(\mathbb P^3\), a generic common level is a smooth complete intersection of two quadrics, hence a genus-one curve. Singular levels can degenerate, while the affine level is a cone of dimension two.}
\GSfield{Reconstruction}{The projective ambient space and generic-smoothness conditions repair the literal dimension and degeneration omissions.}
\GSfield{Fate}{The elliptic parametrization is the Baxter eight-vertex mechanism \citep{baxter1971eight}; later birational treatments are \citet[§2]{bellonmaillardviallet1991higher}. The genus calculation uses adjunction and requires smoothness.}
\GSfield{Present status}{As of November 2026, the qualified statement above is established; no new open problem is inferred from the omissions in the notebook.}
\GSfield{Assessment}{literal statement: correct with missing hypotheses; publication match: partial; verification: read the result; development: annotated; confidence: high}
\end{GSitem}
\begin{GSitem}{1980-515.4}{F}{Invariant varieties beyond the eight-vertex example}
\GSfield{Kind; source}{question; \GSpage{P000965-P000966}}
\GSfield{The notes}{Belavin asks which varieties are invariant under the inversion/crossing operation, before presenting the eight-vertex answer.}
\GSfield{Statement}{For a fixed matrix size, one may study algebraic subvarieties preserved birationally by the composition of inverse and partial transpose. The notes do not specify a dimension, an equivalence relation, or a requirement to classify all such subvarieties.}
\GSfield{Fate}{\citet{bellonmaillardviallet1991higher} studies higher-dimensional examples and their invariants, providing a thematic continuation. It does not identify a unique comprehensive problem posed by this sentence.}
\GSfield{Present status}{As of November 2026, the eight-vertex example is established; no determinate global classification question is reconstructed from the broader sentence.}
\GSfield{Assessment}{literal statement: not determinable; publication match: thematic only; verification: read the result; development: annotated; confidence: medium}
\end{GSitem}

\GStalk{516}{1981-03-02}{I. V. Cherednik}{P000966-P000968}{Factorized scattering and commuting transfer operators}
The continuation derives the Yang--Baxter relation from three-particle reordering, rewrites it in tensor notation, gives an elliptic Pauli-matrix solution, and states the commuting-transfer-operator mechanism. The index in the component relation and the tensor leg traced in the last theorem are faithfully recorded mathematical errors, not transcription mistakes.

\begin{GSitem}{1980-516.1}{E}{From factorized scattering to spin chains and spin fields}
\GSfield{Kind; source}{remark; \GSpage{P000966}}
\GSfield{The notes}{The programme runs from a factorized particle theory through the eight-vertex model and the XYZ chain to the Landau--Lifshitz equation.}
\GSfield{Statement}{This is a sequence of model-dependent identifications and limits: vertex weights define transfer matrices, logarithmic derivatives at a regular spectral value give spin-chain Hamiltonians, and suitable long-wave classical limits give spin-field equations.}
\GSfield{Fate}{The eight-vertex and quantum inverse-scattering ingredients predate the lecture \citep{baxter1971eight,sklyaninfaddeev1978quantum}. The note expressly omits the Baxter step and provides no general limit theorem.}
\GSfield{Present status}{As of November 2026, the qualified statement above is established; no new open problem is inferred from the omissions in the notebook.}
\GSfield{Assessment}{literal statement: correct with missing hypotheses; publication match: partial; verification: publisher metadata checked; development: annotated; confidence: high}
\end{GSitem}
\begin{GSitem}{1980-516.2}{E}{Two-body amplitudes and factorization}
\GSfield{Kind; source}{definition; \GSpage{P000966-P000967}}
\GSfield{The notes}{Particles have equal mass and \(d\) internal types. Each crossing contributes \(S_{ij}^{kl}\), and a many-particle amplitude is the product summed over internal labels.}
\GSfield{Statement}{In elastic factorized scattering in \(1+1\) dimensions, the two-body amplitude acts on \(V\otimes V\), \(\dim V=d\). Lorentz invariance makes its spectral variable a rapidity difference. Summing internal states in an ordered product gives the many-body amplitude; unitarity is an additional condition.}
\GSfield{Fate}{Compare the rapidity formulation in \citet[§1]{belavin1980discrete}. An ordinary angle in a Euclidean spacetime sketch is not the invariant rapidity.}
\GSfield{Present status}{As of November 2026, the qualified statement above is established; no new open problem is inferred from the omissions in the notebook.}
\GSfield{Assessment}{literal statement: correct with missing hypotheses; publication match: exact; verification: read the result; development: annotated; confidence: high}
\end{GSitem}
\begin{GSitem}{1980-516.3}{E}{Three-line consistency gives the Yang–Baxter equation}
\GSfield{Kind; source}{theorem; \GSpage{P000967}}
\GSfield{The notes}{It is enough to slide one vertex past the intersection of two other lines. Equation (12) gives the component relation.}
\GSfield{Statement}{Consistency of the two resolutions of a generic three-particle collision requires \(S_{12}(\alpha)S_{13}(\alpha+\beta)S_{23}(\beta)=S_{23}(\beta)S_{13}(\alpha+\beta)S_{12}(\alpha)\). For generic line arrangements with fixed endpoints, triangle moves together with commuting disjoint crossings generate the relevant reorderings.}
\GSfield{Fate}{The factorization relation is explicitly developed in \citet[§1]{belavin1980discrete}; it recurs from 1978-454.3. The notebook’s component equation has a mismatched free lower index on its right-hand side.}
\GSfield{Present status}{As of November 2026, the qualified statement above is established; no new open problem is inferred from the omissions in the notebook.}
\GSfield{Assessment}{literal statement: false as written; publication match: exact; verification: read the result; development: annotated; confidence: high}
\end{GSitem}
\begin{GSitem}{1980-516.4}{E}{Finding a nontrivial two-body Yang–Baxter matrix}
\GSfield{Kind; source}{problem; \GSpage{P000967}}
\GSfield{The notes}{The stated problem is to find \(S\); the note contrasts \(d^6\) component relations with \(d^4\) unknown functions.}
\GSfield{Statement}{For a chosen internal space and spectral-parameter convention, construct non-scalar meromorphic matrices satisfying the Yang--Baxter relation. The counts \(d^6,d^4\) count written components and entries; they do not prove algebraic independence or nonexistence.}
\GSfield{Fate}{The elliptic two-state answer is given on the next page and derives from the Baxter family \citep{baxter1971eight}. The sentence does not demand a classification of all meromorphic solutions.}
\GSfield{Present status}{As of November 2026, the qualified statement above is established; no new open problem is inferred from the omissions in the notebook.}
\GSfield{Assessment}{literal statement: correct; publication match: partial; verification: publisher metadata checked; development: annotated; confidence: high}
\end{GSitem}
\begin{GSitem}{1980-516.5}{E}{Tensor-leg notation with noncommutative coefficients}
\GSfield{Kind; source}{construction; \GSpage{P000967-P000968}}
\GSfield{The notes}{\(S\) is rewritten in matrix-unit notation and the relation generalized to \(R_{12}S_{13}S^{\prime}_{23}=S^{\prime}_{23}S_{13}R_{12}\), with coefficients in an associative algebra \(D\).}
\GSfield{Statement}{Let \(V_1,V_2\) be finite-dimensional complex auxiliary spaces and \(\mathcal A\) a unital associative algebra. An exchange relation is an identity in \(\operatorname{End}(V_1\otimes V_2)\otimes\mathcal A\), with numeric \(R\) on the two auxiliary legs and the entries of the two monodromy matrices in \(\mathcal A\).}
\GSfield{Fate}{This is the algebraic transfer-matrix mechanism of quantum inverse scattering \citep{sklyaninfaddeev1978quantum}. The reconstruction distinguishes an algebra of coefficients from an additional finite-dimensional matrix factor.}
\GSfield{Present status}{As of November 2026, the qualified statement above is established; no new open problem is inferred from the omissions in the notebook.}
\GSfield{Assessment}{literal statement: correct with missing hypotheses; publication match: partial; verification: publisher metadata checked; development: annotated; confidence: high}
\end{GSitem}
\begin{GSitem}{1980-516.6}{E}{The commuting-transfer-operator lemma}
\GSfield{Kind; source}{theorem; \GSpage{P000968}}
\GSfield{The notes}{The theorem claims commutativity after writing \(\operatorname{tr}_3 S_{13}\) and \(\operatorname{tr}_3 S^{\prime}_{13}\), described as a commutator in the first matrix space.}
\GSfield{Statement}{Assume \(R\) is an invertible numeric operator and \(R_{12}T_1T_2=T_2T_1R_{12}\). Then \([\operatorname{tr}_{V_1}T_1,\operatorname{tr}_{V_2}T_2]=0\) in \(\mathcal A\). This follows by conjugating with \(R\) and taking both auxiliary traces; tracing the common quantum leg is not this theorem.}
\GSfield{Reconstruction}{The auxiliary trace legs and invertibility of R are supplied explicitly.}
\GSfield{Fate}{The standard mechanism underlies \citet{sklyaninfaddeev1978quantum}. Its elementary proof uses cyclicity only for the numeric \(R\), so it remains valid for noncommuting coefficient entries.}
\GSfield{Present status}{As of November 2026, the qualified statement above is established; no new open problem is inferred from the omissions in the notebook.}
\GSfield{Assessment}{literal statement: false as written; publication match: partial; verification: publisher metadata checked; development: annotated; confidence: high}
\end{GSitem}
\begin{GSitem}{1980-516.7}{E}{The elliptic Pauli-matrix solution}
\GSfield{Kind; source}{construction; \GSpage{P000968}}
\GSfield{The notes}{\(S(\alpha)=\sum_{j=0}^3w_j\sigma_j\otimes\sigma_j\), with \(w_0=1\), \(w_1=\operatorname{sn}(x)/\operatorname{sn}(x+\alpha)\), \(w_2=[\operatorname{dn}(x+\alpha)/\operatorname{sn}(x+\alpha)][\operatorname{sn}(x)/\operatorname{dn}(x)]\), and \(w_3=[\operatorname{cn}(x+\alpha)/\operatorname{sn}(x+\alpha)][\operatorname{sn}(x)/\operatorname{cn}(x)]\). The last denominator is overwritten in the autograph; \(\operatorname{cn}\) is the tentative reading.}
\GSfield{Statement}{The recorded two-state matrix is the elliptic eight-vertex family in a Pauli basis, defined meromorphically away from the poles of the ratios. A common elliptic modulus and a fixed crossing parameter are understood; scalar normalization and a permutation of Pauli axes do not change the exchange mechanism.}
\GSfield{Fate}{This is an explicit special-family construction associated with \citet{baxter1971eight}; the reference verifies the model and elliptic solution, not every normalization in the notebook.}
\GSfield{Present status}{As of November 2026, the elliptic family is standard. The displayed normalization is retained as a source formula, without a claim that all its meromorphic or degenerating parameter values are regular.}
\GSfield{Assessment}{literal statement: correct with missing hypotheses; publication match: partial; verification: publisher metadata checked; development: annotated; confidence: medium}
\end{GSitem}
\begin{GSitem}{1980-516.8}{E}{A parameter family yields commuting transfer matrices}
\GSfield{Kind; source}{method; \GSpage{P000968}}
\GSfield{The notes}{For a one-parameter family \(S\), an \(R(\alpha,\alpha^{\prime})\) intertwines \(S(\alpha)\) and \(S(\alpha^{\prime})\), yielding commuting trace operators.}
\GSfield{Statement}{If the invertible auxiliary exchange relation holds for every pair of regular parameters, the correctly traced transfer operators \(t(\alpha)\) satisfy \([t(\alpha),t(\alpha^{\prime})]=0\). Taylor or logarithmic coefficients commute wherever those operations are defined.}
\GSfield{Fate}{This is the family form of the preceding lemma and the construction used in \citet{sklyaninfaddeev1978quantum}. No independence or completeness of the resulting integrals follows merely from their commutativity.}
\GSfield{Present status}{As of November 2026, the qualified statement above is established; no new open problem is inferred from the omissions in the notebook.}
\GSfield{Assessment}{literal statement: correct with missing hypotheses; publication match: partial; verification: publisher metadata checked; development: annotated; confidence: high}
\end{GSitem}

\GStalk{517}{1981-03-23}{I. M. Gelfand}{P000969}{Moser’s isospectral formulation of ellipsoid geodesics}
Gelfand discusses Moser’s contribution to the Chern Symposium, giving the constrained geodesic equation and a projected matrix with constant spectrum. Jacobi’s earlier integration is mentioned; the displayed matrix construction is the modern Moser formulation.

\begin{GSitem}{1980-517.1}{E}{The constrained geodesic equation on an ellipsoid}
\GSfield{Kind; source}{construction; \GSpage{P000969}}
\GSfield{The notes}{For \(\langle A^{-1}x,x\rangle=1\), the notes give \(x_{tt}=-\nu A^{-1}x\), \(y=x_t\), and \(\nu=\langle A^{-1}y,y\rangle/|A^{-1}x|^2\).}
\GSfield{Statement}{Let \(A\) be a positive definite symmetric real matrix. For initial data satisfying \(\langle A^{-1}x,x\rangle=1\) and \(\langle A^{-1}x,y\rangle=0\), the displayed equation is the affinely parametrized geodesic equation of the induced Euclidean metric on the ellipsoid. The multiplier follows by differentiating the constraint twice.}
\GSfield{Fate}{The topic and source are identified exactly by \citet{moser1980quadrics}, a contribution to the 1979 Chern Symposium published in 1980.}
\GSfield{Present status}{As of November 2026, the qualified statement above is established; no new open problem is inferred from the omissions in the notebook.}
\GSfield{Assessment}{literal statement: correct with missing hypotheses; publication match: exact; verification: publisher metadata checked; development: annotated; confidence: high}
\end{GSitem}
\begin{GSitem}{1980-517.2}{E}{Isospectrality of Moser’s projected matrix}
\GSfield{Kind; source}{theorem; \GSpage{P000969}}
\GSfield{The notes}{With \(P_y=I-y\otimes y/|y|^2\), the matrix \(L=P_y(A-x\otimes x)P_y\) has constant eigenvalues along the geodesic. A parenthesis attributes this to Jacobi, December 1838.}
\GSfield{Statement}{For a nonconstant geodesic in the preceding ellipsoid, \(y\ne0\), and the spectrum, with multiplicities, of this symmetric matrix is independent of time. The zero eigenvalue belonging to \(y\) is included; the remaining spectral data encode the classical geodesic integrals.}
\GSfield{Reconstruction}{Positive definiteness, nonzero velocity, and the distinction between Jacobi’s integrability and Moser’s matrix formulation are explicit qualifications.}
\GSfield{Fate}{The matrix formulation is the Moser presentation \citep[pp.~147--188]{moser1980quadrics}; Jacobi is the antecedent for integration of the ellipsoid geodesic problem. The notebook alone does not establish that Jacobi used this matrix or the precise December date.}
\GSfield{Present status}{As of November 2026, the qualified statement above is established; no new open problem is inferred from the omissions in the notebook.}
\GSfield{Assessment}{literal statement: correct with missing hypotheses; publication match: partial; verification: publisher metadata checked; development: annotated; confidence: high}
\end{GSitem}

\GStalk{518}{1981-03-23}{P. G. Grinevich}{P000969-P000970}{Rational commuting differential operators of rank two}
Grinevich explains the elliptic rank-two Krichever--Novikov construction in Tyurin parameters and its rational specializations. The Dixmier case is a choice of rational parameter, not a consequence of setting the centre parameter to zero alone.

\begin{GSitem}{1980-518.1}{E}{Elliptic rank-two operators in Tyurin coordinates}
\GSfield{Kind; source}{construction; \GSpage{P000969}}
\GSfield{The notes}{The notes give \(L_6^2=L_4^3+g_1L+g_0\), and \(L_4=(D^2-u)^2+2c_x(\wp(\gamma_2)-\wp(\gamma_1))D+[c_x(\wp(\gamma_2)-\wp(\gamma_1))]_x-\wp(\gamma_1)-\wp(\gamma_2)\), where \(\gamma_1=\gamma_0+c\), \(\gamma_2=\gamma_0-c\).}
\GSfield{Statement}{The Krichever--Novikov construction associates a commuting pair of orders four and six and generic common-eigenspace dimension two to elliptic spectral data and a locally regular Tyurin deformation. The bare \(L\) in the cubic relation means \(L_4\). Elliptic discriminant nonvanishing and avoidance of poles and of \(c_x=0\) are needed for the displayed local formula.}
\GSfield{Reconstruction}{The commuting relation and nonsingularity assumptions are made explicit. The notebook’s unreconciled potential normalization is not asserted to be a verified transcription of the published formula.}
\GSfield{Fate}{The construction predates this talk; \citet[Thms.~3.2--3.3]{mironov2012higher} reproduces the Krichever--Novikov and Dixmier formulas with explicit conventions. The notebook uses \(D^2-u\) where that account uses \(D^2+u\), so the potential formula must be compared with its sign convention.}
\GSfield{Present status}{As of November 2026, the qualified statement above is established; no new open problem is inferred from the omissions in the notebook.}
\GSfield{Assessment}{literal statement: correct with missing hypotheses; publication match: partial; verification: read the result; development: annotated; confidence: medium}
\end{GSitem}
\begin{GSitem}{1980-518.2}{E}{Tyurin deformation equations and the potential formula}
\GSfield{Kind; source}{method; \GSpage{P000969-P000970}}
\GSfield{The notes}{The note records \(\gamma_{1,x}=1/(\alpha_1-\alpha_2)=-\gamma_{2,x}\), \(\alpha_{1,x}=\alpha_1^2-u-\varphi\), \(\alpha_{2,x}=\alpha_2^2-u+\varphi\), and \(u=1/(4c_x^2)-c_{xx}^2/(4c_x^2)-2\varphi c_{xx}-c_x^2(\varphi_c-\varphi^2)+c_{xxx}/(2c_x)\). Here \(\varphi=\zeta(\gamma_2-\gamma_1)+\zeta(\gamma_1)-\zeta(\gamma_2)\), \(\zeta_x=-\wp\), and \(\zeta\) is odd.}
\GSfield{Statement}{The first pair of equations forces \(\gamma_1+\gamma_2\) to be constant. These are local deformation formulas for the Tyurin data of the elliptic rank-two construction, valid away from coincident Tyurin points and singular denominators. The \(u\)-formula and the signs in the Riccati equations require a consistent choice of potential and elliptic normalization; the excerpt does not fix all of them.}
\GSfield{Fate}{Compare \citet[§3, Thm.~3.2]{mironov2012higher}. This is a substantive identification of the method, not a certification that every sign in the source formula is correct.}
\GSfield{Present status}{As of November 2026, the Tyurin deformation method and elliptic construction are established; the exact notebook potential normalization remains unresolved in this companion.}
\GSfield{Assessment}{literal statement: not determinable; publication match: partial; verification: read the result; development: annotated; confidence: medium}
\end{GSitem}
\begin{GSitem}{1980-518.3}{A}{Grinevich’s rationality criterion and the Dixmier specialization}
\GSfield{Kind; source}{theorem; \GSpage{P000970}}
\GSfield{The notes}{For rational \(q(x)\), the integral \(c(x)=\int_{q(x)}^\infty dt/\sqrt{t^3+g_1t+g_0}\) gives rational coefficients. The next line writes \(\gamma_0=0\Rightarrow q(x)=x\) and calls this the Dixmier case.}
\GSfield{Statement}{In the regular elliptic Krichever--Novikov family, rational coefficients are obtained by choosing the elliptic Abel parameter as the integral from a rational function \(q(x)\) to infinity; the conventional Weierstrass model uses \(4t^3+g_2t+g_3\). The centre \(\gamma_0\) is constant. The joint choices \(\gamma_0=0\) and \(q(x)=x\) yield the Dixmier specialization in the corresponding normalization; the first choice does not imply the second.}
\GSfield{Reconstruction}{The arrow is changed to a joint specialization, and the harmless cubic scaling is separated from the substantive rationality criterion.}
\GSfield{Fate}{\citet{grinevich1982rational} was received 24 November 1980 and published in 1982. The criterion is stated explicitly as \citet[Thm.~3.4]{mironov2012higher}; the latter supports the mathematical identification while the original publisher record verifies chronology.}
\GSfield{Present status}{As of November 2026, the qualified statement above is established; no new open problem is inferred from the omissions in the notebook.}
\GSfield{Assessment}{literal statement: false as written; publication match: exact; verification: read the result; development: annotated; confidence: high}
\end{GSitem}

\GStalk{519}{1981-03-23}{A. A. Kirillov}{P000970}{A Dixmier formula and algebraic curves in the Weyl division ring}
Kirillov gives a formal pseudodifferential description, followed by the question of rational commuting operators with a prescribed genus-two spectral curve. Gelfand proposes an embedding of the function field of any algebraic curve in the rank-one Weyl division ring. The sixth-order formula needs emendation, and the universal curve assertion remains distinct from results for generic genus-two curves or some curves of every genus.

\begin{GSitem}{1980-519.1}{E}{Kirillov’s formal pseudodifferential formula for the Dixmier pair}
\GSfield{Kind; source}{construction; \GSpage{P000970}}
\GSfield{The notes}{With \(L=D^2+x^3\), \(A=L+xL^{-1}\), the notes give \(L_4=(A^2)_+\), \(L_6=(A^3)_+\), then \(L_4=L^2+2x\), \(L_6=L^3+3x^2L+3Lx^2\).}
\GSfield{Statement}{Here positive part means the differential part of a formal pseudodifferential operator. The commuting pair is \(L_4=L^2+2x\), \(L_6=L^3+\tfrac32(xL+Lx)=L^3+3xL+3D\), with \([L_4,L_6]=0\) and \(L_6^2=L_4^3\). The two-term \(A=L+xL^{-1}\) in the notes instead gives \((A^3)_+=L^3+3xL+6D\), so this formula also requires correction. The actual formal square root of \(L_4\) begins \(L+xL^{-1}-DL^{-2}+\cdots\). Replacing \(L\) by \(D^2+x^3+a\) gives the nonsingular elliptic Dixmier cubic when \(a\ne0\).}
\GSfield{Fate}{The construction is due to \citet{dixmier1968weyl}. Both the displayed sixth-order alternative and the two-term formal-root description require correction; this chapter verifies the emendations directly. The parameter-free source specialization has the singular cubic \(w^2=z^3\), whose normalization has genus zero.}
\GSfield{Present status}{As of November 2026, the qualified statement above is established; no new open problem is inferred from the omissions in the notebook.}
\GSfield{Assessment}{literal statement: false as written; publication match: exact; verification: publisher metadata checked; development: developed in this chapter; confidence: high}
\end{GSitem}
\begin{GSitem}{1980-519.2}{B}{Rational differential operators for a genus-two spectral curve}
\GSfield{Kind; source}{question; \GSpage{P000970}}
\GSfield{The notes}{For a genus-two curve \(P(x,y)=0\), is there a pair of operators with rational coefficients satisfying \(P(L_1,L_2)=0\)?}
\GSfield{Statement}{For a smooth projective genus-two curve over \(\mathbb C\), choose a hyperelliptic model \(w^2=z^5+\beta_4z^4+\cdots+\beta_0\). There exist commuting rational-coefficient ordinary differential operators \(L,M\) whose spectral curve is this model and whose generic common eigenspace has dimension two. The coordinate map is faithful. If instead one insists on every given affine equation \(P\), including a presentation with a nonconstant unit, the analogous ordinary-operator assertion is false; the notebook does not distinguish these readings.}
\GSfield{Reconstruction}{Commutativity, faithfulness, and the complex base field are supplied. The primary theorem settles the intended abstract-curve or standard-quintic reading; the stronger literal arbitrary-affine-model reading is separated explicitly.}
\GSfield{Fate}{\citet[preprint Thm.~3(2)]{oganesyan2016rational}, first circulated 18 August 2016 and published in 2018, proves the assertion for every quintic coefficient tuple by using the union of two rational-operator families. Thus the abstract genus-two realization question is resolved. Earlier \citet[Thm.~3]{mironov2011selfadjoint} supplies explicit higher-genus families (2011 preprint, online 2013, volume 2014), while \citet[Thm.~1]{davletshinamironov2016orbits} proves the stronger infinite-orbit assertion for generic genus-two curves (2016 preprint, 2017 publication). The chapter’s affine-unit obstruction qualifies the literal prescribed-presentation reading.}
\GSfield{Present status}{As of November 2026, every smooth projective complex genus-two curve, and every standard monic quintic hyperelliptic model, is realized by rational commuting operators. This is an established result. Universal realization of every affine presentation is false, as the chapter demonstrates using \(zw^2-z^4-1=0\); this distinction is not promoted to an artificial residual open problem.}
\GSfield{Assessment}{literal statement: correct with missing hypotheses; publication match: exact; verification: read the result; development: developed in this chapter; confidence: high; problem \hyperref[prob:GP-1980-40]{GP-1980-40}}
\end{GSitem}
\begin{GSitem}{1980-519.3}{C}{Curve function fields in the Weyl division ring}
\GSfield{Kind; source}{conjecture; \GSpage{P000970}; I. M. Gelfand}
\GSfield{The notes}{Gelfand proposes that any curve embeds in the division ring of fractions of the Weyl algebra. A marginal gloss specifies the field of rational functions on the curve, and an inserted word changes a possible theorem heading to conjecture.}
\GSfield{Statement}{For every smooth projective integral curve \(C/\mathbb C\), does there exist a unital \(\mathbb C\)-embedding \(\mathbb C(C)\hookrightarrow\operatorname{Frac}(A_1(\mathbb C))\), where \(A_1=\mathbb C\langle x,D\rangle/(Dx-xD-1)\)? Singular integral curves give the same question after normalization.}
\GSfield{Fate}{The Dixmier and Grinevich constructions give genus one; \citet[preprint Thm.~3(2)]{oganesyan2016rational} supplies all complex genus-two curves (2016 preprint, 2018 publication). \citet[Thm.~3]{mironov2011selfadjoint} supplies some hyperelliptic curves of every genus, while \citet[Thm.~1]{davletshinamironov2016orbits} proves an infinite-orbit refinement for generic genus-two curves. These positive subcases do not establish the universal statement for all curves. The chapter proves the field embedding for its displayed family and separates it from affine-coordinate ordinary-operator realizations. The earlier 1977-434.6 is a thematic precursor, without an exact recurring problem ID.}
\GSfield{Present status}{As of November 2026, all complex curves of genus at most two and explicit higher-genus families have positive realizations. No source proving or disproving the universal curve-function-field conjecture was located in the documented search. This is a search limitation, not a claim that no such theorem exists.}
\GSfield{Assessment}{literal statement: not determinable; publication match: partial; verification: read the result; development: developed in this chapter; confidence: medium; problem \hyperref[prob:GP-1980-41]{GP-1980-41}}
\GSfield{Worth developing}{Worth developing because the division ring allows inverses and negative differential orders, and the universal birational question is strictly broader than the known explicit spectral families.}
\end{GSitem}
\begin{GSitem}{1980-519.4}{B}{The remark excluding coprime differential orders}
\GSfield{Kind; source}{claim; \GSpage{P000970}}
\GSfield{The notes}{The sentence following the genus-two question says that there are no coprime orders.}
\GSfield{Statement}{The verified theorem for polynomial coefficients states the following. Over a field of characteristic zero, if the centralizer of a nonconstant element of the first Weyl algebra has differential-order rank one, it embeds in a polynomial ring in one variable and has genus zero. Thus polynomial commuting operators of coprime positive orders cannot faithfully generate a positive-genus spectral field. The sentence in context concerns rational coefficients, for which this cited theorem alone is insufficient.}
\GSfield{Fate}{\citet[Thm.~1.1]{makarlimanov2021centralizers}, received January and published May 2021, proves the polynomial centralizer statement. Rank is the gcd of orders of all nonconstant elements of the commutative centralizer; it is not generally inferred from an arbitrary displayed pair.}
\GSfield{Present status}{As of November 2026, the polynomial-coefficient version is established. The unrestricted rational-coefficient wording in this notebook is not certified by the source checked here.}
\GSfield{Assessment}{literal statement: correct with missing hypotheses; publication match: partial; verification: read the result; development: annotated; confidence: medium}
\end{GSitem}
\begin{GSitem}{1980-519.5}{E}{Existence of the Weyl skew field of fractions}
\GSfield{Kind; source}{problem; \GSpage{P000970}}
\GSfield{The notes}{The algebra generated by \(q,p\) with \(qp-pq=1\) embeds in a division ring, offered as an exercise.}
\GSfield{Statement}{The first Weyl algebra over a characteristic-zero field is a noetherian domain and satisfies the left and right Ore conditions for its nonzero elements. Its classical ring of fractions is consequently a division ring containing it injectively.}
\GSfield{Fate}{This is standard background to \citet{dixmier1968weyl}; \citet[§2]{makarlimanov2021centralizers} explicitly recalls the Ore embedding. The order of the generators only determines the sign convention for the defining commutator.}
\GSfield{Present status}{As of November 2026, the qualified statement above is established; no new open problem is inferred from the omissions in the notebook.}
\GSfield{Assessment}{literal statement: correct; publication match: exact; verification: read the result; development: annotated; confidence: high}
\end{GSitem}

\GStalk{520}{1981-03-23}{J.-L. Brylinski}{P000970}{Verma modules and geometric realization}
The uncertain handwritten surname is resolved by the explicit Kashiwara collaboration and the subject. The short record defines Verma modules and gestures toward realizations on the flag variety, with an SL(2) example. It does not state a new proof of the Kazhdan--Lusztig formula.

\begin{GSitem}{1980-520.1}{E}{Verma modules}
\GSfield{Kind; source}{definition; \GSpage{P000970}}
\GSfield{The notes}{For a complex semisimple \(\mathfrak g\), a Borel \(\mathfrak b\), and a character \(\lambda:\mathfrak b\to\mathbb C\), the notes introduce \(M_\lambda\).}
\GSfield{Statement}{The Verma module in the stated induction convention is \(M_\lambda=U(\mathfrak g)\otimes_{U(\mathfrak b)}\mathbb C_\lambda\), with the nilpotent radical of \(\mathfrak b\) acting trivially on \(\mathbb C_\lambda\). Shifts by \(\rho\), opposite Borels, and dual modules belong to alternative conventions and must be specified in a geometric realization.}
\GSfield{Fate}{The definition is established representation theory. \citet[§1]{brylinski1981differential}, presented at Toruń in August--September 1980 and published in 1981, gives the relevant geometric setting.}
\GSfield{Present status}{As of November 2026, the qualified statement above is established; no new open problem is inferred from the omissions in the notebook.}
\GSfield{Assessment}{literal statement: correct; publication match: exact; verification: read the result; development: annotated; confidence: high}
\end{GSitem}
\begin{GSitem}{1980-520.2}{A}{Flag-variety and projective-line realizations}
\GSfield{Kind; source}{construction; \GSpage{P000970}}
\GSfield{The notes}{The module is represented by geometric objects on the flag space. For \(SL(2)\), the sketch labels holomorphic functions on \(\mathbb P^1\setminus\{\infty\}\) meromorphic near infinity, and germs of meromorphic functions near infinity.}
\GSfield{Statement}{Flag-variety realizations use twisted differential operators and appropriate standard, costandard, or local-cohomology modules. For \(SL_2\), global functions on the affine line that are meromorphic at infinity are polynomials; local principal parts arise from a quotient of meromorphic germs by holomorphic germs. A weight-dependent action and the relevant duality convention are necessary to identify either space with a specified Verma module.}
\GSfield{Reconstruction}{The twist, the principal-parts quotient, and the warning about dual Verma modules supply data missing from the diagram.}
\GSfield{Fate}{\citet[§§1--2]{brylinski1981differential} supplies the geometric framework. The note’s mention of the Kashiwara collaboration identifies the speaker; it does not state the Kazhdan--Lusztig theorem anew, so GP-1979-41 is not assigned to this abbreviated construction.}
\GSfield{Present status}{As of November 2026, the qualified statement above is established; no new open problem is inferred from the omissions in the notebook.}
\GSfield{Assessment}{literal statement: correct with missing hypotheses; publication match: partial; verification: read the result; development: annotated; confidence: medium}
\end{GSitem}

\GStalk{521}{1981-03-23}{E. Viehweg}{P000971}{Classification of threefolds after Esnault’s Bourbaki report}
The talk explicitly follows Hélène Esnault’s February 1981 Bourbaki exposition. Kodaira dimension, irregularity, abelian varieties, and Albanese fibrations of threefolds are summarized. A question mark in the irregularity-zero case is preserved as an incomplete classification entry, not promoted to a newly formulated conjecture.

\begin{GSitem}{1980-521.1}{E}{Kodaira dimension and irregularity}
\GSfield{Kind; source}{definition; \GSpage{P000971}}
\GSfield{The notes}{Viehweg defines \(\kappa\) using the graded algebra \(\bigoplus_{m\ge0}H^0(X,\Omega_X^{n\otimes m})\) and \(q=\dim H^0(X,\Omega_X^1)\).}
\GSfield{Statement}{For a smooth connected projective complex \(n\)-fold, let \(R(X,K_X)=\bigoplus_{m\ge0}H^0(X,K_X^{\otimes m})\). If every positive-degree piece vanishes, \(\kappa(X)=-\infty\); otherwise \(\kappa(X)=\operatorname{trdeg}_{\mathbb C}R(X,K_X)-1\). Irregularity is \(q(X)=h^0(X,\Omega_X^1)=h^1(X,\mathcal O_X)\).}
\GSfield{Reconstruction}{Smoothness, projectivity, the zero-ring exception, and transcendence degree rather than an unspecified degree are supplied.}
\GSfield{Fate}{These conventions are given at the start of \citet[pp.~111--112]{esnault1981classification}, Bourbaki exposé 568, February 1981, explicitly named by the lecture.}
\GSfield{Present status}{As of November 2026, the qualified statement above is established; no new open problem is inferred from the omissions in the notebook.}
\GSfield{Assessment}{literal statement: correct with missing hypotheses; publication match: exact; verification: read the result; development: annotated; confidence: high}
\end{GSitem}
\begin{GSitem}{1980-521.2}{E}{Maximal irregularity at Kodaira dimension zero}
\GSfield{Kind; source}{theorem; \GSpage{P000971}}
\GSfield{The notes}{The table relates \(\kappa=0\), \(q=n\), and being birationally equivalent to an abelian variety; the direction of an arrow is unclear.}
\GSfield{Statement}{If \(X\) is smooth projective over \(\mathbb C\), \(\kappa(X)=0\), and \(q(X)=\dim X\), then its Albanese map is birational onto an abelian variety. Conversely a variety birational to an abelian variety has these invariants.}
\GSfield{Reconstruction}{The complete equivalence is the modern reconstruction; the manuscript arrow itself remains a palaeographic ambiguity.}
\GSfield{Fate}{\citet[Thm.~1 and the adjacent table]{esnault1981classification} reports the Kawamata result before the seminar date. This is an exact subject match, although the notebook arrow is not silently read as an equivalence.}
\GSfield{Present status}{As of November 2026, the qualified statement above is established; no new open problem is inferred from the omissions in the notebook.}
\GSfield{Assessment}{literal statement: correct with missing hypotheses; publication match: exact; verification: read the result; development: annotated; confidence: high}
\end{GSitem}
\begin{GSitem}{1980-521.3}{E}{Irregular threefolds of Kodaira dimension zero}
\GSfield{Kind; source}{theorem; \GSpage{P000971}}
\GSfield{The notes}{For \(n=3\), \(\kappa=0\), and \(q=2\) or \(1\), the notes say there is an étale fibration and the fibre has \(\kappa=0\).}
\GSfield{Statement}{For the smooth projective complex threefolds in the cited classification, the Albanese description becomes a product after a suitable finite étale base change, up to the birational modifications allowed in the theorem. The general fibre has dimension \(3-q\) and Kodaira dimension zero. Calling the map itself étale is impossible when its fibres have positive dimension.}
\GSfield{Reconstruction}{Étale refers to the base change and local trivialization, not to a positive-dimensional-fibre morphism.}
\GSfield{Fate}{\citet[Thm.~2 and table, pp.~112--114]{esnault1981classification} is the explicitly identified source. The note compresses birational and finite-cover qualifications into the phrase about étaleness.}
\GSfield{Present status}{As of November 2026, the qualified statement above is established; no new open problem is inferred from the omissions in the notebook.}
\GSfield{Assessment}{literal statement: correct with missing hypotheses; publication match: partial; verification: read the result; development: annotated; confidence: high}
\end{GSitem}
\begin{GSitem}{1980-521.4}{F}{The irregularity-zero entry in the classification table}
\GSfield{Kind; source}{question; \GSpage{P000971}}
\GSfield{The notes}{The line \(\kappa=0,q=0\) ends in a question mark.}
\GSfield{Statement}{The source does not specify whether it asks for a birational classification, a list of deformation types, the geometry of the Albanese map, or another invariant of threefolds with \(\kappa=q=0\). These are inequivalent problems.}
\GSfield{Fate}{The contemporary limitations and cases are surveyed by \citet[§1]{esnault1981classification}. That report identifies the context but cannot turn the isolated mark into one precise conjecture.}
\GSfield{Present status}{As of November 2026, the mark remains an incomplete classification entry in this record; no new globally indexed open problem is inferred.}
\GSfield{Assessment}{literal statement: not determinable; publication match: thematic only; verification: read the result; development: annotated; confidence: high}
\end{GSitem}
\begin{GSitem}{1980-521.5}{F}{The negative-infinite Kodaira-dimension line}
\GSfield{Kind; source}{remark; \GSpage{P000971}}
\GSfield{The notes}{The \(\kappa=-\infty\) line records only an indefinite description.}
\GSfield{Statement}{No definite classification assertion, dimension restriction beyond the preceding discussion, or birational conclusion can be recovered from this fragment.}
\GSfield{Fate}{The surrounding report is \citet{esnault1981classification}; its results cannot be substituted for words that are absent in the notebook.}
\GSfield{Present status}{As of November 2026, this fragment has no determinate mathematical claim to certify.}
\GSfield{Assessment}{literal statement: not determinable; publication match: thematic only; verification: read the result; development: annotated; confidence: high}
\end{GSitem}

\GStalk{522}{1981-03-23}{C. Kassel}{P000971-P000972}{K-theory of spaces and arithmetic K-groups}
The manuscript has separate headings for Christian Kassel and Christophe Soulé, although the catalogue supplies only the Kassel record. Kassel reports Waldhausen’s programme and homotopy-ring constructions. Soulé’s contribution concerns orders of K-groups of the integers and zeta values; its bibliographic shorthand is not an exact identification of an Inventiones paper dated February 1980.

\begin{GSitem}{1980-522.1}{E}{Waldhausen’s K-theory programme for diffeomorphism spaces}
\GSfield{Kind; source}{method; \GSpage{P000971-P000972}}
\GSfield{The notes}{Kassel relates the homotopy type of a diffeomorphism group to the algebraic K-theory of a loop-space ring, invoking Waldhausen and homotopy rings.}
\GSfield{Statement}{For a connected pointed CW space \(X\), algebraic K-theory of spaces, \(A(X)\), is a homotopy-theoretic construction involving retractive spaces over \(X\), or equivalently appropriate modules over the suspension ring spectrum of \(\Omega X\). Its relation to stable pseudoisotopy and stable smooth h-cobordisms is not an identity between an arbitrary diffeomorphism group and the ordinary K-theory of \(\mathbb Z[\pi_1X]\).}
\GSfield{Reconstruction}{The distinction between A-theory and ordinary group-ring K-theory is explicit; no unstated stable-range theorem is supplied.}
\GSfield{Fate}{\citet[pp.~35--60]{waldhausen1978topological} appeared in the proceedings of the 1976 Stanford symposium, published 1978. \citet[Introduction and §1]{kassel1983stabilisation} develops the homotopy-ring and stabilization context, received after this lecture and published in 1983. No exact diffeomorphism-space equivalence is stated in this notebook.}
\GSfield{Present status}{As of November 2026, the qualified statement above is established; no new open problem is inferred from the omissions in the notebook.}
\GSfield{Assessment}{literal statement: correct with missing hypotheses; publication match: partial; verification: read the result; development: annotated; confidence: high}
\end{GSitem}
\begin{GSitem}{1980-522.2}{E}{Quillen’s plus construction and homotopy rings}
\GSfield{Kind; source}{construction; \GSpage{P000972}}
\GSfield{The notes}{The notes compare Quillen’s plus construction for rings with a homotopy-ring construction and mention maps between the corresponding classifying spaces.}
\GSfield{Statement}{For a unital discrete ring \(R\), the positive K-groups are \(K_i(R)=\pi_i(BGL(R)^+)\), \(i\ge1\), with the plus construction killing the stable elementary subgroup. Passing to a ring structure defined only up to homotopy requires coherent multiplication and a compatible homotopical K-theory construction; a bare nonassociative operation up to an unspecified homotopy does not suffice.}
\GSfield{Fate}{The ordinary construction is Quillen’s earlier theory; \citet[§1]{kassel1983stabilisation} provides the relevant homotopy-theoretic refinement. Its 1983 publication is a partial identification of the discussion, not a proof that every diagram in the shorthand is the same theorem.}
\GSfield{Present status}{As of November 2026, the qualified statement above is established; no new open problem is inferred from the omissions in the notebook.}
\GSfield{Assessment}{literal statement: correct with missing hypotheses; publication match: partial; verification: read the result; development: annotated; confidence: high}
\end{GSitem}
\begin{GSitem}{1980-522.3}{B}{Soulé’s zeta-value bounds for the orders of K-groups}
\GSfield{Kind; source}{claim; \GSpage{P000972}; C. Soulé}
\GSfield{The notes}{Under a separate Soulé heading, for even \(i\ge2\) the notes bound \(|K_{2i-2}(\mathbb Z)|\) and \(|K_{2i-1}(\mathbb Z)|\) below by the numerator and denominator of \(\zeta(1-i)\). They cite Inventiones, February 1980, and a method involving topoi.}
\GSfield{Statement}{Write \(|\zeta(1-i)|=a_i/b_i\) in lowest terms, \(a_i,b_i>0\), with even \(i\ge2\). The modern finite-order formulas imply \(|K_{2i-2}(\mathbb Z)|\ge a_i\) and \(|K_{2i-1}(\mathbb Z)|\ge b_i\). Absolute numerator is necessary when the zeta value is negative. Étale Chern classes and arithmetic cohomology are part of the source context; the cryptic date does not identify an exact theorem by itself.}
\GSfield{Reconstruction}{The sign convention for numerator and the separation between present validity and historical attribution are supplied.}
\GSfield{Fate}{Soulé’s arithmetic paper was published in October 1979, not February 1980 \citep{soule1979k}. His 1978 thesis gives related lower bounds with prime restrictions \citep[Introduction]{soule1978thesis}. Complete formulas, including the factors at 2, are recorded in \citet[Cor.~95 and Table~5.1]{weibel2005rings}. The cited later formulas verify the inequalities but do not establish that their full strength was proved in the unidentified February reference.}
\GSfield{Present status}{As of November 2026, these numerical lower bounds follow from established arithmetic K-theory formulas. Exact attribution of the full displayed bounds to the manuscript’s February 1980 reference remains unresolved.}
\GSfield{Assessment}{literal statement: correct with missing hypotheses; publication match: partial; verification: read the result; development: annotated; confidence: medium}
\end{GSitem}

\GStalk{523}{1981-03-30}{I. M. Gelfand}{P000973}{The collisionless Vlasov description}
A brief introduction explains a self-consistent field generated by many charged particles whose binary collisions are neglected. The record states a modelling principle, not a quantified mean-field limit theorem.

\begin{GSitem}{1980-523.1}{E}{The collisionless self-consistent-field model}
\GSfield{Kind; source}{method; \GSpage{P000973}}
\GSfield{The notes}{A large collection of charged particles creates a potential determined by its density; direct binary collisions are neglected.}
\GSfield{Statement}{A representative electrostatic formulation is \(f_t+v\cdot\nabla_xf-\nabla_x\phi\cdot\nabla_vf=0\), with a Poisson equation relating \(\phi\) to \(\rho(x)=\int f(x,v)\,dv\), including a background and boundary conditions when appropriate. This is a collisionless mean-field model; the notebook does not state a scaling limit from a finite particle system.}
\GSfield{Reconstruction}{The displayed PDE is supplied to make the modelling description precise; electrostatic force and boundary conventions are not fixed by the notes.}
\GSfield{Fate}{The model goes back to Vlasov’s 1938 work, reprinted in \citet{vlasov1968vibrational}. The opened publisher page verifies that identification and the topic, not a modern derivation theorem.}
\GSfield{Present status}{As of November 2026, the qualified statement above is established; no new open problem is inferred from the omissions in the notebook.}
\GSfield{Assessment}{literal statement: correct; publication match: exact; verification: publisher metadata checked; development: annotated; confidence: high}
\end{GSitem}

\GStalk{524}{1981-03-30}{V. I. Arnold}{P000973-P000974}{Weinstein’s Hamiltonian formulations of continua}
Arnold begins with the geodesic interpretation of incompressible Euler flow, then describes reduction to coadjoint orbits and the extension to Vlasov and gauge-coupled kinetic equations. The density variable required for compressible fluid mechanics is mentioned but no full bracket is written.

\begin{GSitem}{1980-524.1}{E}{Incompressible Euler flow as geodesic motion}
\GSfield{Kind; source}{theorem; \GSpage{P000973}}
\GSfield{The notes}{The quadratic form \(\frac12\int|v|^2\) on divergence-free vector fields is extended by left translations to volume-preserving diffeomorphisms; its geodesics are ideal-fluid flows.}
\GSfield{Statement}{For a compact oriented Riemannian manifold without boundary, or with the usual impermeable boundary condition, equip volume-preserving diffeomorphisms with the right-invariant weak \(L^2\) metric in the usual material-to-spatial convention. Its smooth geodesics have Eulerian velocity satisfying \(u_t+\nabla_u u=-\nabla p\), \(\operatorname{div}u=0\). An inverse-map convention interchanges left and right.}
\GSfield{Reconstruction}{The manifold, boundary condition, smooth lifespan, and right-translation convention are supplied.}
\GSfield{Fate}{\citet[§§1--3]{arnold1966geometrie} supplies the formulation cited explicitly in the notes. The translation side must be reconciled with the velocity convention.}
\GSfield{Present status}{As of November 2026, the qualified statement above is established; no new open problem is inferred from the omissions in the notebook.}
\GSfield{Assessment}{literal statement: correct with missing hypotheses; publication match: exact; verification: read the result; development: annotated; confidence: high}
\end{GSitem}
\begin{GSitem}{1980-524.2}{E}{Hamiltonian reduction and coadjoint-orbit symplectic forms}
\GSfield{Kind; source}{method; \GSpage{P000973-P000974}}
\GSfield{The notes}{The invariant kinetic energy on \(T^*G\) projects to \(\mathfrak g^*\), whose coadjoint orbits carry the Kirillov form.}
\GSfield{Statement}{For a finite-dimensional Lie group, an invariant Hamiltonian on its cotangent bundle reduces under the lifted translation action to a Lie--Poisson Hamiltonian on \(\mathfrak g^*\). Its symplectic leaves are coadjoint orbits with the Kirillov--Kostant--Souriau form. Infinite-dimensional fluid versions require a chosen regular dual, functional setting, and smoothness domain.}
\GSfield{Reconstruction}{The note’s phrase that the Kirillov form descends is interpreted through reduction; it is not a globally nondegenerate form on all of the dual.}
\GSfield{Fate}{The Euler-group reduction is part of \citet{arnold1966geometrie}. The Vlasov application is developed by \citet{marsdenweinstein1982hamiltonian}, received August 1981 and published March 1982.}
\GSfield{Present status}{As of November 2026, the qualified statement above is established; no new open problem is inferred from the omissions in the notebook.}
\GSfield{Assessment}{literal statement: correct with missing hypotheses; publication match: partial; verification: read the result; development: annotated; confidence: high}
\end{GSitem}
\begin{GSitem}{1980-524.3}{E}{The Vlasov Lie–Poisson bracket}
\GSfield{Kind; source}{construction; \GSpage{P000973}}
\GSfield{The notes}{The group of canonical transformations is said to produce the Vlasov equation, with the parenthetical phrase cold plasma.}
\GSfield{Statement}{On canonical phase space \(T^*\mathbb R^d\), identified with \((x,v)\) in unit-mass coordinates, the Lie--Poisson bracket is \(\{F,G\}(f)=\int f\{\delta F/\delta f,\delta G/\delta f\}_{x,v}\,dx\,dv\), with a consistent overall sign convention. A rigorous algebraic core consists of smooth compactly supported densities and cylindrical polynomial functionals built from pairings with compactly supported smooth test functions. For more general functionals, smooth variational derivatives and decay sufficient for all integrals and integrations by parts are assumed. The kinetic-plus-self-field Hamiltonian formally yields Vlasov evolution within its smooth domain; this construction alone is not an analytic existence theorem and does not require a monokinetic ansatz.}
\GSfield{Fate}{The pure Vlasov functional bracket is already explicit in \citet[p.~385, Eq.~(12)]{morrison1980maxwell}, received August and published December 1980, before this lecture. In canonical momentum variables or unit-mass normalization it is the displayed formula. \citet[pp.~235--236]{weinsteinmorrison1981comments} corrects the mass factors and the separate magnetic coupling, whose original form failed Jacobi. \citet{marsdenweinstein1982hamiltonian} subsequently gives the geometric coadjoint/gauge construction; that later implementation is not the first occurrence of the pure bracket.}
\GSfield{Present status}{As of November 2026, the qualified statement above is established; no new open problem is inferred from the omissions in the notebook.}
\GSfield{Assessment}{literal statement: correct with missing hypotheses; publication match: exact; verification: read the result; development: annotated; confidence: high}
\end{GSitem}
\begin{GSitem}{1980-524.4}{B}{Gauge-coupled kinetic systems}
\GSfield{Kind; source}{construction; \GSpage{P000973}}
\GSfield{The notes}{Adding local degrees of freedom and a gauge group is said to lead to further equations, including Yang--Mills.}
\GSfield{Statement}{The identifiable extension couples a distribution on particle phase space with internal gauge charge to a gauge field, using a reduced Hamiltonian bracket and constraints. The note does not specify a gauge group, charge orbit, field variables, or an exact Yang--Mills--Vlasov bracket.}
\GSfield{Fate}{\citet[abstract]{marsdenweinstein1982hamiltonian} explicitly states applicability to the Yang--Mills--Vlasov system. This verifies the direction of development, not a formula absent from the manuscript.}
\GSfield{Present status}{As of November 2026, the qualified statement above is established; no new open problem is inferred from the omissions in the notebook.}
\GSfield{Assessment}{literal statement: correct with missing hypotheses; publication match: partial; verification: publisher metadata checked; development: annotated; confidence: medium}
\end{GSitem}
\begin{GSitem}{1980-524.5}{F}{Density and the incomplete compressible-fluid Hamiltonian claim}
\GSfield{Kind; source}{remark; \GSpage{P000974}}
\GSfield{The notes}{The note says vorticity must be specified before a Hamiltonian formalism exists, then refers to bundles over the group carrying density.}
\GSfield{Statement}{Compressible-fluid Hamiltonian descriptions require the advected mass density in addition to momentum; a semidirect-product or equivalent cotangent construction provides this. Fixing vorticity is not a general prerequisite for a Hamiltonian formulation. No particular bracket or bundle is specified here.}
\GSfield{Fate}{The contextual fluid geometry is \citet{arnold1966geometrie}; the nearby compressibility work is discussed more explicitly in record 527. The incomplete vorticity phrase is not identified with a published theorem.}
\GSfield{Present status}{As of November 2026, Hamiltonian descriptions of compressible flow are established, but the notebook’s claimed prerequisite is neither necessary in general nor sufficiently specified for a narrower interpretation.}
\GSfield{Assessment}{literal statement: false as written; publication match: thematic only; verification: read the result; development: annotated; confidence: medium}
\end{GSitem}

\GStalk{525}{1981-03-30}{A. V. Bobylev}{P000974-P000975}{Maxwell molecules and Fourier methods for Boltzmann’s equation}
Bobylev presents the spatially homogeneous collision equation and singles out inverse-fourth-power potentials. Fourier methods, Gaussian convolution, linearization, and the isotropic nonlinear hierarchy appear, together with remarks on exact solutions and hydrodynamic limits. Angular cutoff, function-space, and normalization assumptions must be supplied when interpreting the compact-operator claim.

\begin{GSitem}{1980-525.1}{E}{Elastic collisions and the homogeneous Boltzmann equation}
\GSfield{Kind; source}{definition; \GSpage{P000974}}
\GSfield{The notes}{The collision integral is quadratic in \(f\), with \(v^{\prime},w^{\prime}=(v+w\pm|v-w|n)/2\), \(|n|=1\), and kernel \(g=|v-w|\,d\sigma/d\Omega\).}
\GSfield{Statement}{For equal-mass elastic collisions in three dimensions, the displayed postcollision velocities preserve total momentum and kinetic energy. The spatially homogeneous Boltzmann equation is \(\partial_tf=Q(f,f)\), where the angular and relative-speed kernel is nonnegative and the integral is interpreted under the integrability or cancellation hypotheses appropriate to it. Molecular chaos is the modelling assumption behind the product of one-particle densities.}
\GSfield{Fate}{\citet[§1]{bobylev1984exact} gives the collision equation and conventions. This is standard background, not a theorem identifying a deterministic hard-sphere limit for arbitrary times.}
\GSfield{Present status}{As of November 2026, the qualified statement above is established; no new open problem is inferred from the omissions in the notebook.}
\GSfield{Assessment}{literal statement: correct with missing hypotheses; publication match: exact; verification: read the result; development: annotated; confidence: high}
\end{GSitem}
\begin{GSitem}{1980-525.2}{E}{Inverse-power potentials and Maxwell molecules}
\GSfield{Kind; source}{claim; \GSpage{P000974}}
\GSfield{The notes}{For \(V(r)=\alpha r^{-s}\), the collision kernel is \(\widetilde g(\mu)|v-w|^{1-4/s}\). The case \(s=4\) depends only on the angle.}
\GSfield{Statement}{For repulsive inverse-power pair potentials, homogeneity yields the displayed relative-speed exponent in the usual classical scattering convention in three dimensions. The inverse-fourth-power potential is the Maxwell-molecule case. This speed independence does not imply that the angular kernel is integrable; an angular cutoff is a separate assumption.}
\GSfield{Fate}{\citet[§1]{bobylev1984exact} distinguishes true Maxwell molecules from their cutoff or pseudo-Maxwell models. The distinction matters to the later compactness and normalization statements.}
\GSfield{Present status}{As of November 2026, the qualified statement above is established; no new open problem is inferred from the omissions in the notebook.}
\GSfield{Assessment}{literal statement: correct with missing hypotheses; publication match: exact; verification: read the result; development: annotated; confidence: high}
\end{GSitem}
\begin{GSitem}{1980-525.3}{E}{The Bobylev Fourier identity}
\GSfield{Kind; source}{method; \GSpage{P000974-P000975}}
\GSfield{The notes}{When the kernel depends only on angle, Fourier transformation greatly simplifies the equation and gives exact solutions.}
\GSfield{Statement}{For a homogeneous Maxwell model with integrable angular kernel \(b\), let \(\phi(k)=\int e^{-ik\cdot v}f(v)\,dv\). Then \(\partial_t\phi(k)=\int_{S^2}b(\widehat k\cdot\sigma)[\phi(k^+)\phi(k^-)-\phi(k)\phi(0)]\,d\sigma\), where \(k^\pm=(k\pm|k|\sigma)/2\). Noncutoff versions require the corresponding cancellation interpretation.}
\GSfield{Reconstruction}{The Fourier variables and angular-cutoff hypothesis make the shorthand precise.}
\GSfield{Fate}{Bobylev’s Fourier method and exact-solution papers appeared in 1975, before this seminar. \citet[§1 and references~4--8]{bobylev1984exact} reproduces the identity and its chronology.}
\GSfield{Present status}{As of November 2026, the qualified statement above is established; no new open problem is inferred from the omissions in the notebook.}
\GSfield{Assessment}{literal statement: correct with missing hypotheses; publication match: exact; verification: read the result; development: annotated; confidence: high}
\end{GSitem}
\begin{GSitem}{1980-525.4}{E}{Gaussian convolution symmetry}
\GSfield{Kind; source}{theorem; \GSpage{P000975}}
\GSfield{The notes}{Convolving with a Gaussian density does not change the equation.}
\GSfield{Statement}{If \(f(t)\) solves the homogeneous Maxwell-molecule equation and \(G_a\) is a fixed centred normalized isotropic Gaussian, then \(G_a*f(t)\) solves the same equation, whenever both are in the relevant solution class. In Fourier variables this follows from \(|k^+|^2+|k^-|^2=|k|^2\). The variance is constant in time.}
\GSfield{Fate}{This symmetry belongs to Bobylev’s Maxwell Fourier method \citep[§1]{bobylev1984exact}. The statement is not a symmetry of an arbitrary velocity-dependent collision kernel or of an arbitrary time-dependent convolution.}
\GSfield{Present status}{As of November 2026, the qualified statement above is established; no new open problem is inferred from the omissions in the notebook.}
\GSfield{Assessment}{literal statement: correct with missing hypotheses; publication match: exact; verification: read the result; development: annotated; confidence: high}
\end{GSitem}
\begin{GSitem}{1980-525.5}{E}{Maxwellians and linearization}
\GSfield{Kind; source}{construction; \GSpage{P000975}}
\GSfield{The notes}{The stationary family is \(f_M=Ae^{-\alpha|v-v_0|^2}\), called three-parameter for the blocks \(A,\alpha,v_0\). Substitution \(f=f_M(1+F)\) gives the displayed linear collision operator.}
\GSfield{Statement}{For \(A,\alpha>0\) and \(v_0\in\mathbb R^3\), these equilibria have five scalar parameters. Linearizing gives \(\mathcal LF=f_M^{-1}[Q(f_MF,f_M)+Q(f_M,f_MF)]\). With detailed balance and the appropriate domain, \(\mathcal L\) is nonpositive and self-adjoint in \(L^2(f_M\,dv)\); its nullspace is spanned by \(1,v_1,v_2,v_3,|v|^2\) for a nondegenerate collision kernel.}
\GSfield{Reconstruction}{The three named blocks contain five real scalar parameters; this repairs a counting ambiguity rather than introducing a new equilibrium family.}
\GSfield{Fate}{\citet[§§1,3]{bobylev1984exact} treats the Maxwellian and its spectral linearization. The notebook’s integral has specialized the equilibrium scale, mean, and normalization.}
\GSfield{Present status}{As of November 2026, the qualified statement above is established; no new open problem is inferred from the omissions in the notebook.}
\GSfield{Assessment}{literal statement: correct with missing hypotheses; publication match: exact; verification: read the result; development: annotated; confidence: high}
\end{GSitem}
\begin{GSitem}{1980-525.6}{E}{The constant-loss plus compact-gain decomposition}
\GSfield{Kind; source}{claim; \GSpage{P000975}}
\GSfield{The notes}{The linearized operator is \(CI+T\), with compact \(T\) and \(C=-1\).}
\GSfield{Statement}{For an angular-cutoff Maxwell model in the Maxwellian weighted Hilbert space, under the standard integrability assumptions giving compact gain, the linearized collision operator has the form \(-\nu I+K\), with \(K\) compact and constant collision frequency \(\nu>0\). Rescaling time so that \(\nu=1\) yields the notebook formula. Without cutoff the operator is generally unbounded, and for other potentials the loss frequency depends on velocity.}
\GSfield{Fate}{The spectral setting is discussed in \citet[§3]{bobylev1984exact}. The phrase \(C=-1\) is a normalization, not an intrinsic constant of every inverse-fourth-power gas.}
\GSfield{Present status}{As of November 2026, the qualified statement above is established; no new open problem is inferred from the omissions in the notebook.}
\GSfield{Assessment}{literal statement: correct with missing hypotheses; publication match: partial; verification: read the result; development: annotated; confidence: high}
\end{GSitem}
\begin{GSitem}{1980-525.7}{F}{Unspecified new eigenfunctions}
\GSfield{Kind; source}{claim; \GSpage{P000975}}
\GSfield{The notes}{The notes say new eigenfunctions of the linearized equation were obtained, without listing them.}
\GSfield{Statement}{No eigenvalue, angular sector, orthogonality convention, or new eigenfunction is given, so this novelty claim cannot be matched to a particular formula independently of the general Maxwell spectral discussion.}
\GSfield{Fate}{\citet[§3]{bobylev1984exact} supplies the detailed spectral framework, but thematic overlap does not identify which functions this sentence calls new.}
\GSfield{Present status}{As of November 2026, the general spectral theory is established; the particular novelty claimed by this sentence remains unidentified.}
\GSfield{Assessment}{literal statement: not determinable; publication match: thematic only; verification: read the result; development: annotated; confidence: medium}
\end{GSitem}
\begin{GSitem}{1980-525.8}{A}{Triangular linearization of the isotropic Maxwell hierarchy}
\GSfield{Kind; source}{theorem; \GSpage{P000975}}
\GSfield{The notes}{For isotropic \(f(|v|^2)\) and \(s=4\), an infinite system of ordinary differential equations can be transformed into a linear system. The surrounding text mentions general solutions among rapidly decaying functions and says other cases are unknown.}
\GSfield{Statement}{In the normalized isotropic cutoff Maxwell model, the Fourier or moment coefficients satisfy a triangular nonlinear hierarchy. On Bobylev’s specified analytic coefficient class, a triangular polynomial change of coefficients conjugates it to \(\dot y_n=-\lambda_n y_n\), \(n\ge2\), with \(\lambda_n=\int_0^1p(s)[1-s^n-(1-s)^n]\,ds\) and nonnegative integrable angular weight \(p\). The analytic class and invertibility estimates are essential; the result is not a blanket linearization for every Schwartz density or every anisotropic model.}
\GSfield{Reconstruction}{The result is stated for the analytic class used by the primary proof; a larger unrestricted rapidly decaying class is not asserted.}
\GSfield{Fate}{\citet{bobylev1981poincare} was received 4 August 1980 and appeared in 1981. A full explicit proof is in \citet[§4, theorem and Eqs.~(4.8)--(4.14)]{bobylev1984exact}, received May 1984. Its nonresonance and function-space restrictions qualify the notebook’s broad wording.}
\GSfield{Present status}{As of November 2026, this specified isotropic linearization is established. The manuscript’s phrase unknown in other cases is historical and too unspecific to define an additional open problem.}
\GSfield{Assessment}{literal statement: correct with missing hypotheses; publication match: exact; verification: read the result; development: annotated; confidence: high}
\end{GSitem}
\begin{GSitem}{1980-525.9}{A}{Countably many nonlinear invariants of relaxation}
\GSfield{Kind; source}{theorem; \GSpage{P000975}}
\GSfield{The notes}{The notes claim infinitely many conservation laws and exact solutions.}
\GSfield{Statement}{For non-equilibrium solutions in the preceding normal-coordinate class, fix an index \(m\) with \(y_m\ne0\). Quantities \(|y_n|^{1/\lambda_n}/|y_m|^{1/\lambda_m}\) are constant wherever defined. Finite sets of nonzero normal coordinates also generate explicit special-solution families. These nonlinear functionals differ from the five linear collision invariants expressing mass, momentum, and energy.}
\GSfield{Fate}{\citet[§4, Corollaries~1--2, Eqs.~(4.15)--(4.21)]{bobylev1984exact} gives these consequences of the hierarchy already announced in \citet{bobylev1981poincare}.}
\GSfield{Present status}{As of November 2026, the qualified statement above is established; no new open problem is inferred from the omissions in the notebook.}
\GSfield{Assessment}{literal statement: correct with missing hypotheses; publication match: exact; verification: read the result; development: annotated; confidence: high}
\end{GSitem}
\begin{GSitem}{1980-525.10}{F}{The spatially inhomogeneous and hydrodynamic-limit programme}
\GSfield{Kind; source}{remark; \GSpage{P000975}}
\GSfield{The notes}{The spatially inhomogeneous equation is important, and a limiting mean-free-path regime gives Navier--Stokes.}
\GSfield{Statement}{The schematic kinetic-to-fluid statement requires a specified Knudsen scaling, collision model, boundary conditions, preparation of initial data, and solution class. Depending on the scaling, limiting equations can instead be compressible or incompressible Euler, acoustic, or diffusive systems. The notebook gives none of these choices.}
\GSfield{Fate}{\citet[§5]{bobylev1984exact} discusses the limitations of formal hydrodynamic expansions. It does not identify a single rigorous limit theorem with the brief seminar sentence.}
\GSfield{Present status}{As of November 2026, many specific hydrodynamic limits are established; no definite unresolved limit problem is reconstructed from this unscaled statement.}
\GSfield{Assessment}{literal statement: not determinable; publication match: thematic only; verification: read the result; development: annotated; confidence: high}
\end{GSitem}

\GStalk{526}{1981-03-30}{B. V. Lidskii}{P000975-P000976}{Spectra of sums of Hermitian matrices}
Lidskii announces the recursive Horn inequalities and a connection with tensor-product multiplicities. His two-page announcement appeared in 1982, but the later literature distinguishes it from a published proof. Klyachko’s work and the Knutson--Tao saturation theorem provide the established resolution.

\begin{GSitem}{1980-526.1}{B}{The Hermitian sum spectral problem}
\GSfield{Kind; source}{problem; \GSpage{P000975}}
\GSfield{The notes}{Three decreasing lists \(x,y,z\) satisfy \(\sum_i(x_i+y_i-z_i)=0\). The problem is to characterize when they are spectra of Hermitian \(X,Y,Z\) with \(X+Y=Z\).}
\GSfield{Statement}{For real weakly decreasing \(n\)-tuples \(x,y,z\), characterize the existence of complex Hermitian \(n\times n\) matrices \(X,Y\) whose eigenvalue lists are \(x,y\) and whose sum has eigenvalue list \(z\). Trace equality is necessary and is retained in the final characterization.}
\GSfield{Fate}{Lidskii announced the recursive answer in \citet{lidskii1982spectral}, received October 1981 and published April 1982. A complete resolution follows from the work connected by \citet[§1]{knutsontao1999honeycomb}, first circulated July 1998 and published in 1999.}
\GSfield{Present status}{As of November 2026, the Horn problem is solved by the recursive inequalities, equivalently by the honeycomb and representation-theoretic criteria.}
\GSfield{Assessment}{literal statement: correct; publication match: exact; verification: read the result; development: annotated; confidence: high; problem \hyperref[prob:GP-1980-42]{GP-1980-42}}
\end{GSitem}
\begin{GSitem}{1980-526.2}{B}{The recursive Horn inequalities}
\GSfield{Kind; source}{theorem; \GSpage{P000976}}
\GSfield{The notes}{Triples of increasing index lists are declared consistent by a trace-of-indices equality and inequalities for smaller consistent triples; the final answer is \(\sum z_{c_j}\le\sum x_{a_j}+\sum y_{b_j}\).}
\GSfield{Statement}{For \(1\le r\le n\), let \(U_r^n\) consist of triples \(I,J,K\) of \(r\)-subsets of \(\{1,\ldots,n\}\) with \(\sum I+\sum J-\sum K=r(r+1)/2\). Set \(T_1^n=U_1^n\); for \(r>1\), retain \((I,J,K)\in U_r^n\) exactly when every \(p<r\) and \((F,G,H)\in T_p^r\) satisfy \(\sum_{f\in F}i_f+\sum_{g\in G}j_g\le\sum_{h\in H}k_h+p(p+1)/2\). Hermitian realizability is equivalent to trace equality and \(\sum_{k\in K}z_k\le\sum_{i\in I}x_i+\sum_{j\in J}y_j\) for all \(r<n\) and \((I,J,K)\in T_r^n\).}
\GSfield{Reconstruction}{The recursion requires \(p<r\) (the notebook writes only \(\ell<n\)), and the final answer also requires the trace equality already stated on the preceding page.}
\GSfield{Fate}{\citet{lidskii1982spectral} is the exact short announcement. \citet{dayso1998spectrum} documents the absence of a published supporting proof at that stage; \citet[§1 and the saturation theorem]{knutsontao1999honeycomb} gives the later route completing Horn’s conjecture.}
\GSfield{Present status}{As of November 2026, the qualified statement above is established; no new open problem is inferred from the omissions in the notebook.}
\GSfield{Assessment}{literal statement: correct with missing hypotheses; publication match: exact; verification: read the result; development: annotated; confidence: high; problem \hyperref[prob:GP-1980-42]{GP-1980-42}}
\end{GSitem}
\begin{GSitem}{1980-526.3}{B}{Tensor-product multiplicities and saturation}
\GSfield{Kind; source}{theorem; \GSpage{P000976}}
\GSfield{The notes}{The substitution \(a_i\mapsto a_i-i\) simplifies the conditions, and a dual theorem determines which representations occur in a tensor product.}
\GSfield{Statement}{An increasing subset \(I=\{i_1<\cdots<i_r\}\) corresponds to the partition \((i_r-r,\ldots,i_1-1)\). For dominant integral \(GL_n\) weights \(\lambda,\mu,\nu\), the Hermitian eigenvalue inequalities with the trace condition characterize positivity of the Littlewood--Richardson multiplicity \(c_{\lambda\mu}^{\nu}\). Equivalence at integral scale uses saturation: \(c_{N\lambda,N\mu}^{N\nu}>0\) for some positive integer \(N\) implies \(c_{\lambda\mu}^{\nu}>0\).}
\GSfield{Fate}{\citet[Introduction and saturation theorem]{knutsontao1999honeycomb}, preprint July 1998, proves the saturation step and connects it with Klyachko’s 1990s work. The notebook anticipates the relation but contains no proof of this full equivalence.}
\GSfield{Present status}{As of November 2026, the qualified statement above is established; no new open problem is inferred from the omissions in the notebook.}
\GSfield{Assessment}{literal statement: correct with missing hypotheses; publication match: exact; verification: read the result; development: annotated; confidence: high; problem \hyperref[prob:GP-1980-42]{GP-1980-42}}
\end{GSitem}
\begin{GSitem}{1980-526.4}{E}{Convolution of orbital measures}
\GSfield{Kind; source}{construction; \GSpage{P000976}}
\GSfield{The notes}{The notes add that probability can also be written as a convolution.}
\GSfield{Statement}{For fixed spectra \(x,y\), take independent Haar-distributed unitary conjugates of their diagonal matrices. The law of the matrix sum is the convolution of the corresponding orbital probability measures on Hermitian matrices; its spectral pushforward is supported on the Horn polytope. The source gives no density or probability formula beyond this construction.}
\GSfield{Fate}{The spectral-sum context is \citet{dayso1998spectrum}; identifying a convolution is elementary once the orbital measures are defined. No explicit multiplicity or density algorithm is attributed to the one-line remark.}
\GSfield{Present status}{As of November 2026, the qualified statement above is established; no new open problem is inferred from the omissions in the notebook.}
\GSfield{Assessment}{literal statement: correct with missing hypotheses; publication match: partial; verification: publisher metadata checked; development: annotated; confidence: medium}
\end{GSitem}

\GStalk{527}{1981-04-06}{S. G. Krein}{P000977-P000978}{Compressible fluids, shocks, and dispersive step evolution}
Krein reports the diffeomorphism-group approach to compressible flow and the incompressible limit. Scalar characteristics and shock formation prompt interventions by Gelfand and Zakharov. Gelfand’s final question concerns step-like data for a polynomial regularization of Burgers’ equation; the polynomial is not specified, so distinct dispersive and dissipative cases must be separated.

\begin{GSitem}{1980-527.1}{E}{Barotropic compressible Euler equations}
\GSfield{Kind; source}{definition; \GSpage{P000977}}
\GSfield{The notes}{The notes give \(u_t+(u\cdot\nabla)u=-\rho^{-1}\nabla p\), \(\rho_t+\operatorname{div}(\rho u)=0\), \(p=p(\rho)\), exemplified by \(p=A\rho^\gamma\), and separately say \(T\) is constant.}
\GSfield{Statement}{For smooth positive density and a smooth barotropic pressure law, these are the compressible inviscid mass and momentum equations. Strict hyperbolicity of the acoustic part requires \(p^{\prime}(\rho)>0\). Barotropy does not require constant temperature: the polytropic ideal-gas law generally models an isentropic, nonisothermal regime, while the isothermal ideal-gas case has \(\gamma=1\).}
\GSfield{Reconstruction}{Positivity, smoothness, hyperbolicity, and the separation of barotropy from isothermality are supplied.}
\GSfield{Fate}{This is the continuum model underlying \citet{ebin1977slightly}. Gelfand’s heating objection identifies a limitation of the simplified thermodynamic model, not a mathematical contradiction in barotropic Euler theory.}
\GSfield{Present status}{As of November 2026, the qualified statement above is established; no new open problem is inferred from the omissions in the notebook.}
\GSfield{Assessment}{literal statement: correct with missing hypotheses; publication match: partial; verification: publisher metadata checked; development: annotated; confidence: high}
\end{GSitem}
\begin{GSitem}{1980-527.2}{E}{Characteristics and the onset of a scalar shock}
\GSfield{Kind; source}{method; \GSpage{P000977}}
\GSfield{The notes}{For \(u_t+cu_x=0\), \(u=f(x-ct)\); for \(u_t+uu_x=0\), \(u=f(x-ut)\). The drawing shows a falling flank overturning, and says shocks necessarily occur.}
\GSfield{Statement}{For scalar Burgers data \(u(0,x)=f(x)\in C^1\) with bounded derivative, the classical solution is parametrized by \(x=\xi+tf(\xi)\), \(u=f(\xi)\), while \(1+tf^{\prime}(\xi)>0\). If \(\inf f^{\prime}<0\), its first gradient catastrophe is at \(-1/\inf f^{\prime}\); nondecreasing data do not form this forward-time compression shock. Constant data also remain smooth.}
\GSfield{Fate}{The characteristic calculation is classical background to the entropy theory of \citet{kruzhkov1970first}, received April 1969 and published 1970. The unconditional assertion that nonlinear hyperbolic equations always require shocks is false.}
\GSfield{Present status}{As of November 2026, the qualified statement above is established; no new open problem is inferred from the omissions in the notebook.}
\GSfield{Assessment}{literal statement: false as written; publication match: partial; verification: publisher metadata checked; development: annotated; confidence: high}
\end{GSitem}
\begin{GSitem}{1980-527.3}{E}{Conservation form and entropy selection}
\GSfield{Kind; source}{method; \GSpage{P000977-P000978}}
\GSfield{The notes}{After characteristic crossing one should write \(u_t+\frac12(u^2)_x=0\) and define discontinuous solutions.}
\GSfield{Statement}{Bounded weak solutions satisfy the distributional conservation law, but that condition alone does not select a unique evolution. The scalar entropy solution obeys the Kruzhkov inequalities; a shock between constant states has Rankine--Hugoniot speed \((u_-+u_+)/2\), and for convex Burgers flux a decreasing jump is admissible.}
\GSfield{Fate}{\citet{kruzhkov1970first} proves existence, uniqueness, and stability in the bounded entropy class. The source’s request to understand discontinuities is thus an exposition of an established theory, not an unresolved uniqueness question.}
\GSfield{Present status}{As of November 2026, the qualified statement above is established; no new open problem is inferred from the omissions in the notebook.}
\GSfield{Assessment}{literal statement: correct with missing hypotheses; publication match: exact; verification: publisher metadata checked; development: annotated; confidence: high}
\end{GSitem}
\begin{GSitem}{1980-527.4}{F}{Zakharov’s matrix Hopf equation}
\GSfield{Kind; source}{claim; \GSpage{P000978}; V. E. Zakharov}
\GSfield{The notes}{Zakharov says \(U_t+UU_x=0\), with matrix \(U\), can also be integrated completely.}
\GSfield{Statement}{For matrix-valued \(U\), multiplication order matters: \(UU_x\) is generally not \(\frac12(U^2)_x\). The source provides no spectral assumptions, initial-data class, solution formula, or definition of complete integration. Diagonal or mutually commuting reductions reduce to scalar characteristic equations, but that is only a restricted case.}
\GSfield{Fate}{The nearby scalar theory does not certify the unqualified matrix claim. Searches for matrix Hopf and Zakharov noncommutative hydrodynamic systems did not identify a primary theorem with exactly these unspecified hypotheses. No general integrability theorem is inferred.}
\GSfield{Present status}{As of November 2026, the exact general matrix-integration assertion in this note remains unidentified; scalar commuting reductions are elementary and do not settle it.}
\GSfield{Assessment}{literal statement: not determinable; publication match: none found; verification: unresolved lead; development: annotated; confidence: medium}
\end{GSitem}
\begin{GSitem}{1980-527.5}{F}{Expansion of gas into vacuum}
\GSfield{Kind; source}{question; \GSpage{P000978}}
\GSfield{The notes}{For a vessel initially divided into gas and vacuum, removing the partition is said to produce a system of shocks moving back and forth.}
\GSfield{Statement}{The initial expansion of a positive-pressure gas into vacuum in the standard one-dimensional isentropic Riemann problem is described by a rarefaction toward vacuum. Later wall reflections depend on the equation of state, boundary conditions, and geometry. The diagram alone does not determine a universal shock-train evolution or equilibration theorem.}
\GSfield{Fate}{The displayed compressible equations permit both compression and expansion waves; the scalar entropy context is \citet{kruzhkov1970first}. No exact vessel problem or solution theorem was found that warrants the notebook’s unqualified description.}
\GSfield{Present status}{As of November 2026, the recorded vessel question is underspecified; its blanket initial shock description is not valid for expansion into vacuum.}
\GSfield{Assessment}{literal statement: false as written; publication match: thematic only; verification: unresolved lead; development: annotated; confidence: medium}
\end{GSitem}
\begin{GSitem}{1980-527.6}{E}{The hyperbolic-plane nonimmersion analogy}
\GSfield{Kind; source}{claim; \GSpage{P000978}}
\GSfield{The notes}{The notes connect the preceding shock discussion to the assertion that the Lobachevsky plane does not embed in \(\mathbb R^3\).}
\GSfield{Statement}{The classical geometric obstruction is that the complete constant-negative-curvature hyperbolic plane has no smooth isometric immersion into Euclidean three-space. Topological or smooth embeddings without preservation of the metric exist. The notebook supplies no argument connecting this theorem to the gas example.}
\GSfield{Fate}{The intended geometric result is \citet{hilbert1901flachen}; the primary AMS access attempt was blocked, so the bibliography records that limitation. The isometric and completeness qualifications are indispensable and the asserted causal link to shocks is not established.}
\GSfield{Present status}{As of November 2026, the qualified Hilbert nonimmersion theorem is classical; the unqualified embedding sentence and the suggested gas analogy are not valid as a proof.}
\GSfield{Assessment}{literal statement: false as written; publication match: partial; verification: indirect; development: annotated; confidence: medium}
\end{GSitem}
\begin{GSitem}{1980-527.7}{E}{The slightly compressible to incompressible limit}
\GSfield{Kind; source}{theorem; \GSpage{P000978}}
\GSfield{The notes}{Krein says that with density close to one, a parameter limit gives incompressible-fluid solutions.}
\GSfield{Statement}{For the smooth slightly compressible models and compatible, suitably prepared initial data of Ebin’s theorem, solutions approach the incompressible motion on a common time interval where the latter remains smooth. A small Mach parameter, appropriate control of acoustic oscillations, and boundary compatibility are essential; arbitrary data and passage through shock formation are not covered.}
\GSfield{Reconstruction}{This is a calibrated description of the source theorem, not a new quantified convergence statement without its original hypotheses.}
\GSfield{Fate}{The identifiable result is \citet{ebin1977slightly}, published in January 1977, before this lecture. The opened Annals record verifies the source and its strong-constraint interpretation; the notebook does not state the exact function spaces or convergence norm.}
\GSfield{Present status}{As of November 2026, the qualified statement above is established; no new open problem is inferred from the omissions in the notebook.}
\GSfield{Assessment}{literal statement: correct with missing hypotheses; publication match: partial; verification: publisher metadata checked; development: annotated; confidence: medium}
\end{GSitem}
\begin{GSitem}{1980-527.8}{E}{Internal energy in material coordinates}
\GSfield{Kind; source}{construction; \GSpage{P000978}}
\GSfield{The notes}{For the Jacobian \(g(x)\) of the flow map, the potential energy is \[\int_\Omega\int_1^{g^{-1}(x)}p(\rho)\rho^{-2}\,d\rho\,dx.\]}
\GSfield{Statement}{With unit reference mass density, write the material flow as \(\eta(X)\), its positive Jacobian as \(J(X)\), and \(\rho(\eta(X))=J(X)^{-1}\). If specific internal energy obeys \(e^{\prime}(\rho)=p(\rho)/\rho^2\), then \(\int e(J^{-1})\,dX\) is the potential energy up to an additive constant. In Eulerian coordinates the same energy is \(\int \rho e(\rho)\,dx\).}
\GSfield{Reconstruction}{The reference measure and distinction between specific energy and energy density are supplied explicitly.}
\GSfield{Fate}{This is the mechanical energy interpretation used in the strongly constrained flow setting of \citet{ebin1977slightly}. The formula is consistent once the integral is read in material coordinates.}
\GSfield{Present status}{As of November 2026, the qualified statement above is established; no new open problem is inferred from the omissions in the notebook.}
\GSfield{Assessment}{literal statement: correct with missing hypotheses; publication match: partial; verification: publisher metadata checked; development: annotated; confidence: high}
\end{GSitem}
\begin{GSitem}{1980-527.9}{F}{The unspecific viscous-fluid extension}
\GSfield{Kind; source}{remark; \GSpage{P000978}}
\GSfield{The notes}{The notes say the same exists for a viscous fluid, following Arnold.}
\GSfield{Statement}{Viscosity adds a dissipative stress and destroys conservation of the inviscid Hamiltonian in general. A geometric formulation can still be given, but the sentence does not specify whether it refers to a manifold description, a limit theorem, or a Hamiltonian structure with extra variables.}
\GSfield{Fate}{\citet{arnold1966geometrie} supplies the ideal-fluid geometric context; it does not make a viscous equation an energy-conserving geodesic flow by this sentence alone.}
\GSfield{Present status}{As of November 2026, the intended viscous analogue remains unspecified in this record.}
\GSfield{Assessment}{literal statement: not determinable; publication match: thematic only; verification: read the result; development: annotated; confidence: medium}
\end{GSitem}
\begin{GSitem}{1980-527.10}{F}{A descending step for a polynomially regularized Burgers equation}
\GSfield{Kind; source}{problem; \GSpage{P000978}; I. M. Gelfand}
\GSfield{The notes}{Gelfand asks what happens to a descending step profile for \(u_t+uu_x+p(\partial_x)u=0\), with \(p\) a polynomial. The final sentence says KdV continually sheds solitons.}
\GSfield{Statement}{Long-time behaviour depends on the polynomial and on the solution concept: \(p=0\) gives entropy shocks, \(p(D)=-\nu D^2\), \(\nu>0\), gives viscous smoothing and travelling shocks, and a real cubic dispersive term gives the KdV dispersive-shock problem. Some polynomial symbols cause linear ill-posedness. The notes prescribe neither coefficients nor admissible class, so there is no single well-posed universal evolution question.}
\GSfield{Fate}{For the KdV subcase, \citet{gurevichpitaevskii1974shock}, submitted April 1973, Russian publication 1973 and English 1974, describes the expanding modulated oscillatory zone with a soliton edge. This is a concrete prior result, not an answer for every polynomial.}
\GSfield{Present status}{As of November 2026, several specified subcases are established; the unrestricted polynomial question remains an underspecified programme, not a newly indexed open conjecture.}
\GSfield{Assessment}{literal statement: not determinable; publication match: partial; verification: read the result; development: annotated; confidence: high}
\end{GSitem}

\GStalk{528}{1981-04-06}{G. I. Olshanskii}{P000978-P000980}{Invariant cones, Lie semigroups, and highest-weight representations}
Olshanskii develops semigroups of contractions, invariant cones in Hermitian Lie algebras, and holomorphic extension of unitary highest-weight representations. The primary paper was received in December 1980 and published in 1981. Properness of cones, a choice of sign, and the analytic representation category are essential qualifications absent from the shorthand notes.

\begin{GSitem}{1980-528.1}{E}{The semigroup of indefinite contractions}
\GSfield{Kind; source}{construction; \GSpage{P000978}}
\GSfield{The notes}{For \(J=\operatorname{diag}(-I_p,I_q)\), a matrix \(A\) is a J-contraction when \(\langle Ax,Ax\rangle_J\le\langle x,x\rangle_J\). Such matrices form a semigroup.}
\GSfield{Statement}{The inequality is \(A^*JA\le J\). If both \(A,B\) satisfy it, then \((AB)^*J(AB)=B^*A^*JAB\le B^*JB\le J\). The full matrix semigroup can include singular matrices; its invertible part is the part naturally compared with semigroups inside a complex Lie group.}
\GSfield{Fate}{This is an elementary model for the constructions in \citet[§§3,6]{olshanskii1981cones}, received December 1980 and published in 1981. The singular boundary must not be conflated with the group-contained semigroup \(G\exp(iC)\).}
\GSfield{Present status}{As of November 2026, the qualified statement above is established; no new open problem is inferred from the omissions in the notebook.}
\GSfield{Assessment}{literal statement: correct; publication match: exact; verification: read the result; development: annotated; confidence: high}
\end{GSitem}
\begin{GSitem}{1980-528.2}{A}{Compression semigroups of domains}
\GSfield{Kind; source}{construction; \GSpage{P000978-P000979}}
\GSfield{The notes}{Examples are hyperbolic motions taking a half-plane into itself, and transformations of the compact dual of a bounded symmetric domain \(D\) taking \(D\) into itself. The matrix ball gives J-contractions.}
\GSfield{Statement}{For a fixed domain \(D\) in a homogeneous space and an ambient transformation group, the set \(\{g:gD\subseteq D\}\) is a semigroup. For a Hermitian symmetric domain embedded in its compact dual, the connected complex group supplies the compression semigroup identified with the maximal-cone Olshanskii semigroup in the cited theorem. Matrix-ball realizations connect this geometry with appropriately normalized indefinite contractions.}
\GSfield{Reconstruction}{The ambient group, connected component, and distinction from all holomorphic self-maps are made explicit.}
\GSfield{Fate}{\citet[§6, Thm.~6.2]{olshanskii1981cones} proves the symmetric-domain identification. The elementary closure property applies also to the real half-plane example, but does not classify arbitrary domain-preserving maps.}
\GSfield{Present status}{As of November 2026, the qualified statement above is established; no new open problem is inferred from the omissions in the notebook.}
\GSfield{Assessment}{literal statement: correct with missing hypotheses; publication match: exact; verification: read the result; development: annotated; confidence: high}
\end{GSitem}
\begin{GSitem}{1980-528.3}{A}{Contractive involutive-semigroup representations}
\GSfield{Kind; source}{definition; \GSpage{P000979}}
\GSfield{The notes}{The notes require \(T(ab)=T(a)T(b)\), \(T(a^\#)=T(a)^*\), and \(\|T(a)\|\le1\), where \((ab)^\#=b^\#a^\#\).}
\GSfield{Statement}{A representation in the category used here is a continuous nondegenerate homomorphism of a topological involutive semigroup into bounded Hilbert-space contractions, respecting the involution. For a unital semigroup one takes \(T(1)=I\); weak operator continuity is sufficient in the primary definition and yields strong continuity in the setting considered.}
\GSfield{Fate}{The exact definition appears in \citet[§4.2]{olshanskii1981cones}. Continuity and nondegeneracy are necessary category qualifications absent from the abbreviated notebook list.}
\GSfield{Present status}{As of November 2026, the qualified statement above is established; no new open problem is inferred from the omissions in the notebook.}
\GSfield{Assessment}{literal statement: correct with missing hypotheses; publication match: exact; verification: read the result; development: annotated; confidence: high}
\end{GSitem}
\begin{GSitem}{1980-528.4}{E}{Tensor-power representations of ordinary contractions}
\GSfield{Kind; source}{construction; \GSpage{P000979}}
\GSfield{The notes}{For contractions on \(V=\mathbb C^n\), tensor powers give analytic representations indexed by sequences; anti-analytic representations also occur.}
\GSfield{Statement}{The tensor-power action \(A\mapsto A^{\otimes m}\) restricts to the Schur modules \(S_\lambda V\), indexed by partitions \(\lambda\vdash m\) with at most \(n\) rows, and remains contractive. Conjugating gives antiholomorphic examples. These are polynomial holomorphic representations; the notebook does not justify classifying every continuous involutive-semigroup representation by these partitions alone.}
\GSfield{Reconstruction}{The sequences are interpreted as partitions for polynomial representations, without extending the claim to an unrestricted continuous category.}
\GSfield{Fate}{The examples motivate the analytic category in \citet[Introduction and §4]{olshanskii1981cones}. Ordinary matrix-contraction representations and the later \(\Gamma(C_{\min})\) theorem must be kept distinct.}
\GSfield{Present status}{As of November 2026, the qualified statement above is established; no new open problem is inferred from the omissions in the notebook.}
\GSfield{Assessment}{literal statement: correct with missing hypotheses; publication match: partial; verification: read the result; development: annotated; confidence: high}
\end{GSitem}
\begin{GSitem}{1980-528.5}{E}{The Hermitian criterion for invariant proper cones}
\GSfield{Kind; source}{theorem; \GSpage{P000979-P000980}}
\GSfield{The notes}{A convex adjoint-invariant cone with nonempty interior exists in a simple real Lie algebra exactly when it is of Hermitian type.}
\GSfield{Statement}{For a noncompact simple real Lie algebra \(\mathfrak g\), the existence of a nonzero closed convex pointed generating adjoint-invariant cone is equivalent to \(\mathfrak g\) being of Hermitian type. Nonempty interior alone is insufficient: \(\mathfrak g\) itself is an invariant convex cone. For \(\mathfrak g=\mathfrak k\oplus\mathfrak p\), the Hermitian condition is \(\dim Z(\mathfrak k)=1\).}
\GSfield{Reconstruction}{Closedness, pointedness, generating property, and noncompactness are supplied.}
\GSfield{Fate}{\citet[Thm.~1.2]{olshanskii1981cones} states the Vinberg--Kostant criterion and specifies the nondegenerate cone convention. This is explicitly attributed earlier work, even though its use in the speaker’s paper is new.}
\GSfield{Present status}{As of November 2026, the qualified statement above is established; no new open problem is inferred from the omissions in the notebook.}
\GSfield{Assessment}{literal statement: false as written; publication match: exact; verification: read the result; development: annotated; confidence: high}
\end{GSitem}
\begin{GSitem}{1980-528.6}{E}{The list of simple Hermitian real forms}
\GSfield{Kind; source}{theorem; \GSpage{P000979}}
\GSfield{The notes}{The list contains \(\mathfrak{su}(p,q)\), \(\mathfrak{sp}(2n,\mathbb R)\), \(\mathfrak{so}^*(2n)\), \(\mathfrak{so}(2,n)\), and two exceptional forms.}
\GSfield{Statement}{The noncompact simple Hermitian families are \(\mathfrak{su}(p,q)\) for \(p,q\ge1\), \(\mathfrak{sp}(2n,\mathbb R)\) for \(n\ge1\), \(\mathfrak{so}^*(2n)\) for \(n\ge3\), and \(\mathfrak{so}(2,n)\) for \(n=1\) or \(n\ge3\), together with \(\mathfrak e_{6(-14)}\) and \(\mathfrak e_{7(-25)}\), with low-rank isomorphisms identified. In particular \(\mathfrak{so}(2,2)\) is semisimple rather than simple.}
\GSfield{Reconstruction}{Exceptional signatures and the low-rank simplicity qualification are supplied.}
\GSfield{Fate}{The criterion and root-theoretic conventions are in \citet[§1]{olshanskii1981cones}; the exceptional names expand the source’s shorthand.}
\GSfield{Present status}{As of November 2026, the qualified statement above is established; no new open problem is inferred from the omissions in the notebook.}
\GSfield{Assessment}{literal statement: correct with missing hypotheses; publication match: exact; verification: read the result; development: annotated; confidence: high}
\end{GSitem}
\begin{GSitem}{1980-528.7}{A}{Minimal and maximal cones after choosing an orientation}
\GSfield{Kind; source}{theorem; \GSpage{P000979-P000980}}
\GSfield{The notes}{There are unique \(C_{\min}\subset C_{\max}\), and a Cartan-subalgebra description, with attribution to Vinberg, Olshanskii, and Paneitz.}
\GSfield{Statement}{After choosing a positive orientation in the one-dimensional centre of the maximal compact algebra, the positive invariant proper cones have a unique minimal and maximal member. Their negatives give the opposite orientation, so uniqueness without a sign choice is false. Intersection with a compact Cartan subspace identifies the intermediate cones with Weyl-invariant cones between the corresponding root-theoretic extrema.}
\GSfield{Reconstruction}{A positive central direction, and hence an orientation of the cone, is fixed before stating uniqueness.}
\GSfield{Fate}{\citet[Thms.~1.5,1.7,1.9 and Cor.~1.6]{olshanskii1981cones} gives the precise bounds and classification, distinguishing Vinberg’s earlier maximal cone from the new Cartan-cone description. The paper was received 4 December 1980.}
\GSfield{Present status}{As of November 2026, the qualified statement above is established; no new open problem is inferred from the omissions in the notebook.}
\GSfield{Assessment}{literal statement: correct with missing hypotheses; publication match: exact; verification: read the result; development: annotated; confidence: high}
\end{GSitem}
\begin{GSitem}{1980-528.8}{A}{Global semigroups from Hermitian invariant cones}
\GSfield{Kind; source}{theorem; \GSpage{P000980}}
\GSfield{The notes}{For every cone under discussion, \(\Gamma(C)=G\exp(iC)\) is a semigroup inside \(G_{\mathbb C}\).}
\GSfield{Statement}{Let \(\mathfrak g\) be a noncompact simple Hermitian real Lie algebra, \(G_{\mathbb C}\) a connected complex Lie group with Lie algebra \(\mathfrak g_{\mathbb C}\), and \(G\) its connected analytic subgroup with Lie algebra \(\mathfrak g\). For a closed pointed generating adjoint-invariant cone \(C\subset\mathfrak g\), \(G\exp(iC)\) is a closed semigroup, and polar multiplication gives a homeomorphism \(G\times C\to\Gamma(C)\). Appropriate universal covers are used for representations that do not descend to this \(G\).}
\GSfield{Reconstruction}{The Lie-group, cone, and covering hypotheses are stated explicitly.}
\GSfield{Fate}{This is \citet[Thm.~3.5 and §4.1]{olshanskii1981cones}, after identifying the paper’s cone in \(i\mathfrak g\) with \(iC\). The theorem does not claim this global polar description for arbitrary Lie algebras or arbitrary convex cones.}
\GSfield{Present status}{As of November 2026, the qualified statement above is established; no new open problem is inferred from the omissions in the notebook.}
\GSfield{Assessment}{literal statement: correct with missing hypotheses; publication match: exact; verification: read the result; development: annotated; confidence: high}
\end{GSitem}
\begin{GSitem}{1980-528.9}{A}{Local semigroups for general invariant cones}
\GSfield{Kind; source}{theorem; \GSpage{P000980}}
\GSfield{The notes}{A weaker local version is said to hold for every Lie algebra.}
\GSfield{Statement}{If \(M\) is a connected finite-dimensional Lie group and \(C\) a closed pointed generating invariant convex cone in its Lie algebra, the germ at the identity of \(\exp(C\cap U)\), with \(U\) shrinking to zero, is a local semigroup. This is a statement about products sufficiently near the identity, not automatic global closure of an exponential image.}
\GSfield{Fate}{\citet[Thm.~3.3]{olshanskii1981cones} proves this assertion independently of the Hermitian global theorem.}
\GSfield{Present status}{As of November 2026, the qualified statement above is established; no new open problem is inferred from the omissions in the notebook.}
\GSfield{Assessment}{literal statement: correct with missing hypotheses; publication match: exact; verification: read the result; development: annotated; confidence: high}
\end{GSitem}
\begin{GSitem}{1980-528.10}{A}{The representation cone and highest-weight condition}
\GSfield{Kind; source}{theorem; \GSpage{P000980}}
\GSfield{The notes}{The note defines \(C(T)\) by a nonpositivity condition on infinitesimal generators and says it is nontrivial exactly for highest-weight representations.}
\GSfield{Statement}{For a strongly continuous unitary representation \(\pi\), use the self-adjoint closures to define \(C_\pi=\{Y\in\mathfrak g:i\,d\pi(Y)\le0\}\). For a nontrivial irreducible unitary representation of the simply connected Hermitian group, a nonzero such cone, with its orientation fixed, characterizes a unitary highest-weight representation; the opposite orientation interchanges highest and lowest weight. The trivial representation has \(C_\pi=\mathfrak g\) and is a separate exception.}
\GSfield{Reconstruction}{The infinitesimal generator, operator closure, irreducibility, exclusion of the trivial case, and orientation are supplied.}
\GSfield{Fate}{\citet[Thms.~2.2,2.5]{olshanskii1981cones} gives the exact nontriviality, irreducibility, and sign qualifications. The paper works in \(i\mathfrak g\); the formula here pulls that convention back to \(\mathfrak g\).}
\GSfield{Present status}{As of November 2026, the qualified statement above is established; no new open problem is inferred from the omissions in the notebook.}
\GSfield{Assessment}{literal statement: correct with missing hypotheses; publication match: exact; verification: read the result; development: annotated; confidence: high}
\end{GSitem}
\begin{GSitem}{1980-528.11}{A}{Holomorphic extension and the contractive representation category}
\GSfield{Kind; source}{theorem; \GSpage{P000980}}
\GSfield{The notes}{A highest-weight representation extends holomorphically to \(\Gamma(C(T))\), called the exact extension domain, and there is a category equivalence.}
\GSfield{Statement}{For a fixed positive invariant proper cone \(C\), unitary representations satisfying \(C\subset C_\pi\) extend uniquely to contractive involutive representations on the appropriate covering semigroup, holomorphic on its interior, with \(T(g\exp(iY))=\pi(g)\exp(i,d\pi(Y))\). Restriction is an equivalence of the corresponding categories. The cone \(C_\pi\) is maximal for contractive extension in this setting, not necessarily for all analytic continuations of every individual matrix coefficient.}
\GSfield{Reconstruction}{The cone condition and the word contractive make precise both the category and the maximality assertion.}
\GSfield{Fate}{\citet[Thms.~4.5--4.6 and Thm.~2.9]{olshanskii1981cones} provides the extension, uniqueness, and direct-integral/category formulation. The source’s phrase all highest-weight representations requires the compatibility condition with the chosen cone.}
\GSfield{Present status}{As of November 2026, the qualified statement above is established; no new open problem is inferred from the omissions in the notebook.}
\GSfield{Assessment}{literal statement: correct with missing hypotheses; publication match: exact; verification: read the result; development: annotated; confidence: high}
\end{GSitem}
\begin{GSitem}{1980-528.12}{A}{Holomorphic–antiholomorphic factorization for the minimal cone}
\GSfield{Kind; source}{theorem; \GSpage{P000980}}
\GSfield{The notes}{Irreducible representations are said to be tensor products of holomorphic and antiholomorphic irreducible representations.}
\GSfield{Statement}{For the involutive semigroup associated with the minimal positive cone, in the continuous contractive category and on the cover specified by the theorem, every irreducible representation is the tensor product of a holomorphic highest-weight extension and the conjugate of another such extension. The theorem as checked does not assert the same unrestricted formulation for every intermediate cone.}
\GSfield{Reconstruction}{The minimal-cone hypothesis is restored from the primary theorem.}
\GSfield{Fate}{This is the exact scope of \citet[Cor.~5.8]{olshanskii1981cones}, whose notation explicitly has \(C_{\min}\). The notebook has dropped that qualifier.}
\GSfield{Present status}{As of November 2026, the qualified statement above is established; no new open problem is inferred from the omissions in the notebook.}
\GSfield{Assessment}{literal statement: correct with missing hypotheses; publication match: exact; verification: read the result; development: annotated; confidence: high}
\end{GSitem}
\begin{GSitem}{1980-528.13}{F}{Gelfand’s request to go beyond semisimple groups}
\GSfield{Kind; source}{remark; \GSpage{P000980}; I. M. Gelfand}
\GSfield{The notes}{Gelfand objects to restricting the theory to simple or semisimple algebras, especially for infinite-dimensional representation theory.}
\GSfield{Statement}{This is a programme to extend invariant-cone and semigroup methods beyond the finite-dimensional semisimple setting. No class of nonsemisimple or infinite-dimensional groups, topology, cone condition, or target theorem is specified.}
\GSfield{Fate}{\citet[Introduction]{olshanskii1981cones} already discusses an oscillator example as motivation. That overlap is thematic and does not solve an unspecified universal extension problem.}
\GSfield{Present status}{As of November 2026, this remark remains a broad direction in the record, without a definite question suitable for a new problem-index entry.}
\GSfield{Assessment}{literal statement: not determinable; publication match: thematic only; verification: read the result; development: annotated; confidence: high}
\end{GSitem}

\section{Developments}\label{dev:1980}
\subsection{The resolved genus-two subcase: a coordinate obstruction and an explicit certificate}
\label{dev:genus2}

This section audits the now-resolved genus-two subcase, item 1980-519.2 and problem GP-1980-40, as preparation for the still-unresolved universal function-field question in the next subsection. The genus-two item is class B, not an open problem manufactured to justify a development. On P000970 the notebook asks whether a genus-two equation \(P(x,y)=0\) can be satisfied by commuting operators with rational coefficients. It does not specify a normal form for the curve, a nontriviality condition, or permission to change its affine coordinates. These omissions matter. Constant operators corresponding to a complex point satisfy the equation trivially. For a prescribed affine model, the substantive requirement is an \emph{injective unital} map from its integral coordinate ring to
\[
 R=\mathbb C(x)[D],\qquad Df=fD+f',\quad D=d/dx.
\]
The letter \(x\) here is the differential variable; \(z,w\) will denote coordinates of a spectral curve. Injectivity rules out a constant-point representation. This qualification, and the distinction between a prescribed affine model and its function field, are ours.

\paragraph{The decisive quantifier distinction.}
For every smooth genus-two curve, one can choose its standard hyperelliptic quintic model and realize it by commuting ordinary differential operators with rational coefficients. This is a verified later theorem, not merely a generic existence result. The stronger assertion that every prescribed affine model works is false, by the proposition below. We do not convert the difference into a new historical open classification problem.

\paragraph{Checked literature and chronology.}
\citet[Lemma~5.4, Proposition~5.5 and Remark~5.6, pp.~221--223]{dixmier1968weyl} gives the elliptic polynomial-operator example. \citet[Theorems~2--3 and Example~(a), preprint p.~11]{mironov2011selfadjoint} constructs higher-genus hyperelliptic examples; the preprint appeared on 18 July 2011, the article online on 19 October 2013 and in the August 2014 issue. Its genus-two family is the source of the explicit pair verified below. \citet[Theorem~1 and its parameter-map proof]{davletshinamironov2016orbits}, first circulated on 4 June 2016 and published in 2017, realizes a fixed \emph{generic} genus-two spectral curve by polynomial operators and proves infinitely many automorphism orbits. Dominance of that parameter map alone does not settle all special curves. Crucially, \citet[Theorem~3(2), using the union of Theorems~1 and~2]{oganesyan2016rational}, first circulated on 18 August 2016 and published in 2018, realizes \emph{every} monic quintic \(w^2=z^5+\beta_4z^4+\cdots+\beta_0\) by rational-coefficient operators of orders four and ten. Both operator families are needed for the stated coverage. The bounded contribution here is a coordinate obstruction, a correction of the notebook's formula and an exact reproducible certificate for a known family. No novelty is claimed for the known existence theorems.

\begin{lemma}[Units and a rational multiplication operator]\ours\quad
\label{lem:units}
The units of \(R\) are precisely \(\mathbb C(x)^*\). If \(f\in\mathbb C(x)\setminus\mathbb C\), then its centralizer in \(R\) is \(\mathbb C(x)\).
\end{lemma}
\begin{proof}
For nonzero ordinary differential operators the highest derivative coefficient of a product is the product of the highest coefficients. Consequently
\(\operatorname{ord}(AB)=\operatorname{ord}(A)+\operatorname{ord}(B)\), with both orders nonnegative. An invertible operator and its inverse therefore have order zero. Conversely, every nonzero rational function is invertible in \(R\).

Write \(A=a_m(x)D^m+\cdots+a_0(x)\), with \(a_m\ne0\). For \(m>0\), the coefficient of \(D^{m-1}\) in \([A,f]\) is \(m a_m f'\). Characteristic zero and nonconstant rational \(f\) imply \(f'\ne0\), so this coefficient cannot vanish. Thus an operator commuting with \(f\) has order zero, and all order-zero operators commute with it.
\end{proof}

\begin{proposition}[A unit obstruction for an affine curve]\ours\quad
\label{prop:unitobstruction}
Let \(C^\circ\) be a smooth irreducible affine curve over \(\mathbb C\) whose smooth projective model \(C\) has positive genus. If \(\mathbb C[C^\circ]\) has a nonconstant unit, there is no injective unital \(\mathbb C\)-algebra homomorphism
\[
 \mathbb C[C^\circ]\longrightarrow \mathbb C(x)[D].
\]
\end{proposition}
\begin{proof}
Suppose \(\iota\) is such a map and \(u\) is a nonconstant unit. Lemma~\ref{lem:units} gives \(\iota(u)=f(x)\in\mathbb C(x)^*\). If \(f=c\) were constant, \(u-c\ne0\) would belong to the kernel. Hence \(f\) is nonconstant. The commutativity of the source forces the entire image into the centralizer of \(f\), which is \(\mathbb C(x)\). Taking ordinary commutative fractions gives a field embedding
\(\mathbb C(C)\hookrightarrow\mathbb C(x)\).

This embedding induces a nonconstant rational map \(\mathbb P^1\dashrightarrow C\), which extends to a finite morphism because \(C\) is smooth and projective. If its degree is \(d\ge1\), the characteristic-zero Riemann--Hurwitz formula gives
\[
 -2=d(2g(C)-2)+\sum_{p\in\mathbb P^1}(e_p-1)\ge0,
\]
since \(g(C)\ge1\). This contradiction proves the assertion.
\end{proof}

\begin{example}[A smooth genus-two equation that cannot be realized in the prescribed coordinates]\ours\quad
\label{ex:affineobstruction}
Consider
\[
 C^\circ:\quad z w^2-z^4-1=0.
\]
In its coordinate ring, \(z(w^2-z^3)=1\), so \(z\) is a nonconstant unit. Set \(Y=zw\). On \(z\ne0\), the equation becomes
\[
 Y^2=z^5+z.
\]
The polynomial \(z^5+z\) is squarefree: at its zero \(z=0\) its derivative is \(1\), and at each nonzero zero, where \(z^4=-1\), its derivative is \(-4\). The double cover of \(\mathbb P^1_z\) is branched at these five roots and at infinity, so its smooth projective model has genus two. The affine equation is irreducible because \(z^5+z\) is not a square in \(\mathbb C(z)\). It is smooth: since \(z\ne0\), vanishing of the derivative with respect to \(w\) would imply \(w=0\); then the derivative with respect to \(z\) is \(-4z^3\ne0\). Proposition~\ref{prop:unitobstruction} applies. Thus the universal prescribed-affine-equation reading of the notebook question has a negative answer, even for a smooth genus-two affine curve. This does not give a counterexample to Gelfand's function-field conjecture.
\end{example}

\paragraph{Correcting the preliminary elliptic example.}
The literal P000970 expression has \(L=D^2+x^3+a\), \(L_4=L^2+2x\) and \(L_6=L^3+3x^2L+3Lx^2\). With the convention \([D,x]=1\), the coefficient of \(D^5\) in the commutator of these last two operators is
\[
 48x-12\ne0.\qquad\oursComputed
\]
The formula is therefore false, including at \(a=0\). A corrected form of the classical pair is
\[
 P_1=L^2+2x,\qquad
 M_1=L^3+\frac32(xL+Lx)=L^3+3xL+3D.
\]
Our exact computation gives
\[
 [P_1,M_1]=0,\qquad M_1^2=P_1^3-a.\qquad\oursComputed
\]
For \(a\ne0\) the spectral cubic is smooth; for \(a=0\) it is singular. This is the normalization of \citet[Proposition~5.5]{dixmier1968weyl} with \(p=x,q=-D\). The distinction is recorded as a faithful mathematical error in the notebook, rather than an unsubstantiated change of handwriting.

The second notebook description also needs qualification. Its two-term expression \(A=L+xL^{-1}\) has
\[
 (A^3)_+=L^3+3xL+6D,\qquad\oursComputed
\]
where \((\cdot)_+\) denotes the differential part of a formal pseudodifferential operator. Thus it does not reproduce \(M_1\). The actual formal square root of \(P_1\) begins \(L+xL^{-1}-DL^{-2}+\cdots\); including this next term gives the correct differential part of the cube. The exact truncated calculation in \texttt{code/check\_formal.py} verifies both formulas with a cutoff that is sufficient for every nonnegative derivative order. This is a correction of an abbreviated formal expression, not evidence against the classical commuting pair.

\begin{proposition}[An exact genus-two operator certificate]\oursComputed\quad
\label{prop:certificate}
Let \(a\in\mathbb C\), \(L=D^2+x^3+a\), and set
\begin{align*}
 P&=L^2+6x,\\
 M&=LP^2+(3xL-3D)P+9x^2L-18xD+9.
\end{align*}
These are polynomial-coefficient differential operators of orders \(4\) and \(10\). They satisfy
\[
 [P,M]=0,\qquad M^2=P^5+27aP^2+81.
\]
The polynomial \(F_a(z)=z^5+27az^2+81\) is squarefree precisely when
\(2916a^5+3125\ne0\). In particular \(a=0\) gives a smooth genus-two spectral curve.
\end{proposition}
\begin{proof}[Exact computation and smoothness]
All operator calculations are performed over \(\mathbb Q[a,x]\), not at floating-point samples. In left normal form \(A=\sum_i a_iD^i\), \(B=\sum_j b_jD^j\), the complete multiplication rule is
\[
 (AB)_n=\sum_{\substack{i,j\ge0,\ 0\le k\le i\\i+j-k=n}}
 \binom{i}{k}a_i\,\partial_x^k b_j.
\]
It follows by induction from \(Df=fD+f'\). The included script \texttt{code/check\_weyl.py} implements exactly this finite formula and expands every coefficient as a polynomial. The two residual coefficient dictionaries are empty: every coefficient of \([P,M]\) and of \(M^2-P^5-27aP^2-81\) is identically zero. The output contains the full expanded \(P\) and \(M\), the nonzero literal notebook commutator, and the zero residuals, allowing the finite identities to be reproduced without interpreting an unspecified numerical tolerance.

An additional identity check compares the constructed \(M\) with the displayed preprint form and detects a scalar normalization:
\[
 M=L^5+\frac{15}{2}(xL^3+L^3x)+45(x^2L+Lx^2)-9.
\]
Example~(a) of \citet{mironov2011selfadjoint}, in the checked 2012 preprint, displays the expression without \(-9\). That expression still commutes with \(P\), but it equals \(M+9\); its spectral relation must be centered as \((M_{\rm displayed}-9)^2=F_a(P)\). The independent comparison caught this difference before delivery. No claim is made here about whether it was corrected in the final journal typesetting. A second independent algebraic route uses
\(Q=z^2+3xz+9x^2\), \(V=x^3+a\), \(W=6x\). Substitution in Mironov's equation, with primes meaning \(x\)-derivatives,
\begin{align*}
 4F={}&4(z-W)Q^2-4V(Q')^2+(Q'')^2-2Q'Q'''\\
      &+2Q(2V'Q'+4VQ''+Q'''')
\end{align*}
gives \(4(z^5+27az^2+81)\). Here \(Q'=3z+18x\), \(Q''=18\) and the higher derivatives vanish. This verifies the spectral polynomial against the published construction, separately from operator multiplication.

For the squarefreeness claim, a multiple root cannot be \(z=0\), since \(F_a(0)=81\). The equations \(F_a=F_a'=0\) at a nonzero root reduce to
\[
 z^3=-54a/5,\qquad az^2=-5.
\]
They imply \(a\ne0\), \(z=54a^2/25\), and \(2916a^5+3125=0\). Conversely these formulas produce a common root whenever that parameter equation holds. This proves the exact exceptional locus without trusting a discriminant command. The separately computed discriminant is
\[
 \operatorname{disc}_z(F_a)=3^{16}(2916a^5+3125).
\]
A squarefree quintic gives five finite branch points and one at infinity, hence genus two by Riemann--Hurwitz.
\end{proof}

\paragraph{Independent checks and their scope.}
The second operator check computes direct action on polynomials, without multiplying coefficient dictionaries, at \(a=-2,0,1,3\). For every \(n=0,\ldots,20\), it verifies both residuals on \(x^n\). An ordinary differential operator of order at most \(r\) that annihilates \(1,x,\ldots,x^r\) is zero: evaluation on these monomials successively forces its coefficients to vanish, starting with the constant coefficient. Since the residual orders are at most \(13\) and \(20\), respectively, these action checks are complete at each selected parameter. They are a second check, not a replacement for the all-parameter polynomial identities. A separate conceptual review of the unit obstruction checked injectivity, irreducibility, the order filtration, extension of the map to projective models, and the genus calculation; it found no counterexample or missing hypothesis after those qualifications were made explicit. Both check records are included.

\paragraph{What this calculation does and a feasible further check.}
The unit obstruction rules out certain prescribed affine models, while the certificate illustrates a known family in suitable coordinates. General genus-two existence is already supplied by Oganesyan and must not be advertised as work still to do. A useful bounded follow-up would implement his two-family parameter inversion for a specified quintic, such as \(z^5+z\), print the resulting rational operators and verify their identities exactly. That would make an existing theorem executable. No historical reason for the problem having been dropped is established by the sources. The obstruction relevant to the notebook's wording is the distinction between affine pole data and the abstract curve field.

\subsection{From a commuting pair to a curve field in the first Weyl division ring}
\label{dev:fieldembedding}

This development concerns item 1980-519.3 and problem GP-1980-41. The notebook explicitly adds \ru{т.е. поле рац. ф-й на любой алг. кривой}. Its strongest defensible modern formulation over the complex numbers is:

\begin{conjecture}[Gelfand's function-field question as recorded on P000970]
For every smooth projective irreducible complex algebraic curve \(C\), there is an injective unital \(\mathbb C\)-algebra map
\[
 \mathbb C(C)\hookrightarrow Q(A_1),\qquad
 A_1=\mathbb C\langle x,D\rangle/(Dx-xD-1),
\]
where \(Q(A_1)\) is the Ore division ring of fractions.
\end{conjecture}

The notebook does not specify a base field; \(\mathbb C\), smooth projective normalization and unital base-field-linear embeddings are explicit reconstruction choices. The statement concerns a \emph{commutative subfield} of the rank-one Weyl division ring. It is distinct from the Dixmier endomorphism conjecture, from isomorphisms between Weyl fields of different ranks, and from the Gelfand--Kirillov conjecture for enveloping algebras. The unit obstruction above applies to ordinary differential operators and does not apply to a division ring, in which elements of negative differential order are allowed.

\paragraph{Present scope of the checked results.}
Dixmier's example and the higher-genus polynomial constructions give subfields after an injectivity argument. More strongly, \citet[Theorem~3(1)--(2)]{oganesyan2016rational} settles all curve fields of genus one and two; genus zero is immediate. \citet[Theorem~1.1]{makarlimanov2021centralizers} proves that a rank-one centralizer in \(A_1\), over a characteristic-zero field, embeds in a polynomial ring. That is a genus-zero obstruction within the specified polynomial-operator class and does not classify commutative subfields of \(Q(A_1)\). A recent check also located \citet[abstract]{ivlev2026trigonal}, which announces particular rank-two trigonal genus-three examples, not a theorem for every curve; only its publisher metadata and abstract were accessible. Searches of the primary literature and author bibliographies through November 2026 did not identify a proof or counterexample for the exact universal function-field statement. This negative search result is a limitation of the evidence, not a theorem about the completeness of the literature. We retain C with that qualification.

\begin{lemma}[An order test for injectivity]\ours\quad
\label{lem:orderembedding}
Let \(P,M\in A_1\) commute, with \(\operatorname{ord}(P)=r>0\) and \(\operatorname{ord}(M)=s>0\). Suppose \(r\nmid s\), and \(M^2=F(P)\), where \(F\in\mathbb C[z]\) is not a square in \(\mathbb C(z)\). Then the map
\[
 \mathbb C[z,w]/(w^2-F(z))\longrightarrow A_1,
 \qquad z\longmapsto P,\quad w\longmapsto M
\]
is injective.
\end{lemma}
\begin{proof}
Division by the monic polynomial in \(w\) gives a unique representative \(f(z)+g(z)w\). For nonzero \(f\), the order of \(f(P)\) is \(r\deg f\): the highest power of \(P\) has strictly greater order than all lower powers. Likewise, for nonzero \(g\), the order of \(g(P)M\) is \(r\deg g+s\). These orders are in different residue classes modulo \(r\), so the two summands cannot cancel. If only one summand is present it is nonzero, since \(A_1\) is a domain and powers of a positive-order operator are linearly independent. Thus the kernel is zero. The nonsquare condition ensures that the source is the integral coordinate ring of an irreducible curve, as needed for the field conclusion below.
\end{proof}

\begin{corollary}[An explicit family of genus-two field embeddings]\ours\quad
\label{cor:fieldembedding}
For every \(a\in\mathbb C\) with \(2916a^5+3125\ne0\), the curve
\[
 C_a:\quad w^2=z^5+27az^2+81
\]
has smooth projective genus two, and \(\mathbb C(C_a)\) embeds in \(Q(A_1)\) by \(z\mapsto P\), \(w\mapsto M\) from Proposition~\ref{prop:certificate}.
\end{corollary}
\begin{proof}
A squarefree polynomial of odd degree five is not a square in \(\mathbb C(z)\), because its simple zeros have odd valuations. Apply Lemma~\ref{lem:orderembedding} with \(r=4\) and \(s=10\), since \(4\nmid10\). Every nonzero element of the embedded commutative coordinate ring becomes invertible in the Ore division ring. Define a fraction \(b/c\) to map to \(\iota(b)\iota(c)^{-1}\). Images of the coordinate-ring elements and their inverses commute, so the usual fraction equivalence proves this map well defined; injectivity follows because a nonzero numerator remains nonzero in a division ring. The genus statement was proved in Proposition~\ref{prop:certificate}.
\end{proof}

This corollary supplies the elementary injectivity step that a mere polynomial identity does not provide. It uses a known operator family and makes no claim to resolve Gelfand's universal conjecture. The same argument applies to the corrected elliptic pair with \(r=4,s=6\) and \(a\ne0\). The greatest common divisor of two displayed orders alone is not a proof about the rank of their entire centralizer, and no such inference is used.

\begin{corollary}[The obstructed affine model nevertheless has an embeddable function field]\ours\quad
The function field of \(zw^2-z^4-1=0\) embeds in \(Q(A_1)\), even though its affine coordinate ring does not embed in \(\mathbb C(x)[D]\).
\end{corollary}
\begin{proof}
Use \citet[Theorem~3(2)]{oganesyan2016rational} for \(Y^2=z^5+z\), obtaining commuting rational-coefficient operators \(P_0,M_0\) of orders four and ten. The proof of Lemma~\ref{lem:orderembedding} applies unchanged to \(\mathbb C(x)[D]\), whose order is still additive. This ring is contained in \(Q(A_1)\), since nonzero polynomial coefficients can be inverted there. The resulting field embedding sends the original coordinates to \(z\mapsto P_0\), \(w\mapsto M_0P_0^{-1}\).

In fact \(M_0P_0^{-1}\) cannot be an ordinary differential operator. If it were, the spectral relation would imply
\[
 P_0\bigl((M_0P_0^{-1})^2-P_0^3\bigr)=1
\]
inside \(\mathbb C(x)[D]\), making the positive-order \(P_0\) a unit, contrary to Lemma~\ref{lem:units}. The required quotient is exactly where the two questions diverge.
\end{proof}

\paragraph{Remaining obstruction and a feasible next question.}
The missing quantifier is \emph{every prescribed function field in arbitrary genus}. The order argument is sufficient once appropriate operators have been constructed; it does not construct them for an arbitrary curve. One further bounded run could take a specified squarefree degree-seven hyperelliptic equation, determine its membership in the two rational-operator families of \citet[Theorem~3(3)]{oganesyan2016rational}, and produce an exact certificate when successful. Failure of those families or of a chosen degree bound would not prove nonexistence of a division-ring embedding. The affine unit obstruction likewise cannot be promoted to a negative result for that conjecture. These boundaries are part of the mathematical conclusion.

\section{Problem index}\label{problems:1980}
Earlier-year occurrence identifiers are retained as text; current-year identifiers link to their entries.

\subsection{GP-1967-21: The unextended semisimple Gelfand--Kirillov conjecture}\label{prob:GP-1967-21}
For a finite-dimensional semisimple Lie algebra over an algebraically closed field of characteristic zero, is its enveloping quotient division ring a Weyl division ring over a purely transcendental central field? The P929 formulation also gestures toward the wider algebraic-Lie-algebra setting.
\par\noindent\textit{Occurrences:} 1967-144.7, \hyperref[item:1980-497.2]{1980-497.2}.
\par\noindent\textit{Status (2026-11):} As of November 2026, The general statement is false; type A is established, and the reviewed literature still leaves simple types \(C_n\), \(n\ge2\), and \(G_2\) unresolved. Reuse the existing problem identity GP-1967-21; this row supplies the additional occurrence, not a replacement of earlier occurrences. A finite central extension changes the question.

\subsection{GP-1980-01: Projective structures on the real line}\label{prob:GP-1980-01}
Classify connected real projective one-manifolds diffeomorphic to \(\mathbb R\), up to projective diffeomorphism.
\par\noindent\textit{Occurrences:} \hyperref[item:1980-494.4]{1980-494.4}.
\par\noindent\textit{Status (2026-11):} As of November 2026, Resolved: the full universal projective line, a half-line, and the positive half-integral winding classes of bounded intervals; see \citet[Corollary 4.16]{belgunmoroianu2025projective}, which attributes the open-curve result to Kuiper (1954).

\subsection{GP-1980-02: Levi-flat filling of real two-spheres}\label{prob:GP-1980-02}
For a specified smooth embedded sphere \(S\subset\mathbb C^2\), determine whether it bounds a Levi-flat hypersurface foliated by analytic discs.
\par\noindent\textit{Occurrences:} \hyperref[item:1980-495.1]{1980-495.1}.
\par\noindent\textit{Status (2026-11):} As of November 2026, The two-elliptic-point sphere in a strictly pseudoconvex boundary is fillable by Bedford--Gaveau (1983); arbitrary spheres are not universally fillable. \citet[Theorem 1.1 and Introduction]{gaussiersukhov2011filling} supplies a precise later extension and records the counterexamples. This entry does not assert an exhaustive classification of all embeddings.

\subsection{GP-1980-03: Categorical characterization of D-affinity}\label{prob:GP-1980-03}
For a smooth separated finite-type complex variety, does equivalence of the natural global-sections/localization adjunction for quasi-coherent left D-modules imply vanishing of all higher sheaf cohomology and generation by global sections?
\par\noindent\textit{Occurrences:} \hyperref[item:1980-496.8]{1980-496.8}, \hyperref[item:1980-498.2]{1980-498.2}, \hyperref[item:1980-498.3]{1980-498.3}.
\par\noindent\textit{Status (2026-11):} As of November 2026, Yes. Beilinson’s 10 November continuation, 1980-498.2--3, gives the projective-generator and \v{C}ech argument. This is a resolved question, not a present open problem.

\subsection{GP-1980-04: Infinite-dimensional algebras with a multiplicative norm}\label{prob:GP-1980-04}
Does a real infinite-dimensional normed vector space admit a continuous bilinear multiplication satisfying \(\|xy\|=\|x\|\|y\|\)? Does the answer change when a two-sided identity is required?
\par\noindent\textit{Occurrences:} \hyperref[item:1980-499.4]{1980-499.4}.
\par\noindent\textit{Status (2026-11):} As of November 2026, Yes without a unit, by the Hilbert tensor-coordinate construction; no with a unit, by Urbanik--Wright (1960). The notebook does not choose the unit convention, so this is a resolved disambiguation rather than a currently open problem.

\subsection{GP-1980-20: Continuity and differentiability of Thomae functions}\label{prob:GP-1980-20}
For \(T_\alpha(p/q)=q^{-\alpha}\) in lowest terms and \(T_\alpha(x)=0\) for irrational \(x\), determine continuity and differentiability, including \(\alpha=1\).
\par\noindent\textit{Occurrences:} \hyperref[item:1980-501.1]{1980-501.1}, \hyperref[item:1980-501.2]{1980-501.2}.
\par\noindent\textit{Status (2026-11):} As of November 2026, this is a solved classical exercise. For \(\alpha>0\), continuity holds exactly at irrational points; for \(\alpha\leq0\) there are no continuity points. For \(0<\alpha\leq2\) there are no differentiability points. For \(\alpha>2\), at an irrational point with convergent denominators \(q_k\), differentiability is equivalent to \(q_{k+1}/q_k^{\alpha-1}\to0\), and the derivative is zero. The criterion is supplied by the elementary convergent estimate in the companion; \citet{varona2009differentiability} documents the classical general-power problem and its earlier literature.

\subsection{GP-1980-21: A single continuous-function limit for the Dirichlet function}\label{prob:GP-1980-21}
Can the indicator of the rational numbers on \(\mathbb R\), already expressed by an iterated pointwise limit in the notes, be the pointwise limit of one sequence of continuous real-valued functions?
\par\noindent\textit{Occurrences:} \hyperref[item:1980-501.4]{1980-501.4}.
\par\noindent\textit{Status (2026-11):} As of November 2026, the answer is no: Baire-one real functions on \(\mathbb R\) have a dense set of continuity points, whereas this indicator is nowhere continuous. This is an application of the classical Baire criterion, read in \citet[p.173]{baire1900nouvelle}.

\subsection{GP-1980-22: Large-frequency asymptotics of Gaussian, quartic, and sextic Fourier integrals}\label{prob:GP-1980-22}
Determine the asymptotics as \(\lambda\to+\infty\) of \(\int_{\mathbb R}e^{-x^p}\cos(\lambda x)\,dx\) for \(p=2,4,6\).
\par\noindent\textit{Occurrences:} \hyperref[item:1980-501.5]{1980-501.5}.
\par\noindent\textit{Status (2026-11):} As of November 2026, this is a solved steepest-descent exercise. The Gaussian is \(\sqrt\pi e^{-\lambda^2/4}\); the higher powers have exponentially damped oscillatory leading terms from a conjugate pair of dominant saddles, explicitly reconstructed in the occurrence. The method and branch convention are checked against \citet[\S2.4(iv)]{nist2026dlmf}. No new asymptotic theorem is claimed.

\subsection{GP-1980-23: Isomorphic cohomology without a direct chain map}\label{prob:GP-1980-23}
Find complexes with isomorphic cohomology for which a prescribed cohomology isomorphism cannot be induced by a chain map in the stated direction.
\par\noindent\textit{Occurrences:} \hyperref[item:1980-502.1]{1980-502.1}.
\par\noindent\textit{Status (2026-11):} As of November 2026, this is a solved exercise: \(A=(\mathbb Z/2)[0]\) and \(B=[\mathbb Z\xrightarrow{2}\mathbb Z]\) in degrees \(-1,0\) have the same cohomology, but every map \(A\to B\) is zero. Their derived-category isomorphism is represented by the reverse quasi-isomorphism, inverted by localization. Earlier occurrence 1971-266.3 concerns the general derived-category universal property, not the exact counterexample question, so its problem identifier is not reused.

\subsection{GP-1980-24: Genus from a Newton polygon}\label{prob:GP-1980-24}
Determine the genus of the smooth projective model of a generic irreducible plane Laurent curve with prescribed two-dimensional Newton polygon.
\par\noindent\textit{Occurrences:} \hyperref[item:1980-509.4]{1980-509.4}.
\par\noindent\textit{Status (2026-11):} As of November 2026, the problem is solved under Newton nondegeneracy: the genus is the number of interior lattice points. Khovanskii's 1978 publication predates this exposition; the question is a motivating problem, not an open 1981 conjecture.

\subsection{GP-1980-25: Root counts from prescribed Newton polytopes}\label{prob:GP-1980-25}
For a generic square system of Laurent polynomials with prescribed Newton polytopes, determine the number of roots in the algebraic torus, with multiplicity.
\par\noindent\textit{Occurrences:} \hyperref[item:1980-509.9]{1980-509.9}.
\par\noindent\textit{Status (2026-11):} As of November 2026, Bernshtein's 1975 theorem gives \(n!\) times ordinary mixed volume, with equality under the appropriate initial-system nondegeneracy. The notebook presents a problem already solved before this lecture. Earlier local Newton-polyhedron questions are related but not exact recurrences.

\subsection{GP-1980-26: Newton-type asymptotic root count for complex exponential sums}\label{prob:GP-1980-26}
Extend the Laurent-polynomial root count to the density of roots of a regular square system of finite exponential sums with complex spectra.
\par\noindent\textit{Occurrences:} \hyperref[item:1980-513.2]{1980-513.2}.
\par\noindent\textit{Status (2026-11):} As of November 2026, this formulation is solved by Kazarnovskii's 1981 regular-system theorem, with leading coefficient \(\kappa_n n!P_n/(2\pi)^n\) and remainder \(O(R^{n-1})\). The theorem requires its face-system regularity; arbitrary fixed systems are outside the assertion.

\subsection{GP-1980-27: Relating exponential zero density to the growth indicatrix}\label{prob:GP-1980-27}
Explain how the exponential-sum zero-density theorem is expressed through the support functions of the spectral convex bodies, as asked by I. M. Gelfand on 2 March 1981.
\par\noindent\textit{Occurrences:} \hyperref[item:1980-513.6]{1980-513.6}.
\par\noindent\textit{Status (2026-11):} As of November 2026, the finite-spectrum connection is established through the mixed Monge--Ampère integral already in \citet[\S1(c)]{kazarnovskii1981zeros}. \citet{kazarnovskii2024distribution} extends this framework to averaged zero currents of regular finite-dimensional spaces of entire functions of exponential growth. The later result is not a universal fixed-system convergence theorem, and the notebook does not specify such a conjecture.

\subsection{GP-1980-40: Genus-two spectral curves of rational commuting differential operators}\label{prob:GP-1980-40}
For a smooth projective genus-two curve over \(\mathbb C\), choose a hyperelliptic model \(w^2=z^5+\beta_4z^4+\cdots+\beta_0\). There exist commuting rational-coefficient ordinary differential operators \(L,M\) whose spectral curve is this model and whose generic common eigenspace has dimension two. The coordinate map is faithful. If instead one insists on every given affine equation \(P\), including a presentation with a nonconstant unit, the analogous ordinary-operator assertion is false; the notebook does not distinguish these readings.
\par\noindent\textit{Occurrences:} \hyperref[item:1980-519.2]{1980-519.2}.
\par\noindent\textit{Status (2026-11):} As of November 2026, every smooth projective complex genus-two curve, and every standard monic quintic hyperelliptic model, is realized by rational commuting operators. This is an established result. Universal realization of every affine presentation is false, as the chapter demonstrates using \(zw^2-z^4-1=0\); this distinction is not promoted to an artificial residual open problem.

\subsection{GP-1980-41: Curve function fields in the first Weyl division ring}\label{prob:GP-1980-41}
For every smooth projective integral curve \(C/\mathbb C\), does there exist a unital \(\mathbb C\)-embedding \(\mathbb C(C)\hookrightarrow\operatorname{Frac}(A_1(\mathbb C))\), where \(A_1=\mathbb C\langle x,D\rangle/(Dx-xD-1)\)? Singular integral curves give the same question after normalization.
\par\noindent\textit{Occurrences:} \hyperref[item:1980-519.3]{1980-519.3}.
\par\noindent\textit{Status (2026-11):} As of November 2026, all complex curves of genus at most two and explicit higher-genus families have positive realizations. No source proving or disproving the universal curve-function-field conjecture was located in the documented search. This is a search limitation, not a claim that no such theorem exists.

\subsection{GP-1980-42: Spectra of sums of Hermitian matrices and tensor-product multiplicities}\label{prob:GP-1980-42}
For real weakly decreasing \(n\)-tuples \(x,y,z\), characterize the existence of complex Hermitian \(n\times n\) matrices \(X,Y\) whose eigenvalue lists are \(x,y\) and whose sum has eigenvalue list \(z\). Trace equality is necessary and is retained in the final characterization.
\par\noindent\textit{Occurrences:} \hyperref[item:1980-526.1]{1980-526.1}, \hyperref[item:1980-526.2]{1980-526.2}, \hyperref[item:1980-526.3]{1980-526.3}.
\par\noindent\textit{Status (2026-11):} As of November 2026, the Horn problem is solved by the recursive inequalities, equivalently by the honeycomb and representation-theoretic criteria.

\section{Corrections to the transcription and translation}\label{corr:1980}
The evidence distinguishes a transcription correction from a mathematical emendation of wording faithfully preserved in the notebook.
\subsection*{P000916 (both; provenance): Academic-year boundary}
\noindent\textit{Transcription reads.} The catalogue includes this page in records 490--491.\par
\noindent\textit{Critical reading.} Treat the Benney/Zakharov material as the continuation already indexed in 1979-488.1--7, not as a statement by Fuchs or Graev.\par
\noindent\textit{Evidence.} The first explicit date and Fuchs heading occur on P000917; the supplied previous index identifies every mathematical theme on P000916.\par
\subsection*{P000917 (both; faithful error): Filtration and formula (2)}
\noindent\textit{Transcription reads.} \(X_{n-3},X_{n-4}\) after \(X_{2n-2}\); \(\dim(\sigma_j\cap X_{2n-i})=j-2i\).\par
\noindent\textit{Critical reading.} Use consistently real-dimensional skeleta and \(\dim(|\xi|\cap X_{2n-c})\le j-c+\bar m(c)\), also for \(\partial\xi\).\par
\noindent\textit{Evidence.} Scan inspected: the anomalous indices, equality sign and \(j-2i\) really are present. This is a mathematical emendation, not a mistranscription.\par
\subsection*{P000918 (both; faithful error): Curve coordinates and incidence equations}
\noindent\textit{Transcription reads.} \(\xi=(\xi_1,\ldots,\xi_n)\), then \((\xi_1,\ldots,\xi_{n+k-1})\); \(f(x_1,\varphi_1,\ldots,\varphi_n)\); \(x_i^0=\varphi_i(x_i^0,\xi)\).\par
\noindent\textit{Critical reading.} Use \(\xi\in\mathbb R^{n+k-1}\), \(\varphi_2,\ldots,\varphi_n\), and \(x_i^0=\varphi_i(x_1^0,\xi)\), \(2\le i\le n\).\par
\noindent\textit{Evidence.} The scan confirms the incompatible indices. The corrected dimensional bookkeeping is corroborated by Gelfand--Graev, equations (2)--(4).\par
\subsection*{P000918 (both; provenance): FA references}
\noindent\textit{Transcription reads.} Both issue references are placed after Gelfand--Gindikin--Shapiro.\par
\noindent\textit{Critical reading.} FA 13(2) is Gelfand--Gindikin--Shapiro (1979); FA 14(4) is Gelfand--Graev (1980).\par
\noindent\textit{Evidence.} The abbreviated author line and issue numbers are visible; both primary Math-Net records and full articles were opened.\par
\subsection*{P000920 (both; faithful error): Homogeneous incidence example}
\noindent\textit{Transcription reads.} \(A=G\), with joint injectivity of \(G\to H\backslash G\times K\backslash G\).\par
\noindent\textit{Critical reading.} Replace \(A\) by \((H\cap K)\backslash G\), or impose \(H\cap K=\{1\}\).\par
\noindent\textit{Evidence.} The two coset projections identify precisely the left cosets of the intersection; the notebook has not imposed trivial intersection.\par
\subsection*{P000920 (both; faithful error): Lower-index coframe transformation (18)}
\noindent\textit{Transcription reads.} The lower-index transformation has incompatible repeated indices and no visible final \(\theta\) factor.\par
\noindent\textit{Critical reading.} Use \(\widehat\theta_\alpha=B_\alpha^\beta\theta_\beta+c_\alpha\theta\).\par
\noindent\textit{Evidence.} Scan P000920 inspected; the malformed expression is faithfully reproduced. This is the type-correct form of the block action.\par
\subsection*{P000921 (both; faithful error): Coframe transformations (21)}
\noindent\textit{Transcription reads.} \(\theta^*=u\theta^*\) and a \(u^\alpha\theta^\alpha\) term.\par
\noindent\textit{Critical reading.} Use one common rescaled form and a multiple of \(\theta\) in each off-diagonal term, with invertible diagonal blocks.\par
\noindent\textit{Evidence.} Scan inspected; the self-starred expression and repeated upper-index term are present. The adjacent matrix resolves the intended block transformation.\par
\subsection*{P000922 (both; faithful error): Allowability formula (3)}
\noindent\textit{Transcription reads.} \(\dim(\sigma^q\cap X_{n-i})\le q-i\), with \(\operatorname{codim}_{\mathbb R}X_{n-i}=2i\).\par
\noindent\textit{Critical reading.} For middle perversity use \(q-i-1\), and impose the corresponding boundary restriction.\par
\noindent\textit{Evidence.} Scan inspected: the right side is \(q-i\) and the real-codimension \(2i\) parenthesis is explicit. This is not an OCR correction.\par
\subsection*{P000923 (both; faithful error): Unique irreducible submodule, assertion (6)}
\noindent\textit{Transcription reads.} A unique irreducible submodule for the extension of the regular locus, with a cross as example.\par
\noindent\textit{Critical reading.} Require irreducible support and an irreducible local system; for a cross use componentwise intermediate extensions.\par
\noindent\textit{Evidence.} The scan explicitly draws the cross and asserts uniqueness/irreducibility. Its regular locus has two components.\par
\subsection*{P000924 (both; faithful error): General case, finite-length statement}
\noindent\textit{Transcription reads.} Every module has a finite filtration with irreducible factors.\par
\noindent\textit{Critical reading.} Use coherent holonomic modules, locally in the analytic setting or globally in the algebraic finite-type setting.\par
\noindent\textit{Evidence.} The scan says every module and finite length. Infinite direct sums already refute the unrestricted statement.\par
\subsection*{P000924 (both; faithful error): Torsion condition (13)}
\noindent\textit{Transcription reads.} \(f\in\mathcal D\), \(f\ne0\), \(x\ne0\Rightarrow fx\ne0\), hence a locally constant system.\par
\noindent\textit{Critical reading.} Replace the criterion by coherence over \(\mathcal O\) for a coherent D-module; D-torsion-freeness is false as a characterization.\par
\noindent\textit{Evidence.} The original scan clearly has \(f\in\mathcal D\). \(\mathcal O\) has the nonzero annihilator \(\partial\) acting on \(1\), and \(\mathcal D\) itself supplies the opposite counterexample.\par
\subsection*{P000926 (both; ambiguity): Countability and atlas equivalence}
\noindent\textit{Transcription reads.} Countably many locally projective structures after defining equivalent atlases by compatibility.\par
\noindent\textit{Critical reading.} Countability holds for isomorphism classes under projective diffeomorphisms, not equality of structures on a fixed line.\par
\noindent\textit{Evidence.} The scan’s compatibility definition and countability conclusion concern different equivalence levels unless a diffeomorphism quotient is supplied.\par
\subsection*{P000926 (both; faithful error): Covering-group pairs}
\noindent\textit{Transcription reads.} \(g_j\in SL(2,\mathbb C)\), followed by signs of entries and a covering of \(SL(2,\mathbb R)\).\par
\noindent\textit{Critical reading.} Read the intended base group as \(SL(2,\mathbb R)\); retain the unresolved normalization of the triple-line cocycle explicitly.\par
\noindent\textit{Evidence.} The scan clearly has \(\mathbb C\), while the preceding question, inequalities and real-line geometry all require \(\mathbb R\).\par
\subsection*{P000926 (both; faithful error): Parenthetical torus assertion}
\noindent\textit{Transcription reads.} This problem cannot be solved for a torus.\par
\noindent\textit{Critical reading.} No topology-only obstruction exists: the product solid torus in \(\mathbb C^2\) is Levi-flat and disc-foliated; further boundary or hull restrictions would be needed.\par
\noindent\textit{Evidence.} Scan inspected. The parenthetical claim is explicit, and the elementary product example refutes its unrestricted reading.\par
\subsection*{P000928 (both; faithful error): Exponential estimate (13)--(14)}
\noindent\textit{Transcription reads.} \(|\Delta^{-1}V|\le C<1\).\par
\noindent\textit{Critical reading.} Use \(\sup_x\mathbb E_x\int_0^\infty V_-(B_s)ds\le C<1\), with the Green operator normalized to the process.\par
\noindent\textit{Evidence.} The scan contains an absolute value around the signed Green potential, not the negative-part potential. Khasminskii’s estimate needs the latter.\par
\subsection*{P000928 (both; ambiguity): Large-time formula (12)}
\noindent\textit{Transcription reads.} \(\pm\lambda_1\).\par
\noindent\textit{Critical reading.} For \(H=-\tfrac12\Delta+V\), the displayed logarithmic expectation limit is \(-\lambda_0(H)\), under the required spectral hypotheses.\par
\noindent\textit{Evidence.} Scan inspected: the plus-or-minus sign is genuine. The semigroup convention in (10) determines the negative sign.\par
\subsection*{P000928 (both; faithful error): Stochastic equation (9)}
\noindent\textit{Transcription reads.} \(d\Pi_t=e^{-\tau\Box_\tau}d\beta_t\).\par
\noindent\textit{Critical reading.} The operator depends on \(\Pi_t\); distinguish the heat time in a noise coefficient from the heat time in its covariance, and retain the It\^o factor \(1/2\).\par
\noindent\textit{Evidence.} The scan has the tau subscript. Multiplying the heat operator by its adjoint doubles its time parameter.\par
\subsection*{P000930 (both; faithful error): Commutative ring equals affine algebraic variety}
\noindent\textit{Transcription reads.} A commutative ring is an affine algebraic variety.\par
\noindent\textit{Critical reading.} Use affine scheme for arbitrary rings; classical complex affine varieties require finitely generated reduced complex algebras.\par
\noindent\textit{Evidence.} The manuscript explicitly moves from varieties to Spec of an unrestricted ring.\par
\subsection*{P000931 (both; faithful error): Formal adjoint (16)}
\noindent\textit{Transcription reads.} \(f\partial_x\mapsto\partial_x f\).\par
\noindent\textit{Critical reading.} \((f\partial_x)^*=-\partial_x f\) for the usual formal transpose against constant volume.\par
\noindent\textit{Evidence.} Scan inspected: no minus sign is visible. The Weyl commutation relation forces the sign for an anti-involution fixing functions.\par
\subsection*{P000931 (both; faithful error): Negative-twist cohomology (13)}
\noindent\textit{Transcription reads.} Every \(\mathcal O(-i)\) on \(\mathbb P^N\) has higher cohomology.\par
\noindent\textit{Critical reading.} For \(i>0\), nonzero higher cohomology occurs only when \(i\ge N+1\); twists \(1\le i\le N\) are acyclic.\par
\noindent\textit{Evidence.} Scan inspected: the symbol is \(i\) and no range beyond negativity is stated.\par
\subsection*{P000932 (both; faithful error): Canonical-class obstruction (20)--(22)}
\noindent\textit{Transcription reads.} Vanishing first Chern class implies non-D-affinity, without properness assumptions.\par
\noindent\textit{Critical reading.} Supply an appropriate projective/compact hypothesis and specify which characteristic class; affine space refutes the unrestricted statement.\par
\noindent\textit{Evidence.} The scan has no projectivity hypothesis at this point; the prior definition explicitly includes affine varieties.\par
\subsection*{P000932 (both; faithful error): Interlinear line-bundle assertion}
\noindent\textit{Transcription reads.} Invertible sheaves are isomorphic if and only if their first Chern classes agree.\par
\noindent\textit{Critical reading.} Equal first Chern classes need not give holomorphically or algebraically isomorphic line bundles.\par
\noindent\textit{Evidence.} The interlinear assertion is visible; nontrivial Picard-zero line bundles on an elliptic curve give counterexamples.\par
\subsection*{P000932 (both; faithful error): Last projective-line sentence}
\noindent\textit{Transcription reads.} There is a unique operator commuting with all three vector fields.\par
\noindent\textit{Critical reading.} The enveloping algebra has one polynomial central generator; after the trivial-character quotient the commuting center is the scalar field.\par
\noindent\textit{Evidence.} The source’s uniqueness wording is explicit and cannot hold literally, since every scalar commutes.\par
\subsection*{P000933 (both; ambiguity): Finite and infinite absolute-valued algebras}
\noindent\textit{Transcription reads.} Only \(\mathbb R\), \(\mathbb C\), \(\mathbb H\) and Cayley in finite dimensions; an infinite-dimensional example exists.\par
\noindent\textit{Critical reading.} Distinguish the unit axiom. Unital algebras are precisely \(\mathbb R\), \(\mathbb C\), \(\mathbb H\), \(\mathbb O\) even without a dimension hypothesis; nonunital algebras admit infinite-dimensional examples and extra finite-dimensional isotopes.\par
\noindent\textit{Evidence.} Scan inspected: only the multiplicative norm condition is written, with no identity or division-surjectivity axiom. The ambiguity changes the mathematical answer.\par
\subsection*{P000934 (both; faithful error): Constant sheaf in localization (2)}
\noindent\textit{Transcription reads.} \(M\) is a constant sheaf of \(\mathcal A\)-modules.\par
\noindent\textit{Critical reading.} \(M\) defines a constant sheaf of modules over \(\underline A\), where \(A=\Gamma\mathcal A\); extension of scalars gives the \(\mathcal A\)-module.\par
\noindent\textit{Evidence.} Scan inspected: the module-ring symbol is calligraphic \(\mathcal A\). A general \(A\)-module has no pre-existing action by local sections of the sheaf ring.\par
\subsection*{P000935 (both; faithful error): Endomorphism ring (11)--(12)}
\noindent\textit{Transcription reads.} Left \(\operatorname{End}(P)\)-modules and \(\operatorname{End}_{\mathcal A}(\mathcal A)=A\).\par
\noindent\textit{Critical reading.} Use right \(\operatorname{End}(P)\)-modules, or left modules over its opposite; for a left regular module the endomorphism ring is \(A^{\mathrm{op}}\).\par
\noindent\textit{Evidence.} The left-module convention is explicit on P933--934, while the scan of (12) has no opposite-ring symbol.\par
\subsection*{P000936 (both; ambiguity): Last line after (25)}
\noindent\textit{Transcription reads.} \(df(v)\) acts on \(f^{-1}\mathcal D_Y\).\par
\noindent\textit{Critical reading.} The coefficients of \(df(v)\) lie in \(\mathcal O_X\); the induced action is on the scalar-extended transfer module.\par
\noindent\textit{Evidence.} The scan double-underlines the inverse-image expression without the \(\mathcal O_X\) tensor factor. The preceding correct formula (24) supplies the needed module.\par
\subsection*{P000938 (both; faithful error): delta-module pullback example}
\noindent\textit{Transcription reads.} \(\delta(y)=y\delta^{\prime}(y)\)\par
\noindent\textit{Critical reading.} \(\delta(y)=-y\delta^{\prime}(y)\)\par
\noindent\textit{Evidence.} The scan has no minus sign. Differentiating \(y\delta=0\) gives \(\delta+y\delta^{\prime}=0\). Faithful mathematical emendation.\par
\subsection*{P000940 (both; faithful error): Fuchs Whitehead-product target}
\noindent\textit{Transcription reads.} \(\pi_{n+m}(X)\)\par
\noindent\textit{Critical reading.} \(\pi_{n+m-1}(X)\)\par
\noindent\textit{Evidence.} The actual scan has \(n+m\). The boundary of \(I^n\times I^m\) has dimension \(n+m-1\); Hatcher section4.2 p381.\par
\subsection*{P000940 (both; ambiguity): Lomonosov second commutator and spectral subspaces}
\noindent\textit{Transcription reads.} \([B,A^{\prime}]=0\); \(\mathcal H_{ac}(B),\mathcal H_{ac}(B^{\prime})\)\par
\noindent\textit{Critical reading.} Leave the literal text; annotate the possible repair \([B^{\prime},A^{\prime}]=0\), \(\mathcal H_{ac}(A),\mathcal H_{ac}(A^{\prime})\).\par
\noindent\textit{Evidence.} Scan explicitly shows \(B\) without a prime before \(A^{\prime}\) and \(B\) in the spectral-subspace labels. No textual authority for silently replacing them; item 1980-505.1 distinguishes the reconstruction.\par
\subsection*{P000941 (both; ambiguity): de Rham shift after equation10}
\noindent\textit{Transcription reads.} shifted to the left by \(\mathcal M\)\par
\noindent\textit{Critical reading.} shifted to the left by \(\dim X\), as a mathematical reconstruction\par
\noindent\textit{Evidence.} The symbol resembles the module \(M\) in the scan; the next displayed equation fixes the intended dimensional shift.\par
\subsection*{P000942 (both; faithful error): exactness for an embedding}
\noindent\textit{Transcription reads.} \(f\) an embedding\par
\noindent\textit{Critical reading.} \(f\) a closed embedding of smooth varieties\par
\noindent\textit{Evidence.} Both exactness claims are the closed-embedding theorem. They do not hold for arbitrary open embeddings in the same form.\par
\subsection*{P000943 (both; faithful error): final irreducibility line}
\noindent\textit{Transcription reads.} \(Y\) smooth implies already irreducible\par
\noindent\textit{Critical reading.} For a closed embedding, a simple module stays simple; otherwise use intermediate extension.\par
\noindent\textit{Evidence.} The scan contains no closedness condition. The open immersion \(\mathbb G_m\) into \(\mathbb A^1\) gives \(\mathcal O(*0)\), with the proper submodule \(\mathcal O\).\par
\subsection*{P000944-P000948 (both; provenance): Catalog range for record 506}
\noindent\textit{Transcription reads.} Supplied manifest: P000942-P000950.\par
\noindent\textit{Critical reading.} The Novikov heading starts P000944, and the talk and concluding assignment end P000948.\par
\noindent\textit{Evidence.} P000942-P000943 continue Beilinson; P000949 starts MacPherson on 9 February 1981. The overlap is retained as provenance, not attributed to Novikov.\par
\subsection*{P000944; P000946; P000948 (both; faithful error): Curvature sign in (2), (27), and the irrational-case formula}
\noindent\textit{Transcription reads.} \(H_{12}=\partial_2A_1-\partial_1A_2\), but \(A_1=-\varphi_y,A_2=\varphi_x,\Delta\varphi=H\).\par
\noindent\textit{Critical reading.} With the displayed gauge, \(\partial_2A_1-\partial_1A_2=-\Delta\varphi\). Either reverse the definition of \(H_{12}\) or write \(\Delta\varphi=-H_{12}\), adjusting the spin and flux signs consistently.\par
\noindent\textit{Evidence.} The scan confirms all signs. Direct differentiation detects the mismatch; it is not resolved by changing the exponential parameter \(a\).\par
\subsection*{P000944 (both; faithful error): Last line, rational-flux reduction}
\noindent\textit{Transcription reads.} \ru{Рац. свод. к нему укрупнением решётки}; English: rational reduces to it, immediately after irrational.\par
\noindent\textit{Critical reading.} Rational flux reduces to integral flux by passing to a larger lattice cell.\par
\noindent\textit{Evidence.} The scan faithfully contains the pronoun. Multiplication by a lattice index cannot turn a rational into an irrational number; the cited paper explicitly says integral (\citealp[Sec.~I]{dubrovinnovikov1980ground}).\par
\subsection*{P000945 (both; faithful error): Cell-area estimates and superconducting integrality}
\noindent\textit{Transcription reads.} Crystal: \(10^{-10}\,\mathrm{cm}^2\); superconducting lattice: \(10^{-16}\,\mathrm{cm}^2\); flux automatically integral.\par
\noindent\textit{Critical reading.} The examples are reversed: the source gives crystal area about \(10^{-16}\,\mathrm{cm}^2\) and superconducting area \(10^4\)--\(10^6\) times larger. Integrality is for the effective charge \(2e_0\); an electron sees a half quantum per such cell.\par
\noindent\textit{Evidence.} Both exponents are clear in the scan; \citet[Sec.~I, Remark]{dubrovinnovikov1980ground} supplies the contrary values and the charge qualification.\par
\subsection*{P000946 (both; faithful error): Equation (25), Chern class of the reciprocal-torus bundle}
\noindent\textit{Transcription reads.} \(N\)-dimensional bundle, \(c_1=N\).\par
\noindent\textit{Critical reading.} Rank \(N\), \(c_1=1\) in Novikov's convention; absolute Chern number \(1\).\par
\noindent\textit{Evidence.} Scan P000946 unequivocally has \(c_1=N\). \citet[Sec.~II, Lemma~1]{novikov1981magnetic}, received 2 December 1980, states \(c_1=1\). The physical line bundle and the reciprocal-torus eigenbundle were conflated.\par
\subsection*{P000949 (both; faithful error): Grassmannian inclusion table, second row third entry}
\noindent\textit{Transcription reads.} \(G_1^3\)\par
\noindent\textit{Critical reading.} \(G_2^3\)\par
\noindent\textit{Evidence.} The scan and both transcriptions agree on \(G_1^3\). The second row parametrizes two-planes, so this is a faithful index error.\par
\subsection*{P000950 (both; faithful error): Theorem 5 global regularity}
\noindent\textit{Transcription reads.} diffeomorphic\par
\noindent\textit{Critical reading.} homeomorphic, with a diffeomorphism on each orbit/face interior\par
\noindent\textit{Evidence.} P000953 explicitly qualifies the global statement and mentions possible infinite curvature at vertices; \citet[Theorem~2']{atiyah1982convexity}.\par
\subsection*{P000951 (both; faithful error): volume formula for \(0<c_1<c_2<1\)}
\noindent\textit{Transcription reads.} \(\Phi(c_1)-\Phi(c_2)\)\par
\noindent\textit{Critical reading.} For unsigned positive volume use the opposite difference, or specify an oriented integral and its normalization.\par
\noindent\textit{Evidence.} The scan confirms this order. \(\Phi^{\prime}\) is strictly positive on \((0,1)\), so the recorded difference is negative.\par
\subsection*{P000952 (both; faithful error): harmonic-form vanishing theorem 13}
\noindent\textit{Transcription reads.} every harmonic form on \(G_n^m\)\par
\noindent\textit{Critical reading.} positive-degree harmonic forms on the unoriented real Grassmannian, with \(nm\) even and the invariant metric\par
\noindent\textit{Evidence.} \citet[pp.~13--14]{atiyah1982convexity} explicitly states the parity hypothesis and failures for odd \(nm\) and oriented Grassmannians. Constants also refute the all-degrees reading.\par
\subsection*{P000952 (both; ambiguity): incomplete hypersimplex boundary formula}
\noindent\textit{Transcription reads.} second sum starts at -1 and stops after the hypersimplex label\par
\noindent\textit{Critical reading.} Reconstruct the second family as \(x_i=1\), \(i=1,\ldots,n+m\); do not silently supply orientation signs.\par
\noindent\textit{Evidence.} Actual formula is incomplete. The inequalities give the two facet families; no full signed formula survives.\par
\subsection*{P000952 (both; faithful error): one harmonic four-form; tautological two- or three-bundle}
\noindent\textit{Transcription reads.} exactly one form; either bundle without a sign distinction\par
\noindent\textit{Critical reading.} a one-dimensional space; \(p_1\) of the complementary bundles differs by a sign\par
\noindent\textit{Evidence.} Forms can be rescaled; Whitney product for the trivial rank-five sum gives \(p_1(E)+p_1(Q)=0\).\par
\subsection*{P000953 (both; faithful error): Atiyah theorem 15 hypotheses and orbit diagram}
\noindent\textit{Transcription reads.} commuting Hamiltonians plus a common critical point; Hamiltonian orbit maps to a polyhedron\par
\noindent\textit{Critical reading.} compact connected symplectic \(M\) with almost-periodic Hamiltonian flows; \(f(M)\) is convex. Positive complexified orbits belong to the Kähler theorem.\par
\noindent\textit{Evidence.} Actual scan labels \(H_{f_i}\). Commuting Hamiltonians are constant along their Hamiltonian orbits; \((\cos x,\cos^2 x)\) on a symplectic torus gives a nonconvex image under the written hypotheses.\par
\subsection*{P000953 (both; faithful error): torus-orbit polytope assertion}
\noindent\textit{Transcription reads.} every positive-real orbit is homeomorphic to a polyhedron\par
\noindent\textit{Critical reading.} every positive-real orbit closure is homeomorphic to its moment polytope; the orbit maps to the relative interior\par
\noindent\textit{Evidence.} The missing closure is substantive; \citet[Theorem~2']{atiyah1982convexity} concerns the closure.\par
\subsection*{P000954 (both; faithful error): reference to Toroidal Embeddings I, Lecture Notes 339}
\noindent\textit{Transcription reads.} 1972\par
\noindent\textit{Critical reading.} 1973\par
\noindent\textit{Evidence.} \citet[bibliography, no.~26]{danilov1978toric} explicitly lists Lecture Notes in Mathematics 339 (1973); this corrects the seminar bibliographic year.\par
\subsection*{P000955 (both; faithful error): comparison after plane genus formula}
\noindent\textit{Transcription reads.} \(y^2=P_n(x)\) is not an equation of degree \(n\)\par
\noindent\textit{Critical reading.} For \(n\ge2\) it has total degree \(n\), but is a special plane model, generally singular at infinity.\par
\noindent\textit{Evidence.} Actual scan confirms the literal wording; its total-degree assertion is false. The intended Newton-support distinction remains valid.\par
\subsection*{P000956 (both; faithful error): final edge count}
\noindent\textit{Transcription reads.} number of integer points on this segment\par
\noindent\textit{Critical reading.} lattice length = number of lattice points minus one\par
\noindent\textit{Evidence.} The same scan earlier states length correctly. A binomial edge of primitive length1 already has one root and two endpoints.\par
\subsection*{P000957 (both; faithful error): divisor of \(Q\) restricted to the compact curve}
\noindent\textit{Transcription reads.} its divisor is zero\par
\noindent\textit{Critical reading.} its divisor has degree zero\par
\noindent\textit{Evidence.} The scan uses the word divisor without degree; principal divisors need not vanish as divisors.\par
\subsection*{P000957 (both; faithful error): equation11 local expansion}
\noindent\textit{Transcription reads.} \(\sum c_{ab}x^a y^b=c_{ab}t^{an+bm}\)\par
\noindent\textit{Critical reading.} Use a full sum on the right and nonvanishing unit factors from \(x=t^n u(t)\), \(y=t^m v(t)\).\par
\noindent\textit{Evidence.} Both transcriptions preserve the missing summation and schematic expansion; exact order uses noncancellation.\par
\subsection*{P000957 (both; faithful error): equation13 boundary sum equals root count}
\noindent\textit{Transcription reads.} sum of minimum boundary orders\par
\noindent\textit{Critical reading.} negative sum of boundary orders\par
\noindent\textit{Evidence.} Principal-divisor degree zero and absence of torus poles give \(N=-\sum\operatorname{ord}_{\mathrm{boundary}}(Q)\); the scan confirms no minus sign.\par
\subsection*{P000958 (both; faithful error): mixed-volume definition followed by theorem 14}
\noindent\textit{Transcription reads.} \(V(\Delta,\ldots,\Delta)=\operatorname{Vol}(\Delta)\), then number of roots \(=V\)\par
\noindent\textit{Critical reading.} number of roots \(=n!V\); in the planar formulation \(2V=S\)\par
\noindent\textit{Evidence.} The scan moves between coefficient normalization \(S\) and diagonal normalization \(V\). \citet[\S1]{khovanskii1978genus} gives the factorial explicitly.\par
\subsection*{P000958 (both; faithful error): surface intersection inequality19}
\noindent\textit{Transcription reads.} any two curves on any algebraic surface\par
\noindent\textit{Critical reading.} nef divisor classes on a smooth projective surface suffice; arbitrary curves need not satisfy the inequality\par
\noindent\textit{Evidence.} Two disjoint \((-1)\)-exceptional curves give \(0\ge1\) in the literal assertion. Hodge index supplies the positive-cone statement.\par
\subsection*{P000958 (both; faithful error): theorem 15 after curve is changed to hypersurface}
\noindent\textit{Transcription reads.} Euler characteristic \(=-n!\operatorname{Vol}(\Delta)\)\par
\noindent\textit{Critical reading.} Euler characteristic \(=(-1)^{n-1}n!\operatorname{Vol}(\Delta)\), under Newton nondegeneracy\par
\noindent\textit{Evidence.} Actual scan preserves the minus sign after generalizing the word; \citet[\S3, Theorem~2]{khovanskii1978genus} supplies the dimension-dependent sign.\par
\subsection*{P000960 (both; faithful error): Stage I Lax pair}
\noindent\textit{Transcription reads.} \(L=D^2-u\); \(A=D^3-uD-Du\); \([L,A]=L_t\)\par
\noindent\textit{Critical reading.} One consistent choice for the printed PDE: \(L=D^2+u\), \(A=-4D^3-3(uD+Du)\), \(L_t=[L,A]\)\par
\noindent\textit{Evidence.} Autograph confirms formulas; Weyl multiplication gives the corrected identity.\par
\subsection*{P000960 (both; provenance): date heading for recovered records 510 and 511}
\noindent\textit{Transcription reads.} MCCME catalogue: 23-Feb-91\par
\noindent\textit{Critical reading.} 1981-02-23\par
\noindent\textit{Evidence.} Autograph 23/II/81 and uninterrupted surrounding 1981 entries; catalogue decade error.\par
\subsection*{P000961 (both; transcription): elliptic auxiliary system, second line}
\noindent\textit{Transcription reads.} \(E_t\,(J(\lambda)U)\times E\)\par
\noindent\textit{Critical reading.} \(E_t=(J(\lambda)U)\times E\)\par
\noindent\textit{Evidence.} Missing equality sign is visible in the autograph; mathematical normalization still requires separate checking.\par
\subsection*{P000961 (both; ambiguity): nonlinear two-field system, second equation}
\noindent\textit{Transcription reads.} \(U_t=(JU)\times V\)\par
\noindent\textit{Critical reading.} Preserve \(U_t\) with an uncertainty note; \(V_t\) is a plausible emendation, not established by the autograph.\par
\noindent\textit{Evidence.} Both readings are mathematically consequential; the manuscript writes \(U\).\par
\subsection*{P000961-P000962 (both; provenance): speaker and date provenance}
\noindent\textit{Transcription reads.} catalogue continuation included within record 509 or omitted because dated 1991\par
\noindent\textit{Critical reading.} record 512, A. G. Khovanskii, 23 February 1981, P961 lower through P962 upper\par
\noindent\textit{Evidence.} The actual Khovanskii heading is onP961 below the integrable-system discussion and above Kazarnovskii onP962; the meeting date onP960 is23/II/81.\par
\subsection*{P000962 (both; faithful error): exponential-system zero-count theorem}
\noindent\textit{Transcription reads.} \(N_F(R)=cVR^n+O(R^{n-1})\), without conditions\par
\noindent\textit{Critical reading.} Add Kazarnovskii regularity, multiplicities, and \(c=\kappa_n n!/(2\pi)^n\) in the real-volume normalization.\par
\noindent\textit{Evidence.} \citet[Definition~3 and Theorem~1]{kazarnovskii1981zeros}; \(e^z-1\) gives the independent calibration \(R/\pi+O(1)\).\par
\subsection*{P000962 (both; ambiguity): sigma-process chart}
\noindent\textit{Transcription reads.} the displayed rational inverse is the sigma-process\par
\noindent\textit{Critical reading.} Specify the blow-up morphism direction and include the second chart for the full blow-up.\par
\noindent\textit{Evidence.} The two displayed equations describe one standard affine blow-up chart; the inverse rational map fails at \(u_2=0\).\par
\subsection*{P000963 (both; faithful error): last Levi-eigenvalue integral labelled \(S(\Delta)\)}
\noindent\textit{Transcription reads.} \(S(\Delta)=\int\prod_i\kappa_i\,dS\)\par
\noindent\textit{Critical reading.} This is a body pseudovolume formula up to an explicit convention-dependent factor, not the dimensionless face projection factor \(S(\Delta)\).\par
\noindent\textit{Evidence.} The scan uses \(S(\Delta)\) both for projection distortion and for a dimensional boundary integral; the published formulas distinguish these quantities.\par
\subsection*{P000964 (both; provenance): question relating Kazarnovskii to growth indicatrix}
\noindent\textit{Transcription reads.} placed adjacent to the next meeting without attribution to its actual date\par
\noindent\textit{Critical reading.} I. M. Gelfand, 2 March 1981; a follow-up to record 513 rather than part of the 23 February speech\par
\noindent\textit{Evidence.} The EN/RU page begins2March and labels Gelfand. This is kept under513 only as an explicitly dated cross-record follow-up.\par
\subsection*{P000965 (both; faithful error): simultaneous \(\mathbb Z_2\) conjugation}
\noindent\textit{Transcription reads.} \(R\mapsto-(g^{-1}\otimes g^{-1})R(g\otimes g)\)\par
\noindent\textit{Critical reading.} Remove the leading minus for literal invariance; alternatively explicitly work projectively.\par
\noindent\textit{Evidence.} Autograph contains a leading minus; conjugation fixes the displayed nonzero eight-vertex matrix without that sign.\par
\subsection*{P000966 (both; faithful error): invariant level sets}
\noindent\textit{Transcription reads.} elliptic curves\par
\noindent\textit{Critical reading.} generic smooth projective common levels are genus-one curves\par
\noindent\textit{Evidence.} The projective intersection is defined by two homogeneous quadrics in \(\mathbb P^3\). Affine common levels are cones of dimension two, and special fibres may be singular.\par
\subsection*{P000966-P000967 (both; faithful error): relativistic scattering angle}
\noindent\textit{Transcription reads.} ordinary angle / hyperbolic cosine of momentum\par
\noindent\textit{Critical reading.} rapidity difference (with rapidity parametrizing the mass shell)\par
\noindent\textit{Evidence.} The diagram is Euclidean; Lorentz invariance uses a hyperbolic rapidity parameter, as in Belavin 1980 §1.\par
\subsection*{P000967 (both; faithful error): component Yang–Baxter relation (12), second factor on RHS}
\noindent\textit{Transcription reads.} \(S_{j_1k_3}^{k_1j_3}\)\par
\noindent\textit{Critical reading.} \(S_{i_1k_3}^{k_1j_3}\)\par
\noindent\textit{Evidence.} Autograph and both transcriptions repeat \(j_1\); free-index balance and the adjacent tensor relation require \(i_1\).\par
\subsection*{P000968 (both; faithful error): transfer-matrix theorem}
\noindent\textit{Transcription reads.} \([\operatorname{tr}_3 S_{13},\operatorname{tr}_3 S^{\prime}_{13}]=0\)\par
\noindent\textit{Critical reading.} \([\operatorname{tr}_1 S_{13},\operatorname{tr}_2 S^{\prime}_{23}]=0\), with invertible numeric \(R_{12}\)\par
\noindent\textit{Evidence.} Scan confirms \(\operatorname{tr}_3\); the exchange relation proves commutativity only after auxiliary traces.\par
\subsection*{P000969 (both; ambiguity): Grinevich cubic relation}
\noindent\textit{Transcription reads.} \(L_6^2=L_4^3+g_1L+g_0\)\par
\noindent\textit{Critical reading.} \(L_6^2=L_4^3+g_1L_4+g_0\)\par
\noindent\textit{Evidence.} The unscripted \(L\) is literal in the autograph; the displayed orders and elliptic spectral relation identify \(L_4\).\par
\subsection*{P000970 (both; faithful error): Dixmier sixth-order operator}
\noindent\textit{Transcription reads.} \(L_6=L^3+3x^2L+3Lx^2\)\par
\noindent\textit{Critical reading.} \(L_6=L^3+\tfrac32(xL+Lx)=L^3+3xL+3D\)\par
\noindent\textit{Evidence.} Both x-squares are clear in the scan. Formal positive-part calculation and direct Weyl multiplication verify the emendation.\par
\subsection*{P000970 (both; faithful error): Dixmier specialization}
\noindent\textit{Transcription reads.} \(\gamma_0=0\Rightarrow q(x)=x\)\par
\noindent\textit{Critical reading.} \(\gamma_0=0\) and \(q(x)=x\)\par
\noindent\textit{Evidence.} The arbitrary rational \(q\) is a separate choice; \citet[Theorem~3.4]{mironov2012higher} reproduces the rationality criterion and identifies the joint specialization.\par
\subsection*{P000970 (both; faithful error): formal pseudodifferential cube}
\noindent\textit{Transcription reads.} \(A=L+xL^{-1}\); \(L_6=(A^3)_+\)\par
\noindent\textit{Critical reading.} \((A^3)_+=L^3+3xL+6D\); use the full square root of \(L_4\), beginning \(L+xL^{-1}-DL^{-2}\), to obtain the commuting \(L_6\)\par
\noindent\textit{Evidence.} Independent formal positive-part expansion and the chapter’s executable Weyl certificate agree.\par
\subsection*{P000970 (both; provenance): speaker surname}
\noindent\textit{Transcription reads.} Berlinskii / Brelinskii / Berlinski (?)\par
\noindent\textit{Critical reading.} Jean-Luc Brylinski\par
\noindent\textit{Evidence.} The Kashiwara collaboration and the geometric Verma-module topic identify Brylinski, whose 1981 primary exposition is exact bibliographic evidence.\par
\subsection*{P000971 (both; ambiguity): definition of Kodaira dimension}
\noindent\textit{Transcription reads.} degree of the canonical graded algebra minus one\par
\noindent\textit{Critical reading.} transcendence degree minus one, with the no-positive-section case assigned minus infinity\par
\noindent\textit{Evidence.} The source abbreviation is imprecise; Esnault 1981 p111 supplies the exact definition.\par
\subsection*{P000971 (both; faithful error): positive-dimensional Albanese fibres}
\noindent\textit{Transcription reads.} étale fibration\par
\noindent\textit{Critical reading.} fibration trivialized by finite étale base change, with the theorem’s birational qualifications\par
\noindent\textit{Evidence.} An étale map has relative dimension zero; Esnault 1981 Thm2 is the explicitly named source.\par
\subsection*{P000972 (both; provenance): Soulé bibliographic shorthand}
\noindent\textit{Transcription reads.} Inventiones, February 1980\par
\noindent\textit{Critical reading.} Unresolved shorthand; the identifiable arithmetic K-theory paper is Inventiones55 (1979),251–295\par
\noindent\textit{Evidence.} Soulé’s publication list and the publisher’s October 1979 record were opened; no exact February 1980 paper matching the line was identified.\par
\subsection*{P000972 (both; ambiguity): numerator of zeta value}
\noindent\textit{Transcription reads.} numerator of \(\zeta(1-i)\)\par
\noindent\textit{Critical reading.} absolute numerator, if the inequality is intended as a positive lower bound\par
\noindent\textit{Evidence.} For example \(\zeta(-1)=-1/12\); the group order is nonnegative.\par
\subsection*{P000973 (both; ambiguity): invariant metric convention}
\noindent\textit{Transcription reads.} extend by left translations\par
\noindent\textit{Critical reading.} right translations for the usual Eulerian material-to-spatial convention; left translations after taking inverse maps\par
\noindent\textit{Evidence.} \citet{arnold1966geometrie} defines the fluid metric with the appropriate right invariance; the convention matters to the Euler equation.\par
\subsection*{P000975 (both; ambiguity): Maxwellian parameter count}
\noindent\textit{Transcription reads.} three-parameter family\par
\noindent\textit{Critical reading.} three parameter blocks, comprising five real scalar parameters in three dimensions\par
\noindent\textit{Evidence.} \(A\) and \(\alpha\) are scalar and \(v_0\) has three components; the scan confirms the shorthand count.\par
\subsection*{P000975 (both; faithful error): linearized collision operator}
\noindent\textit{Transcription reads.} \(L=-I+\text{compact operator}\) without qualifications\par
\noindent\textit{Critical reading.} angular-cutoff Maxwell model, suitable weighted space, and collision-frequency normalization \(\nu=1\)\par
\noindent\textit{Evidence.} Without angular cutoff the loss term is not a bounded constant multiple of identity; the collision model and space are omitted.\par
\subsection*{P000976 (both; ambiguity): final answer}
\noindent\textit{Transcription reads.} ordering plus subset inequalities\par
\noindent\textit{Critical reading.} ordering, subset inequalities, and the trace equality (1) on P000975\par
\noindent\textit{Evidence.} Trace is not repeated in the last formula but is indispensable; the source states it on the preceding page.\par
\subsection*{P000976 (both; faithful error): recursive consistency condition}
\noindent\textit{Transcription reads.} \(\ell<n\) below the smaller lists\par
\noindent\textit{Critical reading.} \(k<\ell\) for the smaller-list recursion (and \(\ell<n\) for the final spectral inequalities)\par
\noindent\textit{Evidence.} The scanned inequality \(\ell<n\) does not make the definition inductive; the smaller triples belong to \(T_k^{\ell}\).\par
\subsection*{P000977 (both; ambiguity): barotropic assumption}
\noindent\textit{Transcription reads.} \(T\) is constant\par
\noindent\textit{Critical reading.} Barotropy \(p=p(\rho)\) does not by itself impose constant temperature\par
\noindent\textit{Evidence.} The separate line is literal; polytropic isentropic and isothermal laws are different closures.\par
\subsection*{P000977-P000978 (both; faithful error): shock inevitability}
\noindent\textit{Transcription reads.} shock waves necessarily arise / discontinuities always needed\par
\noindent\textit{Critical reading.} compression can produce shocks; constant and suitable expanding smooth data do not\par
\noindent\textit{Evidence.} The characteristic map \(x=\xi+t f(\xi)\) fails only when \(1+t f^{\prime}(\xi)=0\).\par
\subsection*{P000978 (both; faithful error): Lobachevsky-plane assertion}
\noindent\textit{Transcription reads.} does not embed in \(\mathbb R^3\)\par
\noindent\textit{Critical reading.} has no complete smooth isometric immersion in Euclidean \(\mathbb R^3\)\par
\noindent\textit{Evidence.} Without isometry, the hyperbolic plane is diffeomorphic to a disc and embeds.\par
\subsection*{P000979-P000980 (both; faithful error): invariant-cone existence criterion}
\noindent\textit{Transcription reads.} convex invariant cone with nonempty interior\par
\noindent\textit{Critical reading.} closed pointed generating nonzero invariant convex cone, in a noncompact simple real Lie algebra\par
\noindent\textit{Evidence.} Otherwise the whole Lie algebra is a counterexample to the claimed Hermitian criterion; \citet[\S1]{olshanskii1981cones} specifies nondegeneracy.\par
\subsection*{P000980 (both; faithful error): irreducible factorization}
\noindent\textit{Transcription reads.} irreducible equals holomorphic tensor antiholomorphic for \(\Gamma(C)\)\par
\noindent\textit{Critical reading.} the checked theorem is for \(\Gamma(C_{\min})\), with its covering and contractive representation category\par
\noindent\textit{Evidence.} \citet[Corollary~5.8]{olshanskii1981cones} explicitly uses the minimal cone.\par
\subsection*{P000980 (both; ambiguity): maximal domain of extension}
\noindent\textit{Transcription reads.} exact domain of extendibility\par
\noindent\textit{Critical reading.} maximal domain of contractive holomorphic semigroup extension in the stated category\par
\noindent\textit{Evidence.} The generator inequality encodes contractivity; it does not prohibit every other analytic continuation of matrix coefficients.\par
\subsection*{P000980 (both; faithful error): unique extremal cones}
\noindent\textit{Transcription reads.} unique \(C_{\min}\) and \(C_{\max}\)\par
\noindent\textit{Critical reading.} unique positive extremal cones after choosing a central orientation\par
\noindent\textit{Evidence.} Their negatives give the opposite pair; \citet[\S1]{olshanskii1981cones} fixes positivity.\par
